% DeskFlow - listing kodu źródłowego backendu, testów i konfiguracji Docker
% Auto-generated. Do not edit manually.


% Ustawienia wymagane dla czytelnego listingu kodu.
% Wariant oparty o fvextra lepiej zachowuje wciecia zawinietych linii
% niz standardowe lstlisting z postbreak. Zawiniete fragmenty nie dostaja
% strzalek ani dodatkowych znakow i wracaja do wciecia poczatku linii.
% W preambule glownego pliku .tex dodaj: \usepackage{fvextra} oraz \usepackage{pdflscape}
% oraz zostaw/dodaj: \usepackage{xcolor}, \usepackage{hyperref}
\providecommand{\phantomsection}{}
\makeatletter
\@ifundefined{c@lstlisting}{%
  \newcounter{lstlisting}[chapter]
  \renewcommand{\thelstlisting}{\thechapter.\arabic{lstlisting}}%
}{}
\providecommand{\lstlistingname}{Listing}
\newcommand{\DeskFlowListingCaption}[2]{%
  \refstepcounter{lstlisting}%
  \label{#2}%
  \noindent\begin{minipage}{\linewidth}%
  \textbf{\lstlistingname~\thelstlisting:}\quad\raggedright #1%
  \end{minipage}\par\nobreak\vspace{2pt}%
}
\makeatother
\DefineVerbatimEnvironment{DeskFlowCode}{Verbatim}{%
  numbers=left,%
  numbersep=8pt,%
  fontsize=\scriptsize,%
  breaklines=true,%
  breakanywhere=false,%
  breakautoindent=true,%
  breakindentnchars=0,%
  breaksymbolleft={},%
  breaksymbolright={},%
  frame=single,%
  framesep=3pt,%
  tabsize=4,%
  obeytabs=true%
}

% Wpisy klas/listingow beda dodawane do spisu tresci jako subsubsection.
% Ta linia sprawia, ze sa widoczne w spisie tresci po kolejnej kompilacji.
\addtocontents{toc}{\protect\setcounter{tocdepth}{3}}

\Chapter{Listing kodu źródłowego, testów i skryptów wdrożeniowych}\label{chapter:listing}


\begin{landscape}
\section{Kod źródłowy backendu}

\phantomsection
\addcontentsline{toc}{subsubsection}{DeskflowApplication.java}
\DeskFlowListingCaption{Plik główny: DeskflowApplication.java}{lst:DeskflowApplication-java}
\begin{DeskFlowCode}
package com.deskflow;

import org.springframework.boot.SpringApplication;
import org.springframework.boot.autoconfigure.SpringBootApplication;
import org.springframework.scheduling.annotation.EnableScheduling;

@SpringBootApplication
@EnableScheduling
public class DeskflowApplication {
    public static void main(String[] args) {
        SpringApplication.run(DeskflowApplication.class, args);
    }
}
\end{DeskFlowCode}

\subsection{Kontekst Allocation - alokacja zasobów}

\phantomsection
\addcontentsline{toc}{subsubsection}{AllocationRequest}
\DeskFlowListingCaption{allocation/application/domain/model/AllocationRequest.java - AllocationRequest}{lst:allocation-application-domain-model-AllocationRequest-java}
\begin{DeskFlowCode}
package com.deskflow.allocation.application.domain.model;

import com.deskflow.common.application.domain.model.vo.EmployeeId;
import com.deskflow.common.application.domain.DomainEvent;
import com.deskflow.allocation.application.domain.model.events.*;
import com.deskflow.allocation.application.domain.model.vo.AllocationRequestId;
import com.deskflow.allocation.application.domain.model.vo.DateRange;
import com.deskflow.allocation.application.domain.model.vo.RequestStatus;
import com.deskflow.common.application.domain.model.vo.ResourceId;

import java.time.LocalDateTime;
import java.util.ArrayList;
import java.util.List;

public class AllocationRequest {

    private final AllocationRequestId id;
    private final EmployeeId employeeId;
    private final ResourceId resourceId;
    private final DateRange dateRange;
    private final String justification;
    private final LocalDateTime submittedAt;
    private RequestStatus status;
    private LocalDateTime issuedAt;
    private LocalDateTime returnedAt;
    private String rejectionReason;
    private final List<DomainEvent> events = new ArrayList<>();

    private AllocationRequest(AllocationRequestId id, EmployeeId employeeId, ResourceId resourceId,
                              DateRange dateRange, String justification, LocalDateTime submittedAt,
                              RequestStatus status, LocalDateTime issuedAt,
                              LocalDateTime returnedAt, String rejectionReason) {
        this.id = id;
        this.employeeId = employeeId;
        this.resourceId = resourceId;
        this.dateRange = dateRange;
        this.justification = justification;
        this.submittedAt = submittedAt;
        this.status = status;
        this.issuedAt = issuedAt;
        this.returnedAt = returnedAt;
        this.rejectionReason = rejectionReason;
    }

    public static AllocationRequest create(AllocationRequestId id, EmployeeId employeeId,
                                            ResourceId resourceId, DateRange dateRange,
                                            String justification, LocalDateTime submittedAt) {
        var r = new AllocationRequest(id, employeeId, resourceId, dateRange, justification, submittedAt,
                RequestStatus.SUBMITTED, null, null, null);
        r.events.add(new AllocationRequestSubmitted(id, employeeId, resourceId, dateRange, submittedAt));
        return r;
    }

    public static AllocationRequest reconstitute(AllocationRequestId id, EmployeeId employeeId,
                                                  ResourceId resourceId, DateRange dateRange,
                                                  String justification, LocalDateTime submittedAt,
                                                  RequestStatus status, LocalDateTime issuedAt,
                                                  LocalDateTime returnedAt, String rejectionReason) {
        return new AllocationRequest(id, employeeId, resourceId, dateRange, justification, submittedAt,
                status, issuedAt, returnedAt, rejectionReason);
    }

    public void approve() {
        requireStatus(RequestStatus.SUBMITTED, "approve");
        this.status = RequestStatus.APPROVED;
        events.add(new AllocationRequestApproved(id, employeeId, LocalDateTime.now()));
    }

    public void reject(String reason) {
        requireStatus(RequestStatus.SUBMITTED, "reject");
        this.status = RequestStatus.REJECTED;
        this.rejectionReason = reason;
        events.add(new AllocationRequestRejected(id, employeeId, reason, LocalDateTime.now()));
    }

    public void cancel() {
        if (status != RequestStatus.SUBMITTED && status != RequestStatus.APPROVED) {
            throw new IllegalStateException(
                    "Can only cancel a SUBMITTED or APPROVED request, current: " + status);
        }
        this.status = RequestStatus.CANCELLED;
        events.add(new AllocationRequestCancelled(id, employeeId, LocalDateTime.now()));
    }

    public void issueResource() {
        requireStatus(RequestStatus.APPROVED, "issue resource");
        this.status = RequestStatus.ACTIVE;
        this.issuedAt = LocalDateTime.now();
        events.add(new ResourceIssued(id, employeeId, resourceId, issuedAt));
    }

    public void returnResource() {
        if (status != RequestStatus.ACTIVE && status != RequestStatus.OVERDUE) {
            throw new IllegalStateException(
                    "Can only return resource for ACTIVE or OVERDUE request, current: " + status);
        }
        this.status = RequestStatus.COMPLETED;
        this.returnedAt = LocalDateTime.now();
        events.add(new ResourceReturned(id, employeeId, resourceId, returnedAt));
    }

    public void markAsOverdue() {
        requireStatus(RequestStatus.ACTIVE, "mark as overdue");
        this.status = RequestStatus.OVERDUE;
        events.add(new AllocationOverdue(id, employeeId, resourceId, LocalDateTime.now()));
    }

    private void requireStatus(RequestStatus required, String operation) {
        if (status != required) {
            throw new IllegalStateException(
                    "Operation '" + operation + "' requires status " + required + ", current: " + status);
        }
    }

    public List<DomainEvent> flushEvents() {
        var copy = List.copyOf(events);
        events.clear();
        return copy;
    }

    public AllocationRequestId getId() { return id; }
    public EmployeeId getEmployeeId() { return employeeId; }
    public ResourceId getResourceId() { return resourceId; }
    public DateRange getDateRange() { return dateRange; }
    public String getJustification() { return justification; }
    public LocalDateTime getSubmittedAt() { return submittedAt; }
    public RequestStatus getStatus() { return status; }
    public LocalDateTime getIssuedAt() { return issuedAt; }
    public LocalDateTime getReturnedAt() { return returnedAt; }
    public String getRejectionReason() { return rejectionReason; }
}
\end{DeskFlowCode}

\phantomsection
\addcontentsline{toc}{subsubsection}{EmployeeAllocationProfile}
\DeskFlowListingCaption{allocation/application/domain/model/EmployeeAllocationProfile.java - EmployeeAllocationProfile}{lst:allocation-application-domain-model-EmployeeAllocationProfile-java}
\begin{DeskFlowCode}
package com.deskflow.allocation.application.domain.model;

import java.util.Map;

public record EmployeeAllocationProfile(Map<String, Integer> limitsByResourceType) {
    public int limitForType(String resourceType) {
        return limitsByResourceType.getOrDefault(resourceType, 1);
    }
}
\end{DeskFlowCode}

\phantomsection
\addcontentsline{toc}{subsubsection}{ResourceAllocationProfile}
\DeskFlowListingCaption{allocation/application/domain/model/ResourceAllocationProfile.java - ResourceAllocationProfile}{lst:allocation-application-domain-model-ResourceAllocationProfile-java}
\begin{DeskFlowCode}
package com.deskflow.allocation.application.domain.model;

import com.deskflow.common.application.domain.model.vo.ResourceId;

public record ResourceAllocationProfile(
        ResourceId id,
        String type,
        String name,
        String building,
        int floor,
        String number,
        boolean active,
        boolean available
) {
    public boolean isActive() {
        return active;
    }

    public boolean isAvailable() {
        return available;
    }
}
\end{DeskFlowCode}

\phantomsection
\addcontentsline{toc}{subsubsection}{AllocationOverdue}
\DeskFlowListingCaption{allocation/application/domain/model/events/AllocationOverdue.java - AllocationOverdue}{lst:allocation-application-domain-model-events-AllocationOverdue-java}
\begin{DeskFlowCode}
package com.deskflow.allocation.application.domain.model.events;

import com.deskflow.common.application.domain.model.vo.EmployeeId;
import com.deskflow.common.application.domain.DomainEvent;
import com.deskflow.allocation.application.domain.model.vo.AllocationRequestId;
import com.deskflow.common.application.domain.model.vo.ResourceId;

import java.time.LocalDateTime;
import java.util.UUID;

public record AllocationOverdue(
        AllocationRequestId requestId,
        EmployeeId employeeId,
        ResourceId resourceId,
        LocalDateTime occurredAt
) implements DomainEvent {
    @Override
    public UUID eventId() { return UUID.randomUUID(); }
    @Override
    public String eventType() { return "AllocationOverdue"; }
    @Override
    public UUID aggregateId() { return requestId.getValue(); }
}
\end{DeskFlowCode}

\phantomsection
\addcontentsline{toc}{subsubsection}{AllocationRequestApproved}
\DeskFlowListingCaption{allocation/application/domain/model/events/AllocationRequestApproved.java - AllocationRequestApproved}{lst:allocation-application-domain-model-events-AllocationRequestApproved-java}
\begin{DeskFlowCode}
package com.deskflow.allocation.application.domain.model.events;

import com.deskflow.common.application.domain.model.vo.EmployeeId;
import com.deskflow.common.application.domain.DomainEvent;
import com.deskflow.allocation.application.domain.model.vo.AllocationRequestId;

import java.time.LocalDateTime;
import java.util.UUID;

public record AllocationRequestApproved(
        AllocationRequestId requestId,
        EmployeeId employeeId,
        LocalDateTime occurredAt
) implements DomainEvent {
    @Override
    public UUID eventId() { return UUID.randomUUID(); }
    @Override
    public String eventType() { return "AllocationRequestApproved"; }
    @Override
    public UUID aggregateId() { return requestId.getValue(); }
}
\end{DeskFlowCode}

\phantomsection
\addcontentsline{toc}{subsubsection}{AllocationRequestCancelled}
\DeskFlowListingCaption{allocation/application/domain/model/events/AllocationRequestCancelled.java - AllocationRequestCancelled}{lst:allocation-application-domain-model-events-AllocationRequestCancelled-java}
\begin{DeskFlowCode}
package com.deskflow.allocation.application.domain.model.events;

import com.deskflow.common.application.domain.model.vo.EmployeeId;
import com.deskflow.common.application.domain.DomainEvent;
import com.deskflow.allocation.application.domain.model.vo.AllocationRequestId;

import java.time.LocalDateTime;
import java.util.UUID;

public record AllocationRequestCancelled(
        AllocationRequestId requestId,
        EmployeeId employeeId,
        LocalDateTime occurredAt
) implements DomainEvent {
    @Override
    public UUID eventId() { return UUID.randomUUID(); }
    @Override
    public String eventType() { return "AllocationRequestCancelled"; }
    @Override
    public UUID aggregateId() { return requestId.getValue(); }
}
\end{DeskFlowCode}

\phantomsection
\addcontentsline{toc}{subsubsection}{AllocationRequestRejected}
\DeskFlowListingCaption{allocation/application/domain/model/events/AllocationRequestRejected.java - AllocationRequestRejected}{lst:allocation-application-domain-model-events-AllocationRequestRejected-java}
\begin{DeskFlowCode}
package com.deskflow.allocation.application.domain.model.events;

import com.deskflow.common.application.domain.model.vo.EmployeeId;
import com.deskflow.common.application.domain.DomainEvent;
import com.deskflow.allocation.application.domain.model.vo.AllocationRequestId;

import java.time.LocalDateTime;
import java.util.UUID;

public record AllocationRequestRejected(
        AllocationRequestId requestId,
        EmployeeId employeeId,
        String reason,
        LocalDateTime occurredAt
) implements DomainEvent {
    @Override
    public UUID eventId() { return UUID.randomUUID(); }
    @Override
    public String eventType() { return "AllocationRequestRejected"; }
    @Override
    public UUID aggregateId() { return requestId.getValue(); }
}
\end{DeskFlowCode}

\phantomsection
\addcontentsline{toc}{subsubsection}{AllocationRequestSubmitted}
\DeskFlowListingCaption{allocation/application/domain/model/events/AllocationRequestSubmitted.java - AllocationRequestSubmitted}{lst:allocation-application-domain-model-events-AllocationRequestSubmitted-java}
\begin{DeskFlowCode}
package com.deskflow.allocation.application.domain.model.events;

import com.deskflow.common.application.domain.model.vo.EmployeeId;
import com.deskflow.common.application.domain.DomainEvent;
import com.deskflow.allocation.application.domain.model.vo.AllocationRequestId;
import com.deskflow.allocation.application.domain.model.vo.DateRange;
import com.deskflow.common.application.domain.model.vo.ResourceId;

import java.time.LocalDateTime;
import java.util.UUID;

public record AllocationRequestSubmitted(
        AllocationRequestId requestId,
        EmployeeId employeeId,
        ResourceId resourceId,
        DateRange dateRange,
        LocalDateTime occurredAt
) implements DomainEvent {
    @Override
    public UUID eventId() { return UUID.randomUUID(); }
    @Override
    public String eventType() { return "AllocationRequestSubmitted"; }
    @Override
    public UUID aggregateId() { return requestId.getValue(); }
}
\end{DeskFlowCode}

\phantomsection
\addcontentsline{toc}{subsubsection}{ResourceIssued}
\DeskFlowListingCaption{allocation/application/domain/model/events/ResourceIssued.java - ResourceIssued}{lst:allocation-application-domain-model-events-ResourceIssued-java}
\begin{DeskFlowCode}
package com.deskflow.allocation.application.domain.model.events;

import com.deskflow.common.application.domain.model.vo.EmployeeId;
import com.deskflow.common.application.domain.DomainEvent;
import com.deskflow.allocation.application.domain.model.vo.AllocationRequestId;
import com.deskflow.common.application.domain.model.vo.ResourceId;

import java.time.LocalDateTime;
import java.util.UUID;

public record ResourceIssued(
        AllocationRequestId requestId,
        EmployeeId employeeId,
        ResourceId resourceId,
        LocalDateTime occurredAt
) implements DomainEvent {
    @Override
    public UUID eventId() { return UUID.randomUUID(); }
    @Override
    public String eventType() { return "ResourceIssued"; }
    @Override
    public UUID aggregateId() { return requestId.getValue(); }
}
\end{DeskFlowCode}

\phantomsection
\addcontentsline{toc}{subsubsection}{ResourceReturned}
\DeskFlowListingCaption{allocation/application/domain/model/events/ResourceReturned.java - ResourceReturned}{lst:allocation-application-domain-model-events-ResourceReturned-java}
\begin{DeskFlowCode}
package com.deskflow.allocation.application.domain.model.events;

import com.deskflow.common.application.domain.model.vo.EmployeeId;
import com.deskflow.common.application.domain.DomainEvent;
import com.deskflow.allocation.application.domain.model.vo.AllocationRequestId;
import com.deskflow.common.application.domain.model.vo.ResourceId;

import java.time.LocalDateTime;
import java.util.UUID;

public record ResourceReturned(
        AllocationRequestId requestId,
        EmployeeId employeeId,
        ResourceId resourceId,
        LocalDateTime occurredAt
) implements DomainEvent {
    @Override
    public UUID eventId() { return UUID.randomUUID(); }
    @Override
    public String eventType() { return "ResourceReturned"; }
    @Override
    public UUID aggregateId() { return requestId.getValue(); }
}
\end{DeskFlowCode}

\phantomsection
\addcontentsline{toc}{subsubsection}{AllocationRequestId}
\DeskFlowListingCaption{allocation/application/domain/model/vo/AllocationRequestId.java - AllocationRequestId}{lst:allocation-application-domain-model-vo-AllocationRequestId-java}
\begin{DeskFlowCode}
package com.deskflow.allocation.application.domain.model.vo;

import java.util.Objects;
import java.util.UUID;

public final class AllocationRequestId {
    private final UUID value;

    public AllocationRequestId(UUID value) {
        this.value = Objects.requireNonNull(value);
    }

    public static AllocationRequestId of(UUID value) {
        return new AllocationRequestId(value);
    }

    public static AllocationRequestId of(String value) {
        return new AllocationRequestId(UUID.fromString(value));
    }

    public static AllocationRequestId newId() {
        return new AllocationRequestId(UUID.randomUUID());
    }

    public UUID getValue() {
        return value;
    }

    @Override
    public boolean equals(Object o) {
        if (this == o) return true;
        if (!(o instanceof AllocationRequestId that)) return false;
        return value.equals(that.value);
    }

    @Override
    public int hashCode() {
        return value.hashCode();
    }

    @Override
    public String toString() {
        return value.toString();
    }
}
\end{DeskFlowCode}

\phantomsection
\addcontentsline{toc}{subsubsection}{DateRange}
\DeskFlowListingCaption{allocation/application/domain/model/vo/DateRange.java - DateRange}{lst:allocation-application-domain-model-vo-DateRange-java}
\begin{DeskFlowCode}
package com.deskflow.allocation.application.domain.model.vo;

import java.time.LocalDate;
import java.util.Objects;

public final class DateRange {
    private final LocalDate dateFrom;
    private final LocalDate dateTo;

    public DateRange(LocalDate dateFrom, LocalDate dateTo) {
        Objects.requireNonNull(dateFrom, "dateFrom cannot be null");
        Objects.requireNonNull(dateTo, "dateTo cannot be null");
        if (!dateFrom.isBefore(dateTo)) {
            throw new IllegalArgumentException("dateFrom must be before dateTo");
        }
        this.dateFrom = dateFrom;
        this.dateTo = dateTo;
    }

    public LocalDate getDateFrom() { return dateFrom; }
    public LocalDate getDateTo() { return dateTo; }

    @Override
    public boolean equals(Object o) {
        if (this == o) return true;
        if (!(o instanceof DateRange that)) return false;
        return dateFrom.equals(that.dateFrom) && dateTo.equals(that.dateTo);
    }

    @Override
    public int hashCode() {
        return Objects.hash(dateFrom, dateTo);
    }

    @Override
    public String toString() {
        return dateFrom + " — " + dateTo;
    }
}
\end{DeskFlowCode}

\phantomsection
\addcontentsline{toc}{subsubsection}{RequestStatus}
\DeskFlowListingCaption{allocation/application/domain/model/vo/RequestStatus.java - RequestStatus}{lst:allocation-application-domain-model-vo-RequestStatus-java}
\begin{DeskFlowCode}
package com.deskflow.allocation.application.domain.model.vo;

public enum RequestStatus {
    SUBMITTED,
    APPROVED,
    ACTIVE,
    COMPLETED,
    REJECTED,
    CANCELLED,
    OVERDUE
}
\end{DeskFlowCode}

\phantomsection
\addcontentsline{toc}{subsubsection}{AllocationRequestFactory}
\DeskFlowListingCaption{allocation/application/domain/service/AllocationRequestFactory.java - AllocationRequestFactory}{lst:allocation-application-domain-service-AllocationRequestFactory-java}
\begin{DeskFlowCode}
package com.deskflow.allocation.application.domain.service;

import com.deskflow.allocation.application.domain.model.AllocationRequest;
import com.deskflow.common.application.domain.model.vo.EmployeeId;
import com.deskflow.allocation.application.domain.model.vo.AllocationRequestId;
import com.deskflow.allocation.application.domain.model.vo.DateRange;
import com.deskflow.common.application.domain.model.vo.ResourceId;

import java.time.LocalDateTime;

public class AllocationRequestFactory {

    public AllocationRequest create(EmployeeId employeeId, ResourceId resourceId,
                                    DateRange dateRange, String justification) {
        return AllocationRequest.create(
                AllocationRequestId.newId(),
                employeeId,
                resourceId,
                dateRange,
                justification,
                LocalDateTime.now()
        );
    }
}
\end{DeskFlowCode}

\phantomsection
\addcontentsline{toc}{subsubsection}{ApproveAllocationRequestDomainService}
\DeskFlowListingCaption{allocation/application/domain/service/ApproveAllocationRequestDomainService.java - ApproveAllocationRequestDomainService}{lst:allocation-application-domain-service-ApproveAllocationRequestDomainService-java}
\begin{DeskFlowCode}
package com.deskflow.allocation.application.domain.service;

import com.deskflow.allocation.application.domain.model.AllocationRequest;

public class ApproveAllocationRequestDomainService {

    public void approveAllocationRequest(AllocationRequest request) {
        request.approve();
    }
}
\end{DeskFlowCode}

\phantomsection
\addcontentsline{toc}{subsubsection}{CancelAllocationRequestDomainService}
\DeskFlowListingCaption{allocation/application/domain/service/CancelAllocationRequestDomainService.java - CancelAllocationRequestDomainService}{lst:allocation-application-domain-service-CancelAllocationRequestDomainService-java}
\begin{DeskFlowCode}
package com.deskflow.allocation.application.domain.service;

import com.deskflow.allocation.application.domain.model.AllocationRequest;
import com.deskflow.common.application.domain.model.vo.EmployeeId;

public class CancelAllocationRequestDomainService {

    public void cancelAllocationRequest(AllocationRequest request, EmployeeId employeeId) {
        if (!request.getEmployeeId().equals(employeeId)) {
            throw new IllegalStateException("Cannot cancel another employee's request");
        }
        request.cancel();
    }
}
\end{DeskFlowCode}

\phantomsection
\addcontentsline{toc}{subsubsection}{EnsureEmployeeWithinLimitDomainService}
\DeskFlowListingCaption{allocation/application/domain/service/EnsureEmployeeWithinLimitDomainService.java - EnsureEmployeeWithinLimitDomainService}{lst:allocation-application-domain-service-EnsureEmployeeWithinLimitDomainService-java}
\begin{DeskFlowCode}
package com.deskflow.allocation.application.domain.service;

import com.deskflow.allocation.application.domain.model.EmployeeAllocationProfile;

public class EnsureEmployeeWithinLimitDomainService {

    public void ensureEmployeeWithinLimit(EmployeeAllocationProfile employee, String resourceType, long currentActiveForType) {
        int limit = employee.limitForType(resourceType);
        if (currentActiveForType >= limit) {
            throw new IllegalStateException("Employee has reached the allocation limit for this resource type");
        }
    }
}
\end{DeskFlowCode}

\phantomsection
\addcontentsline{toc}{subsubsection}{EnsureResourceAvailableInPeriodDomainService}
\DeskFlowListingCaption{allocation/application/domain/service/EnsureResourceAvailableInPeriodDomainService.java - EnsureResourceAvailableInPeriodDomainService}{lst:allocation-application-domain-service-EnsureResourceAvailableInPeriodDomainService-java}
\begin{DeskFlowCode}
package com.deskflow.allocation.application.domain.service;

import com.deskflow.allocation.application.domain.model.AllocationRequest;

import java.util.List;

public class EnsureResourceAvailableInPeriodDomainService {

    public void ensureResourceAvailableInPeriod(List<AllocationRequest> overlappingRequests) {
        if (!overlappingRequests.isEmpty()) {
            throw new IllegalStateException("Resource is not available for the requested period");
        }
    }
}
\end{DeskFlowCode}

\phantomsection
\addcontentsline{toc}{subsubsection}{EnsureResourceEligibleForAllocationDomainService}
\DeskFlowListingCaption{allocation/application/domain/service/EnsureResourceEligibleForAllocationDomainService.java - EnsureResourceEligibleForAllocationDomainService}{lst:allocation-application-domain-service-EnsureResourceEligibleForAllocationDomainService-java}
\begin{DeskFlowCode}
package com.deskflow.allocation.application.domain.service;

import com.deskflow.allocation.application.domain.model.ResourceAllocationProfile;

public class EnsureResourceEligibleForAllocationDomainService {

    public void ensureResourceEligibleForAllocation(ResourceAllocationProfile resource) {
        if (!resource.isActive() || !resource.isAvailable()) {
            throw new IllegalStateException("Resource is not available for allocation");
        }
    }
}
\end{DeskFlowCode}

\phantomsection
\addcontentsline{toc}{subsubsection}{IssueAllocationResourceDomainService}
\DeskFlowListingCaption{allocation/application/domain/service/IssueAllocationResourceDomainService.java - IssueAllocationResourceDomainService}{lst:allocation-application-domain-service-IssueAllocationResourceDomainService-java}
\begin{DeskFlowCode}
package com.deskflow.allocation.application.domain.service;

import com.deskflow.allocation.application.domain.model.AllocationRequest;

public class IssueAllocationResourceDomainService {

    public void issueAllocationResource(AllocationRequest request) {
        request.issueResource();
    }
}
\end{DeskFlowCode}

\phantomsection
\addcontentsline{toc}{subsubsection}{MarkAllocationRequestOverdueDomainService}
\DeskFlowListingCaption{allocation/application/domain/service/MarkAllocationRequestOverdueDomainService.java - MarkAllocationRequestOverdueDomainService}{lst:allocation-application-domain-service-MarkAllocationRequestOverdueDomainService-java}
\begin{DeskFlowCode}
package com.deskflow.allocation.application.domain.service;

import com.deskflow.allocation.application.domain.model.AllocationRequest;

public class MarkAllocationRequestOverdueDomainService {

    public void markAllocationRequestOverdue(AllocationRequest request) {
        request.markAsOverdue();
    }
}
\end{DeskFlowCode}

\phantomsection
\addcontentsline{toc}{subsubsection}{RejectAllocationRequestDomainService}
\DeskFlowListingCaption{allocation/application/domain/service/RejectAllocationRequestDomainService.java - RejectAllocationRequestDomainService}{lst:allocation-application-domain-service-RejectAllocationRequestDomainService-java}
\begin{DeskFlowCode}
package com.deskflow.allocation.application.domain.service;

import com.deskflow.allocation.application.domain.model.AllocationRequest;

public class RejectAllocationRequestDomainService {

    public void rejectAllocationRequest(AllocationRequest request, String reason) {
        if (reason == null || reason.isBlank()) {
            throw new IllegalArgumentException("Rejection reason cannot be blank");
        }
        request.reject(reason);
    }
}
\end{DeskFlowCode}

\phantomsection
\addcontentsline{toc}{subsubsection}{ReturnAllocationResourceDomainService}
\DeskFlowListingCaption{allocation/application/domain/service/ReturnAllocationResourceDomainService.java - ReturnAllocationResourceDomainService}{lst:allocation-application-domain-service-ReturnAllocationResourceDomainService-java}
\begin{DeskFlowCode}
package com.deskflow.allocation.application.domain.service;

import com.deskflow.allocation.application.domain.model.AllocationRequest;

public class ReturnAllocationResourceDomainService {

    public void returnAllocationResource(AllocationRequest request) {
        request.returnResource();
    }
}
\end{DeskFlowCode}

\phantomsection
\addcontentsline{toc}{subsubsection}{AllocationRequestView}
\DeskFlowListingCaption{allocation/application/model/AllocationRequestView.java - AllocationRequestView}{lst:allocation-application-model-AllocationRequestView-java}
\begin{DeskFlowCode}
package com.deskflow.allocation.application.model;

import com.deskflow.allocation.application.domain.model.AllocationRequest;

import java.time.LocalDate;
import java.time.LocalDateTime;
import java.util.UUID;

public record AllocationRequestView(
        UUID id,
        UUID employeeId,
        UUID resourceId,
        String status,
        LocalDate dateFrom,
        LocalDate dateTo,
        String justification,
        LocalDateTime submittedAt,
        LocalDateTime issuedAt,
        LocalDateTime returnedAt,
        String rejectionReason
) {
    public static AllocationRequestView from(AllocationRequest request) {
        return new AllocationRequestView(
                request.getId().getValue(),
                request.getEmployeeId().getValue(),
                request.getResourceId().getValue(),
                request.getStatus().name(),
                request.getDateRange().getDateFrom(),
                request.getDateRange().getDateTo(),
                request.getJustification(),
                request.getSubmittedAt(),
                request.getIssuedAt(),
                request.getReturnedAt(),
                request.getRejectionReason()
        );
    }
}
\end{DeskFlowCode}

\phantomsection
\addcontentsline{toc}{subsubsection}{ResourceCatalogItem}
\DeskFlowListingCaption{allocation/application/model/ResourceCatalogItem.java - ResourceCatalogItem}{lst:allocation-application-model-ResourceCatalogItem-java}
\begin{DeskFlowCode}
package com.deskflow.allocation.application.model;

import com.deskflow.allocation.application.domain.model.ResourceAllocationProfile;

import java.util.UUID;

public record ResourceCatalogItem(
        UUID id,
        String type,
        String name,
        String building,
        int floor,
        String number,
        boolean available
) {
    public static ResourceCatalogItem from(ResourceAllocationProfile resource) {
        return new ResourceCatalogItem(
                resource.id().getValue(),
                resource.type(),
                resource.name(),
                resource.building(),
                resource.floor(),
                resource.number(),
                resource.available()
        );
    }
}
\end{DeskFlowCode}

\phantomsection
\addcontentsline{toc}{subsubsection}{AllocationAdministrationInputPort}
\DeskFlowListingCaption{allocation/application/port/in/AllocationAdministrationInputPort.java - AllocationAdministrationInputPort}{lst:allocation-application-port-in-AllocationAdministrationInputPort-java}
\begin{DeskFlowCode}
package com.deskflow.allocation.application.port.in;

import com.deskflow.allocation.application.model.AllocationRequestView;
import com.deskflow.allocation.application.domain.model.vo.RequestStatus;
import com.deskflow.allocation.application.domain.model.vo.AllocationRequestId;

import java.util.List;
import java.util.Optional;

public interface AllocationAdministrationInputPort {
    Optional<AllocationRequestView> getRequest(AllocationRequestId requestId);

    List<AllocationRequestView> listRequests(RequestStatus status);

    void approveRequest(AllocationRequestId requestId);

    void rejectRequest(AllocationRequestId requestId, String reason);

    void issueResource(AllocationRequestId requestId);

    void returnResource(AllocationRequestId requestId);

    List<AllocationRequestView> monitorActiveAllocations();

    void detectOverdueAllocations();
}
\end{DeskFlowCode}

\phantomsection
\addcontentsline{toc}{subsubsection}{AllocationRequestInputPort}
\DeskFlowListingCaption{allocation/application/port/in/AllocationRequestInputPort.java - AllocationRequestInputPort}{lst:allocation-application-port-in-AllocationRequestInputPort-java}
\begin{DeskFlowCode}
package com.deskflow.allocation.application.port.in;

import com.deskflow.allocation.application.model.AllocationRequestView;
import com.deskflow.allocation.application.model.ResourceCatalogItem;
import com.deskflow.allocation.application.domain.model.vo.AllocationRequestId;
import com.deskflow.common.application.domain.model.vo.EmployeeId;
import com.deskflow.common.application.domain.model.vo.ResourceId;

import java.time.LocalDate;
import java.util.List;

public interface AllocationRequestInputPort {

    AllocationRequestId submitRequest(EmployeeId employeeId, ResourceId resourceId,
                                      LocalDate dateFrom, LocalDate dateTo, String justification);

    void cancelRequest(AllocationRequestId requestId, EmployeeId employeeId);

    List<ResourceCatalogItem> browseResources();

    List<AllocationRequestView> browseHistory(EmployeeId employeeId);

    List<AllocationRequestView> findActiveForResource(ResourceId resourceId);
}
\end{DeskFlowCode}

\phantomsection
\addcontentsline{toc}{subsubsection}{AllocationRequestRepositoryPort}
\DeskFlowListingCaption{allocation/application/port/out/AllocationRequestRepositoryPort.java - AllocationRequestRepositoryPort}{lst:allocation-application-port-out-AllocationRequestRepositoryPort-java}
\begin{DeskFlowCode}
package com.deskflow.allocation.application.port.out;

import com.deskflow.common.application.domain.model.vo.EmployeeId;
import com.deskflow.allocation.application.domain.model.AllocationRequest;
import com.deskflow.allocation.application.domain.model.vo.AllocationRequestId;
import com.deskflow.allocation.application.domain.model.vo.DateRange;
import com.deskflow.allocation.application.domain.model.vo.RequestStatus;
import com.deskflow.common.application.domain.model.vo.ResourceId;

import java.time.LocalDateTime;
import java.util.List;
import java.util.Optional;

public interface AllocationRequestRepositoryPort {

    void save(AllocationRequest request);

    Optional<AllocationRequest> find(AllocationRequestId id);

    List<AllocationRequest> findActiveForResource(ResourceId resourceId, DateRange dateRange);

    List<AllocationRequest> findActiveForResource(ResourceId resourceId);

    List<AllocationRequest> findActiveForEmployee(EmployeeId employeeId);

    List<AllocationRequest> findOverdue(LocalDateTime now);

    List<AllocationRequest> findForEmployee(EmployeeId employeeId);

    List<AllocationRequest> findAll(RequestStatus status);

    List<AllocationRequest> findAll();
}
\end{DeskFlowCode}

\phantomsection
\addcontentsline{toc}{subsubsection}{EmployeeProfilePort}
\DeskFlowListingCaption{allocation/application/port/out/EmployeeProfilePort.java - EmployeeProfilePort}{lst:allocation-application-port-out-EmployeeProfilePort-java}
\begin{DeskFlowCode}
package com.deskflow.allocation.application.port.out;

import com.deskflow.allocation.application.domain.model.EmployeeAllocationProfile;
import com.deskflow.common.application.domain.model.vo.EmployeeId;

import java.util.Optional;

public interface EmployeeProfilePort {
    Optional<EmployeeAllocationProfile> findEmployee(EmployeeId employeeId);
}
\end{DeskFlowCode}

\phantomsection
\addcontentsline{toc}{subsubsection}{ResourceCatalogPort}
\DeskFlowListingCaption{allocation/application/port/out/ResourceCatalogPort.java - ResourceCatalogPort}{lst:allocation-application-port-out-ResourceCatalogPort-java}
\begin{DeskFlowCode}
package com.deskflow.allocation.application.port.out;

import com.deskflow.allocation.application.domain.model.ResourceAllocationProfile;
import com.deskflow.common.application.domain.model.vo.ResourceId;

import java.util.List;
import java.util.Optional;

public interface ResourceCatalogPort {
    Optional<ResourceAllocationProfile> findResource(ResourceId resourceId);

    List<ResourceAllocationProfile> findActiveResources();
}
\end{DeskFlowCode}

\phantomsection
\addcontentsline{toc}{subsubsection}{AllocationAdministrationAppService}
\DeskFlowListingCaption{allocation/application/service/AllocationAdministrationAppService.java - AllocationAdministrationAppService}{lst:allocation-application-service-AllocationAdministrationAppService-java}
\begin{DeskFlowCode}
package com.deskflow.allocation.application.service;

import com.deskflow.allocation.application.model.AllocationRequestView;
import com.deskflow.allocation.application.port.in.AllocationAdministrationInputPort;
import com.deskflow.allocation.application.port.out.AllocationRequestRepositoryPort;
import com.deskflow.allocation.application.domain.model.AllocationRequest;
import com.deskflow.allocation.application.domain.model.vo.RequestStatus;
import com.deskflow.allocation.application.domain.service.ApproveAllocationRequestDomainService;
import com.deskflow.allocation.application.domain.service.IssueAllocationResourceDomainService;
import com.deskflow.allocation.application.domain.service.MarkAllocationRequestOverdueDomainService;
import com.deskflow.allocation.application.domain.service.RejectAllocationRequestDomainService;
import com.deskflow.allocation.application.domain.service.ReturnAllocationResourceDomainService;
import com.deskflow.common.application.port.out.DomainEventPublisherPort;
import com.deskflow.common.application.domain.DomainEvent;
import com.deskflow.allocation.application.domain.model.vo.AllocationRequestId;
import org.springframework.stereotype.Service;
import org.springframework.transaction.annotation.Transactional;

import java.time.LocalDateTime;
import java.util.List;
import java.util.NoSuchElementException;
import java.util.Optional;

@Service
@Transactional
public class AllocationAdministrationAppService implements AllocationAdministrationInputPort {

    private final AllocationRequestRepositoryPort requestRepository;
    private final DomainEventPublisherPort eventPublisher;
    private final ApproveAllocationRequestDomainService approveAllocationRequestDomainService;
    private final RejectAllocationRequestDomainService rejectAllocationRequestDomainService;
    private final IssueAllocationResourceDomainService issueAllocationResourceDomainService;
    private final ReturnAllocationResourceDomainService returnAllocationResourceDomainService;
    private final MarkAllocationRequestOverdueDomainService markAllocationRequestOverdueDomainService;

    public AllocationAdministrationAppService(AllocationRequestRepositoryPort requestRepository,
                                              DomainEventPublisherPort eventPublisher,
                                              ApproveAllocationRequestDomainService approveAllocationRequestDomainService,
                                              RejectAllocationRequestDomainService rejectAllocationRequestDomainService,
                                              IssueAllocationResourceDomainService issueAllocationResourceDomainService,
                                              ReturnAllocationResourceDomainService returnAllocationResourceDomainService,
                                              MarkAllocationRequestOverdueDomainService markAllocationRequestOverdueDomainService) {
        this.requestRepository = requestRepository;
        this.eventPublisher = eventPublisher;
        this.approveAllocationRequestDomainService = approveAllocationRequestDomainService;
        this.rejectAllocationRequestDomainService = rejectAllocationRequestDomainService;
        this.issueAllocationResourceDomainService = issueAllocationResourceDomainService;
        this.returnAllocationResourceDomainService = returnAllocationResourceDomainService;
        this.markAllocationRequestOverdueDomainService = markAllocationRequestOverdueDomainService;
    }

    @Override
    @Transactional(readOnly = true)
    public Optional<AllocationRequestView> getRequest(AllocationRequestId requestId) {
        return requestRepository.find(requestId).map(AllocationRequestView::from);
    }

    @Override
    @Transactional(readOnly = true)
    public List<AllocationRequestView> listRequests(RequestStatus status) {
        List<AllocationRequest> requests = status == null ? requestRepository.findAll() : requestRepository.findAll(status);
        return requests.stream().map(AllocationRequestView::from).toList();
    }

    @Override
    public void approveRequest(AllocationRequestId requestId) {
        mutate(requestId, approveAllocationRequestDomainService::approveAllocationRequest);
    }

    @Override
    public void rejectRequest(AllocationRequestId requestId, String reason) {
        mutate(requestId, request -> rejectAllocationRequestDomainService.rejectAllocationRequest(request, reason));
    }

    @Override
    public void issueResource(AllocationRequestId requestId) {
        mutate(requestId, issueAllocationResourceDomainService::issueAllocationResource);
    }

    @Override
    public void returnResource(AllocationRequestId requestId) {
        mutate(requestId, returnAllocationResourceDomainService::returnAllocationResource);
    }

    @Override
    @Transactional(readOnly = true)
    public List<AllocationRequestView> monitorActiveAllocations() {
        return requestRepository.findAll(RequestStatus.ACTIVE).stream()
                .map(AllocationRequestView::from)
                .toList();
    }

    @Override
    public void detectOverdueAllocations() {
        List<AllocationRequest> overdue = requestRepository.findOverdue(LocalDateTime.now());
        for (AllocationRequest request : overdue) {
            markAllocationRequestOverdueDomainService.markAllocationRequestOverdue(request);
            requestRepository.save(request);
            publishAndClear(request);
        }
    }

    private void mutate(AllocationRequestId requestId, java.util.function.Consumer<AllocationRequest> mutation) {
        AllocationRequest request = requestRepository.find(requestId)
                .orElseThrow(() -> new NoSuchElementException("Allocation request not found: " + requestId));
        mutation.accept(request);
        requestRepository.save(request);
        publishAndClear(request);
    }

    private void publishAndClear(AllocationRequest request) {
        for (DomainEvent event : request.flushEvents()) {
            eventPublisher.publish(event);
        }
    }
}
\end{DeskFlowCode}

\phantomsection
\addcontentsline{toc}{subsubsection}{AllocationRequestAppService}
\DeskFlowListingCaption{allocation/application/service/AllocationRequestAppService.java - AllocationRequestAppService}{lst:allocation-application-service-AllocationRequestAppService-java}
\begin{DeskFlowCode}
package com.deskflow.allocation.application.service;

import com.deskflow.allocation.application.port.in.AllocationRequestInputPort;
import com.deskflow.allocation.application.model.AllocationRequestView;
import com.deskflow.allocation.application.model.ResourceCatalogItem;
import com.deskflow.allocation.application.port.out.AllocationRequestRepositoryPort;
import com.deskflow.allocation.application.port.out.EmployeeProfilePort;
import com.deskflow.allocation.application.port.out.ResourceCatalogPort;
import com.deskflow.common.application.domain.model.vo.EmployeeId;
import com.deskflow.common.application.domain.DomainEvent;
import com.deskflow.allocation.application.domain.model.AllocationRequest;
import com.deskflow.allocation.application.domain.model.ResourceAllocationProfile;
import com.deskflow.allocation.application.domain.service.AllocationRequestFactory;
import com.deskflow.allocation.application.domain.service.CancelAllocationRequestDomainService;
import com.deskflow.allocation.application.domain.service.EnsureEmployeeWithinLimitDomainService;
import com.deskflow.allocation.application.domain.service.EnsureResourceAvailableInPeriodDomainService;
import com.deskflow.allocation.application.domain.service.EnsureResourceEligibleForAllocationDomainService;
import com.deskflow.allocation.application.domain.model.vo.AllocationRequestId;
import com.deskflow.allocation.application.domain.model.vo.DateRange;
import com.deskflow.common.application.port.out.DomainEventPublisherPort;
import com.deskflow.common.application.domain.model.vo.ResourceId;
import org.springframework.stereotype.Service;
import org.springframework.transaction.annotation.Transactional;

import java.time.LocalDate;
import java.util.List;
import java.util.NoSuchElementException;

@Service
@Transactional
public class AllocationRequestAppService implements AllocationRequestInputPort {

    private final AllocationRequestRepositoryPort requestRepository;
    private final ResourceCatalogPort resourceCatalogPort;
    private final EmployeeProfilePort employeeProfilePort;
    private final AllocationRequestFactory factory;
    private final CancelAllocationRequestDomainService cancelAllocationRequestDomainService;
    private final EnsureResourceEligibleForAllocationDomainService ensureResourceEligibleForAllocationDomainService;
    private final EnsureResourceAvailableInPeriodDomainService ensureResourceAvailableInPeriodDomainService;
    private final EnsureEmployeeWithinLimitDomainService ensureEmployeeWithinLimitDomainService;
    private final DomainEventPublisherPort eventPublisher;

    public AllocationRequestAppService(AllocationRequestRepositoryPort requestRepository,
                                      ResourceCatalogPort resourceCatalogPort,
                                      EmployeeProfilePort employeeProfilePort,
                                      AllocationRequestFactory factory,
                                      CancelAllocationRequestDomainService cancelAllocationRequestDomainService,
                                      EnsureResourceEligibleForAllocationDomainService ensureResourceEligibleForAllocationDomainService,
                                      EnsureResourceAvailableInPeriodDomainService ensureResourceAvailableInPeriodDomainService,
                                      EnsureEmployeeWithinLimitDomainService ensureEmployeeWithinLimitDomainService,
                                      DomainEventPublisherPort eventPublisher) {
        this.requestRepository = requestRepository;
        this.resourceCatalogPort = resourceCatalogPort;
        this.employeeProfilePort = employeeProfilePort;
        this.factory = factory;
        this.cancelAllocationRequestDomainService = cancelAllocationRequestDomainService;
        this.ensureResourceEligibleForAllocationDomainService = ensureResourceEligibleForAllocationDomainService;
        this.ensureResourceAvailableInPeriodDomainService = ensureResourceAvailableInPeriodDomainService;
        this.ensureEmployeeWithinLimitDomainService = ensureEmployeeWithinLimitDomainService;
        this.eventPublisher = eventPublisher;
    }

    @Override
    public AllocationRequestId submitRequest(EmployeeId employeeId, ResourceId resourceId,
                                              LocalDate dateFrom, LocalDate dateTo, String justification) {
        var employee = employeeProfilePort.findEmployee(employeeId)
                .orElseThrow(() -> new IllegalArgumentException("Employee not found: " + employeeId));
        ResourceAllocationProfile requestedResource = resourceCatalogPort.findResource(resourceId)
                .orElseThrow(() -> new IllegalArgumentException("Resource not found: " + resourceId));

        ensureResourceEligibleForAllocationDomainService.ensureResourceEligibleForAllocation(requestedResource);

        String requestedType = requestedResource.type();
        long currentActiveForType = requestRepository.findActiveForEmployee(employeeId).stream()
                .map(AllocationRequest::getResourceId)
                .map(resourceCatalogPort::findResource)
                .flatMap(java.util.Optional::stream)
                .filter(resource -> resource.type().equals(requestedType))
                .count();

        ensureEmployeeWithinLimitDomainService.ensureEmployeeWithinLimit(employee, requestedType, currentActiveForType);

        DateRange dateRange = new DateRange(dateFrom, dateTo);

        List<AllocationRequest> overlappingRequests = requestRepository.findActiveForResource(resourceId, dateRange);
        ensureResourceAvailableInPeriodDomainService.ensureResourceAvailableInPeriod(overlappingRequests);

        AllocationRequest request = factory.create(employeeId, resourceId, dateRange, justification);
        requestRepository.save(request);
        publishAndClear(request);
        return request.getId();
    }

    @Override
    public void cancelRequest(AllocationRequestId requestId, EmployeeId employeeId) {
        AllocationRequest request = findRequest(requestId);
        cancelAllocationRequestDomainService.cancelAllocationRequest(request, employeeId);
        requestRepository.save(request);
        publishAndClear(request);
    }

    @Override
    @Transactional(readOnly = true)
    public List<ResourceCatalogItem> browseResources() {
        return resourceCatalogPort.findActiveResources().stream()
                .map(ResourceCatalogItem::from)
                .toList();
    }

    @Override
    @Transactional(readOnly = true)
    public List<AllocationRequestView> browseHistory(EmployeeId employeeId) {
        return requestRepository.findForEmployee(employeeId).stream()
                .map(AllocationRequestView::from)
                .toList();
    }

    @Override
    @Transactional(readOnly = true)
    public List<AllocationRequestView> findActiveForResource(ResourceId resourceId) {
        return requestRepository.findActiveForResource(resourceId).stream()
                .map(AllocationRequestView::from)
                .toList();
    }

    private AllocationRequest findRequest(AllocationRequestId id) {
        return requestRepository.find(id)
                .orElseThrow(() -> new NoSuchElementException("Allocation request not found: " + id));
    }

    private void publishAndClear(AllocationRequest request) {
        for (DomainEvent event : request.flushEvents()) {
            eventPublisher.publish(event);
        }
    }
}
\end{DeskFlowCode}

\phantomsection
\addcontentsline{toc}{subsubsection}{AllocationRequestController}
\DeskFlowListingCaption{allocation/infrastructure/in/web/AllocationRequestController.java - AllocationRequestController}{lst:allocation-infrastructure-in-web-AllocationRequestController-java}
\begin{DeskFlowCode}
package com.deskflow.allocation.infrastructure.in.web;

import com.deskflow.allocation.application.port.out.EmployeeProfilePort;
import com.deskflow.allocation.infrastructure.in.web.dto.AllocationRequestDTO;
import com.deskflow.allocation.infrastructure.in.web.dto.AllocationRequestResponseDTO;
import com.deskflow.allocation.application.port.in.AllocationRequestInputPort;
import com.deskflow.common.application.domain.model.vo.EmployeeId;
import com.deskflow.allocation.application.domain.model.vo.AllocationRequestId;
import com.deskflow.common.application.domain.model.vo.ResourceId;
import com.deskflow.employees.application.port.in.EmployeeManagementInputPort;
import com.deskflow.employees.application.domain.model.Employee;
import jakarta.validation.Valid;
import org.springframework.http.ResponseEntity;
import org.springframework.security.core.Authentication;
import org.springframework.web.bind.annotation.*;

import java.net.URI;
import java.util.List;
import java.util.UUID;

@RestController
@RequestMapping("/api")
public class AllocationRequestController {

    private final AllocationRequestInputPort allocationRequestInputPort;
    private final CurrentUserResolver userResolver;
    private final EmployeeManagementInputPort employeeManagementInputPort;

    public AllocationRequestController(AllocationRequestInputPort allocationRequestInputPort,
                                       CurrentUserResolver userResolver,
                                       EmployeeManagementInputPort employeeManagementInputPort) {
        this.allocationRequestInputPort = allocationRequestInputPort;
        this.userResolver = userResolver;
        this.employeeManagementInputPort = employeeManagementInputPort;
    }

    @GetMapping("/resources")
    public ResponseEntity<?> browseResources() {
        return ResponseEntity.ok(allocationRequestInputPort.browseResources().stream()
                .map(r -> new java.util.LinkedHashMap<>(java.util.Map.of(
                        "id", r.id(),
                        "type", r.type(),
                        "name", r.name(),
                        "building", r.building(),
                        "floor", r.floor(),
                        "number", r.number(),
                        "available", r.available()
                )))
                .toList());
    }

    @PostMapping("/requests")
    public ResponseEntity<?> submitRequest(@Valid @RequestBody AllocationRequestDTO dto, Authentication auth) {
        EmployeeId employeeId = userResolver.resolveEmployeeId(auth);
        AllocationRequestId id = allocationRequestInputPort.submitRequest(
                employeeId, ResourceId.of(dto.resourceId()),
                dto.dateFrom(), dto.dateTo(), dto.justification());
        return ResponseEntity.created(URI.create("/api/requests/" + id))
                .body(java.util.Map.of("requestId", id.getValue()));
    }

    @PutMapping("/requests/{id}/cancel")
    public ResponseEntity<Void> cancelRequest(@PathVariable UUID id, Authentication auth) {
        EmployeeId employeeId = userResolver.resolveEmployeeId(auth);
        allocationRequestInputPort.cancelRequest(AllocationRequestId.of(id), employeeId);
        return ResponseEntity.ok().build();
    }

    @GetMapping("/requests/my")
    public ResponseEntity<List<AllocationRequestResponseDTO>> browseHistory(Authentication auth) {
        EmployeeId employeeId = userResolver.resolveEmployeeId(auth);
        Employee employee = employeeManagementInputPort.getEmployee(employeeId).orElse(null);
        List<AllocationRequestResponseDTO> history = allocationRequestInputPort.browseHistory(employeeId)
                .stream()
                .map(r -> AllocationRequestRestMapper.toResponse(r,
                        employee != null ? employee.getFirstName() : null,
                        employee != null ? employee.getLastName() : null))
                .toList();
        return ResponseEntity.ok(history);
    }
}
\end{DeskFlowCode}

\phantomsection
\addcontentsline{toc}{subsubsection}{AllocationRequestRestMapper}
\DeskFlowListingCaption{allocation/infrastructure/in/web/AllocationRequestRestMapper.java - AllocationRequestRestMapper}{lst:allocation-infrastructure-in-web-AllocationRequestRestMapper-java}
\begin{DeskFlowCode}
package com.deskflow.allocation.infrastructure.in.web;

import com.deskflow.allocation.application.model.AllocationRequestView;
import com.deskflow.allocation.infrastructure.in.web.dto.AllocationRequestResponseDTO;

public class AllocationRequestRestMapper {

    public static AllocationRequestResponseDTO toResponse(AllocationRequestView r) {
        return toResponse(r, null, null);
    }

    public static AllocationRequestResponseDTO toResponse(AllocationRequestView r, String firstName, String lastName) {
        return new AllocationRequestResponseDTO(
                r.id(),
                r.employeeId(),
                firstName,
                lastName,
                r.resourceId(),
                r.status(),
                r.dateFrom(),
                r.dateTo(),
                r.justification(),
                r.submittedAt(),
                r.issuedAt(),
                r.returnedAt(),
                r.rejectionReason()
        );
    }
}
\end{DeskFlowCode}

\phantomsection
\addcontentsline{toc}{subsubsection}{CurrentUserResolver}
\DeskFlowListingCaption{allocation/infrastructure/in/web/CurrentUserResolver.java - CurrentUserResolver}{lst:allocation-infrastructure-in-web-CurrentUserResolver-java}
\begin{DeskFlowCode}
package com.deskflow.allocation.infrastructure.in.web;

import com.deskflow.common.application.domain.model.vo.EmployeeId;
import com.deskflow.employees.application.port.in.EmployeeManagementInputPort;
import org.springframework.security.core.Authentication;
import org.springframework.stereotype.Component;

@Component
public class CurrentUserResolver {

    private final EmployeeManagementInputPort employeeManagementInputPort;

    public CurrentUserResolver(EmployeeManagementInputPort employeeManagementInputPort) {
        this.employeeManagementInputPort = employeeManagementInputPort;
    }

    public EmployeeId resolveEmployeeId(Authentication auth) {
        return employeeManagementInputPort.getEmployeeByEmail(auth.getName())
                .map(e -> e.getId())
                .orElseThrow(() -> new IllegalStateException("Employee not found for: " + auth.getName()));
    }
}
\end{DeskFlowCode}

\phantomsection
\addcontentsline{toc}{subsubsection}{AllocationRequestDTO}
\DeskFlowListingCaption{allocation/infrastructure/in/web/dto/AllocationRequestDTO.java - AllocationRequestDTO}{lst:allocation-infrastructure-in-web-dto-AllocationRequestDTO-java}
\begin{DeskFlowCode}
package com.deskflow.allocation.infrastructure.in.web.dto;

import jakarta.validation.constraints.FutureOrPresent;
import jakarta.validation.constraints.NotNull;

import java.time.LocalDate;
import java.util.UUID;

public record AllocationRequestDTO(
        @NotNull UUID resourceId,
        @NotNull @FutureOrPresent LocalDate dateFrom,
        @NotNull LocalDate dateTo,
        String justification
) {}
\end{DeskFlowCode}

\phantomsection
\addcontentsline{toc}{subsubsection}{AllocationRequestResponseDTO}
\DeskFlowListingCaption{allocation/infrastructure/in/web/dto/AllocationRequestResponseDTO.java - AllocationRequestResponseDTO}{lst:allocation-infrastructure-in-web-dto-AllocationRequestResponseDTO-java}
\begin{DeskFlowCode}
package com.deskflow.allocation.infrastructure.in.web.dto;

import java.time.LocalDate;
import java.time.LocalDateTime;
import java.util.UUID;

public record AllocationRequestResponseDTO(
        UUID id,
        UUID employeeId,
        String employeeFirstName,
        String employeeLastName,
        UUID resourceId,
        String status,
        LocalDate dateFrom,
        LocalDate dateTo,
        String justification,
        LocalDateTime submittedAt,
        LocalDateTime issuedAt,
        LocalDateTime returnedAt,
        String rejectionReason
) {}
\end{DeskFlowCode}

\phantomsection
\addcontentsline{toc}{subsubsection}{EmployeeProfileAdapter}
\DeskFlowListingCaption{allocation/infrastructure/out/employees/EmployeeProfileAdapter.java - EmployeeProfileAdapter}{lst:allocation-infrastructure-out-employees-EmployeeProfileAdapter-java}
\begin{DeskFlowCode}
package com.deskflow.allocation.infrastructure.out.employees;

import com.deskflow.allocation.application.domain.model.EmployeeAllocationProfile;
import com.deskflow.allocation.application.port.out.EmployeeProfilePort;
import com.deskflow.employees.application.port.in.EmployeeManagementInputPort;
import com.deskflow.common.application.domain.model.vo.EmployeeId;
import org.springframework.stereotype.Component;

import java.util.Optional;
import java.util.stream.Collectors;

@Component
public class EmployeeProfileAdapter implements EmployeeProfilePort {

    private final EmployeeManagementInputPort employeeManagementInputPort;

    public EmployeeProfileAdapter(EmployeeManagementInputPort employeeManagementInputPort) {
        this.employeeManagementInputPort = employeeManagementInputPort;
    }

    @Override
    public Optional<EmployeeAllocationProfile> findEmployee(EmployeeId employeeId) {
        return employeeManagementInputPort.getEmployee(employeeId)
                .map(employee -> new EmployeeAllocationProfile(
                        employee.getLimits().entrySet().stream()
                                .collect(Collectors.toMap(entry -> entry.getKey().name(), java.util.Map.Entry::getValue))
                ));
    }
}
\end{DeskFlowCode}

\phantomsection
\addcontentsline{toc}{subsubsection}{AllocationRequestEntity}
\DeskFlowListingCaption{allocation/infrastructure/out/persistence/AllocationRequestEntity.java - AllocationRequestEntity}{lst:allocation-infrastructure-out-persistence-AllocationRequestEntity-java}
\begin{DeskFlowCode}
package com.deskflow.allocation.infrastructure.out.persistence;

import jakarta.persistence.*;

import java.time.LocalDate;
import java.time.LocalDateTime;
import java.util.UUID;

@Entity
@Table(name = "allocation_requests")
class AllocationRequestEntity {

    @Id
    UUID id;

    @Column(name = "employee_id", nullable = false)
    UUID employeeId;

    @Column(name = "resource_id", nullable = false)
    UUID resourceId;

    @Column(nullable = false)
    String status;

    @Column(name = "date_from", nullable = false)
    LocalDate dateFrom;

    @Column(name = "date_to", nullable = false)
    LocalDate dateTo;

    @Column(name = "justification")
    String justification;

    @Column(name = "submitted_at", nullable = false)
    LocalDateTime submittedAt;

    @Column(name = "issued_at")
    LocalDateTime issuedAt;

    @Column(name = "returned_at")
    LocalDateTime returnedAt;

    @Column(name = "rejection_reason")
    String rejectionReason;
}
\end{DeskFlowCode}

\phantomsection
\addcontentsline{toc}{subsubsection}{AllocationRequestMapper}
\DeskFlowListingCaption{allocation/infrastructure/out/persistence/AllocationRequestMapper.java - AllocationRequestMapper}{lst:allocation-infrastructure-out-persistence-AllocationRequestMapper-java}
\begin{DeskFlowCode}
package com.deskflow.allocation.infrastructure.out.persistence;

import com.deskflow.common.application.domain.model.vo.EmployeeId;
import com.deskflow.allocation.application.domain.model.AllocationRequest;
import com.deskflow.allocation.application.domain.model.vo.AllocationRequestId;
import com.deskflow.allocation.application.domain.model.vo.DateRange;
import com.deskflow.allocation.application.domain.model.vo.RequestStatus;
import com.deskflow.common.application.domain.model.vo.ResourceId;
import org.springframework.stereotype.Component;

@Component
public class AllocationRequestMapper {

    AllocationRequestEntity toEntity(AllocationRequest request) {
        AllocationRequestEntity e = new AllocationRequestEntity();
        updateEntity(request, e);
        return e;
    }

    void updateEntity(AllocationRequest request, AllocationRequestEntity e) {
        e.id = request.getId().getValue();
        e.employeeId = request.getEmployeeId().getValue();
        e.resourceId = request.getResourceId().getValue();
        e.status = request.getStatus().name();
        e.dateFrom = request.getDateRange().getDateFrom();
        e.dateTo = request.getDateRange().getDateTo();
        e.justification = request.getJustification();
        e.submittedAt = request.getSubmittedAt();
        e.issuedAt = request.getIssuedAt();
        e.returnedAt = request.getReturnedAt();
        e.rejectionReason = request.getRejectionReason();
    }

    AllocationRequest toDomain(AllocationRequestEntity e) {
        return AllocationRequest.reconstitute(
                AllocationRequestId.of(e.id),
                EmployeeId.of(e.employeeId),
                ResourceId.of(e.resourceId),
                new DateRange(e.dateFrom, e.dateTo),
                e.justification,
                e.submittedAt,
                RequestStatus.valueOf(e.status),
                e.issuedAt,
                e.returnedAt,
                e.rejectionReason
        );
    }
}
\end{DeskFlowCode}

\phantomsection
\addcontentsline{toc}{subsubsection}{HibernateAllocationRequestRepositoryAdapter}
\DeskFlowListingCaption{allocation/infrastructure/out/persistence/HibernateAllocationRequestRepositoryAdapter.java - HibernateAllocationRequestRepositoryAdapter}{lst:allocation-infrastructure-out-persistence-HibernateAllocationRequestRepositoryAdapter-java}
\begin{DeskFlowCode}
package com.deskflow.allocation.infrastructure.out.persistence;

import com.deskflow.allocation.application.port.out.AllocationRequestRepositoryPort;
import com.deskflow.allocation.application.domain.model.AllocationRequest;
import com.deskflow.allocation.application.domain.model.vo.AllocationRequestId;
import com.deskflow.allocation.application.domain.model.vo.DateRange;
import com.deskflow.allocation.application.domain.model.vo.RequestStatus;
import com.deskflow.common.application.domain.model.vo.EmployeeId;
import com.deskflow.common.application.domain.model.vo.ResourceId;
import org.springframework.stereotype.Repository;

import java.time.LocalDateTime;
import java.util.List;
import java.util.Optional;

@Repository
public class HibernateAllocationRequestRepositoryAdapter implements AllocationRequestRepositoryPort {

    private final SpringDataAllocationRequestRepository jpa;
    private final AllocationRequestMapper mapper;

    public HibernateAllocationRequestRepositoryAdapter(SpringDataAllocationRequestRepository jpa, AllocationRequestMapper mapper) {
        this.jpa = jpa;
        this.mapper = mapper;
    }

    @Override
    public void save(AllocationRequest request) {
        jpa.findById(request.getId().getValue()).ifPresentOrElse(
                managed -> mapper.updateEntity(request, managed),
                () -> jpa.save(mapper.toEntity(request))
        );
    }

    @Override
    public Optional<AllocationRequest> find(AllocationRequestId id) {
        return jpa.findById(id.getValue()).map(mapper::toDomain);
    }

    @Override
    public List<AllocationRequest> findActiveForResource(ResourceId resourceId, DateRange dateRange) {
        return jpa.findActiveOverlapping(resourceId.getValue(), dateRange.getDateFrom(), dateRange.getDateTo())
                .stream().map(mapper::toDomain).toList();
    }

    @Override
    public List<AllocationRequest> findActiveForResource(ResourceId resourceId) {
        return jpa.findActiveByResourceId(resourceId.getValue())
                .stream().map(mapper::toDomain).toList();
    }

    @Override
    public List<AllocationRequest> findActiveForEmployee(EmployeeId employeeId) {
        return jpa.findActiveByEmployeeId(employeeId.getValue())
                .stream().map(mapper::toDomain).toList();
    }

    @Override
    public List<AllocationRequest> findOverdue(LocalDateTime now) {
        return jpa.findOverdue(now.toLocalDate())
                .stream().map(mapper::toDomain).toList();
    }

    @Override
    public List<AllocationRequest> findForEmployee(EmployeeId employeeId) {
        return jpa.findByEmployeeId(employeeId.getValue())
                .stream().map(mapper::toDomain).toList();
    }

    @Override
    public List<AllocationRequest> findAll(RequestStatus status) {
        return jpa.findByStatus(status.name())
                .stream().map(mapper::toDomain).toList();
    }

    @Override
    public List<AllocationRequest> findAll() {
        return jpa.findAll().stream().map(mapper::toDomain).toList();
    }
}
\end{DeskFlowCode}

\phantomsection
\addcontentsline{toc}{subsubsection}{SpringDataAllocationRequestRepository}
\DeskFlowListingCaption{allocation/infrastructure/out/persistence/SpringDataAllocationRequestRepository.java - SpringDataAllocationRequestRepository}{lst:allocation-infrastructure-out-persistence-SpringDataAllocationRequestRepository-java}
\begin{DeskFlowCode}
package com.deskflow.allocation.infrastructure.out.persistence;

import org.springframework.data.jpa.repository.JpaRepository;
import org.springframework.data.jpa.repository.Query;
import org.springframework.data.repository.query.Param;

import java.time.LocalDate;
import java.util.List;
import java.util.UUID;

interface SpringDataAllocationRequestRepository extends JpaRepository<AllocationRequestEntity, UUID> {

    List<AllocationRequestEntity> findByEmployeeId(UUID employeeId);

    List<AllocationRequestEntity> findByStatus(String status);

    @Query("""
            SELECT w FROM AllocationRequestEntity w
            WHERE w.resourceId = :resourceId
            AND w.status IN ('SUBMITTED', 'APPROVED', 'ACTIVE', 'OVERDUE')
            AND w.dateFrom <= :dateTo
            AND w.dateTo >= :dateFrom
            """)
    List<AllocationRequestEntity> findActiveOverlapping(
            @Param("resourceId") UUID resourceId,
            @Param("dateFrom") LocalDate dateFrom,
            @Param("dateTo") LocalDate dateTo
    );

    @Query("""
            SELECT w FROM AllocationRequestEntity w
            WHERE w.employeeId = :employeeId
            AND w.status IN ('SUBMITTED', 'APPROVED', 'ACTIVE', 'OVERDUE')
            """)
    List<AllocationRequestEntity> findActiveByEmployeeId(@Param("employeeId") UUID employeeId);

    @Query("""
            SELECT w FROM AllocationRequestEntity w
            WHERE w.resourceId = :resourceId
            AND w.status IN ('SUBMITTED', 'APPROVED', 'ACTIVE', 'OVERDUE')
            """)
    List<AllocationRequestEntity> findActiveByResourceId(@Param("resourceId") UUID resourceId);

    @Query("""
            SELECT w FROM AllocationRequestEntity w
            WHERE w.status = 'ACTIVE'
            AND w.dateTo < :now
            """)
    List<AllocationRequestEntity> findOverdue(@Param("now") LocalDate now);
}
\end{DeskFlowCode}

\phantomsection
\addcontentsline{toc}{subsubsection}{ResourceCatalogAdapter}
\DeskFlowListingCaption{allocation/infrastructure/out/resources/ResourceCatalogAdapter.java - ResourceCatalogAdapter}{lst:allocation-infrastructure-out-resources-ResourceCatalogAdapter-java}
\begin{DeskFlowCode}
package com.deskflow.allocation.infrastructure.out.resources;

import com.deskflow.allocation.application.domain.model.ResourceAllocationProfile;
import com.deskflow.allocation.application.port.out.ResourceCatalogPort;
import com.deskflow.resources.application.model.ResourceView;
import com.deskflow.resources.application.port.in.ResourceManagementInputPort;
import com.deskflow.common.application.domain.model.vo.ResourceId;
import org.springframework.stereotype.Component;

import java.util.List;
import java.util.Optional;

@Component
public class ResourceCatalogAdapter implements ResourceCatalogPort {

    private final ResourceManagementInputPort resourceManagementInputPort;

    public ResourceCatalogAdapter(ResourceManagementInputPort resourceManagementInputPort) {
        this.resourceManagementInputPort = resourceManagementInputPort;
    }

    @Override
    public Optional<ResourceAllocationProfile> findResource(ResourceId resourceId) {
        return resourceManagementInputPort.getResource(resourceId).map(this::toProfile);
    }

    @Override
    public List<ResourceAllocationProfile> findActiveResources() {
        return resourceManagementInputPort.getActiveResources().stream()
                .map(this::toProfile)
                .toList();
    }

    private ResourceAllocationProfile toProfile(ResourceView resource) {
        return new ResourceAllocationProfile(
                ResourceId.of(resource.id()),
                resource.type(),
                resource.name(),
                resource.building(),
                resource.floor(),
                resource.number(),
                resource.active(),
                resource.available()
        );
    }
}
\end{DeskFlowCode}

\subsection{Kontekst Resources - zasoby biurowe}

\phantomsection
\addcontentsline{toc}{subsubsection}{OfficeResource}
\DeskFlowListingCaption{resources/application/domain/model/OfficeResource.java - OfficeResource}{lst:resources-application-domain-model-OfficeResource-java}
\begin{DeskFlowCode}
package com.deskflow.resources.application.domain.model;

import com.deskflow.common.application.domain.DomainEvent;
import com.deskflow.resources.application.domain.model.events.ResourceDeactivated;
import com.deskflow.resources.application.domain.model.events.ResourceMadeAvailable;
import com.deskflow.resources.application.domain.model.events.ResourceMadeUnavailable;
import com.deskflow.resources.application.domain.model.events.ResourceUpdated;
import com.deskflow.resources.application.domain.model.vo.Location;
import com.deskflow.common.application.domain.model.vo.ResourceType;
import com.deskflow.common.application.domain.model.vo.ResourceId;

import java.time.LocalDateTime;
import java.util.ArrayList;
import java.util.List;

public class OfficeResource {

    private final ResourceId id;
    private final ResourceType type;
    private String name;
    private Location location;
    private boolean active;
    private boolean available;
    private final List<DomainEvent> events = new ArrayList<>();

    private OfficeResource(ResourceId id, ResourceType type, String name, Location location,
                           boolean active, boolean available) {
        this.id = id;
        this.type = type;
        this.name = validateName(name);
        this.location = location;
        this.active = active;
        this.available = available;
    }

    public static OfficeResource create(ResourceId id, ResourceType type, String name, Location location) {
        return new OfficeResource(id, type, name, location, true, true);
    }

    public static OfficeResource reconstitute(ResourceId id, ResourceType type, String name,
                                               Location location, boolean active, boolean available) {
        return new OfficeResource(id, type, name, location, active, available);
    }

    public void update(String newName, Location newLocation) {
        if (!active) throw new IllegalStateException("Cannot edit an inactive resource");
        String prevName = this.name;
        this.name = validateName(newName);
        this.location = newLocation;
        events.add(new ResourceUpdated(id, prevName, newName, LocalDateTime.now()));
    }

    public void markAsUnavailable() {
        if (!active) throw new IllegalStateException("Resource is not active");
        this.available = false;
        events.add(new ResourceMadeUnavailable(id, LocalDateTime.now()));
    }

    public void markAsAvailable() {
        if (!active) throw new IllegalStateException("Resource is not active");
        this.available = true;
        events.add(new ResourceMadeAvailable(id, LocalDateTime.now()));
    }

    public void deactivate() {
        if (!active) throw new IllegalStateException("Resource is already deactivated");
        this.active = false;
        this.available = false;
        events.add(new ResourceDeactivated(id, LocalDateTime.now()));
    }

    private String validateName(String name) {
        if (name == null || name.isBlank()) throw new IllegalArgumentException("Resource name cannot be blank");
        return name;
    }

    public List<DomainEvent> flushEvents() {
        var copy = List.copyOf(events);
        events.clear();
        return copy;
    }

    public ResourceId getId() { return id; }
    public ResourceType getType() { return type; }
    public String getName() { return name; }
    public Location getLocation() { return location; }
    public boolean isActive() { return active; }
    public boolean isAvailable() { return available; }
}
\end{DeskFlowCode}

\phantomsection
\addcontentsline{toc}{subsubsection}{ResourceDeactivated}
\DeskFlowListingCaption{resources/application/domain/model/events/ResourceDeactivated.java - ResourceDeactivated}{lst:resources-application-domain-model-events-ResourceDeactivated-java}
\begin{DeskFlowCode}
package com.deskflow.resources.application.domain.model.events;

import com.deskflow.common.application.domain.DomainEvent;
import com.deskflow.common.application.domain.model.vo.ResourceId;

import java.time.LocalDateTime;
import java.util.UUID;

public record ResourceDeactivated(
        ResourceId resourceId,
        LocalDateTime occurredAt
) implements DomainEvent {
    @Override
    public UUID eventId() { return UUID.randomUUID(); }
    @Override
    public String eventType() { return "ResourceDeactivated"; }
    @Override
    public UUID aggregateId() { return resourceId.getValue(); }
}
\end{DeskFlowCode}

\phantomsection
\addcontentsline{toc}{subsubsection}{ResourceMadeAvailable}
\DeskFlowListingCaption{resources/application/domain/model/events/ResourceMadeAvailable.java - ResourceMadeAvailable}{lst:resources-application-domain-model-events-ResourceMadeAvailable-java}
\begin{DeskFlowCode}
package com.deskflow.resources.application.domain.model.events;

import com.deskflow.common.application.domain.DomainEvent;
import com.deskflow.common.application.domain.model.vo.ResourceId;

import java.time.LocalDateTime;
import java.util.UUID;

public record ResourceMadeAvailable(
        ResourceId resourceId,
        LocalDateTime occurredAt
) implements DomainEvent {
    @Override
    public UUID eventId() { return UUID.randomUUID(); }
    @Override
    public String eventType() { return "ResourceMadeAvailable"; }
    @Override
    public UUID aggregateId() { return resourceId.getValue(); }
}
\end{DeskFlowCode}

\phantomsection
\addcontentsline{toc}{subsubsection}{ResourceMadeUnavailable}
\DeskFlowListingCaption{resources/application/domain/model/events/ResourceMadeUnavailable.java - ResourceMadeUnavailable}{lst:resources-application-domain-model-events-ResourceMadeUnavailable-java}
\begin{DeskFlowCode}
package com.deskflow.resources.application.domain.model.events;

import com.deskflow.common.application.domain.DomainEvent;
import com.deskflow.common.application.domain.model.vo.ResourceId;

import java.time.LocalDateTime;
import java.util.UUID;

public record ResourceMadeUnavailable(
        ResourceId resourceId,
        LocalDateTime occurredAt
) implements DomainEvent {
    @Override
    public UUID eventId() { return UUID.randomUUID(); }
    @Override
    public String eventType() { return "ResourceMadeUnavailable"; }
    @Override
    public UUID aggregateId() { return resourceId.getValue(); }
}
\end{DeskFlowCode}

\phantomsection
\addcontentsline{toc}{subsubsection}{ResourceUpdated}
\DeskFlowListingCaption{resources/application/domain/model/events/ResourceUpdated.java - ResourceUpdated}{lst:resources-application-domain-model-events-ResourceUpdated-java}
\begin{DeskFlowCode}
package com.deskflow.resources.application.domain.model.events;

import com.deskflow.common.application.domain.DomainEvent;
import com.deskflow.common.application.domain.model.vo.ResourceId;

import java.time.LocalDateTime;
import java.util.UUID;

public record ResourceUpdated(
        ResourceId resourceId,
        String previousName,
        String newName,
        LocalDateTime occurredAt
) implements DomainEvent {
    @Override
    public UUID eventId() { return UUID.randomUUID(); }
    @Override
    public String eventType() { return "ResourceUpdated"; }
    @Override
    public UUID aggregateId() { return resourceId.getValue(); }
}
\end{DeskFlowCode}

\phantomsection
\addcontentsline{toc}{subsubsection}{AllocationUsageStatus}
\DeskFlowListingCaption{resources/application/domain/model/vo/AllocationUsageStatus.java - AllocationUsageStatus}{lst:resources-application-domain-model-vo-AllocationUsageStatus-java}
\begin{DeskFlowCode}
package com.deskflow.resources.application.domain.model.vo;

public enum AllocationUsageStatus {
    REQUESTED,
    APPROVED,
    IN_USE,
    OVERDUE,
    CLOSED
}
\end{DeskFlowCode}

\phantomsection
\addcontentsline{toc}{subsubsection}{Location}
\DeskFlowListingCaption{resources/application/domain/model/vo/Location.java - Location}{lst:resources-application-domain-model-vo-Location-java}
\begin{DeskFlowCode}
package com.deskflow.resources.application.domain.model.vo;

import java.util.Objects;

public final class Location {
    private final String building;
    private final int floor;
    private final String number;

    public Location(String building, int floor, String number) {
        if (building == null || building.isBlank()) throw new IllegalArgumentException("Building cannot be blank");
        if (number == null || number.isBlank()) throw new IllegalArgumentException("Number cannot be blank");
        this.building = building;
        this.floor = floor;
        this.number = number;
    }

    public String getBuilding() { return building; }
    public int getFloor() { return floor; }
    public String getNumber() { return number; }

    @Override
    public boolean equals(Object o) {
        if (this == o) return true;
        if (!(o instanceof Location that)) return false;
        return floor == that.floor && building.equals(that.building) && number.equals(that.number);
    }

    @Override
    public int hashCode() {
        return Objects.hash(building, floor, number);
    }

    @Override
    public String toString() {
        return building + ", floor " + floor + ", no. " + number;
    }
}
\end{DeskFlowCode}

\phantomsection
\addcontentsline{toc}{subsubsection}{DeactivateOfficeResourceDomainService}
\DeskFlowListingCaption{resources/application/domain/service/DeactivateOfficeResourceDomainService.java - DeactivateOfficeResourceDomainService}{lst:resources-application-domain-service-DeactivateOfficeResourceDomainService-java}
\begin{DeskFlowCode}
package com.deskflow.resources.application.domain.service;

import com.deskflow.resources.application.domain.model.OfficeResource;
import com.deskflow.resources.application.domain.model.vo.AllocationUsageStatus;

import java.util.List;

public class DeactivateOfficeResourceDomainService {

    public void deactivateOfficeResource(OfficeResource resource, List<AllocationUsageStatus> relatedRequestStatuses) {
        if (!canModifyResource(relatedRequestStatuses)) {
            throw new IllegalStateException("Cannot deactivate a resource with active allocation requests");
        }
        resource.deactivate();
    }

    private boolean canModifyResource(List<AllocationUsageStatus> relatedRequestStatuses) {
        return relatedRequestStatuses.stream()
                .noneMatch(status -> status == AllocationUsageStatus.IN_USE
                        || status == AllocationUsageStatus.APPROVED
                        || status == AllocationUsageStatus.REQUESTED
                        || status == AllocationUsageStatus.OVERDUE);
    }
}
\end{DeskFlowCode}

\phantomsection
\addcontentsline{toc}{subsubsection}{MarkOfficeResourceAvailableDomainService}
\DeskFlowListingCaption{resources/application/domain/service/MarkOfficeResourceAvailableDomainService.java - MarkOfficeResourceAvailableDomainService}{lst:resources-application-domain-service-MarkOfficeResourceAvailableDomainService-java}
\begin{DeskFlowCode}
package com.deskflow.resources.application.domain.service;

import com.deskflow.resources.application.domain.model.OfficeResource;

public class MarkOfficeResourceAvailableDomainService {

    public void markOfficeResourceAvailable(OfficeResource resource) {
        resource.markAsAvailable();
    }
}
\end{DeskFlowCode}

\phantomsection
\addcontentsline{toc}{subsubsection}{MarkOfficeResourceUnavailableDomainService}
\DeskFlowListingCaption{resources/application/domain/service/MarkOfficeResourceUnavailableDomainService.java - MarkOfficeResourceUnavailableDomainService}{lst:resources-application-domain-service-MarkOfficeResourceUnavailableDomainService-java}
\begin{DeskFlowCode}
package com.deskflow.resources.application.domain.service;

import com.deskflow.resources.application.domain.model.OfficeResource;

public class MarkOfficeResourceUnavailableDomainService {

    public void markOfficeResourceUnavailable(OfficeResource resource) {
        resource.markAsUnavailable();
    }
}
\end{DeskFlowCode}

\phantomsection
\addcontentsline{toc}{subsubsection}{OfficeResourceFactory}
\DeskFlowListingCaption{resources/application/domain/service/OfficeResourceFactory.java - OfficeResourceFactory}{lst:resources-application-domain-service-OfficeResourceFactory-java}
\begin{DeskFlowCode}
package com.deskflow.resources.application.domain.service;

import com.deskflow.resources.application.domain.model.OfficeResource;
import com.deskflow.resources.application.domain.model.vo.Location;
import com.deskflow.common.application.domain.model.vo.ResourceId;
import com.deskflow.common.application.domain.model.vo.ResourceType;

import java.util.UUID;

public class OfficeResourceFactory {

    public OfficeResource create(ResourceType type, String name, Location location) {
        return OfficeResource.create(
                ResourceId.of(UUID.randomUUID()),
                type,
                name,
                location
        );
    }
}
\end{DeskFlowCode}

\phantomsection
\addcontentsline{toc}{subsubsection}{UpdateOfficeResourceDomainService}
\DeskFlowListingCaption{resources/application/domain/service/UpdateOfficeResourceDomainService.java - UpdateOfficeResourceDomainService}{lst:resources-application-domain-service-UpdateOfficeResourceDomainService-java}
\begin{DeskFlowCode}
package com.deskflow.resources.application.domain.service;

import com.deskflow.resources.application.domain.model.OfficeResource;
import com.deskflow.resources.application.domain.model.vo.Location;

public class UpdateOfficeResourceDomainService {

    public void updateOfficeResource(OfficeResource resource, String name, Location location) {
        resource.update(name, location);
    }
}
\end{DeskFlowCode}

\phantomsection
\addcontentsline{toc}{subsubsection}{ResourceView}
\DeskFlowListingCaption{resources/application/model/ResourceView.java - ResourceView}{lst:resources-application-model-ResourceView-java}
\begin{DeskFlowCode}
package com.deskflow.resources.application.model;

import com.deskflow.common.application.domain.model.vo.ResourceId;
import com.deskflow.common.application.domain.model.vo.ResourceType;
import com.deskflow.resources.application.domain.model.vo.Location;

import java.util.UUID;

public record ResourceView(
        UUID id,
        String type,
        String name,
        String building,
        int floor,
        String number,
        boolean active,
        boolean available
) {
    public static ResourceView from(ResourceId id, ResourceType type, String name, Location location, boolean active, boolean available) {
        return new ResourceView(
                id.getValue(),
                type.name(),
                name,
                location.getBuilding(),
                location.getFloor(),
                location.getNumber(),
                active,
                available
        );
    }
}
\end{DeskFlowCode}

\phantomsection
\addcontentsline{toc}{subsubsection}{ResourceManagementInputPort}
\DeskFlowListingCaption{resources/application/port/in/ResourceManagementInputPort.java - ResourceManagementInputPort}{lst:resources-application-port-in-ResourceManagementInputPort-java}
\begin{DeskFlowCode}
package com.deskflow.resources.application.port.in;

import com.deskflow.common.application.domain.model.vo.ResourceId;
import com.deskflow.common.application.domain.model.vo.ResourceType;
import com.deskflow.resources.application.model.ResourceView;

import java.util.List;
import java.util.Optional;

public interface ResourceManagementInputPort {

    ResourceId addResource(ResourceType type, String name, String building, int floor, String number);

    void updateResource(ResourceId resourceId, String name, String building, int floor, String number);

    void markUnavailable(ResourceId resourceId);

    void markAvailable(ResourceId resourceId);

    void deactivateResource(ResourceId resourceId);

    Optional<ResourceView> getResource(ResourceId resourceId);

    List<ResourceView> getActiveResources();
}
\end{DeskFlowCode}

\phantomsection
\addcontentsline{toc}{subsubsection}{OfficeResourceRepositoryPort}
\DeskFlowListingCaption{resources/application/port/out/OfficeResourceRepositoryPort.java - OfficeResourceRepositoryPort}{lst:resources-application-port-out-OfficeResourceRepositoryPort-java}
\begin{DeskFlowCode}
package com.deskflow.resources.application.port.out;

import com.deskflow.resources.application.domain.model.OfficeResource;
import com.deskflow.common.application.domain.model.vo.ResourceId;

import java.util.List;
import java.util.Optional;

public interface OfficeResourceRepositoryPort {

    void save(OfficeResource resource);

    Optional<OfficeResource> find(ResourceId id);

    List<OfficeResource> findActive();
}
\end{DeskFlowCode}

\phantomsection
\addcontentsline{toc}{subsubsection}{ResourceAllocationRequestPort}
\DeskFlowListingCaption{resources/application/port/out/ResourceAllocationRequestPort.java - ResourceAllocationRequestPort}{lst:resources-application-port-out-ResourceAllocationRequestPort-java}
\begin{DeskFlowCode}
package com.deskflow.resources.application.port.out;

import com.deskflow.resources.application.domain.model.vo.AllocationUsageStatus;
import com.deskflow.common.application.domain.model.vo.ResourceId;

import java.util.List;

public interface ResourceAllocationRequestPort {
    List<AllocationUsageStatus> findStatusesForResource(ResourceId resourceId);
}
\end{DeskFlowCode}

\phantomsection
\addcontentsline{toc}{subsubsection}{ResourceAppService}
\DeskFlowListingCaption{resources/application/service/ResourceAppService.java - ResourceAppService}{lst:resources-application-service-ResourceAppService-java}
\begin{DeskFlowCode}
package com.deskflow.resources.application.service;

import com.deskflow.common.application.port.out.DomainEventPublisherPort;
import com.deskflow.common.application.domain.DomainEvent;
import com.deskflow.resources.application.model.ResourceView;
import com.deskflow.resources.application.port.in.ResourceManagementInputPort;
import com.deskflow.resources.application.port.out.OfficeResourceRepositoryPort;
import com.deskflow.resources.application.port.out.ResourceAllocationRequestPort;
import com.deskflow.resources.application.domain.model.OfficeResource;
import com.deskflow.resources.application.domain.service.OfficeResourceFactory;
import com.deskflow.resources.application.domain.service.DeactivateOfficeResourceDomainService;
import com.deskflow.resources.application.domain.service.MarkOfficeResourceAvailableDomainService;
import com.deskflow.resources.application.domain.service.MarkOfficeResourceUnavailableDomainService;
import com.deskflow.resources.application.domain.service.UpdateOfficeResourceDomainService;
import com.deskflow.resources.application.domain.model.vo.Location;
import com.deskflow.common.application.domain.model.vo.ResourceId;
import com.deskflow.common.application.domain.model.vo.ResourceType;
import org.springframework.stereotype.Service;
import org.springframework.transaction.annotation.Transactional;

import java.util.List;
import java.util.NoSuchElementException;
import java.util.Optional;

@Service
@Transactional
public class ResourceAppService implements ResourceManagementInputPort {

    private final OfficeResourceRepositoryPort resourceRepository;
    private final ResourceAllocationRequestPort allocationRequestPort;
    private final OfficeResourceFactory factory;
    private final UpdateOfficeResourceDomainService updateOfficeResourceDomainService;
    private final MarkOfficeResourceUnavailableDomainService markOfficeResourceUnavailableDomainService;
    private final MarkOfficeResourceAvailableDomainService markOfficeResourceAvailableDomainService;
    private final DeactivateOfficeResourceDomainService deactivateOfficeResourceDomainService;
    private final DomainEventPublisherPort eventPublisher;

    public ResourceAppService(OfficeResourceRepositoryPort resourceRepository,
                            ResourceAllocationRequestPort allocationRequestPort,
                            OfficeResourceFactory factory,
                            UpdateOfficeResourceDomainService updateOfficeResourceDomainService,
                            MarkOfficeResourceUnavailableDomainService markOfficeResourceUnavailableDomainService,
                            MarkOfficeResourceAvailableDomainService markOfficeResourceAvailableDomainService,
                            DeactivateOfficeResourceDomainService deactivateOfficeResourceDomainService,
                            DomainEventPublisherPort eventPublisher) {
        this.resourceRepository = resourceRepository;
        this.allocationRequestPort = allocationRequestPort;
        this.factory = factory;
        this.updateOfficeResourceDomainService = updateOfficeResourceDomainService;
        this.markOfficeResourceUnavailableDomainService = markOfficeResourceUnavailableDomainService;
        this.markOfficeResourceAvailableDomainService = markOfficeResourceAvailableDomainService;
        this.deactivateOfficeResourceDomainService = deactivateOfficeResourceDomainService;
        this.eventPublisher = eventPublisher;
    }

    @Override
    public ResourceId addResource(ResourceType type, String name, String building, int floor, String number) {
        Location location = new Location(building, floor, number);
        OfficeResource resource = factory.create(type, name, location);
        resourceRepository.save(resource);
        return resource.getId();
    }

    @Override
    public void updateResource(ResourceId resourceId, String name, String building, int floor, String number) {
        OfficeResource resource = findResource(resourceId);
        Location location = new Location(building, floor, number);
        updateOfficeResourceDomainService.updateOfficeResource(resource, name, location);
        resourceRepository.save(resource);
        publishAndClear(resource);
    }

    @Override
    public void markUnavailable(ResourceId resourceId) {
        OfficeResource resource = findResource(resourceId);
        markOfficeResourceUnavailableDomainService.markOfficeResourceUnavailable(resource);
        resourceRepository.save(resource);
        publishAndClear(resource);
    }

    @Override
    public void markAvailable(ResourceId resourceId) {
        OfficeResource resource = findResource(resourceId);
        markOfficeResourceAvailableDomainService.markOfficeResourceAvailable(resource);
        resourceRepository.save(resource);
        publishAndClear(resource);
    }

    @Override
    public void deactivateResource(ResourceId resourceId) {
        OfficeResource resource = findResource(resourceId);
        deactivateOfficeResourceDomainService.deactivateOfficeResource(resource, allocationRequestPort.findStatusesForResource(resourceId));
        resourceRepository.save(resource);
        publishAndClear(resource);
    }

    @Override
    @Transactional(readOnly = true)
    public Optional<ResourceView> getResource(ResourceId resourceId) {
        return resourceRepository.find(resourceId).map(this::toView);
    }

    @Override
    @Transactional(readOnly = true)
    public List<ResourceView> getActiveResources() {
        return resourceRepository.findActive().stream().map(this::toView).toList();
    }

    private ResourceView toView(OfficeResource r) {
        return ResourceView.from(r.getId(), r.getType(), r.getName(), r.getLocation(), r.isActive(), r.isAvailable());
    }

    private OfficeResource findResource(ResourceId id) {
        return resourceRepository.find(id)
                .orElseThrow(() -> new NoSuchElementException("Resource not found: " + id));
    }

    private void publishAndClear(OfficeResource resource) {
        for (DomainEvent event : resource.flushEvents()) {
            eventPublisher.publish(event);
        }
    }
}
\end{DeskFlowCode}

\phantomsection
\addcontentsline{toc}{subsubsection}{ResourceAvailabilityEventListener}
\DeskFlowListingCaption{resources/infrastructure/in/events/ResourceAvailabilityEventListener.java - ResourceAvailabilityEventListener}{lst:resources-infrastructure-in-events-ResourceAvailabilityEventListener-java}
\begin{DeskFlowCode}
package com.deskflow.resources.infrastructure.in.events;

import com.deskflow.allocation.application.domain.model.events.ResourceIssued;
import com.deskflow.allocation.application.domain.model.events.ResourceReturned;
import com.deskflow.resources.application.port.in.ResourceManagementInputPort;
import org.springframework.stereotype.Component;
import org.springframework.transaction.event.TransactionPhase;
import org.springframework.transaction.event.TransactionalEventListener;

@Component
public class ResourceAvailabilityEventListener {

    private final ResourceManagementInputPort resourceManagementInputPort;

    public ResourceAvailabilityEventListener(ResourceManagementInputPort resourceManagementInputPort) {
        this.resourceManagementInputPort = resourceManagementInputPort;
    }

    @TransactionalEventListener(phase = TransactionPhase.AFTER_COMMIT)
    public void onResourceIssued(ResourceIssued event) {
        resourceManagementInputPort.markUnavailable(event.resourceId());
    }

    @TransactionalEventListener(phase = TransactionPhase.AFTER_COMMIT)
    public void onResourceReturned(ResourceReturned event) {
        resourceManagementInputPort.markAvailable(event.resourceId());
    }
}
\end{DeskFlowCode}

\phantomsection
\addcontentsline{toc}{subsubsection}{ResourceController}
\DeskFlowListingCaption{resources/infrastructure/in/web/ResourceController.java - ResourceController}{lst:resources-infrastructure-in-web-ResourceController-java}
\begin{DeskFlowCode}
package com.deskflow.resources.infrastructure.in.web;

import com.deskflow.resources.infrastructure.in.web.dto.ResourceRequestDTO;
import com.deskflow.resources.application.port.in.ResourceManagementInputPort;
import com.deskflow.common.application.domain.model.vo.ResourceId;
import com.deskflow.common.application.domain.model.vo.ResourceType;
import jakarta.validation.Valid;
import org.springframework.http.ResponseEntity;
import org.springframework.web.bind.annotation.*;

import java.net.URI;
import java.util.Map;
import java.util.UUID;

@RestController
@RequestMapping("/api/resources")
public class ResourceController {

    private final ResourceManagementInputPort resourceManagementInputPort;

    public ResourceController(ResourceManagementInputPort resourceManagementInputPort) {
        this.resourceManagementInputPort = resourceManagementInputPort;
    }

    @PostMapping
    public ResponseEntity<?> addResource(@Valid @RequestBody ResourceRequestDTO dto) {
        ResourceId id = resourceManagementInputPort.addResource(
                ResourceType.valueOf(dto.type()), dto.name(), dto.building(), dto.floor(), dto.number());
        return ResponseEntity.created(URI.create("/api/resources/" + id))
                .body(Map.of("resourceId", id.getValue()));
    }

    @PutMapping("/{id}")
    public ResponseEntity<Void> updateResource(@PathVariable UUID id, @Valid @RequestBody ResourceRequestDTO dto) {
        resourceManagementInputPort.updateResource(ResourceId.of(id), dto.name(), dto.building(), dto.floor(), dto.number());
        return ResponseEntity.ok().build();
    }

    @PutMapping("/{id}/unavailable")
    public ResponseEntity<Void> markUnavailable(@PathVariable UUID id) {
        resourceManagementInputPort.markUnavailable(ResourceId.of(id));
        return ResponseEntity.ok().build();
    }

    @PutMapping("/{id}/available")
    public ResponseEntity<Void> markAvailable(@PathVariable UUID id) {
        resourceManagementInputPort.markAvailable(ResourceId.of(id));
        return ResponseEntity.ok().build();
    }

    @DeleteMapping("/{id}")
    public ResponseEntity<Void> deactivate(@PathVariable UUID id) {
        resourceManagementInputPort.deactivateResource(ResourceId.of(id));
        return ResponseEntity.ok().build();
    }
}
\end{DeskFlowCode}

\phantomsection
\addcontentsline{toc}{subsubsection}{ResourceRequestDTO}
\DeskFlowListingCaption{resources/infrastructure/in/web/dto/ResourceRequestDTO.java - ResourceRequestDTO}{lst:resources-infrastructure-in-web-dto-ResourceRequestDTO-java}
\begin{DeskFlowCode}
package com.deskflow.resources.infrastructure.in.web.dto;

import jakarta.validation.constraints.NotBlank;
import jakarta.validation.constraints.NotNull;

public record ResourceRequestDTO(
        @NotNull String type,
        @NotBlank String name,
        @NotBlank String building,
        int floor,
        @NotBlank String number
) {}
\end{DeskFlowCode}

\phantomsection
\addcontentsline{toc}{subsubsection}{AllocationRequestAdapter}
\DeskFlowListingCaption{resources/infrastructure/out/allocation/AllocationRequestAdapter.java - AllocationRequestAdapter}{lst:resources-infrastructure-out-allocation-AllocationRequestAdapter-java}
\begin{DeskFlowCode}
package com.deskflow.resources.infrastructure.out.allocation;

import com.deskflow.allocation.application.port.in.AllocationRequestInputPort;
import com.deskflow.resources.application.port.out.ResourceAllocationRequestPort;
import com.deskflow.resources.application.domain.model.vo.AllocationUsageStatus;
import com.deskflow.common.application.domain.model.vo.ResourceId;
import org.springframework.context.annotation.Lazy;
import org.springframework.stereotype.Component;

import java.util.List;

@Component
public class AllocationRequestAdapter implements ResourceAllocationRequestPort {

    private final AllocationRequestInputPort allocationRequestInputPort;

    public AllocationRequestAdapter(@Lazy AllocationRequestInputPort allocationRequestInputPort) {
        this.allocationRequestInputPort = allocationRequestInputPort;
    }

    @Override
    public List<AllocationUsageStatus> findStatusesForResource(ResourceId resourceId) {
        return allocationRequestInputPort.findActiveForResource(resourceId).stream()
                .map(request -> request.status())
                .map(this::toUsageStatus)
                .toList();
    }

    private AllocationUsageStatus toUsageStatus(String status) {
        return switch (status) {
            case "SUBMITTED" -> AllocationUsageStatus.REQUESTED;
            case "APPROVED" -> AllocationUsageStatus.APPROVED;
            case "ACTIVE" -> AllocationUsageStatus.IN_USE;
            case "OVERDUE" -> AllocationUsageStatus.OVERDUE;
            case "REJECTED", "CANCELLED", "COMPLETED" -> AllocationUsageStatus.CLOSED;
            default -> throw new IllegalArgumentException("Unknown allocation status: " + status);
        };
    }
}
\end{DeskFlowCode}

\phantomsection
\addcontentsline{toc}{subsubsection}{HibernateOfficeResourceRepositoryAdapter}
\DeskFlowListingCaption{resources/infrastructure/out/persistence/HibernateOfficeResourceRepositoryAdapter.java - HibernateOfficeResourceRepositoryAdapter}{lst:resources-infrastructure-out-persistence-HibernateOfficeResourceRepositoryAdapter-java}
\begin{DeskFlowCode}
package com.deskflow.resources.infrastructure.out.persistence;

import com.deskflow.resources.application.port.out.OfficeResourceRepositoryPort;
import com.deskflow.resources.application.domain.model.OfficeResource;
import com.deskflow.common.application.domain.model.vo.ResourceId;
import org.springframework.stereotype.Repository;

import java.util.List;
import java.util.Optional;

@Repository
public class HibernateOfficeResourceRepositoryAdapter implements OfficeResourceRepositoryPort {

    private final SpringDataOfficeResourceRepository jpa;
    private final OfficeResourceMapper mapper;

    public HibernateOfficeResourceRepositoryAdapter(SpringDataOfficeResourceRepository jpa, OfficeResourceMapper mapper) {
        this.jpa = jpa;
        this.mapper = mapper;
    }

    @Override
    public void save(OfficeResource resource) {
        jpa.findById(resource.getId().getValue()).ifPresentOrElse(
                managed -> mapper.updateEntity(resource, managed),
                () -> jpa.save(mapper.toEntity(resource))
        );
    }

    @Override
    public Optional<OfficeResource> find(ResourceId id) {
        return jpa.findById(id.getValue()).map(mapper::toDomain);
    }

    @Override
    public List<OfficeResource> findActive() {
        return jpa.findByActiveTrue().stream().map(mapper::toDomain).toList();
    }
}
\end{DeskFlowCode}

\phantomsection
\addcontentsline{toc}{subsubsection}{OfficeResourceEntity}
\DeskFlowListingCaption{resources/infrastructure/out/persistence/OfficeResourceEntity.java - OfficeResourceEntity}{lst:resources-infrastructure-out-persistence-OfficeResourceEntity-java}
\begin{DeskFlowCode}
package com.deskflow.resources.infrastructure.out.persistence;

import jakarta.persistence.*;

import java.util.UUID;

@Entity
@Table(name = "office_resources")
public class OfficeResourceEntity {

    @Id
    private UUID id;

    @Column(name = "type", nullable = false)
    private String type;

    @Column(name = "name", nullable = false)
    private String name;

    @Column(name = "building", nullable = false)
    private String building;

    @Column(name = "floor", nullable = false)
    private int floor;

    @Column(name = "number", nullable = false)
    private String number;

    @Column(name = "active", nullable = false)
    private boolean active;

    @Column(name = "available", nullable = false)
    private boolean available;

    public UUID getId() { return id; }
    public void setId(UUID id) { this.id = id; }
    public String getType() { return type; }
    public void setType(String type) { this.type = type; }
    public String getName() { return name; }
    public void setName(String name) { this.name = name; }
    public String getBuilding() { return building; }
    public void setBuilding(String building) { this.building = building; }
    public int getFloor() { return floor; }
    public void setFloor(int floor) { this.floor = floor; }
    public String getNumber() { return number; }
    public void setNumber(String number) { this.number = number; }
    public boolean isActive() { return active; }
    public void setActive(boolean active) { this.active = active; }
    public boolean isAvailable() { return available; }
    public void setAvailable(boolean available) { this.available = available; }
}
\end{DeskFlowCode}

\phantomsection
\addcontentsline{toc}{subsubsection}{OfficeResourceMapper}
\DeskFlowListingCaption{resources/infrastructure/out/persistence/OfficeResourceMapper.java - OfficeResourceMapper}{lst:resources-infrastructure-out-persistence-OfficeResourceMapper-java}
\begin{DeskFlowCode}
package com.deskflow.resources.infrastructure.out.persistence;

import com.deskflow.resources.application.domain.model.OfficeResource;
import com.deskflow.resources.application.domain.model.vo.Location;
import com.deskflow.common.application.domain.model.vo.ResourceId;
import com.deskflow.common.application.domain.model.vo.ResourceType;
import org.springframework.stereotype.Component;

@Component
class OfficeResourceMapper {

    OfficeResourceEntity toEntity(OfficeResource resource) {
        OfficeResourceEntity e = new OfficeResourceEntity();
        updateEntity(resource, e);
        return e;
    }

    void updateEntity(OfficeResource resource, OfficeResourceEntity e) {
        e.setId(resource.getId().getValue());
        e.setType(resource.getType().name());
        e.setName(resource.getName());
        e.setBuilding(resource.getLocation().getBuilding());
        e.setFloor(resource.getLocation().getFloor());
        e.setNumber(resource.getLocation().getNumber());
        e.setActive(resource.isActive());
        e.setAvailable(resource.isAvailable());
    }

    OfficeResource toDomain(OfficeResourceEntity e) {
        return OfficeResource.reconstitute(
                ResourceId.of(e.getId()),
                ResourceType.valueOf(e.getType()),
                e.getName(),
                new Location(e.getBuilding(), e.getFloor(), e.getNumber()),
                e.isActive(),
                e.isAvailable()
        );
    }
}
\end{DeskFlowCode}

\phantomsection
\addcontentsline{toc}{subsubsection}{SpringDataOfficeResourceRepository}
\DeskFlowListingCaption{resources/infrastructure/out/persistence/SpringDataOfficeResourceRepository.java - SpringDataOfficeResourceRepository}{lst:resources-infrastructure-out-persistence-SpringDataOfficeResourceRepository-java}
\begin{DeskFlowCode}
package com.deskflow.resources.infrastructure.out.persistence;

import org.springframework.data.jpa.repository.JpaRepository;

import java.util.List;
import java.util.UUID;

public interface SpringDataOfficeResourceRepository extends JpaRepository<OfficeResourceEntity, UUID> {
    List<OfficeResourceEntity> findByActiveTrue();
}
\end{DeskFlowCode}

\subsection{Kontekst Administration - administracja}

\phantomsection
\addcontentsline{toc}{subsubsection}{AdministeredAllocationRequest}
\DeskFlowListingCaption{administration/application/domain/model/AdministeredAllocationRequest.java - AdministeredAllocationRequest}{lst:administration-application-domain-model-AdministeredAllocationRequest-java}
\begin{DeskFlowCode}
package com.deskflow.administration.application.domain.model;

public record AdministeredAllocationRequest(
        AdministeredAllocationRequestId id,
        AdministeredAllocationStatus status
) {
    public void ensureCanApprove() {
        requireStatus(AdministeredAllocationStatus.SUBMITTED, "approve");
    }

    public void ensureCanReject() {
        requireStatus(AdministeredAllocationStatus.SUBMITTED, "reject");
    }

    public void ensureCanIssue() {
        requireStatus(AdministeredAllocationStatus.APPROVED, "issue resource");
    }

    public void ensureCanReturn() {
        if (status != AdministeredAllocationStatus.ACTIVE && status != AdministeredAllocationStatus.OVERDUE) {
            throw new IllegalStateException(
                    "Operation 'return resource' requires ACTIVE or OVERDUE status, current: " + status);
        }
    }

    private void requireStatus(AdministeredAllocationStatus required, String operation) {
        if (status != required) {
            throw new IllegalStateException(
                    "Operation '" + operation + "' requires status " + required + ", current: " + status);
        }
    }
}
\end{DeskFlowCode}

\phantomsection
\addcontentsline{toc}{subsubsection}{AdministeredAllocationRequestId}
\DeskFlowListingCaption{administration/application/domain/model/AdministeredAllocationRequestId.java - AdministeredAllocationRequestId}{lst:administration-application-domain-model-AdministeredAllocationRequestId-java}
\begin{DeskFlowCode}
package com.deskflow.administration.application.domain.model;

import java.util.Objects;
import java.util.UUID;

public final class AdministeredAllocationRequestId {
    private final UUID value;

    private AdministeredAllocationRequestId(UUID value) {
        this.value = Objects.requireNonNull(value, "value");
    }

    public static AdministeredAllocationRequestId of(UUID value) {
        return new AdministeredAllocationRequestId(value);
    }

    public UUID getValue() {
        return value;
    }

    @Override
    public boolean equals(Object o) {
        if (this == o) return true;
        if (!(o instanceof AdministeredAllocationRequestId that)) return false;
        return value.equals(that.value);
    }

    @Override
    public int hashCode() {
        return value.hashCode();
    }

    @Override
    public String toString() {
        return value.toString();
    }
}
\end{DeskFlowCode}

\phantomsection
\addcontentsline{toc}{subsubsection}{AdministeredAllocationStatus}
\DeskFlowListingCaption{administration/application/domain/model/AdministeredAllocationStatus.java - AdministeredAllocationStatus}{lst:administration-application-domain-model-AdministeredAllocationStatus-java}
\begin{DeskFlowCode}
package com.deskflow.administration.application.domain.model;

public enum AdministeredAllocationStatus {
    SUBMITTED,
    APPROVED,
    ACTIVE,
    OVERDUE,
    REJECTED,
    CANCELLED,
    COMPLETED;

    public static AdministeredAllocationStatus from(String status) {
        return AdministeredAllocationStatus.valueOf(status);
    }
}
\end{DeskFlowCode}

\phantomsection
\addcontentsline{toc}{subsubsection}{AdministeredAllocationView}
\DeskFlowListingCaption{administration/application/model/AdministeredAllocationView.java - AdministeredAllocationView}{lst:administration-application-model-AdministeredAllocationView-java}
\begin{DeskFlowCode}
package com.deskflow.administration.application.model;

import java.time.LocalDate;
import java.time.LocalDateTime;
import java.util.UUID;

public record AdministeredAllocationView(
        UUID id,
        UUID employeeId,
        UUID resourceId,
        String status,
        LocalDate dateFrom,
        LocalDate dateTo,
        String justification,
        LocalDateTime submittedAt,
        LocalDateTime issuedAt,
        LocalDateTime returnedAt,
        String rejectionReason
) {}
\end{DeskFlowCode}

\phantomsection
\addcontentsline{toc}{subsubsection}{AdministrationInputPort}
\DeskFlowListingCaption{administration/application/port/in/AdministrationInputPort.java - AdministrationInputPort}{lst:administration-application-port-in-AdministrationInputPort-java}
\begin{DeskFlowCode}
package com.deskflow.administration.application.port.in;

import com.deskflow.administration.application.model.AdministeredAllocationView;
import com.deskflow.administration.application.domain.model.AdministeredAllocationRequestId;
import com.deskflow.administration.application.domain.model.AdministeredAllocationStatus;

import java.util.List;

public interface AdministrationInputPort {

    List<AdministeredAllocationView> listRequests(AdministeredAllocationStatus status);

    void approveRequest(AdministeredAllocationRequestId requestId);

    void rejectRequest(AdministeredAllocationRequestId requestId, String reason);

    void issueResource(AdministeredAllocationRequestId requestId);

    void returnResource(AdministeredAllocationRequestId requestId);

    List<AdministeredAllocationView> monitorActiveAllocations();

    void detectOverdueAllocations();
}
\end{DeskFlowCode}

\phantomsection
\addcontentsline{toc}{subsubsection}{AdministrationAllocationPort}
\DeskFlowListingCaption{administration/application/port/out/AdministrationAllocationPort.java - AdministrationAllocationPort}{lst:administration-application-port-out-AdministrationAllocationPort-java}
\begin{DeskFlowCode}
package com.deskflow.administration.application.port.out;

import com.deskflow.administration.application.model.AdministeredAllocationView;
import com.deskflow.administration.application.domain.model.AdministeredAllocationRequestId;
import com.deskflow.administration.application.domain.model.AdministeredAllocationStatus;

import java.util.List;
import java.util.Optional;

public interface AdministrationAllocationPort {
    Optional<AdministeredAllocationView> getRequest(AdministeredAllocationRequestId requestId);

    List<AdministeredAllocationView> listRequests(AdministeredAllocationStatus status);

    void approveRequest(AdministeredAllocationRequestId requestId);

    void rejectRequest(AdministeredAllocationRequestId requestId, String reason);

    void issueResource(AdministeredAllocationRequestId requestId);

    void returnResource(AdministeredAllocationRequestId requestId);

    List<AdministeredAllocationView> monitorActiveAllocations();

    void detectOverdueAllocations();
}
\end{DeskFlowCode}

\phantomsection
\addcontentsline{toc}{subsubsection}{AdministrationAppService}
\DeskFlowListingCaption{administration/application/service/AdministrationAppService.java - AdministrationAppService}{lst:administration-application-service-AdministrationAppService-java}
\begin{DeskFlowCode}
package com.deskflow.administration.application.service;

import com.deskflow.administration.application.model.AdministeredAllocationView;
import com.deskflow.administration.application.port.in.AdministrationInputPort;
import com.deskflow.administration.application.port.out.AdministrationAllocationPort;
import com.deskflow.administration.application.domain.model.AdministeredAllocationRequest;
import com.deskflow.administration.application.domain.model.AdministeredAllocationRequestId;
import com.deskflow.administration.application.domain.model.AdministeredAllocationStatus;
import org.springframework.stereotype.Service;
import org.springframework.transaction.annotation.Transactional;

import java.util.List;
import java.util.NoSuchElementException;

@Service
@Transactional
public class AdministrationAppService implements AdministrationInputPort {

    private final AdministrationAllocationPort allocationPort;

    public AdministrationAppService(AdministrationAllocationPort allocationPort) {
        this.allocationPort = allocationPort;
    }

    @Override
    @Transactional(readOnly = true)
    public List<AdministeredAllocationView> listRequests(AdministeredAllocationStatus status) {
        return allocationPort.listRequests(status);
    }

    @Override
    public void approveRequest(AdministeredAllocationRequestId requestId) {
        loadRequest(requestId).ensureCanApprove();
        allocationPort.approveRequest(requestId);
    }

    @Override
    public void rejectRequest(AdministeredAllocationRequestId requestId, String reason) {
        loadRequest(requestId).ensureCanReject();
        if (reason == null || reason.isBlank()) {
            throw new IllegalArgumentException("Rejection reason cannot be blank");
        }
        allocationPort.rejectRequest(requestId, reason);
    }

    @Override
    public void issueResource(AdministeredAllocationRequestId requestId) {
        loadRequest(requestId).ensureCanIssue();
        allocationPort.issueResource(requestId);
    }

    @Override
    public void returnResource(AdministeredAllocationRequestId requestId) {
        loadRequest(requestId).ensureCanReturn();
        allocationPort.returnResource(requestId);
    }

    @Override
    @Transactional(readOnly = true)
    public List<AdministeredAllocationView> monitorActiveAllocations() {
        return allocationPort.monitorActiveAllocations();
    }

    @Override
    public void detectOverdueAllocations() {
        allocationPort.detectOverdueAllocations();
    }

    private AdministeredAllocationRequest loadRequest(AdministeredAllocationRequestId requestId) {
        AdministeredAllocationView view = allocationPort.getRequest(requestId)
                .orElseThrow(() -> new NoSuchElementException("Allocation request not found: " + requestId));
        return new AdministeredAllocationRequest(
                requestId,
                AdministeredAllocationStatus.from(view.status())
        );
    }
}
\end{DeskFlowCode}

\phantomsection
\addcontentsline{toc}{subsubsection}{OverdueAllocationScheduler}
\DeskFlowListingCaption{administration/application/service/OverdueAllocationScheduler.java - OverdueAllocationScheduler}{lst:administration-application-service-OverdueAllocationScheduler-java}
\begin{DeskFlowCode}
package com.deskflow.administration.application.service;

import org.springframework.scheduling.annotation.Scheduled;
import org.springframework.stereotype.Component;

@Component
public class OverdueAllocationScheduler {

    private final AdministrationAppService administrationService;

    public OverdueAllocationScheduler(AdministrationAppService administrationService) {
        this.administrationService = administrationService;
    }

    @Scheduled(fixedRateString = "${app.scheduler.przeterminowania-rate-ms:900000}")
    public void detectOverdueAllocations() {
        administrationService.detectOverdueAllocations();
    }
}
\end{DeskFlowCode}

\phantomsection
\addcontentsline{toc}{subsubsection}{AdministrationController}
\DeskFlowListingCaption{administration/infrastructure/in/web/AdministrationController.java - AdministrationController}{lst:administration-infrastructure-in-web-AdministrationController-java}
\begin{DeskFlowCode}
package com.deskflow.administration.infrastructure.in.web;

import com.deskflow.administration.application.model.AdministeredAllocationView;
import com.deskflow.administration.infrastructure.in.web.dto.RejectRequestDTO;
import com.deskflow.allocation.infrastructure.in.web.dto.AllocationRequestResponseDTO;
import com.deskflow.administration.application.port.in.AdministrationInputPort;
import com.deskflow.administration.application.domain.model.AdministeredAllocationRequestId;
import com.deskflow.administration.application.domain.model.AdministeredAllocationStatus;
import com.deskflow.common.application.domain.model.vo.EmployeeId;
import com.deskflow.employees.application.port.in.EmployeeManagementInputPort;
import com.deskflow.employees.application.domain.model.Employee;
import jakarta.validation.Valid;
import org.springframework.http.ResponseEntity;
import org.springframework.web.bind.annotation.*;

import java.util.List;
import java.util.Map;
import java.util.UUID;
import java.util.stream.Collectors;

@RestController
@RequestMapping("/api/administration")
public class AdministrationController {

    private final AdministrationInputPort administrationInputPort;
    private final EmployeeManagementInputPort employeeManagementInputPort;

    public AdministrationController(AdministrationInputPort administrationInputPort,
                                    EmployeeManagementInputPort employeeManagementInputPort) {
        this.administrationInputPort = administrationInputPort;
        this.employeeManagementInputPort = employeeManagementInputPort;
    }

    @GetMapping("/requests")
    public ResponseEntity<List<AllocationRequestResponseDTO>> listRequests(
            @RequestParam(required = false) String status) {
        AdministeredAllocationStatus statusEnum = status != null ? AdministeredAllocationStatus.valueOf(status) : null;
        return ResponseEntity.ok(enrich(administrationInputPort.listRequests(statusEnum)));
    }

    @PutMapping("/requests/{id}/approve")
    public ResponseEntity<Void> approve(@PathVariable UUID id) {
        administrationInputPort.approveRequest(AdministeredAllocationRequestId.of(id));
        return ResponseEntity.ok().build();
    }

    @PutMapping("/requests/{id}/reject")
    public ResponseEntity<Void> reject(@PathVariable UUID id, @Valid @RequestBody RejectRequestDTO dto) {
        administrationInputPort.rejectRequest(AdministeredAllocationRequestId.of(id), dto.reason());
        return ResponseEntity.ok().build();
    }

    @PutMapping("/requests/{id}/issue")
    public ResponseEntity<Void> issue(@PathVariable UUID id) {
        administrationInputPort.issueResource(AdministeredAllocationRequestId.of(id));
        return ResponseEntity.ok().build();
    }

    @PutMapping("/requests/{id}/return")
    public ResponseEntity<Void> returnResource(@PathVariable UUID id) {
        administrationInputPort.returnResource(AdministeredAllocationRequestId.of(id));
        return ResponseEntity.ok().build();
    }

    @GetMapping("/allocations")
    public ResponseEntity<List<AllocationRequestResponseDTO>> monitorActive() {
        return ResponseEntity.ok(enrich(administrationInputPort.monitorActiveAllocations()));
    }

    private List<AllocationRequestResponseDTO> enrich(List<AdministeredAllocationView> requests) {
        Map<UUID, Employee> cache = requests.stream()
                .map(AdministeredAllocationView::employeeId)
                .distinct()
                .collect(Collectors.toMap(
                        id -> id,
                        id -> employeeManagementInputPort.getEmployee(EmployeeId.of(id)).orElse(null),
                        (a, b) -> a
                ));
        return requests.stream()
                .map(r -> {
                    Employee emp = cache.get(r.employeeId());
                    return toResponse(r,
                            emp != null ? emp.getFirstName() : null,
                            emp != null ? emp.getLastName() : null);
                })
                .toList();
    }

    private AllocationRequestResponseDTO toResponse(AdministeredAllocationView r, String firstName, String lastName) {
        return new AllocationRequestResponseDTO(
                r.id(),
                r.employeeId(),
                firstName,
                lastName,
                r.resourceId(),
                r.status(),
                r.dateFrom(),
                r.dateTo(),
                r.justification(),
                r.submittedAt(),
                r.issuedAt(),
                r.returnedAt(),
                r.rejectionReason()
        );
    }
}
\end{DeskFlowCode}

\phantomsection
\addcontentsline{toc}{subsubsection}{RejectRequestDTO}
\DeskFlowListingCaption{administration/infrastructure/in/web/dto/RejectRequestDTO.java - RejectRequestDTO}{lst:administration-infrastructure-in-web-dto-RejectRequestDTO-java}
\begin{DeskFlowCode}
package com.deskflow.administration.infrastructure.in.web.dto;

import jakarta.validation.constraints.NotBlank;

public record RejectRequestDTO(@NotBlank String reason) {}
\end{DeskFlowCode}

\phantomsection
\addcontentsline{toc}{subsubsection}{AllocationAdministrationAdapter}
\DeskFlowListingCaption{administration/infrastructure/out/allocation/AllocationAdministrationAdapter.java - AllocationAdministrationAdapter}{lst:administration-infrastructure-out-allocation-AllocationAdministrationAdapter-java}
\begin{DeskFlowCode}
package com.deskflow.administration.infrastructure.out.allocation;

import com.deskflow.administration.application.model.AdministeredAllocationView;
import com.deskflow.administration.application.port.out.AdministrationAllocationPort;
import com.deskflow.administration.application.domain.model.AdministeredAllocationRequestId;
import com.deskflow.administration.application.domain.model.AdministeredAllocationStatus;
import com.deskflow.allocation.application.model.AllocationRequestView;
import com.deskflow.allocation.application.port.in.AllocationAdministrationInputPort;
import com.deskflow.allocation.application.domain.model.vo.AllocationRequestId;
import com.deskflow.allocation.application.domain.model.vo.RequestStatus;
import org.springframework.stereotype.Component;

import java.util.List;
import java.util.Optional;

@Component
public class AllocationAdministrationAdapter implements AdministrationAllocationPort {

    private final AllocationAdministrationInputPort allocationAdministrationInputPort;

    public AllocationAdministrationAdapter(AllocationAdministrationInputPort allocationAdministrationInputPort) {
        this.allocationAdministrationInputPort = allocationAdministrationInputPort;
    }

    @Override
    public Optional<AdministeredAllocationView> getRequest(AdministeredAllocationRequestId requestId) {
        return allocationAdministrationInputPort.getRequest(toAllocationId(requestId)).map(this::toAdminView);
    }

    @Override
    public List<AdministeredAllocationView> listRequests(AdministeredAllocationStatus status) {
        RequestStatus allocationStatus = status != null ? RequestStatus.valueOf(status.name()) : null;
        return allocationAdministrationInputPort.listRequests(allocationStatus).stream()
                .map(this::toAdminView)
                .toList();
    }

    @Override
    public void approveRequest(AdministeredAllocationRequestId requestId) {
        allocationAdministrationInputPort.approveRequest(toAllocationId(requestId));
    }

    @Override
    public void rejectRequest(AdministeredAllocationRequestId requestId, String reason) {
        allocationAdministrationInputPort.rejectRequest(toAllocationId(requestId), reason);
    }

    @Override
    public void issueResource(AdministeredAllocationRequestId requestId) {
        allocationAdministrationInputPort.issueResource(toAllocationId(requestId));
    }

    @Override
    public void returnResource(AdministeredAllocationRequestId requestId) {
        allocationAdministrationInputPort.returnResource(toAllocationId(requestId));
    }

    @Override
    public List<AdministeredAllocationView> monitorActiveAllocations() {
        return allocationAdministrationInputPort.monitorActiveAllocations().stream()
                .map(this::toAdminView)
                .toList();
    }

    @Override
    public void detectOverdueAllocations() {
        allocationAdministrationInputPort.detectOverdueAllocations();
    }

    private AllocationRequestId toAllocationId(AdministeredAllocationRequestId requestId) {
        return AllocationRequestId.of(requestId.getValue());
    }

    private AdministeredAllocationView toAdminView(AllocationRequestView view) {
        return new AdministeredAllocationView(
                view.id(),
                view.employeeId(),
                view.resourceId(),
                view.status(),
                view.dateFrom(),
                view.dateTo(),
                view.justification(),
                view.submittedAt(),
                view.issuedAt(),
                view.returnedAt(),
                view.rejectionReason()
        );
    }
}
\end{DeskFlowCode}

\subsection{Kontekst Employees - pracownicy}

\phantomsection
\addcontentsline{toc}{subsubsection}{Employee}
\DeskFlowListingCaption{employees/application/domain/model/Employee.java - Employee}{lst:employees-application-domain-model-Employee-java}
\begin{DeskFlowCode}
package com.deskflow.employees.application.domain.model;

import com.deskflow.common.application.domain.model.vo.EmployeeId;
import com.deskflow.employees.application.domain.model.vo.Role;
import com.deskflow.common.application.domain.model.vo.ResourceType;

import java.util.EnumMap;
import java.util.Map;

public class Employee {

    private final EmployeeId id;
    private final String email;
    private final String firstName;
    private final String lastName;
    private final Role role;
    private final Map<ResourceType, Integer> limits;

    public Employee(EmployeeId id, String email, String firstName, String lastName,
                    Role role, Map<ResourceType, Integer> limits) {
        this.id = id;
        this.email = email;
        this.firstName = firstName;
        this.lastName = lastName;
        this.role = role;
        this.limits = new EnumMap<>(ResourceType.class);
        this.limits.putAll(limits);
    }

    public int limitForType(ResourceType type) {
        return limits.getOrDefault(type, 1);
    }

    public void changeLimit(ResourceType type, int amount) {
        if (amount < 0) {
            throw new IllegalArgumentException("Allocation limit cannot be negative");
        }
        limits.put(type, amount);
    }

    public EmployeeId getId() { return id; }
    public String getEmail() { return email; }
    public String getFirstName() { return firstName; }
    public String getLastName() { return lastName; }
    public Role getRole() { return role; }
    public Map<ResourceType, Integer> getLimits() { return Map.copyOf(limits); }
}
\end{DeskFlowCode}

\phantomsection
\addcontentsline{toc}{subsubsection}{Role}
\DeskFlowListingCaption{employees/application/domain/model/vo/Role.java - Role}{lst:employees-application-domain-model-vo-Role-java}
\begin{DeskFlowCode}
package com.deskflow.employees.application.domain.model.vo;

public enum Role {
    EMPLOYEE,
    MANAGER,
    IT_ADMIN
}
\end{DeskFlowCode}

\phantomsection
\addcontentsline{toc}{subsubsection}{EmployeeCredentials}
\DeskFlowListingCaption{employees/application/model/EmployeeCredentials.java - EmployeeCredentials}{lst:employees-application-model-EmployeeCredentials-java}
\begin{DeskFlowCode}
package com.deskflow.employees.application.model;

public record EmployeeCredentials(String email, String passwordHash, String role) {}
\end{DeskFlowCode}

\phantomsection
\addcontentsline{toc}{subsubsection}{EmployeeManagementInputPort}
\DeskFlowListingCaption{employees/application/port/in/EmployeeManagementInputPort.java - EmployeeManagementInputPort}{lst:employees-application-port-in-EmployeeManagementInputPort-java}
\begin{DeskFlowCode}
package com.deskflow.employees.application.port.in;

import com.deskflow.employees.application.model.EmployeeCredentials;
import com.deskflow.employees.application.domain.model.Employee;
import com.deskflow.common.application.domain.model.vo.EmployeeId;

import java.util.Optional;

public interface EmployeeManagementInputPort {
    Optional<Employee> getEmployee(EmployeeId id);
    Optional<Employee> getEmployeeByEmail(String email);
    Optional<EmployeeCredentials> getEmployeeCredentialsByEmail(String email);
    java.util.List<Employee> listEmployees();
    void upsertLimit(EmployeeId employeeId, String resourceType, int amount);
}
\end{DeskFlowCode}

\phantomsection
\addcontentsline{toc}{subsubsection}{EmployeeRepositoryPort}
\DeskFlowListingCaption{employees/application/port/out/EmployeeRepositoryPort.java - EmployeeRepositoryPort}{lst:employees-application-port-out-EmployeeRepositoryPort-java}
\begin{DeskFlowCode}
package com.deskflow.employees.application.port.out;

import com.deskflow.employees.application.model.EmployeeCredentials;
import com.deskflow.employees.application.domain.model.Employee;
import com.deskflow.common.application.domain.model.vo.EmployeeId;

import java.util.List;
import java.util.Optional;

public interface EmployeeRepositoryPort {

    Optional<Employee> find(EmployeeId id);

    Optional<Employee> findByEmail(String email);

    Optional<EmployeeCredentials> findCredentialsByEmail(String email);

    void save(Employee employee);

    List<Employee> findAll();
}
\end{DeskFlowCode}

\phantomsection
\addcontentsline{toc}{subsubsection}{EmployeeAppService}
\DeskFlowListingCaption{employees/application/service/EmployeeAppService.java - EmployeeAppService}{lst:employees-application-service-EmployeeAppService-java}
\begin{DeskFlowCode}
package com.deskflow.employees.application.service;

import com.deskflow.employees.application.model.EmployeeCredentials;
import com.deskflow.employees.application.port.in.EmployeeManagementInputPort;
import com.deskflow.employees.application.port.out.EmployeeRepositoryPort;
import com.deskflow.employees.application.domain.model.Employee;
import com.deskflow.common.application.domain.model.vo.EmployeeId;
import com.deskflow.common.application.domain.model.vo.ResourceType;
import org.springframework.stereotype.Service;
import org.springframework.transaction.annotation.Transactional;

import java.util.List;
import java.util.NoSuchElementException;
import java.util.Optional;

@Service
@Transactional(readOnly = true)
public class EmployeeAppService implements EmployeeManagementInputPort {

    private final EmployeeRepositoryPort employeeRepository;

    public EmployeeAppService(EmployeeRepositoryPort employeeRepository) {
        this.employeeRepository = employeeRepository;
    }

    @Override
    public Optional<Employee> getEmployee(EmployeeId id) {
        return employeeRepository.find(id);
    }

    @Override
    public Optional<Employee> getEmployeeByEmail(String email) {
        return employeeRepository.findByEmail(email);
    }

    @Override
    public Optional<EmployeeCredentials> getEmployeeCredentialsByEmail(String email) {
        return employeeRepository.findCredentialsByEmail(email);
    }

    @Override
    public List<Employee> listEmployees() {
        return employeeRepository.findAll();
    }

    @Override
    @Transactional
    public void upsertLimit(EmployeeId employeeId, String resourceType, int amount) {
        Employee employee = employeeRepository.find(employeeId)
                .orElseThrow(() -> new NoSuchElementException("Employee not found: " + employeeId));
        employee.changeLimit(ResourceType.valueOf(resourceType), amount);
        employeeRepository.save(employee);
    }
}
\end{DeskFlowCode}

\phantomsection
\addcontentsline{toc}{subsubsection}{EmployeeAdminController}
\DeskFlowListingCaption{employees/infrastructure/in/web/EmployeeAdminController.java - EmployeeAdminController}{lst:employees-infrastructure-in-web-EmployeeAdminController-java}
\begin{DeskFlowCode}
package com.deskflow.employees.infrastructure.in.web;

import com.deskflow.common.application.domain.model.vo.EmployeeId;
import com.deskflow.common.application.domain.model.vo.ResourceType;
import com.deskflow.employees.application.domain.model.Employee;
import com.deskflow.employees.application.port.in.EmployeeManagementInputPort;
import jakarta.validation.Valid;
import jakarta.validation.constraints.Min;
import jakarta.validation.constraints.NotBlank;
import org.springframework.http.ResponseEntity;
import org.springframework.web.bind.annotation.GetMapping;
import org.springframework.web.bind.annotation.PathVariable;
import org.springframework.web.bind.annotation.PutMapping;
import org.springframework.web.bind.annotation.RequestBody;
import org.springframework.web.bind.annotation.RequestMapping;
import org.springframework.web.bind.annotation.RestController;

import java.util.Arrays;
import java.util.List;
import java.util.Map;
import java.util.UUID;
import java.util.stream.Collectors;

@RestController
@RequestMapping("/api/administration/employees")
public class EmployeeAdminController {

    private final EmployeeManagementInputPort employeeManagementInputPort;

    public EmployeeAdminController(EmployeeManagementInputPort employeeManagementInputPort) {
        this.employeeManagementInputPort = employeeManagementInputPort;
    }

    @GetMapping
    public ResponseEntity<List<Map<String, Object>>> listEmployees() {
        List<Map<String, Object>> result = employeeManagementInputPort.listEmployees().stream()
                .map(e -> Map.<String, Object>of(
                        "id", e.getId().getValue(),
                        "email", e.getEmail(),
                        "firstName", e.getFirstName(),
                        "lastName", e.getLastName(),
                        "role", e.getRole().name(),
                        "limits", buildLimits(e)
                ))
                .collect(Collectors.toList());
        return ResponseEntity.ok(result);
    }

    @PutMapping("/{id}/limits")
    public ResponseEntity<Void> upsertLimit(@PathVariable UUID id, @Valid @RequestBody LimitRequest body) {
        ResourceType.valueOf(body.resourceType());
        employeeManagementInputPort.upsertLimit(EmployeeId.of(id), body.resourceType(), body.amount());
        return ResponseEntity.ok().build();
    }

    private Map<String, Integer> buildLimits(Employee e) {
        return Arrays.stream(ResourceType.values())
                .collect(Collectors.toMap(
                        Enum::name,
                        t -> e.limitForType(t)
                ));
    }

    public record LimitRequest(
            @NotBlank String resourceType,
            @Min(0) int amount
    ) {}
}
\end{DeskFlowCode}

\phantomsection
\addcontentsline{toc}{subsubsection}{AllocationLimitEntity}
\DeskFlowListingCaption{employees/infrastructure/out/persistence/AllocationLimitEntity.java - AllocationLimitEntity}{lst:employees-infrastructure-out-persistence-AllocationLimitEntity-java}
\begin{DeskFlowCode}
package com.deskflow.employees.infrastructure.out.persistence;

import jakarta.persistence.*;

import java.util.UUID;

@Entity
@Table(name = "allocation_limits")
public class AllocationLimitEntity {

    @Id
    private UUID id;

    @ManyToOne(fetch = FetchType.LAZY)
    @JoinColumn(name = "employee_id", nullable = false)
    private EmployeeEntity employee;

    @Column(name = "resource_type", nullable = false)
    private String resourceType;

    @Column(name = "limit_amount", nullable = false)
    private int limitAmount;

    public UUID getId() { return id; }
    public void setId(UUID id) { this.id = id; }
    public void setEmployee(EmployeeEntity employee) { this.employee = employee; }
    public String getResourceType() { return resourceType; }
    public void setResourceType(String resourceType) { this.resourceType = resourceType; }
    public int getLimitAmount() { return limitAmount; }
    public void setLimitAmount(int limitAmount) { this.limitAmount = limitAmount; }
}
\end{DeskFlowCode}

\phantomsection
\addcontentsline{toc}{subsubsection}{EmployeeEntity}
\DeskFlowListingCaption{employees/infrastructure/out/persistence/EmployeeEntity.java - EmployeeEntity}{lst:employees-infrastructure-out-persistence-EmployeeEntity-java}
\begin{DeskFlowCode}
package com.deskflow.employees.infrastructure.out.persistence;

import jakarta.persistence.*;

import java.util.ArrayList;
import java.util.List;
import java.util.UUID;

@Entity
@Table(name = "employees")
public class EmployeeEntity {

    @Id
    private UUID id;

    @Column(nullable = false, unique = true)
    private String email;

    @Column(name = "first_name", nullable = false)
    private String firstName;

    @Column(name = "last_name", nullable = false)
    private String lastName;

    @Column(name = "role", nullable = false)
    private String role;

    @Column(name = "password_hash", nullable = false)
    private String passwordHash;

    @OneToMany(mappedBy = "employee", cascade = CascadeType.ALL, orphanRemoval = true, fetch = FetchType.EAGER)
    private List<AllocationLimitEntity> allocationLimits = new ArrayList<>();

    public UUID getId() { return id; }
    public void setId(UUID id) { this.id = id; }
    public String getEmail() { return email; }
    public void setEmail(String email) { this.email = email; }
    public String getFirstName() { return firstName; }
    public void setFirstName(String firstName) { this.firstName = firstName; }
    public String getLastName() { return lastName; }
    public void setLastName(String lastName) { this.lastName = lastName; }
    public String getRole() { return role; }
    public void setRole(String role) { this.role = role; }
    public String getPasswordHash() { return passwordHash; }
    public void setPasswordHash(String passwordHash) { this.passwordHash = passwordHash; }
    public List<AllocationLimitEntity> getAllocationLimits() { return allocationLimits; }

    public void addAllocationLimit(AllocationLimitEntity limit) {
        limit.setEmployee(this);
        allocationLimits.add(limit);
    }
}
\end{DeskFlowCode}

\phantomsection
\addcontentsline{toc}{subsubsection}{EmployeeMapper}
\DeskFlowListingCaption{employees/infrastructure/out/persistence/EmployeeMapper.java - EmployeeMapper}{lst:employees-infrastructure-out-persistence-EmployeeMapper-java}
\begin{DeskFlowCode}
package com.deskflow.employees.infrastructure.out.persistence;

import com.deskflow.employees.application.domain.model.Employee;
import com.deskflow.common.application.domain.model.vo.EmployeeId;
import com.deskflow.employees.application.domain.model.vo.Role;
import com.deskflow.common.application.domain.model.vo.ResourceType;
import org.springframework.stereotype.Component;

import java.util.Map;
import java.util.stream.Collectors;

@Component
public class EmployeeMapper {

    public Employee toDomain(EmployeeEntity e) {
        Map<ResourceType, Integer> limits = e.getAllocationLimits().stream()
                .collect(Collectors.toMap(
                        l -> ResourceType.valueOf(l.getResourceType()),
                        AllocationLimitEntity::getLimitAmount
                ));
        return new Employee(
                EmployeeId.of(e.getId()),
                e.getEmail(), e.getFirstName(), e.getLastName(),
                Role.valueOf(e.getRole()),
                limits
        );
    }
}
\end{DeskFlowCode}

\phantomsection
\addcontentsline{toc}{subsubsection}{HibernateEmployeeRepositoryAdapter}
\DeskFlowListingCaption{employees/infrastructure/out/persistence/HibernateEmployeeRepositoryAdapter.java - HibernateEmployeeRepositoryAdapter}{lst:employees-infrastructure-out-persistence-HibernateEmployeeRepositoryAdapter-java}
\begin{DeskFlowCode}
package com.deskflow.employees.infrastructure.out.persistence;

import com.deskflow.employees.application.model.EmployeeCredentials;
import com.deskflow.employees.application.port.out.EmployeeRepositoryPort;
import com.deskflow.employees.application.domain.model.Employee;
import com.deskflow.common.application.domain.model.vo.EmployeeId;
import org.springframework.stereotype.Repository;

import java.util.List;
import java.util.Optional;
import java.util.UUID;

@Repository
public class HibernateEmployeeRepositoryAdapter implements EmployeeRepositoryPort {

    private final SpringDataEmployeeRepository jpa;
    private final SpringDataAllocationLimitRepository limitJpa;
    private final EmployeeMapper mapper;

    public HibernateEmployeeRepositoryAdapter(SpringDataEmployeeRepository jpa,
                                              SpringDataAllocationLimitRepository limitJpa,
                                              EmployeeMapper mapper) {
        this.jpa = jpa;
        this.limitJpa = limitJpa;
        this.mapper = mapper;
    }

    @Override
    public Optional<Employee> find(EmployeeId id) {
        return jpa.findById(id.getValue()).map(mapper::toDomain);
    }

    @Override
    public Optional<Employee> findByEmail(String email) {
        return jpa.findByEmail(email).map(mapper::toDomain);
    }

    @Override
    public Optional<EmployeeCredentials> findCredentialsByEmail(String email) {
        return jpa.findByEmail(email)
                .map(employee -> new EmployeeCredentials(employee.getEmail(), employee.getPasswordHash(), employee.getRole()));
    }

    @Override
    public void save(Employee employee) {
        EmployeeEntity entity = jpa.findById(employee.getId().getValue())
                .orElseThrow(() -> new IllegalArgumentException("Employee not found: " + employee.getId()));
        employee.getLimits().forEach((type, amount) -> upsertLimit(entity, type.name(), amount));
        jpa.save(entity);
    }

    @Override
    public List<Employee> findAll() {
        return jpa.findAll().stream().map(mapper::toDomain).toList();
    }

    private void upsertLimit(EmployeeEntity employee, String resourceType, int amount) {
        AllocationLimitEntity limit = limitJpa
                .findByEmployeeIdAndResourceType(employee.getId(), resourceType)
                .orElseGet(() -> {
                    AllocationLimitEntity l = new AllocationLimitEntity();
                    l.setId(UUID.randomUUID());
                    l.setEmployee(employee);
                    l.setResourceType(resourceType);
                    return l;
                });
        limit.setLimitAmount(amount);
        limitJpa.save(limit);
    }
}
\end{DeskFlowCode}

\phantomsection
\addcontentsline{toc}{subsubsection}{SpringDataAllocationLimitRepository}
\DeskFlowListingCaption{employees/infrastructure/out/persistence/SpringDataAllocationLimitRepository.java - SpringDataAllocationLimitRepository}{lst:employees-infrastructure-out-persistence-SpringDataAllocationLimitRepository-java}
\begin{DeskFlowCode}
package com.deskflow.employees.infrastructure.out.persistence;

import org.springframework.data.jpa.repository.JpaRepository;

import java.util.Optional;
import java.util.UUID;

public interface SpringDataAllocationLimitRepository extends JpaRepository<AllocationLimitEntity, UUID> {
    Optional<AllocationLimitEntity> findByEmployeeIdAndResourceType(UUID employeeId, String resourceType);
}
\end{DeskFlowCode}

\phantomsection
\addcontentsline{toc}{subsubsection}{SpringDataEmployeeRepository}
\DeskFlowListingCaption{employees/infrastructure/out/persistence/SpringDataEmployeeRepository.java - SpringDataEmployeeRepository}{lst:employees-infrastructure-out-persistence-SpringDataEmployeeRepository-java}
\begin{DeskFlowCode}
package com.deskflow.employees.infrastructure.out.persistence;

import org.springframework.data.jpa.repository.JpaRepository;

import java.util.Optional;
import java.util.UUID;

public interface SpringDataEmployeeRepository extends JpaRepository<EmployeeEntity, UUID> {
    Optional<EmployeeEntity> findByEmail(String email);
}
\end{DeskFlowCode}

\subsection{Kontekst Identity - tożsamość i uwierzytelnianie}

\phantomsection
\addcontentsline{toc}{subsubsection}{AuthenticatedIdentity}
\DeskFlowListingCaption{identity/application/model/AuthenticatedIdentity.java - AuthenticatedIdentity}{lst:identity-application-model-AuthenticatedIdentity-java}
\begin{DeskFlowCode}
package com.deskflow.identity.application.model;

public record AuthenticatedIdentity(
        String token,
        String email,
        String role,
        String firstName,
        String lastName,
        String employeeId
) {}
\end{DeskFlowCode}

\phantomsection
\addcontentsline{toc}{subsubsection}{IdentityProfile}
\DeskFlowListingCaption{identity/application/model/IdentityProfile.java - IdentityProfile}{lst:identity-application-model-IdentityProfile-java}
\begin{DeskFlowCode}
package com.deskflow.identity.application.model;

public record IdentityProfile(
        String email,
        String role,
        String firstName,
        String lastName,
        String employeeId
) {}
\end{DeskFlowCode}

\phantomsection
\addcontentsline{toc}{subsubsection}{UserCredentials}
\DeskFlowListingCaption{identity/application/model/UserCredentials.java - UserCredentials}{lst:identity-application-model-UserCredentials-java}
\begin{DeskFlowCode}
package com.deskflow.identity.application.model;

public record UserCredentials(String email, String passwordHash, String role) {}
\end{DeskFlowCode}

\phantomsection
\addcontentsline{toc}{subsubsection}{IdentityAccessInputPort}
\DeskFlowListingCaption{identity/application/port/in/IdentityAccessInputPort.java - IdentityAccessInputPort}{lst:identity-application-port-in-IdentityAccessInputPort-java}
\begin{DeskFlowCode}
package com.deskflow.identity.application.port.in;

import com.deskflow.identity.application.model.AuthenticatedIdentity;
import com.deskflow.identity.application.model.IdentityProfile;

import java.util.Optional;

public interface IdentityAccessInputPort {
    AuthenticatedIdentity login(String email, String password);

    Optional<IdentityProfile> currentIdentity(String email);
}
\end{DeskFlowCode}

\phantomsection
\addcontentsline{toc}{subsubsection}{AccessTokenIssuerPort}
\DeskFlowListingCaption{identity/application/port/out/AccessTokenIssuerPort.java - AccessTokenIssuerPort}{lst:identity-application-port-out-AccessTokenIssuerPort-java}
\begin{DeskFlowCode}
package com.deskflow.identity.application.port.out;

public interface AccessTokenIssuerPort {
    String issueAccessToken(String email, String role);
}
\end{DeskFlowCode}

\phantomsection
\addcontentsline{toc}{subsubsection}{CredentialsAuthenticationPort}
\DeskFlowListingCaption{identity/application/port/out/CredentialsAuthenticationPort.java - CredentialsAuthenticationPort}{lst:identity-application-port-out-CredentialsAuthenticationPort-java}
\begin{DeskFlowCode}
package com.deskflow.identity.application.port.out;

public interface CredentialsAuthenticationPort {
    void authenticate(String email, String password);
}
\end{DeskFlowCode}

\phantomsection
\addcontentsline{toc}{subsubsection}{IdentityEmployeeProfilePort}
\DeskFlowListingCaption{identity/application/port/out/IdentityEmployeeProfilePort.java - IdentityEmployeeProfilePort}{lst:identity-application-port-out-IdentityEmployeeProfilePort-java}
\begin{DeskFlowCode}
package com.deskflow.identity.application.port.out;

import com.deskflow.identity.application.model.IdentityProfile;

import java.util.Optional;

public interface IdentityEmployeeProfilePort {
    Optional<IdentityProfile> findByEmail(String email);
}
\end{DeskFlowCode}

\phantomsection
\addcontentsline{toc}{subsubsection}{UserCredentialsPort}
\DeskFlowListingCaption{identity/application/port/out/UserCredentialsPort.java - UserCredentialsPort}{lst:identity-application-port-out-UserCredentialsPort-java}
\begin{DeskFlowCode}
package com.deskflow.identity.application.port.out;

import com.deskflow.identity.application.model.UserCredentials;

import java.util.Optional;

public interface UserCredentialsPort {
    Optional<UserCredentials> findByEmail(String email);
}
\end{DeskFlowCode}

\phantomsection
\addcontentsline{toc}{subsubsection}{IdentityAccessAppService}
\DeskFlowListingCaption{identity/application/service/IdentityAccessAppService.java - IdentityAccessAppService}{lst:identity-application-service-IdentityAccessAppService-java}
\begin{DeskFlowCode}
package com.deskflow.identity.application.service;

import com.deskflow.identity.application.model.AuthenticatedIdentity;
import com.deskflow.identity.application.model.IdentityProfile;
import com.deskflow.identity.application.port.in.IdentityAccessInputPort;
import com.deskflow.identity.application.port.out.AccessTokenIssuerPort;
import com.deskflow.identity.application.port.out.CredentialsAuthenticationPort;
import com.deskflow.identity.application.port.out.IdentityEmployeeProfilePort;
import org.springframework.stereotype.Service;
import org.springframework.transaction.annotation.Transactional;

import java.util.Optional;

@Service
@Transactional(readOnly = true)
public class IdentityAccessAppService implements IdentityAccessInputPort {

    private final CredentialsAuthenticationPort credentialsAuthenticationPort;
    private final AccessTokenIssuerPort accessTokenIssuerPort;
    private final IdentityEmployeeProfilePort employeeProfilePort;

    public IdentityAccessAppService(CredentialsAuthenticationPort credentialsAuthenticationPort,
                                    AccessTokenIssuerPort accessTokenIssuerPort,
                                    IdentityEmployeeProfilePort employeeProfilePort) {
        this.credentialsAuthenticationPort = credentialsAuthenticationPort;
        this.accessTokenIssuerPort = accessTokenIssuerPort;
        this.employeeProfilePort = employeeProfilePort;
    }

    @Override
    public AuthenticatedIdentity login(String email, String password) {
        credentialsAuthenticationPort.authenticate(email, password);
        IdentityProfile profile = employeeProfilePort.findByEmail(email).orElseThrow();
        String token = accessTokenIssuerPort.issueAccessToken(email, profile.role());
        return new AuthenticatedIdentity(
                token, profile.email(), profile.role(),
                profile.firstName(), profile.lastName(), profile.employeeId()
        );
    }

    @Override
    public Optional<IdentityProfile> currentIdentity(String email) {
        return employeeProfilePort.findByEmail(email);
    }
}
\end{DeskFlowCode}

\phantomsection
\addcontentsline{toc}{subsubsection}{AuthController}
\DeskFlowListingCaption{identity/infrastructure/in/web/AuthController.java - AuthController}{lst:identity-infrastructure-in-web-AuthController-java}
\begin{DeskFlowCode}
package com.deskflow.identity.infrastructure.in.web;

import com.deskflow.identity.application.model.AuthenticatedIdentity;
import com.deskflow.identity.application.model.IdentityProfile;
import com.deskflow.identity.application.port.in.IdentityAccessInputPort;
import com.deskflow.identity.infrastructure.in.web.dto.LoginRequest;
import com.deskflow.identity.infrastructure.in.web.dto.LoginResponse;
import jakarta.validation.Valid;
import org.springframework.http.ResponseEntity;
import org.springframework.security.core.Authentication;
import org.springframework.web.bind.annotation.GetMapping;
import org.springframework.web.bind.annotation.PostMapping;
import org.springframework.web.bind.annotation.RequestBody;
import org.springframework.web.bind.annotation.RequestMapping;
import org.springframework.web.bind.annotation.RestController;

@RestController
@RequestMapping("/api/auth")
public class AuthController {

    private final IdentityAccessInputPort identityAccessInputPort;

    public AuthController(IdentityAccessInputPort identityAccessInputPort) {
        this.identityAccessInputPort = identityAccessInputPort;
    }

    @PostMapping("/login")
    public ResponseEntity<LoginResponse> login(@Valid @RequestBody LoginRequest req) {
        return ResponseEntity.ok(toResponse(identityAccessInputPort.login(req.email(), req.password())));
    }

    @GetMapping("/me")
    public ResponseEntity<?> me(Authentication auth) {
        if (auth == null) return ResponseEntity.status(401).build();
        return identityAccessInputPort.currentIdentity(auth.getName())
                .<ResponseEntity<?>>map(profile -> ResponseEntity.ok(toResponse(profile)))
                .orElse(ResponseEntity.notFound().build());
    }

    private LoginResponse toResponse(AuthenticatedIdentity identity) {
        return new LoginResponse(
                identity.token(), identity.email(), identity.role(),
                identity.firstName(), identity.lastName(), identity.employeeId()
        );
    }

    private LoginResponse toResponse(IdentityProfile profile) {
        return new LoginResponse(
                null, profile.email(), profile.role(),
                profile.firstName(), profile.lastName(), profile.employeeId()
        );
    }
}
\end{DeskFlowCode}

\phantomsection
\addcontentsline{toc}{subsubsection}{LoginRequest}
\DeskFlowListingCaption{identity/infrastructure/in/web/dto/LoginRequest.java - LoginRequest}{lst:identity-infrastructure-in-web-dto-LoginRequest-java}
\begin{DeskFlowCode}
package com.deskflow.identity.infrastructure.in.web.dto;

import jakarta.validation.constraints.Email;
import jakarta.validation.constraints.NotBlank;

public record LoginRequest(@Email @NotBlank String email, @NotBlank String password) {}
\end{DeskFlowCode}

\phantomsection
\addcontentsline{toc}{subsubsection}{LoginResponse}
\DeskFlowListingCaption{identity/infrastructure/in/web/dto/LoginResponse.java - LoginResponse}{lst:identity-infrastructure-in-web-dto-LoginResponse-java}
\begin{DeskFlowCode}
package com.deskflow.identity.infrastructure.in.web.dto;

public record LoginResponse(String token, String email, String role, String firstName, String lastName, String employeeId) {}
\end{DeskFlowCode}

\phantomsection
\addcontentsline{toc}{subsubsection}{EmployeeIdentityProfileAdapter}
\DeskFlowListingCaption{identity/infrastructure/out/employees/EmployeeIdentityProfileAdapter.java - EmployeeIdentityProfileAdapter}{lst:identity-infrastructure-out-employees-EmployeeIdentityProfileAdapter-java}
\begin{DeskFlowCode}
package com.deskflow.identity.infrastructure.out.employees;

import com.deskflow.employees.application.port.in.EmployeeManagementInputPort;
import com.deskflow.identity.application.model.IdentityProfile;
import com.deskflow.identity.application.port.out.IdentityEmployeeProfilePort;
import org.springframework.stereotype.Component;

import java.util.Optional;

@Component
public class EmployeeIdentityProfileAdapter implements IdentityEmployeeProfilePort {

    private final EmployeeManagementInputPort employeeManagementInputPort;

    public EmployeeIdentityProfileAdapter(EmployeeManagementInputPort employeeManagementInputPort) {
        this.employeeManagementInputPort = employeeManagementInputPort;
    }

    @Override
    public Optional<IdentityProfile> findByEmail(String email) {
        return employeeManagementInputPort.getEmployeeByEmail(email)
                .map(employee -> new IdentityProfile(
                        employee.getEmail(),
                        employee.getRole().name(),
                        employee.getFirstName(),
                        employee.getLastName(),
                        employee.getId().toString()
                ));
    }
}
\end{DeskFlowCode}

\phantomsection
\addcontentsline{toc}{subsubsection}{EmployeeUserCredentialsAdapter}
\DeskFlowListingCaption{identity/infrastructure/out/employees/EmployeeUserCredentialsAdapter.java - EmployeeUserCredentialsAdapter}{lst:identity-infrastructure-out-employees-EmployeeUserCredentialsAdapter-java}
\begin{DeskFlowCode}
package com.deskflow.identity.infrastructure.out.employees;

import com.deskflow.employees.application.port.in.EmployeeManagementInputPort;
import com.deskflow.identity.application.model.UserCredentials;
import com.deskflow.identity.application.port.out.UserCredentialsPort;
import org.springframework.stereotype.Component;

import java.util.Optional;

@Component
public class EmployeeUserCredentialsAdapter implements UserCredentialsPort {

    private final EmployeeManagementInputPort employeeManagementInputPort;

    public EmployeeUserCredentialsAdapter(EmployeeManagementInputPort employeeManagementInputPort) {
        this.employeeManagementInputPort = employeeManagementInputPort;
    }

    @Override
    public Optional<UserCredentials> findByEmail(String email) {
        return employeeManagementInputPort.getEmployeeCredentialsByEmail(email)
                .map(credentials -> new UserCredentials(credentials.email(), credentials.passwordHash(), credentials.role()));
    }
}
\end{DeskFlowCode}

\phantomsection
\addcontentsline{toc}{subsubsection}{SpringCredentialsAuthenticationAdapter}
\DeskFlowListingCaption{identity/infrastructure/out/security/SpringCredentialsAuthenticationAdapter.java - SpringCredentialsAuthenticationAdapter}{lst:identity-infrastructure-out-security-SpringCredentialsAuthenticationAdapter-java}
\begin{DeskFlowCode}
package com.deskflow.identity.infrastructure.out.security;

import com.deskflow.identity.application.port.out.CredentialsAuthenticationPort;
import org.springframework.security.authentication.AuthenticationManager;
import org.springframework.security.authentication.UsernamePasswordAuthenticationToken;
import org.springframework.stereotype.Component;

@Component
public class SpringCredentialsAuthenticationAdapter implements CredentialsAuthenticationPort {

    private final AuthenticationManager authenticationManager;

    public SpringCredentialsAuthenticationAdapter(AuthenticationManager authenticationManager) {
        this.authenticationManager = authenticationManager;
    }

    @Override
    public void authenticate(String email, String password) {
        authenticationManager.authenticate(new UsernamePasswordAuthenticationToken(email, password));
    }
}
\end{DeskFlowCode}

\phantomsection
\addcontentsline{toc}{subsubsection}{DeskflowUserDetailsService}
\DeskFlowListingCaption{identity/infrastructure/security/DeskflowUserDetailsService.java - DeskflowUserDetailsService}{lst:identity-infrastructure-security-DeskflowUserDetailsService-java}
\begin{DeskFlowCode}
package com.deskflow.identity.infrastructure.security;

import com.deskflow.identity.application.port.out.UserCredentialsPort;
import org.springframework.security.core.authority.SimpleGrantedAuthority;
import org.springframework.security.core.userdetails.User;
import org.springframework.security.core.userdetails.UserDetails;
import org.springframework.security.core.userdetails.UserDetailsService;
import org.springframework.security.core.userdetails.UsernameNotFoundException;
import org.springframework.stereotype.Service;

import java.util.List;

@Service
public class DeskflowUserDetailsService implements UserDetailsService {

    private final UserCredentialsPort userCredentialsPort;

    public DeskflowUserDetailsService(UserCredentialsPort userCredentialsPort) {
        this.userCredentialsPort = userCredentialsPort;
    }

    @Override
    public UserDetails loadUserByUsername(String email) throws UsernameNotFoundException {
        return userCredentialsPort.findByEmail(email)
                .map(credentials -> new User(
                        credentials.email(), credentials.passwordHash(),
                        List.of(new SimpleGrantedAuthority("ROLE_" + credentials.role()))
                ))
                .orElseThrow(() -> new UsernameNotFoundException("User not found: " + email));
    }
}
\end{DeskFlowCode}

\phantomsection
\addcontentsline{toc}{subsubsection}{JwtAuthenticationFilter}
\DeskFlowListingCaption{identity/infrastructure/security/JwtAuthenticationFilter.java - JwtAuthenticationFilter}{lst:identity-infrastructure-security-JwtAuthenticationFilter-java}
\begin{DeskFlowCode}
package com.deskflow.identity.infrastructure.security;

import jakarta.servlet.FilterChain;
import jakarta.servlet.ServletException;
import jakarta.servlet.http.HttpServletRequest;
import jakarta.servlet.http.HttpServletResponse;
import org.springframework.security.authentication.UsernamePasswordAuthenticationToken;
import org.springframework.security.core.authority.SimpleGrantedAuthority;
import org.springframework.security.core.context.SecurityContextHolder;
import org.springframework.stereotype.Component;
import org.springframework.web.filter.OncePerRequestFilter;

import java.io.IOException;
import java.util.List;

@Component
public class JwtAuthenticationFilter extends OncePerRequestFilter {

    private final JwtTokenProvider tokenProvider;

    public JwtAuthenticationFilter(JwtTokenProvider tokenProvider) {
        this.tokenProvider = tokenProvider;
    }

    @Override
    protected void doFilterInternal(HttpServletRequest request, HttpServletResponse response,
                                    FilterChain chain) throws ServletException, IOException {
        String token = resolveToken(request);
        if (token != null && tokenProvider.validateToken(token)) {
            String email = tokenProvider.getEmailFromToken(token);
            String role = tokenProvider.getRoleFromToken(token);
            var auth = new UsernamePasswordAuthenticationToken(
                    email, null, List.of(new SimpleGrantedAuthority("ROLE_" + role))
            );
            var context = SecurityContextHolder.createEmptyContext();
            context.setAuthentication(auth);
            SecurityContextHolder.setContext(context);
        }
        chain.doFilter(request, response);
    }

    private String resolveToken(HttpServletRequest request) {
        String header = request.getHeader("Authorization");
        if (header != null && header.startsWith("Bearer ")) {
            return header.substring(7);
        }
        return null;
    }
}
\end{DeskFlowCode}

\phantomsection
\addcontentsline{toc}{subsubsection}{JwtTokenProvider}
\DeskFlowListingCaption{identity/infrastructure/security/JwtTokenProvider.java - JwtTokenProvider}{lst:identity-infrastructure-security-JwtTokenProvider-java}
\begin{DeskFlowCode}
package com.deskflow.identity.infrastructure.security;

import com.deskflow.identity.application.port.out.AccessTokenIssuerPort;
import io.jsonwebtoken.*;
import io.jsonwebtoken.security.Keys;
import org.springframework.beans.factory.annotation.Value;
import org.springframework.stereotype.Component;

import javax.crypto.SecretKey;
import java.security.MessageDigest;
import java.security.NoSuchAlgorithmException;
import java.nio.charset.StandardCharsets;
import java.util.Date;

@Component
public class JwtTokenProvider implements AccessTokenIssuerPort {

    private final SecretKey secretKey;
    private final long expirationMs;

    public JwtTokenProvider(
            @Value("${app.jwt.secret}") String secret,
            @Value("${app.jwt.expiration-ms:86400000}") long expirationMs) {
        this.secretKey = Keys.hmacShaKeyFor(sha256(secret));
        this.expirationMs = expirationMs;
    }

    private static byte[] sha256(String secret) {
        String value = (secret == null || secret.isBlank())
                ? "deskflow-secret-key-at-least-32-characters-long-for-hs256"
                : secret;
        try {
            return MessageDigest.getInstance("SHA-256").digest(value.getBytes(StandardCharsets.UTF_8));
        } catch (NoSuchAlgorithmException e) {
            throw new IllegalStateException("SHA-256 is not available", e);
        }
    }

    @Override
    public String issueAccessToken(String email, String role) {
        return Jwts.builder()
                .subject(email)
                .claim("role", role)
                .issuedAt(new Date())
                .expiration(new Date(System.currentTimeMillis() + expirationMs))
                .signWith(secretKey)
                .compact();
    }

    public String getEmailFromToken(String token) {
        return parseClaims(token).getSubject();
    }

    public String getRoleFromToken(String token) {
        return parseClaims(token).get("role", String.class);
    }

    public boolean validateToken(String token) {
        try {
            parseClaims(token);
            return true;
        } catch (JwtException | IllegalArgumentException e) {
            return false;
        }
    }

    private Claims parseClaims(String token) {
        return Jwts.parser()
                .verifyWith(secretKey)
                .build()
                .parseSignedClaims(token)
                .getPayload();
    }
}
\end{DeskFlowCode}

\phantomsection
\addcontentsline{toc}{subsubsection}{SecurityConfig}
\DeskFlowListingCaption{identity/infrastructure/security/SecurityConfig.java - SecurityConfig}{lst:identity-infrastructure-security-SecurityConfig-java}
\begin{DeskFlowCode}
package com.deskflow.identity.infrastructure.security;

import org.springframework.beans.factory.annotation.Value;
import org.springframework.context.annotation.Bean;
import org.springframework.context.annotation.Configuration;
import org.springframework.http.HttpMethod;
import org.springframework.security.authentication.AuthenticationManager;
import org.springframework.security.config.annotation.authentication.configuration.AuthenticationConfiguration;
import org.springframework.security.config.annotation.method.configuration.EnableMethodSecurity;
import org.springframework.security.config.annotation.web.builders.HttpSecurity;
import org.springframework.security.config.annotation.web.configuration.EnableWebSecurity;
import org.springframework.security.config.annotation.web.configurers.AbstractHttpConfigurer;
import org.springframework.security.config.http.SessionCreationPolicy;
import org.springframework.security.crypto.bcrypt.BCryptPasswordEncoder;
import org.springframework.security.crypto.password.PasswordEncoder;
import org.springframework.security.web.SecurityFilterChain;
import org.springframework.security.web.authentication.UsernamePasswordAuthenticationFilter;
import org.springframework.web.cors.CorsConfiguration;
import org.springframework.web.cors.CorsConfigurationSource;
import org.springframework.web.cors.UrlBasedCorsConfigurationSource;

import java.util.List;

@Configuration
@EnableWebSecurity
@EnableMethodSecurity
public class SecurityConfig {

    private final JwtAuthenticationFilter jwtFilter;
    private final String allowedOrigin;

    public SecurityConfig(JwtAuthenticationFilter jwtFilter,
                          @Value("${app.cors.allowed-origin:http://localhost:3000}") String allowedOrigin) {
        this.jwtFilter = jwtFilter;
        this.allowedOrigin = allowedOrigin;
    }

    @Bean
    public SecurityFilterChain filterChain(HttpSecurity http) throws Exception {
        http
            .cors(cors -> cors.configurationSource(corsConfigurationSource()))
            .csrf(AbstractHttpConfigurer::disable)
            .sessionManagement(s -> s.sessionCreationPolicy(SessionCreationPolicy.STATELESS))
            .authorizeHttpRequests(auth -> auth
                .requestMatchers("/api/auth/**").permitAll()
                .requestMatchers("/actuator/health").permitAll()
                .requestMatchers(HttpMethod.GET, "/api/resources").authenticated()
                .requestMatchers("/api/administration/**").hasAnyRole("MANAGER", "IT_ADMIN")
                .requestMatchers(HttpMethod.POST, "/api/resources").hasAnyRole("IT_ADMIN", "MANAGER")
                .requestMatchers(HttpMethod.PUT, "/api/resources/**").hasAnyRole("IT_ADMIN", "MANAGER")
                .requestMatchers(HttpMethod.DELETE, "/api/resources/**").hasAnyRole("IT_ADMIN")
                .anyRequest().authenticated()
            )
            .addFilterBefore(jwtFilter, UsernamePasswordAuthenticationFilter.class);
        return http.build();
    }

    @Bean
    public PasswordEncoder passwordEncoder() {
        return new BCryptPasswordEncoder();
    }

    @Bean
    public AuthenticationManager authenticationManager(AuthenticationConfiguration config) throws Exception {
        return config.getAuthenticationManager();
    }

    @Bean
    public CorsConfigurationSource corsConfigurationSource() {
        CorsConfiguration config = new CorsConfiguration();
        config.setAllowedOrigins(List.of(allowedOrigin));
        config.setAllowedMethods(List.of("GET", "POST", "PUT", "DELETE", "OPTIONS"));
        config.setAllowedHeaders(List.of("*"));
        config.setAllowCredentials(true);
        UrlBasedCorsConfigurationSource source = new UrlBasedCorsConfigurationSource();
        source.registerCorsConfiguration("/**", config);
        return source;
    }
}
\end{DeskFlowCode}

\subsection{Kontekst Notifications - powiadomienia}

\phantomsection
\addcontentsline{toc}{subsubsection}{NotificationEventInputPort}
\DeskFlowListingCaption{notifications/application/port/in/NotificationEventInputPort.java - NotificationEventInputPort}{lst:notifications-application-port-in-NotificationEventInputPort-java}
\begin{DeskFlowCode}
package com.deskflow.notifications.application.port.in;

import com.deskflow.allocation.application.domain.model.vo.AllocationRequestId;
import com.deskflow.common.application.domain.model.vo.EmployeeId;

public interface NotificationEventInputPort {
    void allocationRequestSubmitted(AllocationRequestId requestId, EmployeeId employeeId);

    void allocationRequestApproved(AllocationRequestId requestId, EmployeeId employeeId);

    void allocationRequestRejected(AllocationRequestId requestId, EmployeeId employeeId, String reason);

    void allocationRequestCancelled(AllocationRequestId requestId, EmployeeId employeeId);

    void resourceIssued(AllocationRequestId requestId, EmployeeId employeeId);

    void resourceReturned(AllocationRequestId requestId, EmployeeId employeeId);

    void allocationOverdue(AllocationRequestId requestId, EmployeeId employeeId);
}
\end{DeskFlowCode}

\phantomsection
\addcontentsline{toc}{subsubsection}{NotificationPort}
\DeskFlowListingCaption{notifications/application/port/out/NotificationPort.java - NotificationPort}{lst:notifications-application-port-out-NotificationPort-java}
\begin{DeskFlowCode}
package com.deskflow.notifications.application.port.out;

public interface NotificationPort {
    void send(String toEmail, String subject, String body);
}
\end{DeskFlowCode}

\phantomsection
\addcontentsline{toc}{subsubsection}{NotificationAppService}
\DeskFlowListingCaption{notifications/application/service/NotificationAppService.java - NotificationAppService}{lst:notifications-application-service-NotificationAppService-java}
\begin{DeskFlowCode}
package com.deskflow.notifications.application.service;

import com.deskflow.allocation.application.domain.model.vo.AllocationRequestId;
import com.deskflow.employees.application.port.in.EmployeeManagementInputPort;
import com.deskflow.common.application.domain.model.vo.EmployeeId;
import com.deskflow.notifications.application.port.in.NotificationEventInputPort;
import com.deskflow.notifications.application.port.out.NotificationPort;
import org.slf4j.Logger;
import org.slf4j.LoggerFactory;
import org.springframework.stereotype.Service;
import org.springframework.transaction.annotation.Transactional;

@Service
@Transactional(readOnly = true)
public class NotificationAppService implements NotificationEventInputPort {

    private static final Logger log = LoggerFactory.getLogger(NotificationAppService.class);

    private final NotificationPort notifications;
    private final EmployeeManagementInputPort employeeManagementInputPort;

    public NotificationAppService(NotificationPort notifications,
                                  EmployeeManagementInputPort employeeManagementInputPort) {
        this.notifications = notifications;
        this.employeeManagementInputPort = employeeManagementInputPort;
    }

    @Override
    public void allocationRequestSubmitted(AllocationRequestId requestId, EmployeeId employeeId) {
        emailEmployee(employeeId,
                "Request submitted",
                "Your allocation request has been submitted (ID: " + requestId + ").");
    }

    @Override
    public void allocationRequestApproved(AllocationRequestId requestId, EmployeeId employeeId) {
        emailEmployee(employeeId,
                "Request approved",
                "Your request has been approved (ID: " + requestId + "). Please collect the resource.");
    }

    @Override
    public void allocationRequestRejected(AllocationRequestId requestId, EmployeeId employeeId, String reason) {
        emailEmployee(employeeId,
                "Request rejected",
                "Your request (ID: " + requestId + ") was rejected. Reason: " + reason);
    }

    @Override
    public void allocationRequestCancelled(AllocationRequestId requestId, EmployeeId employeeId) {
        emailEmployee(employeeId,
                "Request cancelled",
                "Your request (ID: " + requestId + ") has been cancelled.");
    }

    @Override
    public void resourceIssued(AllocationRequestId requestId, EmployeeId employeeId) {
        emailEmployee(employeeId,
                "Resource issued",
                "The resource from request (ID: " + requestId + ") has been issued to you.");
    }

    @Override
    public void resourceReturned(AllocationRequestId requestId, EmployeeId employeeId) {
        emailEmployee(employeeId,
                "Resource returned",
                "We confirm receipt of the resource from request (ID: " + requestId + ").");
    }

    @Override
    public void allocationOverdue(AllocationRequestId requestId, EmployeeId employeeId) {
        emailEmployee(employeeId,
                "Allocation overdue",
                "Your allocation (ID: " + requestId + ") is past its return date. Please return the resource immediately.");
    }

    private void emailEmployee(EmployeeId employeeId, String subject, String body) {
        try {
            employeeManagementInputPort.getEmployee(employeeId)
                    .ifPresent(employee -> notifications.send(employee.getEmail(), subject, body));
        } catch (Exception e) {
            log.warn("Failed to send notification to employee {}: {}", employeeId, e.getMessage(), e);
        }
    }
}
\end{DeskFlowCode}

\phantomsection
\addcontentsline{toc}{subsubsection}{NotificationEventListener}
\DeskFlowListingCaption{notifications/infrastructure/in/events/NotificationEventListener.java - NotificationEventListener}{lst:notifications-infrastructure-in-events-NotificationEventListener-java}
\begin{DeskFlowCode}
package com.deskflow.notifications.infrastructure.in.events;

import com.deskflow.allocation.application.domain.model.events.AllocationOverdue;
import com.deskflow.allocation.application.domain.model.events.AllocationRequestApproved;
import com.deskflow.allocation.application.domain.model.events.AllocationRequestCancelled;
import com.deskflow.allocation.application.domain.model.events.AllocationRequestRejected;
import com.deskflow.allocation.application.domain.model.events.AllocationRequestSubmitted;
import com.deskflow.allocation.application.domain.model.events.ResourceIssued;
import com.deskflow.allocation.application.domain.model.events.ResourceReturned;
import com.deskflow.notifications.application.port.in.NotificationEventInputPort;
import org.springframework.stereotype.Component;
import org.springframework.transaction.event.TransactionPhase;
import org.springframework.transaction.event.TransactionalEventListener;

@Component
public class NotificationEventListener {

    private final NotificationEventInputPort notificationEventInputPort;

    public NotificationEventListener(NotificationEventInputPort notificationEventInputPort) {
        this.notificationEventInputPort = notificationEventInputPort;
    }

    @TransactionalEventListener(phase = TransactionPhase.AFTER_COMMIT)
    public void onRequestSubmitted(AllocationRequestSubmitted event) {
        notificationEventInputPort.allocationRequestSubmitted(event.requestId(), event.employeeId());
    }

    @TransactionalEventListener(phase = TransactionPhase.AFTER_COMMIT)
    public void onRequestApproved(AllocationRequestApproved event) {
        notificationEventInputPort.allocationRequestApproved(event.requestId(), event.employeeId());
    }

    @TransactionalEventListener(phase = TransactionPhase.AFTER_COMMIT)
    public void onRequestRejected(AllocationRequestRejected event) {
        notificationEventInputPort.allocationRequestRejected(event.requestId(), event.employeeId(), event.reason());
    }

    @TransactionalEventListener(phase = TransactionPhase.AFTER_COMMIT)
    public void onRequestCancelled(AllocationRequestCancelled event) {
        notificationEventInputPort.allocationRequestCancelled(event.requestId(), event.employeeId());
    }

    @TransactionalEventListener(phase = TransactionPhase.AFTER_COMMIT)
    public void onResourceIssued(ResourceIssued event) {
        notificationEventInputPort.resourceIssued(event.requestId(), event.employeeId());
    }

    @TransactionalEventListener(phase = TransactionPhase.AFTER_COMMIT)
    public void onResourceReturned(ResourceReturned event) {
        notificationEventInputPort.resourceReturned(event.requestId(), event.employeeId());
    }

    @TransactionalEventListener(phase = TransactionPhase.AFTER_COMMIT)
    public void onAllocationOverdue(AllocationOverdue event) {
        notificationEventInputPort.allocationOverdue(event.requestId(), event.employeeId());
    }
}
\end{DeskFlowCode}

\phantomsection
\addcontentsline{toc}{subsubsection}{SmtpNotificationAdapter}
\DeskFlowListingCaption{notifications/infrastructure/out/smtp/SmtpNotificationAdapter.java - SmtpNotificationAdapter}{lst:notifications-infrastructure-out-smtp-SmtpNotificationAdapter-java}
\begin{DeskFlowCode}
package com.deskflow.notifications.infrastructure.out.smtp;

import com.deskflow.notifications.application.port.out.NotificationPort;
import org.slf4j.Logger;
import org.slf4j.LoggerFactory;
import org.springframework.mail.MailException;
import org.springframework.mail.SimpleMailMessage;
import org.springframework.mail.javamail.JavaMailSender;
import org.springframework.stereotype.Component;

@Component
public class SmtpNotificationAdapter implements NotificationPort {

    private static final Logger log = LoggerFactory.getLogger(SmtpNotificationAdapter.class);

    private final JavaMailSender mailSender;

    public SmtpNotificationAdapter(JavaMailSender mailSender) {
        this.mailSender = mailSender;
    }

    @Override
    public void send(String toEmail, String subject, String body) {
        try {
            SimpleMailMessage message = new SimpleMailMessage();
            message.setFrom("noreply@deskflow.local");
            message.setTo(toEmail);
            message.setSubject(subject);
            message.setText(body);
            mailSender.send(message);
        } catch (MailException e) {
            log.warn("Failed to send email to {}: {}", toEmail, e.getMessage());
        }
    }
}
\end{DeskFlowCode}

\subsection{Kontekst Common - kod współdzielony}

\phantomsection
\addcontentsline{toc}{subsubsection}{EmployeeId}
\DeskFlowListingCaption{common/application/domain/model/vo/EmployeeId.java - EmployeeId}{lst:common-application-domain-model-vo-EmployeeId-java}
\begin{DeskFlowCode}
package com.deskflow.common.application.domain.model.vo;

import java.util.Objects;
import java.util.UUID;

public final class EmployeeId {
    private final UUID value;

    public EmployeeId(UUID value) {
        this.value = Objects.requireNonNull(value);
    }

    public static EmployeeId of(UUID value) {
        return new EmployeeId(value);
    }

    public static EmployeeId of(String value) {
        return new EmployeeId(UUID.fromString(value));
    }

    public UUID getValue() {
        return value;
    }

    @Override
    public boolean equals(Object o) {
        if (this == o) return true;
        if (!(o instanceof EmployeeId that)) return false;
        return value.equals(that.value);
    }

    @Override
    public int hashCode() {
        return value.hashCode();
    }

    @Override
    public String toString() {
        return value.toString();
    }
}
\end{DeskFlowCode}

\phantomsection
\addcontentsline{toc}{subsubsection}{ResourceId}
\DeskFlowListingCaption{common/application/domain/model/vo/ResourceId.java - ResourceId}{lst:common-application-domain-model-vo-ResourceId-java}
\begin{DeskFlowCode}
package com.deskflow.common.application.domain.model.vo;

import java.util.Objects;
import java.util.UUID;

public final class ResourceId {
    private final UUID value;

    public ResourceId(UUID value) {
        this.value = Objects.requireNonNull(value);
    }

    public static ResourceId of(UUID value) {
        return new ResourceId(value);
    }

    public static ResourceId of(String value) {
        return new ResourceId(UUID.fromString(value));
    }

    public UUID getValue() {
        return value;
    }

    @Override
    public boolean equals(Object o) {
        if (this == o) return true;
        if (!(o instanceof ResourceId that)) return false;
        return value.equals(that.value);
    }

    @Override
    public int hashCode() {
        return value.hashCode();
    }

    @Override
    public String toString() {
        return value.toString();
    }
}
\end{DeskFlowCode}

\phantomsection
\addcontentsline{toc}{subsubsection}{ResourceType}
\DeskFlowListingCaption{common/application/domain/model/vo/ResourceType.java - ResourceType}{lst:common-application-domain-model-vo-ResourceType-java}
\begin{DeskFlowCode}
package com.deskflow.common.application.domain.model.vo;

public enum ResourceType {
    DESK,
    CONFERENCE_ROOM,
    LAPTOP,
    MONITOR,
    PROJECTOR,
    VEHICLE
}
\end{DeskFlowCode}

\phantomsection
\addcontentsline{toc}{subsubsection}{DomainEventPublisherPort}
\DeskFlowListingCaption{common/application/port/out/DomainEventPublisherPort.java - DomainEventPublisherPort}{lst:common-application-port-out-DomainEventPublisherPort-java}
\begin{DeskFlowCode}
package com.deskflow.common.application.port.out;

import com.deskflow.common.application.domain.DomainEvent;

public interface DomainEventPublisherPort {
    void publish(DomainEvent event);
}
\end{DeskFlowCode}

\phantomsection
\addcontentsline{toc}{subsubsection}{DataInitializer}
\DeskFlowListingCaption{common/infrastructure/config/DataInitializer.java - DataInitializer}{lst:common-infrastructure-config-DataInitializer-java}
\begin{DeskFlowCode}
package com.deskflow.common.infrastructure.config;

import com.deskflow.employees.infrastructure.out.persistence.AllocationLimitEntity;
import com.deskflow.employees.infrastructure.out.persistence.EmployeeEntity;
import com.deskflow.employees.infrastructure.out.persistence.SpringDataEmployeeRepository;
import com.deskflow.resources.infrastructure.out.persistence.OfficeResourceEntity;
import com.deskflow.resources.infrastructure.out.persistence.SpringDataOfficeResourceRepository;
import org.springframework.boot.CommandLineRunner;
import org.springframework.context.annotation.Profile;
import org.springframework.security.crypto.password.PasswordEncoder;
import org.springframework.stereotype.Component;

import java.util.UUID;

@Component
@Profile("!test")
public class DataInitializer implements CommandLineRunner {

    private final SpringDataEmployeeRepository employeeRepo;
    private final SpringDataOfficeResourceRepository resourceRepo;
    private final PasswordEncoder passwordEncoder;

    public DataInitializer(SpringDataEmployeeRepository employeeRepo,
                           SpringDataOfficeResourceRepository resourceRepo,
                           PasswordEncoder passwordEncoder) {
        this.employeeRepo = employeeRepo;
        this.resourceRepo = resourceRepo;
        this.passwordEncoder = passwordEncoder;
    }

    @Override
    public void run(String... args) {
        if (employeeRepo.count() > 0) return;

        // Employees (test data — names and emails stay in Polish per project requirements)
        EmployeeEntity employeeOne = new EmployeeEntity();
        employeeOne.setId(UUID.randomUUID());
        employeeOne.setEmail("anna.kowalska@firma.pl");
        employeeOne.setFirstName("Anna");
        employeeOne.setLastName("Kowalska");
        employeeOne.setRole("EMPLOYEE");
        employeeOne.setPasswordHash(passwordEncoder.encode("password123"));
        addLimit(employeeOne, "DESK", 2);
        addLimit(employeeOne, "LAPTOP", 1);

        EmployeeEntity employeeTwo = new EmployeeEntity();
        employeeTwo.setId(UUID.randomUUID());
        employeeTwo.setEmail("jan.nowak@firma.pl");
        employeeTwo.setFirstName("Jan");
        employeeTwo.setLastName("Nowak");
        employeeTwo.setRole("EMPLOYEE");
        employeeTwo.setPasswordHash(passwordEncoder.encode("password123"));
        addLimit(employeeTwo, "DESK", 1);
        addLimit(employeeTwo, "CONFERENCE_ROOM", 3);

        EmployeeEntity manager = new EmployeeEntity();
        manager.setId(UUID.randomUUID());
        manager.setEmail("marek.kierownik@firma.pl");
        manager.setFirstName("Marek");
        manager.setLastName("Kierownik");
        manager.setRole("MANAGER");
        manager.setPasswordHash(passwordEncoder.encode("password123"));

        EmployeeEntity admin = new EmployeeEntity();
        admin.setId(UUID.randomUUID());
        admin.setEmail("admin@firma.pl");
        admin.setFirstName("Admin");
        admin.setLastName("IT");
        admin.setRole("IT_ADMIN");
        admin.setPasswordHash(passwordEncoder.encode("admin123"));
        addLimit(admin, "DESK", 5);
        addLimit(admin, "LAPTOP", 3);
        addLimit(admin, "CONFERENCE_ROOM", 5);

        employeeRepo.save(employeeOne);
        employeeRepo.save(employeeTwo);
        employeeRepo.save(manager);
        employeeRepo.save(admin);

        // Office resources (test data — names stay in Polish per project requirements)
        createResource("DESK", "Biurko A-101", "A", 1, "101");
        createResource("DESK", "Biurko A-102", "A", 1, "102");
        createResource("DESK", "Biurko B-201", "B", 2, "201");
        createResource("CONFERENCE_ROOM", "Sala Alpha", "A", 0, "S1");
        createResource("CONFERENCE_ROOM", "Sala Beta", "B", 1, "S2");
        createResource("LAPTOP", "Laptop Dell XPS", "A", 0, "IT-01");
        createResource("LAPTOP", "Laptop Lenovo ThinkPad", "A", 0, "IT-02");
        createResource("MONITOR", "Monitor Samsung 27\"", "A", 1, "M-01");
        createResource("PROJECTOR", "Projektor Epson", "A", 0, "P-01");
        createResource("VEHICLE", "Samochód VW Golf", "Parking", 0, "GAR-01");
    }

    private void addLimit(EmployeeEntity employee, String resourceType, int amount) {
        AllocationLimitEntity l = new AllocationLimitEntity();
        l.setId(UUID.randomUUID());
        l.setResourceType(resourceType);
        l.setLimitAmount(amount);
        employee.addAllocationLimit(l);
    }

    private void createResource(String type, String name, String building, int floor, String number) {
        OfficeResourceEntity resource = new OfficeResourceEntity();
        resource.setId(UUID.randomUUID());
        resource.setType(type);
        resource.setName(name);
        resource.setBuilding(building);
        resource.setFloor(floor);
        resource.setNumber(number);
        resource.setActive(true);
        resource.setAvailable(true);
        resourceRepo.save(resource);
    }
}
\end{DeskFlowCode}

\phantomsection
\addcontentsline{toc}{subsubsection}{DomainBeanConfig}
\DeskFlowListingCaption{common/infrastructure/config/DomainBeanConfig.java - DomainBeanConfig}{lst:common-infrastructure-config-DomainBeanConfig-java}
\begin{DeskFlowCode}
package com.deskflow.common.infrastructure.config;

import com.deskflow.allocation.application.domain.service.AllocationRequestFactory;
import com.deskflow.allocation.application.domain.service.ApproveAllocationRequestDomainService;
import com.deskflow.allocation.application.domain.service.CancelAllocationRequestDomainService;
import com.deskflow.allocation.application.domain.service.EnsureEmployeeWithinLimitDomainService;
import com.deskflow.allocation.application.domain.service.EnsureResourceAvailableInPeriodDomainService;
import com.deskflow.allocation.application.domain.service.EnsureResourceEligibleForAllocationDomainService;
import com.deskflow.allocation.application.domain.service.IssueAllocationResourceDomainService;
import com.deskflow.allocation.application.domain.service.MarkAllocationRequestOverdueDomainService;
import com.deskflow.allocation.application.domain.service.RejectAllocationRequestDomainService;
import com.deskflow.allocation.application.domain.service.ReturnAllocationResourceDomainService;
import com.deskflow.resources.application.domain.service.OfficeResourceFactory;
import com.deskflow.resources.application.domain.service.DeactivateOfficeResourceDomainService;
import com.deskflow.resources.application.domain.service.MarkOfficeResourceAvailableDomainService;
import com.deskflow.resources.application.domain.service.MarkOfficeResourceUnavailableDomainService;
import com.deskflow.resources.application.domain.service.UpdateOfficeResourceDomainService;
import org.springframework.context.annotation.Bean;
import org.springframework.context.annotation.Configuration;

@Configuration
public class DomainBeanConfig {

    @Bean
    public AllocationRequestFactory allocationRequestFactory() {
        return new AllocationRequestFactory();
    }

    @Bean
    public OfficeResourceFactory officeResourceFactory() {
        return new OfficeResourceFactory();
    }

    @Bean
    public EnsureResourceEligibleForAllocationDomainService ensureResourceEligibleForAllocationDomainService() {
        return new EnsureResourceEligibleForAllocationDomainService();
    }

    @Bean
    public EnsureResourceAvailableInPeriodDomainService ensureResourceAvailableInPeriodDomainService() {
        return new EnsureResourceAvailableInPeriodDomainService();
    }

    @Bean
    public EnsureEmployeeWithinLimitDomainService ensureEmployeeWithinLimitDomainService() {
        return new EnsureEmployeeWithinLimitDomainService();
    }

    @Bean
    public CancelAllocationRequestDomainService cancelAllocationRequestDomainService() {
        return new CancelAllocationRequestDomainService();
    }

    @Bean
    public ApproveAllocationRequestDomainService approveAllocationRequestDomainService() {
        return new ApproveAllocationRequestDomainService();
    }

    @Bean
    public RejectAllocationRequestDomainService rejectAllocationRequestDomainService() {
        return new RejectAllocationRequestDomainService();
    }

    @Bean
    public IssueAllocationResourceDomainService issueAllocationResourceDomainService() {
        return new IssueAllocationResourceDomainService();
    }

    @Bean
    public ReturnAllocationResourceDomainService returnAllocationResourceDomainService() {
        return new ReturnAllocationResourceDomainService();
    }

    @Bean
    public MarkAllocationRequestOverdueDomainService markAllocationRequestOverdueDomainService() {
        return new MarkAllocationRequestOverdueDomainService();
    }

    @Bean
    public UpdateOfficeResourceDomainService updateOfficeResourceDomainService() {
        return new UpdateOfficeResourceDomainService();
    }

    @Bean
    public MarkOfficeResourceUnavailableDomainService markOfficeResourceUnavailableDomainService() {
        return new MarkOfficeResourceUnavailableDomainService();
    }

    @Bean
    public MarkOfficeResourceAvailableDomainService markOfficeResourceAvailableDomainService() {
        return new MarkOfficeResourceAvailableDomainService();
    }

    @Bean
    public DeactivateOfficeResourceDomainService deactivateOfficeResourceDomainService() {
        return new DeactivateOfficeResourceDomainService();
    }
}
\end{DeskFlowCode}

\phantomsection
\addcontentsline{toc}{subsubsection}{GlobalExceptionHandler}
\DeskFlowListingCaption{common/infrastructure/in/web/GlobalExceptionHandler.java - GlobalExceptionHandler}{lst:common-infrastructure-in-web-GlobalExceptionHandler-java}
\begin{DeskFlowCode}
package com.deskflow.common.infrastructure.in.web;

import org.springframework.http.HttpStatus;
import org.springframework.http.ResponseEntity;
import org.springframework.web.bind.MethodArgumentNotValidException;
import org.springframework.web.bind.annotation.ExceptionHandler;
import org.springframework.web.bind.annotation.RestControllerAdvice;

import java.util.Map;
import java.util.NoSuchElementException;

@RestControllerAdvice
public class GlobalExceptionHandler {

    @ExceptionHandler(IllegalStateException.class)
    public ResponseEntity<?> handleBusinessRule(IllegalStateException e) {
        return ResponseEntity.status(HttpStatus.CONFLICT).body(Map.of("error", e.getMessage()));
    }

    @ExceptionHandler(IllegalArgumentException.class)
    public ResponseEntity<?> handleBadArgument(IllegalArgumentException e) {
        return ResponseEntity.status(HttpStatus.BAD_REQUEST).body(Map.of("error", e.getMessage()));
    }

    @ExceptionHandler(NoSuchElementException.class)
    public ResponseEntity<?> handleNotFound(NoSuchElementException e) {
        return ResponseEntity.status(HttpStatus.NOT_FOUND).body(Map.of("error", e.getMessage()));
    }

    @ExceptionHandler(MethodArgumentNotValidException.class)
    public ResponseEntity<?> handleValidation(MethodArgumentNotValidException e) {
        return ResponseEntity.status(HttpStatus.BAD_REQUEST)
                .body(Map.of("error", "Validation error", "details", e.getBindingResult().toString()));
    }
}
\end{DeskFlowCode}

\phantomsection
\addcontentsline{toc}{subsubsection}{SpringDomainEventPublisherAdapter}
\DeskFlowListingCaption{common/infrastructure/out/events/SpringDomainEventPublisherAdapter.java - SpringDomainEventPublisherAdapter}{lst:common-infrastructure-out-events-SpringDomainEventPublisherAdapter-java}
\begin{DeskFlowCode}
package com.deskflow.common.infrastructure.out.events;

import com.deskflow.common.application.port.out.DomainEventPublisherPort;
import com.deskflow.common.application.domain.DomainEvent;
import org.springframework.context.ApplicationEventPublisher;
import org.springframework.stereotype.Component;

@Component
public class SpringDomainEventPublisherAdapter implements DomainEventPublisherPort {

    private final ApplicationEventPublisher publisher;

    public SpringDomainEventPublisherAdapter(ApplicationEventPublisher publisher) {
        this.publisher = publisher;
    }

    @Override
    public void publish(DomainEvent event) {
        publisher.publishEvent(event);
    }
}
\end{DeskFlowCode}

\phantomsection
\addcontentsline{toc}{subsubsection}{DomainEvent}
\DeskFlowListingCaption{common/application/domain/DomainEvent.java - DomainEvent}{lst:common-application-domain-DomainEvent-java}
\begin{DeskFlowCode}
package com.deskflow.common.application.domain;

import java.time.LocalDateTime;
import java.util.UUID;

public interface DomainEvent {
    UUID eventId();
    String eventType();
    UUID aggregateId();
    LocalDateTime occurredAt();
}
\end{DeskFlowCode}

\clearpage
\section{Kod źródłowy testów}

\phantomsection
\addcontentsline{toc}{subsubsection}{CleanArchitectureTest}
\DeskFlowListingCaption{architecture/CleanArchitectureTest.java - CleanArchitectureTest}{lst:architecture-CleanArchitectureTest-java}
\begin{DeskFlowCode}
package com.deskflow.architecture;

import com.tngtech.archunit.core.domain.JavaClasses;
import com.tngtech.archunit.core.importer.ClassFileImporter;
import com.tngtech.archunit.lang.ArchRule;
import org.junit.jupiter.api.BeforeAll;
import org.junit.jupiter.api.Test;

import static com.tngtech.archunit.lang.syntax.ArchRuleDefinition.noClasses;
import static com.tngtech.archunit.library.dependencies.SlicesRuleDefinition.slices;

class CleanArchitectureTest {

    private static JavaClasses classes;

    @BeforeAll
    static void importClasses() {
        classes = new ClassFileImporter().importPackages("com.deskflow");
    }

    @Test
    void domain_must_not_depend_on_infrastructure() {
        ArchRule rule = noClasses().that().resideInAPackage("..domain..")
                .should().dependOnClassesThat().resideInAPackage("..infrastructure..");
        rule.check(classes);
    }

    @Test
    void domain_must_not_depend_on_adapters() {
        ArchRule rule = noClasses().that().resideInAPackage("..domain..")
                .should().dependOnClassesThat().resideInAPackage("..adapters..");
        rule.check(classes);
    }

    @Test
    void domain_must_not_depend_on_application_service() {
        ArchRule rule = noClasses().that().resideInAPackage("..application.domain..")
                .should().dependOnClassesThat().resideInAPackage("..application.service..");
        rule.check(classes);
    }

    @Test
    void domain_must_not_depend_on_application_port() {
        ArchRule rule = noClasses().that().resideInAPackage("..application.domain..")
                .should().dependOnClassesThat().resideInAPackage("..application.port..");
        rule.check(classes);
    }

    @Test
    void application_must_not_depend_on_infrastructure() {
        ArchRule rule = noClasses().that().resideInAPackage("..application..")
                .should().dependOnClassesThat().resideInAPackage("..infrastructure..");
        rule.check(classes);
    }

    @Test
    void application_must_not_depend_on_adapters() {
        ArchRule rule = noClasses().that().resideInAPackage("..application..")
                .should().dependOnClassesThat().resideInAPackage("..adapters..");
        rule.check(classes);
    }

    @Test
    void domain_contexts_have_no_cyclic_dependencies() {
        slices().matching("com.deskflow.(*)..domain..").should().beFreeOfCycles()
                .check(classes);
    }
}
\end{DeskFlowCode}

\phantomsection
\addcontentsline{toc}{subsubsection}{DomainIsolationTest}
\DeskFlowListingCaption{architecture/DomainIsolationTest.java - DomainIsolationTest}{lst:architecture-DomainIsolationTest-java}
\begin{DeskFlowCode}
package com.deskflow.architecture;

import com.tngtech.archunit.core.domain.JavaClasses;
import com.tngtech.archunit.core.importer.ClassFileImporter;
import com.tngtech.archunit.lang.ArchRule;
import org.junit.jupiter.api.BeforeAll;
import org.junit.jupiter.api.Test;

import static com.tngtech.archunit.lang.syntax.ArchRuleDefinition.noClasses;

class DomainIsolationTest {

    private static JavaClasses classes;

    @BeforeAll
    static void importClasses() {
        classes = new ClassFileImporter().importPackages("com.deskflow");
    }

    @Test
    void domain_must_not_have_jpa_annotations() {
        ArchRule rule = noClasses().that().resideInAPackage("..domain..")
                .should().beAnnotatedWith("jakarta.persistence.Entity")
                .orShould().beAnnotatedWith("jakarta.persistence.Table")
                .orShould().beAnnotatedWith("jakarta.persistence.Column");
        rule.check(classes);
    }

    @Test
    void domain_must_not_import_spring() {
        ArchRule rule = noClasses().that().resideInAPackage("..domain..")
                .should().dependOnClassesThat().resideInAPackage("org.springframework..")
                .because("Domain layer must be framework-independent");
        rule.check(classes);
    }

    @Test
    void domain_must_not_import_jakarta_persistence() {
        ArchRule rule = noClasses().that().resideInAPackage("..domain..")
                .should().dependOnClassesThat().resideInAPackage("jakarta.persistence..");
        rule.check(classes);
    }
}
\end{DeskFlowCode}

\subsection{Testy - Allocation}

\phantomsection
\addcontentsline{toc}{subsubsection}{AllocationRequestTest}
\DeskFlowListingCaption{allocation/application/domain/model/AllocationRequestTest.java - AllocationRequestTest}{lst:allocation-application-domain-model-AllocationRequestTest-java}
\begin{DeskFlowCode}
package com.deskflow.allocation.application.domain.model;

import com.deskflow.common.application.domain.model.vo.EmployeeId;
import com.deskflow.allocation.application.domain.model.events.*;
import com.deskflow.allocation.application.domain.service.AllocationRequestFactory;
import com.deskflow.allocation.application.domain.model.vo.AllocationRequestId;
import com.deskflow.allocation.application.domain.model.vo.DateRange;
import com.deskflow.allocation.application.domain.model.vo.RequestStatus;
import com.deskflow.common.application.domain.model.vo.ResourceId;
import org.junit.jupiter.api.BeforeEach;
import org.junit.jupiter.api.Test;

import java.time.LocalDate;
import java.time.LocalDateTime;
import java.util.UUID;

import static org.assertj.core.api.Assertions.*;

class AllocationRequestTest {

    private static final EmployeeId EMPLOYEE_ID = EmployeeId.of(UUID.randomUUID());
    private static final ResourceId RESOURCE_ID = ResourceId.of(UUID.randomUUID());
    private static final DateRange DATE_RANGE = new DateRange(
            LocalDate.of(2025, 6, 10),
            LocalDate.of(2025, 6, 15)
    );

    private AllocationRequest request;

    @BeforeEach
    void setUp() {
        AllocationRequestFactory factory = new AllocationRequestFactory();
        request = factory.create(EMPLOYEE_ID, RESOURCE_ID, DATE_RANGE, "Potrzebuję do projektu");
    }

    @Test
    void new_request_has_submitted_status() {
        assertThat(request.getStatus()).isEqualTo(RequestStatus.SUBMITTED);
    }

    @Test
    void new_request_publishes_submitted_event() {
        assertThat(request.flushEvents())
                .hasSize(1)
                .first()
                .isInstanceOf(AllocationRequestSubmitted.class);
    }

    @Test
    void approve_changes_status_to_approved() {
        request.flushEvents();
        request.approve();
        assertThat(request.getStatus()).isEqualTo(RequestStatus.APPROVED);
    }

    @Test
    void approve_publishes_approved_event() {
        request.flushEvents();
        request.approve();
        assertThat(request.flushEvents())
                .hasSize(1)
                .first()
                .isInstanceOf(AllocationRequestApproved.class);
    }

    @Test
    void reject_changes_status_and_sets_reason() {
        request.reject("Brak dostępności");
        assertThat(request.getStatus()).isEqualTo(RequestStatus.REJECTED);
        assertThat(request.getRejectionReason()).isEqualTo("Brak dostępności");
    }

    @Test
    void reject_publishes_rejected_event() {
        request.flushEvents();
        request.reject("powód");
        assertThat(request.flushEvents())
                .hasSize(1)
                .first()
                .isInstanceOf(AllocationRequestRejected.class);
    }

    @Test
    void cancel_from_submitted_changes_status() {
        request.cancel();
        assertThat(request.getStatus()).isEqualTo(RequestStatus.CANCELLED);
    }

    @Test
    void cancel_from_approved_changes_status() {
        request.approve();
        request.cancel();
        assertThat(request.getStatus()).isEqualTo(RequestStatus.CANCELLED);
    }

    @Test
    void cancel_from_active_throws() {
        request.approve();
        request.issueResource();
        assertThatThrownBy(request::cancel)
                .isInstanceOf(IllegalStateException.class);
    }

    @Test
    void issue_resource_changes_status_to_active() {
        request.approve();
        request.issueResource();
        assertThat(request.getStatus()).isEqualTo(RequestStatus.ACTIVE);
        assertThat(request.getIssuedAt()).isNotNull();
    }

    @Test
    void issue_resource_from_submitted_throws() {
        assertThatThrownBy(request::issueResource)
                .isInstanceOf(IllegalStateException.class);
    }

    @Test
    void return_resource_from_active_changes_status_to_completed() {
        request.approve();
        request.issueResource();
        request.returnResource();
        assertThat(request.getStatus()).isEqualTo(RequestStatus.COMPLETED);
        assertThat(request.getReturnedAt()).isNotNull();
    }

    @Test
    void return_resource_from_overdue_changes_status_to_completed() {
        request.approve();
        request.issueResource();
        request.markAsOverdue();
        request.returnResource();
        assertThat(request.getStatus()).isEqualTo(RequestStatus.COMPLETED);
    }

    @Test
    void mark_as_overdue_from_active_changes_status() {
        request.approve();
        request.issueResource();
        request.markAsOverdue();
        assertThat(request.getStatus()).isEqualTo(RequestStatus.OVERDUE);
    }

    @Test
    void mark_as_overdue_from_submitted_throws() {
        assertThatThrownBy(request::markAsOverdue)
                .isInstanceOf(IllegalStateException.class);
    }

    @Test
    void approve_after_rejection_throws() {
        request.reject("za późno");
        assertThatThrownBy(request::approve)
                .isInstanceOf(IllegalStateException.class);
    }

    @Test
    void approve_after_cancellation_throws() {
        request.cancel();
        assertThatThrownBy(request::approve)
                .isInstanceOf(IllegalStateException.class);
    }

    @Test
    void request_stores_input_data() {
        assertThat(request.getEmployeeId()).isEqualTo(EMPLOYEE_ID);
        assertThat(request.getResourceId()).isEqualTo(RESOURCE_ID);
        assertThat(request.getDateRange()).isEqualTo(DATE_RANGE);
        assertThat(request.getJustification()).isEqualTo("Potrzebuję do projektu");
        assertThat(request.getSubmittedAt()).isNotNull();
    }

    @Test
    void reconstituted_request_has_no_events() {
        AllocationRequest reconstituted = AllocationRequest.reconstitute(
                AllocationRequestId.newId(), EMPLOYEE_ID, RESOURCE_ID, DATE_RANGE,
                "uzasadnienie", LocalDateTime.now(),
                RequestStatus.APPROVED, null, null, null
        );
        assertThat(reconstituted.flushEvents()).isEmpty();
        assertThat(reconstituted.getStatus()).isEqualTo(RequestStatus.APPROVED);
    }
}
\end{DeskFlowCode}

\phantomsection
\addcontentsline{toc}{subsubsection}{DateRangeTest}
\DeskFlowListingCaption{allocation/application/domain/model/DateRangeTest.java - DateRangeTest}{lst:allocation-application-domain-model-DateRangeTest-java}
\begin{DeskFlowCode}
package com.deskflow.allocation.application.domain.model;

import com.deskflow.allocation.application.domain.model.vo.DateRange;
import org.junit.jupiter.api.Test;

import java.time.LocalDate;

import static org.assertj.core.api.Assertions.*;

class DateRangeTest {

    private static final LocalDate MON = LocalDate.of(2025, 6, 2);
    private static final LocalDate WED = LocalDate.of(2025, 6, 4);
    private static final LocalDate FRI = LocalDate.of(2025, 6, 6);
    private static final LocalDate MON2 = LocalDate.of(2025, 6, 9);

    @Test
    void creates_date_range_when_from_before_to() {
        DateRange r = new DateRange(MON, FRI);
        assertThat(r.getDateFrom()).isEqualTo(MON);
        assertThat(r.getDateTo()).isEqualTo(FRI);
    }

    @Test
    void throws_when_date_from_equals_date_to() {
        assertThatThrownBy(() -> new DateRange(MON, MON))
                .isInstanceOf(IllegalArgumentException.class)
                .hasMessageContaining("dateFrom must be before dateTo");
    }

    @Test
    void throws_when_date_from_after_date_to() {
        assertThatThrownBy(() -> new DateRange(FRI, MON))
                .isInstanceOf(IllegalArgumentException.class);
    }

    @Test
    void throws_when_date_from_null() {
        assertThatThrownBy(() -> new DateRange(null, FRI))
                .isInstanceOf(NullPointerException.class);
    }

    @Test
    void value_object_equality() {
        DateRange r1 = new DateRange(MON, FRI);
        DateRange r2 = new DateRange(MON, FRI);
        assertThat(r1).isEqualTo(r2);
        assertThat(r1.hashCode()).isEqualTo(r2.hashCode());
    }

    @Test
    void not_equal_when_different_dates() {
        DateRange r1 = new DateRange(MON, FRI);
        DateRange r2 = new DateRange(WED, MON2);
        assertThat(r1).isNotEqualTo(r2);
    }
}
\end{DeskFlowCode}

\phantomsection
\addcontentsline{toc}{subsubsection}{AllocationLifecycleDomainServicesTest}
\DeskFlowListingCaption{allocation/application/domain/service/AllocationLifecycleDomainServicesTest.java - AllocationLifecycleDomainServicesTest}{lst:allocation-application-domain-service-AllocationLifecycleDomainServicesTest-java}
\begin{DeskFlowCode}
package com.deskflow.allocation.application.domain.service;

import com.deskflow.allocation.application.domain.model.AllocationRequest;
import com.deskflow.allocation.application.domain.model.vo.AllocationRequestId;
import com.deskflow.allocation.application.domain.model.vo.DateRange;
import com.deskflow.allocation.application.domain.model.vo.RequestStatus;
import com.deskflow.common.application.domain.model.vo.EmployeeId;
import com.deskflow.common.application.domain.model.vo.ResourceId;
import org.junit.jupiter.api.Test;

import java.time.LocalDate;
import java.time.LocalDateTime;
import java.util.UUID;

import static org.assertj.core.api.Assertions.assertThat;
import static org.assertj.core.api.Assertions.assertThatThrownBy;

class AllocationLifecycleDomainServicesTest {

    private final ApproveAllocationRequestDomainService approveService = new ApproveAllocationRequestDomainService();
    private final RejectAllocationRequestDomainService rejectService = new RejectAllocationRequestDomainService();
    private final IssueAllocationResourceDomainService issueService = new IssueAllocationResourceDomainService();
    private final ReturnAllocationResourceDomainService returnService = new ReturnAllocationResourceDomainService();
    private final MarkAllocationRequestOverdueDomainService overdueService = new MarkAllocationRequestOverdueDomainService();
    private final CancelAllocationRequestDomainService cancelService = new CancelAllocationRequestDomainService();

    private final EmployeeId employeeId = EmployeeId.of(UUID.randomUUID());

    @Test
    void approve_service_approves_submitted_request() {
        AllocationRequest request = request(RequestStatus.SUBMITTED);

        approveService.approveAllocationRequest(request);

        assertThat(request.getStatus()).isEqualTo(RequestStatus.APPROVED);
    }

    @Test
    void reject_service_rejects_request_and_stores_reason() {
        AllocationRequest request = request(RequestStatus.SUBMITTED);

        rejectService.rejectAllocationRequest(request, "No resource available");

        assertThat(request.getStatus()).isEqualTo(RequestStatus.REJECTED);
        assertThat(request.getRejectionReason()).isEqualTo("No resource available");
    }

    @Test
    void reject_service_rejects_blank_reason() {
        AllocationRequest request = request(RequestStatus.SUBMITTED);

        assertThatThrownBy(() -> rejectService.rejectAllocationRequest(request, " "))
                .isInstanceOf(IllegalArgumentException.class)
                .hasMessageContaining("cannot be blank");
    }

    @Test
    void issue_service_issues_approved_request() {
        AllocationRequest request = request(RequestStatus.APPROVED);

        issueService.issueAllocationResource(request);

        assertThat(request.getStatus()).isEqualTo(RequestStatus.ACTIVE);
        assertThat(request.getIssuedAt()).isNotNull();
    }

    @Test
    void return_service_completes_active_request() {
        AllocationRequest request = request(RequestStatus.ACTIVE);

        returnService.returnAllocationResource(request);

        assertThat(request.getStatus()).isEqualTo(RequestStatus.COMPLETED);
        assertThat(request.getReturnedAt()).isNotNull();
    }

    @Test
    void overdue_service_marks_active_request_as_overdue() {
        AllocationRequest request = request(RequestStatus.ACTIVE);

        overdueService.markAllocationRequestOverdue(request);

        assertThat(request.getStatus()).isEqualTo(RequestStatus.OVERDUE);
    }

    @Test
    void cancel_service_cancels_own_request() {
        AllocationRequest request = request(RequestStatus.SUBMITTED);

        cancelService.cancelAllocationRequest(request, employeeId);

        assertThat(request.getStatus()).isEqualTo(RequestStatus.CANCELLED);
    }

    @Test
    void cancel_service_rejects_cancelling_another_employees_request() {
        AllocationRequest request = request(RequestStatus.SUBMITTED);

        assertThatThrownBy(() -> cancelService.cancelAllocationRequest(request, EmployeeId.of(UUID.randomUUID())))
                .isInstanceOf(IllegalStateException.class)
                .hasMessageContaining("another employee");
    }

    private AllocationRequest request(RequestStatus status) {
        return AllocationRequest.reconstitute(
                AllocationRequestId.newId(), employeeId, ResourceId.of(UUID.randomUUID()),
                new DateRange(LocalDate.of(2025, 7, 1), LocalDate.of(2025, 7, 5)),
                "test", LocalDateTime.now(), status, null, null, null
        );
    }
}
\end{DeskFlowCode}

\phantomsection
\addcontentsline{toc}{subsubsection}{EmployeeLimitDomainServiceTest}
\DeskFlowListingCaption{allocation/application/domain/service/EmployeeLimitDomainServiceTest.java - EmployeeLimitDomainServiceTest}{lst:allocation-application-domain-service-EmployeeLimitDomainServiceTest-java}
\begin{DeskFlowCode}
package com.deskflow.allocation.application.domain.service;

import com.deskflow.allocation.application.domain.model.EmployeeAllocationProfile;
import com.deskflow.common.application.domain.model.vo.ResourceType;
import org.junit.jupiter.api.Test;

import java.util.Map;

import static org.assertj.core.api.Assertions.assertThatCode;
import static org.assertj.core.api.Assertions.assertThatThrownBy;

class EmployeeLimitDomainServiceTest {

    private final EnsureEmployeeWithinLimitDomainService service = new EnsureEmployeeWithinLimitDomainService();

    @Test
    void allows_allocation_when_current_usage_is_below_limit() {
        EmployeeAllocationProfile employee = new EmployeeAllocationProfile(Map.of(ResourceType.DESK.name(), 2));

        assertThatCode(() -> service.ensureEmployeeWithinLimit(employee, ResourceType.DESK.name(), 1))
                .doesNotThrowAnyException();
    }

    @Test
    void rejects_allocation_when_current_usage_equals_limit() {
        EmployeeAllocationProfile employee = new EmployeeAllocationProfile(Map.of(ResourceType.DESK.name(), 2));

        assertThatThrownBy(() -> service.ensureEmployeeWithinLimit(employee, ResourceType.DESK.name(), 2))
                .isInstanceOf(IllegalStateException.class)
                .hasMessageContaining("allocation limit");
    }

    @Test
    void rejects_allocation_when_current_usage_exceeds_limit() {
        EmployeeAllocationProfile employee = new EmployeeAllocationProfile(Map.of(ResourceType.DESK.name(), 2));

        assertThatThrownBy(() -> service.ensureEmployeeWithinLimit(employee, ResourceType.DESK.name(), 3))
                .isInstanceOf(IllegalStateException.class)
                .hasMessageContaining("allocation limit");
    }
}
\end{DeskFlowCode}

\phantomsection
\addcontentsline{toc}{subsubsection}{ResourceAvailabilityDomainServicesTest}
\DeskFlowListingCaption{allocation/application/domain/service/ResourceAvailabilityDomainServicesTest.java - ResourceAvailabilityDomainServicesTest}{lst:allocation-application-domain-service-ResourceAvailabilityDomainServicesTest-java}
\begin{DeskFlowCode}
package com.deskflow.allocation.application.domain.service;

import com.deskflow.allocation.application.domain.model.AllocationRequest;
import com.deskflow.allocation.application.domain.model.ResourceAllocationProfile;
import com.deskflow.allocation.application.domain.model.vo.AllocationRequestId;
import com.deskflow.allocation.application.domain.model.vo.DateRange;
import com.deskflow.allocation.application.domain.model.vo.RequestStatus;
import com.deskflow.common.application.domain.model.vo.EmployeeId;
import com.deskflow.common.application.domain.model.vo.ResourceId;
import com.deskflow.common.application.domain.model.vo.ResourceType;
import org.junit.jupiter.api.Test;

import java.time.LocalDate;
import java.time.LocalDateTime;
import java.util.Collections;
import java.util.List;
import java.util.UUID;

import static org.assertj.core.api.Assertions.assertThatCode;
import static org.assertj.core.api.Assertions.assertThatThrownBy;

class ResourceAvailabilityDomainServicesTest {

    private final EnsureResourceAvailableInPeriodDomainService ensureResourceAvailableInPeriodDomainService = new EnsureResourceAvailableInPeriodDomainService();
    private final EnsureResourceEligibleForAllocationDomainService ensureResourceEligibleForAllocationDomainService = new EnsureResourceEligibleForAllocationDomainService();

    @Test
    void requires_available_when_no_overlapping_requests() {
        assertThatCode(() -> ensureResourceAvailableInPeriodDomainService.ensureResourceAvailableInPeriod(Collections.emptyList()))
                .doesNotThrowAnyException();
    }

    @Test
    void throws_when_overlapping_requests_exist() {
        AllocationRequest existing = AllocationRequest.reconstitute(
                AllocationRequestId.newId(), EmployeeId.of(UUID.randomUUID()), ResourceId.of(UUID.randomUUID()),
                new DateRange(LocalDate.of(2025, 7, 1), LocalDate.of(2025, 7, 10)),
                "test", LocalDateTime.now(), RequestStatus.ACTIVE, null, null, null
        );

        assertThatThrownBy(() -> ensureResourceAvailableInPeriodDomainService.ensureResourceAvailableInPeriod(List.of(existing)))
                .isInstanceOf(IllegalStateException.class)
                .hasMessageContaining("not available for the requested period");
    }
    
    @Test
    void requires_eligible_for_allocation_does_not_throw_when_active_and_available() {
        ResourceAllocationProfile resource = new ResourceAllocationProfile(
                ResourceId.of(UUID.randomUUID()), ResourceType.DESK.name(), "D1", "A", 1, "1", true, true);
        assertThatCode(() -> ensureResourceEligibleForAllocationDomainService.ensureResourceEligibleForAllocation(resource))
                .doesNotThrowAnyException();
    }
    
    @Test
    void requires_eligible_throws_when_unavailable() {
        ResourceAllocationProfile resource = new ResourceAllocationProfile(
                ResourceId.of(UUID.randomUUID()), ResourceType.DESK.name(), "D1", "A", 1, "1", true, false);
        assertThatThrownBy(() -> ensureResourceEligibleForAllocationDomainService.ensureResourceEligibleForAllocation(resource))
                .isInstanceOf(IllegalStateException.class)
                .hasMessageContaining("Resource is not available for allocation");
    }
}
\end{DeskFlowCode}

\phantomsection
\addcontentsline{toc}{subsubsection}{AllocationRequestViewTest}
\DeskFlowListingCaption{allocation/application/model/AllocationRequestViewTest.java - AllocationRequestViewTest}{lst:allocation-application-model-AllocationRequestViewTest-java}
\begin{DeskFlowCode}
package com.deskflow.allocation.application.model;

import com.deskflow.allocation.application.domain.model.AllocationRequest;
import com.deskflow.allocation.application.domain.model.vo.AllocationRequestId;
import com.deskflow.allocation.application.domain.model.vo.DateRange;
import com.deskflow.allocation.application.domain.model.vo.RequestStatus;
import com.deskflow.common.application.domain.model.vo.EmployeeId;
import com.deskflow.common.application.domain.model.vo.ResourceId;
import org.junit.jupiter.api.Test;

import java.time.LocalDate;
import java.time.LocalDateTime;
import java.util.UUID;

import static org.assertj.core.api.Assertions.*;

class AllocationRequestViewTest {

    @Test
    void from_creates_view_with_all_fields() {
        var id = AllocationRequestId.newId();
        var employeeId = EmployeeId.of(UUID.randomUUID());
        var resourceId = ResourceId.of(UUID.randomUUID());
        var dateRange = new DateRange(LocalDate.of(2025, 6, 10), LocalDate.of(2025, 6, 15));
        var submittedAt = LocalDateTime.of(2025, 6, 1, 10, 0);
        var request = AllocationRequest.reconstitute(id, employeeId, resourceId, dateRange,
                "justification", submittedAt, RequestStatus.SUBMITTED, null, null, null);

        var view = AllocationRequestView.from(request);

        assertThat(view.id()).isEqualTo(id.getValue());
        assertThat(view.employeeId()).isEqualTo(employeeId.getValue());
        assertThat(view.resourceId()).isEqualTo(resourceId.getValue());
        assertThat(view.status()).isEqualTo("SUBMITTED");
        assertThat(view.dateFrom()).isEqualTo(LocalDate.of(2025, 6, 10));
        assertThat(view.dateTo()).isEqualTo(LocalDate.of(2025, 6, 15));
        assertThat(view.justification()).isEqualTo("justification");
        assertThat(view.submittedAt()).isEqualTo(submittedAt);
        assertThat(view.issuedAt()).isNull();
        assertThat(view.returnedAt()).isNull();
        assertThat(view.rejectionReason()).isNull();
    }

    @Test
    void from_creates_view_for_completed_request() {
        var id = AllocationRequestId.newId();
        var employeeId = EmployeeId.of(UUID.randomUUID());
        var resourceId = ResourceId.of(UUID.randomUUID());
        var dateRange = new DateRange(LocalDate.of(2025, 6, 10), LocalDate.of(2025, 6, 15));
        var submittedAt = LocalDateTime.of(2025, 6, 1, 10, 0);
        var issuedAt = LocalDateTime.of(2025, 6, 2, 9, 0);
        var returnedAt = LocalDateTime.of(2025, 6, 14, 16, 0);
        var request = AllocationRequest.reconstitute(id, employeeId, resourceId, dateRange,
                "test", submittedAt, RequestStatus.COMPLETED, issuedAt, returnedAt, null);

        var view = AllocationRequestView.from(request);

        assertThat(view.status()).isEqualTo("COMPLETED");
        assertThat(view.issuedAt()).isEqualTo(issuedAt);
        assertThat(view.returnedAt()).isEqualTo(returnedAt);
    }

    @Test
    void from_creates_view_for_rejected_request() {
        var id = AllocationRequestId.newId();
        var employeeId = EmployeeId.of(UUID.randomUUID());
        var resourceId = ResourceId.of(UUID.randomUUID());
        var dateRange = new DateRange(LocalDate.of(2025, 6, 10), LocalDate.of(2025, 6, 15));
        var request = AllocationRequest.reconstitute(id, employeeId, resourceId, dateRange,
                "test", LocalDateTime.now(), RequestStatus.REJECTED, null, null, "Brak miejsca");

        var view = AllocationRequestView.from(request);

        assertThat(view.status()).isEqualTo("REJECTED");
        assertThat(view.rejectionReason()).isEqualTo("Brak miejsca");
    }
}
\end{DeskFlowCode}

\phantomsection
\addcontentsline{toc}{subsubsection}{AllocationAdministrationAppServiceTest}
\DeskFlowListingCaption{allocation/application/service/AllocationAdministrationAppServiceTest.java - AllocationAdministrationAppServiceTest}{lst:allocation-application-service-AllocationAdministrationAppServiceTest-java}
\begin{DeskFlowCode}
package com.deskflow.allocation.application.service;

import com.deskflow.allocation.application.port.out.AllocationRequestRepositoryPort;
import com.deskflow.allocation.application.domain.model.AllocationRequest;
import com.deskflow.allocation.application.domain.model.events.AllocationRequestApproved;
import com.deskflow.allocation.application.domain.model.events.AllocationRequestRejected;
import com.deskflow.allocation.application.domain.model.events.ResourceIssued;
import com.deskflow.allocation.application.domain.model.events.ResourceReturned;
import com.deskflow.allocation.application.domain.model.vo.DateRange;
import com.deskflow.allocation.application.domain.model.vo.RequestStatus;
import com.deskflow.allocation.application.domain.service.ApproveAllocationRequestDomainService;
import com.deskflow.allocation.application.domain.service.IssueAllocationResourceDomainService;
import com.deskflow.allocation.application.domain.service.MarkAllocationRequestOverdueDomainService;
import com.deskflow.allocation.application.domain.service.RejectAllocationRequestDomainService;
import com.deskflow.allocation.application.domain.service.ReturnAllocationResourceDomainService;
import com.deskflow.common.application.port.out.DomainEventPublisherPort;
import com.deskflow.allocation.application.domain.model.vo.AllocationRequestId;
import com.deskflow.common.application.domain.model.vo.EmployeeId;
import com.deskflow.common.application.domain.model.vo.ResourceId;
import org.junit.jupiter.api.BeforeEach;
import org.junit.jupiter.api.Test;
import org.junit.jupiter.api.extension.ExtendWith;
import org.mockito.Mock;
import org.mockito.junit.jupiter.MockitoExtension;

import java.time.LocalDate;
import java.time.LocalDateTime;
import java.util.List;
import java.util.NoSuchElementException;
import java.util.Optional;
import java.util.UUID;

import static org.assertj.core.api.Assertions.assertThat;
import static org.assertj.core.api.Assertions.assertThatThrownBy;
import static org.mockito.Mockito.any;
import static org.mockito.Mockito.verify;
import static org.mockito.Mockito.when;

@ExtendWith(MockitoExtension.class)
class AllocationAdministrationAppServiceTest {

    @Mock private AllocationRequestRepositoryPort requestRepository;
    @Mock private DomainEventPublisherPort eventPublisher;

    private AllocationAdministrationAppService appService;
    private EmployeeId employeeId;
    private ResourceId resourceId;

    @BeforeEach
    void setUp() {
        appService = new AllocationAdministrationAppService(
                requestRepository,
                eventPublisher,
                new ApproveAllocationRequestDomainService(),
                new RejectAllocationRequestDomainService(),
                new IssueAllocationResourceDomainService(),
                new ReturnAllocationResourceDomainService(),
                new MarkAllocationRequestOverdueDomainService());
        employeeId = EmployeeId.of(UUID.randomUUID());
        resourceId = ResourceId.of(UUID.randomUUID());
    }

    private AllocationRequest createRequest(RequestStatus status) {
        return AllocationRequest.reconstitute(
                AllocationRequestId.newId(), employeeId, resourceId,
                new DateRange(LocalDate.of(2025, 7, 1), LocalDate.of(2025, 7, 5)),
                "test", LocalDateTime.now(),
                status, null, null, null
        );
    }

    @Test
    void approve_request_changes_status_and_publishes_event() {
        AllocationRequest request = createRequest(RequestStatus.SUBMITTED);
        when(requestRepository.find(request.getId())).thenReturn(Optional.of(request));

        appService.approveRequest(request.getId());

        assertThat(request.getStatus()).isEqualTo(RequestStatus.APPROVED);
        verify(requestRepository).save(request);
        verify(eventPublisher).publish(any(AllocationRequestApproved.class));
    }

    @Test
    void reject_request_changes_status_and_publishes_event() {
        AllocationRequest request = createRequest(RequestStatus.SUBMITTED);
        when(requestRepository.find(request.getId())).thenReturn(Optional.of(request));

        appService.rejectRequest(request.getId(), "brak");

        assertThat(request.getStatus()).isEqualTo(RequestStatus.REJECTED);
        assertThat(request.getRejectionReason()).isEqualTo("brak");
        verify(eventPublisher).publish(any(AllocationRequestRejected.class));
    }

    @Test
    void issue_request_changes_status_and_publishes_event() {
        AllocationRequest request = createRequest(RequestStatus.APPROVED);
        when(requestRepository.find(request.getId())).thenReturn(Optional.of(request));

        appService.issueResource(request.getId());

        assertThat(request.getStatus()).isEqualTo(RequestStatus.ACTIVE);
        verify(eventPublisher).publish(any(ResourceIssued.class));
    }

    @Test
    void return_resource_changes_status_and_publishes_event() {
        AllocationRequest request = createRequest(RequestStatus.ACTIVE);
        when(requestRepository.find(request.getId())).thenReturn(Optional.of(request));

        appService.returnResource(request.getId());

        assertThat(request.getStatus()).isEqualTo(RequestStatus.COMPLETED);
        verify(eventPublisher).publish(any(ResourceReturned.class));
    }

    @Test
    void detect_overdue_marks_active_past_due_requests() {
        AllocationRequest request = createRequest(RequestStatus.ACTIVE);
        when(requestRepository.findOverdue(any(LocalDateTime.class))).thenReturn(List.of(request));

        appService.detectOverdueAllocations();

        assertThat(request.getStatus()).isEqualTo(RequestStatus.OVERDUE);
        verify(requestRepository).save(request);
    }

    @Test
    void approve_throws_when_request_missing() {
        AllocationRequestId requestId = AllocationRequestId.newId();
        when(requestRepository.find(requestId)).thenReturn(Optional.empty());

        assertThatThrownBy(() -> appService.approveRequest(requestId))
                .isInstanceOf(NoSuchElementException.class)
                .hasMessageContaining("Allocation request not found");
    }
}
\end{DeskFlowCode}

\phantomsection
\addcontentsline{toc}{subsubsection}{AllocationRequestAppServiceMissingDataTest}
\DeskFlowListingCaption{allocation/application/service/AllocationRequestAppServiceMissingDataTest.java - AllocationRequestAppServiceMissingDataTest}{lst:allocation-application-service-AllocationRequestAppServiceMissingDataTest-java}
\begin{DeskFlowCode}
package com.deskflow.allocation.application.service;

import com.deskflow.allocation.application.port.out.AllocationRequestRepositoryPort;
import com.deskflow.allocation.application.port.out.EmployeeProfilePort;
import com.deskflow.allocation.application.port.out.ResourceCatalogPort;
import com.deskflow.allocation.application.domain.model.EmployeeAllocationProfile;
import com.deskflow.allocation.application.domain.model.ResourceAllocationProfile;
import com.deskflow.common.application.domain.model.vo.EmployeeId;
import com.deskflow.allocation.application.domain.service.AllocationRequestFactory;
import com.deskflow.allocation.application.domain.service.CancelAllocationRequestDomainService;
import com.deskflow.allocation.application.domain.service.EnsureEmployeeWithinLimitDomainService;
import com.deskflow.allocation.application.domain.service.EnsureResourceAvailableInPeriodDomainService;
import com.deskflow.allocation.application.domain.service.EnsureResourceEligibleForAllocationDomainService;
import com.deskflow.allocation.application.domain.model.vo.AllocationRequestId;
import com.deskflow.common.application.domain.model.vo.ResourceId;
import com.deskflow.common.application.domain.model.vo.ResourceType;
import com.deskflow.common.application.port.out.DomainEventPublisherPort;
import org.junit.jupiter.api.BeforeEach;
import org.junit.jupiter.api.Test;
import org.junit.jupiter.api.extension.ExtendWith;
import org.mockito.Mock;
import org.mockito.junit.jupiter.MockitoExtension;

import java.time.LocalDate;
import java.util.Map;
import java.util.NoSuchElementException;
import java.util.Optional;
import java.util.UUID;

import static org.assertj.core.api.Assertions.assertThatThrownBy;
import static org.mockito.Mockito.when;

@ExtendWith(MockitoExtension.class)
class AllocationRequestAppServiceMissingDataTest {

    @Mock private AllocationRequestRepositoryPort requestRepository;
    @Mock private ResourceCatalogPort resourceCatalogPort;
    @Mock private EmployeeProfilePort employeeProfilePort;
    @Mock private DomainEventPublisherPort eventPublisher;

    private AllocationRequestAppService appService;
    private EmployeeId employeeId;
    private ResourceId resourceId;
    private ResourceAllocationProfile resource;
    private EmployeeAllocationProfile employee;

    @BeforeEach
    void setUp() {
        employeeId = EmployeeId.of(UUID.randomUUID());
        resourceId = ResourceId.of(UUID.randomUUID());

        resource = new ResourceAllocationProfile(resourceId, ResourceType.DESK.name(), "Test desk", "A", 1, "1", true, true);
        employee = new EmployeeAllocationProfile(Map.of(ResourceType.DESK.name(), 3));

        appService = new AllocationRequestAppService(
                requestRepository, resourceCatalogPort, employeeProfilePort,
                new AllocationRequestFactory(), new CancelAllocationRequestDomainService(),
                new EnsureResourceEligibleForAllocationDomainService(), new EnsureResourceAvailableInPeriodDomainService(),
                new EnsureEmployeeWithinLimitDomainService(), eventPublisher
        );
    }

    @Test
    void submit_throws_when_employee_missing() {
        when(employeeProfilePort.findEmployee(employeeId)).thenReturn(Optional.empty());

        assertThatThrownBy(() -> appService.submitRequest(employeeId, resourceId, LocalDate.of(2025, 7, 1), LocalDate.of(2025, 7, 5), "test"))
                .isInstanceOf(IllegalArgumentException.class)
                .hasMessageContaining("Employee not found");
    }

    @Test
    void submit_throws_when_resource_missing() {
        when(employeeProfilePort.findEmployee(employeeId)).thenReturn(Optional.of(employee));
        when(resourceCatalogPort.findResource(resourceId)).thenReturn(Optional.empty());

        assertThatThrownBy(() -> appService.submitRequest(employeeId, resourceId, LocalDate.of(2025, 7, 1), LocalDate.of(2025, 7, 5), "test"))
                .isInstanceOf(IllegalArgumentException.class)
                .hasMessageContaining("Resource not found");
    }

    @Test
    void cancel_throws_when_request_missing() {
        AllocationRequestId requestId = AllocationRequestId.newId();
        when(requestRepository.find(requestId)).thenReturn(Optional.empty());

        assertThatThrownBy(() -> appService.cancelRequest(requestId, employeeId))
                .isInstanceOf(NoSuchElementException.class)
                .hasMessageContaining("Allocation request not found");
    }
}
\end{DeskFlowCode}

\phantomsection
\addcontentsline{toc}{subsubsection}{AllocationRequestAppServiceTest}
\DeskFlowListingCaption{allocation/application/service/AllocationRequestAppServiceTest.java - AllocationRequestAppServiceTest}{lst:allocation-application-service-AllocationRequestAppServiceTest-java}
\begin{DeskFlowCode}
package com.deskflow.allocation.application.service;

import com.deskflow.allocation.application.port.out.AllocationRequestRepositoryPort;
import com.deskflow.allocation.application.port.out.EmployeeProfilePort;
import com.deskflow.allocation.application.port.out.ResourceCatalogPort;
import com.deskflow.allocation.application.domain.model.EmployeeAllocationProfile;
import com.deskflow.allocation.application.domain.model.ResourceAllocationProfile;
import com.deskflow.common.application.domain.model.vo.EmployeeId;
import com.deskflow.allocation.application.domain.model.AllocationRequest;
import com.deskflow.allocation.application.domain.service.AllocationRequestFactory;
import com.deskflow.allocation.application.domain.service.CancelAllocationRequestDomainService;
import com.deskflow.allocation.application.domain.service.EnsureEmployeeWithinLimitDomainService;
import com.deskflow.allocation.application.domain.service.EnsureResourceAvailableInPeriodDomainService;
import com.deskflow.allocation.application.domain.service.EnsureResourceEligibleForAllocationDomainService;
import com.deskflow.allocation.application.domain.model.vo.AllocationRequestId;
import com.deskflow.allocation.application.domain.model.vo.DateRange;
import com.deskflow.allocation.application.domain.model.vo.RequestStatus;
import com.deskflow.common.application.domain.model.vo.ResourceId;
import com.deskflow.common.application.domain.model.vo.ResourceType;
import com.deskflow.common.application.port.out.DomainEventPublisherPort;
import org.junit.jupiter.api.BeforeEach;
import org.junit.jupiter.api.Test;
import org.junit.jupiter.api.extension.ExtendWith;
import org.mockito.Mock;
import org.mockito.junit.jupiter.MockitoExtension;

import java.time.LocalDate;
import java.time.LocalDateTime;
import java.util.Collections;
import java.util.List;
import java.util.Map;
import java.util.Optional;
import java.util.UUID;

import static org.assertj.core.api.Assertions.*;
import static org.mockito.Mockito.*;

@ExtendWith(MockitoExtension.class)
class AllocationRequestAppServiceTest {

    @Mock private AllocationRequestRepositoryPort requestRepository;
    @Mock private ResourceCatalogPort resourceCatalogPort;
    @Mock private EmployeeProfilePort employeeProfilePort;
    @Mock private DomainEventPublisherPort eventPublisher;

    private AllocationRequestAppService appService;
    private EmployeeId employeeId;
    private ResourceId resourceId;
    private ResourceAllocationProfile resource;
    private EmployeeAllocationProfile employee;

    @BeforeEach
    void setUp() {
        employeeId = EmployeeId.of(UUID.randomUUID());
        resourceId = ResourceId.of(UUID.randomUUID());

        resource = new ResourceAllocationProfile(resourceId, ResourceType.DESK.name(), "Test desk", "A", 1, "1", true, true);
        employee = new EmployeeAllocationProfile(Map.of(ResourceType.DESK.name(), 3));

        appService = new AllocationRequestAppService(
                requestRepository, resourceCatalogPort, employeeProfilePort,
                new AllocationRequestFactory(), new CancelAllocationRequestDomainService(),
                new EnsureResourceEligibleForAllocationDomainService(), new EnsureResourceAvailableInPeriodDomainService(),
                new EnsureEmployeeWithinLimitDomainService(), eventPublisher
        );
    }

    @Test
    void submit_saves_request_when_within_limit_and_resource_available() {
        when(employeeProfilePort.findEmployee(employeeId)).thenReturn(Optional.of(employee));
        when(resourceCatalogPort.findResource(resourceId)).thenReturn(Optional.of(resource));
        when(requestRepository.findActiveForEmployee(employeeId)).thenReturn(Collections.emptyList());
        when(requestRepository.findActiveForResource(eq(resourceId), any(DateRange.class))).thenReturn(Collections.emptyList());

        AllocationRequestId id = appService.submitRequest(employeeId, resourceId, LocalDate.of(2025, 7, 1), LocalDate.of(2025, 7, 5), "test");

        assertThat(id).isNotNull();
        verify(requestRepository).save(any(AllocationRequest.class));
    }

    @Test
    void submit_throws_when_limit_exceeded() {
        AllocationRequest a1 = AllocationRequest.reconstitute(
                AllocationRequestId.newId(), employeeId, resourceId,
                new DateRange(LocalDate.of(2025, 6, 1), LocalDate.of(2025, 6, 10)),
                "x", LocalDateTime.now(), RequestStatus.ACTIVE, LocalDateTime.now(), null, null
        );
        AllocationRequest a2 = AllocationRequest.reconstitute(
                AllocationRequestId.newId(), employeeId, resourceId,
                new DateRange(LocalDate.of(2025, 6, 11), LocalDate.of(2025, 6, 20)),
                "x", LocalDateTime.now(), RequestStatus.ACTIVE, LocalDateTime.now(), null, null
        );
        AllocationRequest a3 = AllocationRequest.reconstitute(
                AllocationRequestId.newId(), employeeId, resourceId,
                new DateRange(LocalDate.of(2025, 6, 21), LocalDate.of(2025, 6, 30)),
                "x", LocalDateTime.now(), RequestStatus.ACTIVE, LocalDateTime.now(), null, null
        );
        when(employeeProfilePort.findEmployee(employeeId)).thenReturn(Optional.of(employee));
        when(resourceCatalogPort.findResource(resourceId)).thenReturn(Optional.of(resource));
        when(requestRepository.findActiveForEmployee(employeeId))
                .thenReturn(List.of(a1, a2, a3));

        assertThatThrownBy(() -> appService.submitRequest(employeeId, resourceId, LocalDate.of(2025, 7, 1), LocalDate.of(2025, 7, 5), "test"))
                .isInstanceOf(IllegalStateException.class)
                .hasMessageContaining("limit");
    }

    @Test
    void submit_throws_when_resource_unavailable() {
        when(employeeProfilePort.findEmployee(employeeId)).thenReturn(Optional.of(employee));
        when(resourceCatalogPort.findResource(resourceId)).thenReturn(Optional.of(resource));
        when(requestRepository.findActiveForEmployee(employeeId)).thenReturn(Collections.emptyList());

        AllocationRequest blocking = AllocationRequest.reconstitute(
                AllocationRequestId.newId(), EmployeeId.of(UUID.randomUUID()), resourceId,
                new DateRange(LocalDate.of(2025, 7, 1), LocalDate.of(2025, 7, 10)),
                "x", LocalDateTime.now(), RequestStatus.ACTIVE, LocalDateTime.now(), null, null
        );
        when(requestRepository.findActiveForResource(eq(resourceId), any(DateRange.class)))
                .thenReturn(List.of(blocking));

        assertThatThrownBy(() -> appService.submitRequest(employeeId, resourceId, LocalDate.of(2025, 7, 1), LocalDate.of(2025, 7, 5), "test"))
                .isInstanceOf(IllegalStateException.class)
                .hasMessageContaining("available");
    }

    @Test
    void cancel_changes_status_and_saves() {
        AllocationRequestId requestId = AllocationRequestId.newId();
        AllocationRequest request = AllocationRequest.reconstitute(
                requestId, employeeId, resourceId,
                new DateRange(LocalDate.of(2025, 7, 1), LocalDate.of(2025, 7, 5)),
                "test", LocalDateTime.now(),
                RequestStatus.SUBMITTED, null, null, null
        );
        when(requestRepository.find(requestId)).thenReturn(Optional.of(request));

        appService.cancelRequest(requestId, employeeId);

        assertThat(request.getStatus()).isEqualTo(RequestStatus.CANCELLED);
        verify(requestRepository).save(request);
    }
}
\end{DeskFlowCode}

\phantomsection
\addcontentsline{toc}{subsubsection}{AllocationRequestRestMapperTest}
\DeskFlowListingCaption{allocation/infrastructure/in/web/AllocationRequestRestMapperTest.java - AllocationRequestRestMapperTest}{lst:allocation-infrastructure-in-web-AllocationRequestRestMapperTest-java}
\begin{DeskFlowCode}
package com.deskflow.allocation.infrastructure.in.web;

import com.deskflow.allocation.application.model.AllocationRequestView;
import com.deskflow.allocation.infrastructure.in.web.dto.AllocationRequestResponseDTO;
import org.junit.jupiter.api.Test;

import java.time.LocalDate;
import java.time.LocalDateTime;
import java.util.UUID;

import static org.assertj.core.api.Assertions.*;

class AllocationRequestRestMapperTest {

    private final UUID id = UUID.randomUUID();
    private final UUID employeeId = UUID.randomUUID();
    private final UUID resourceId = UUID.randomUUID();
    private final LocalDate dateFrom = LocalDate.of(2025, 6, 10);
    private final LocalDate dateTo = LocalDate.of(2025, 6, 15);
    private final LocalDateTime submittedAt = LocalDateTime.of(2025, 6, 1, 10, 0);

    @Test
    void toResponse_maps_all_fields() {
        var view = new AllocationRequestView(id, employeeId, resourceId, "SUBMITTED",
                dateFrom, dateTo, "justification", submittedAt, null, null, null);

        var response = AllocationRequestRestMapper.toResponse(view);

        assertThat(response.id()).isEqualTo(id);
        assertThat(response.employeeId()).isEqualTo(employeeId);
        assertThat(response.resourceId()).isEqualTo(resourceId);
        assertThat(response.status()).isEqualTo("SUBMITTED");
        assertThat(response.dateFrom()).isEqualTo(dateFrom);
        assertThat(response.dateTo()).isEqualTo(dateTo);
        assertThat(response.justification()).isEqualTo("justification");
        assertThat(response.submittedAt()).isEqualTo(submittedAt);
        assertThat(response.issuedAt()).isNull();
        assertThat(response.returnedAt()).isNull();
        assertThat(response.rejectionReason()).isNull();
        assertThat(response.employeeFirstName()).isNull();
        assertThat(response.employeeLastName()).isNull();
    }

    @Test
    void toResponse_with_employee_names() {
        var view = new AllocationRequestView(id, employeeId, resourceId, "APPROVED",
                dateFrom, dateTo, "test", submittedAt, LocalDateTime.now(), null, null);

        var response = AllocationRequestRestMapper.toResponse(view, "Anna", "Kowalska");

        assertThat(response.employeeFirstName()).isEqualTo("Anna");
        assertThat(response.employeeLastName()).isEqualTo("Kowalska");
    }

    @Test
    void toResponse_maps_completed_request_with_all_fields() {
        var issuedAt = LocalDateTime.of(2025, 6, 11, 9, 0);
        var returnedAt = LocalDateTime.of(2025, 6, 14, 16, 0);
        var view = new AllocationRequestView(id, employeeId, resourceId, "COMPLETED",
                dateFrom, dateTo, "projekt", submittedAt, issuedAt, returnedAt, null);

        var response = AllocationRequestRestMapper.toResponse(view, "Jan", "Test");

        assertThat(response.status()).isEqualTo("COMPLETED");
        assertThat(response.issuedAt()).isEqualTo(issuedAt);
        assertThat(response.returnedAt()).isEqualTo(returnedAt);
        assertThat(response.employeeFirstName()).isEqualTo("Jan");
        assertThat(response.employeeLastName()).isEqualTo("Test");
    }

    @Test
    void toResponse_maps_rejected_request() {
        var view = new AllocationRequestView(id, employeeId, resourceId, "REJECTED",
                dateFrom, dateTo, "test", submittedAt, null, null, "Brak miejsca");

        var response = AllocationRequestRestMapper.toResponse(view);

        assertThat(response.status()).isEqualTo("REJECTED");
        assertThat(response.rejectionReason()).isEqualTo("Brak miejsca");
    }
}
\end{DeskFlowCode}

\phantomsection
\addcontentsline{toc}{subsubsection}{EmployeeProfileAdapterTest}
\DeskFlowListingCaption{allocation/infrastructure/out/employees/EmployeeProfileAdapterTest.java - EmployeeProfileAdapterTest}{lst:allocation-infrastructure-out-employees-EmployeeProfileAdapterTest-java}
\begin{DeskFlowCode}
package com.deskflow.allocation.infrastructure.out.employees;

import com.deskflow.allocation.application.domain.model.EmployeeAllocationProfile;
import com.deskflow.common.application.domain.model.vo.EmployeeId;
import com.deskflow.common.application.domain.model.vo.ResourceType;
import com.deskflow.employees.application.port.in.EmployeeManagementInputPort;
import com.deskflow.employees.application.domain.model.Employee;
import com.deskflow.employees.application.domain.model.vo.Role;
import org.junit.jupiter.api.BeforeEach;
import org.junit.jupiter.api.Test;
import org.junit.jupiter.api.extension.ExtendWith;
import org.mockito.Mock;
import org.mockito.junit.jupiter.MockitoExtension;

import java.util.Map;
import java.util.Optional;
import java.util.UUID;

import static org.assertj.core.api.Assertions.*;
import static org.mockito.Mockito.*;

@ExtendWith(MockitoExtension.class)
class EmployeeProfileAdapterTest {

    @Mock private EmployeeManagementInputPort employeeManagementInputPort;

    private EmployeeProfileAdapter adapter;
    private EmployeeId employeeId;

    @BeforeEach
    void setUp() {
        adapter = new EmployeeProfileAdapter(employeeManagementInputPort);
        employeeId = EmployeeId.of(UUID.randomUUID());
    }

    @Test
    void findEmployee_returns_profile_with_limits() {
        var employee = new Employee(employeeId, "test@company.test", "Jan", "Kowalski", Role.EMPLOYEE,
                Map.of(ResourceType.DESK, 3, ResourceType.LAPTOP, 1));
        when(employeeManagementInputPort.getEmployee(employeeId)).thenReturn(Optional.of(employee));

        var result = adapter.findEmployee(employeeId);

        assertThat(result).isPresent();
        assertThat(result.get().limitsByResourceType()).containsEntry("DESK", 3);
        assertThat(result.get().limitsByResourceType()).containsEntry("LAPTOP", 1);
    }

    @Test
    void findEmployee_returns_empty_when_not_found() {
        when(employeeManagementInputPort.getEmployee(employeeId)).thenReturn(Optional.empty());

        var result = adapter.findEmployee(employeeId);

        assertThat(result).isEmpty();
    }

    @Test
    void findEmployee_empty_limits_when_employee_has_none() {
        var employee = new Employee(employeeId, "test@company.test", "Jan", "Kowalski", Role.EMPLOYEE, Map.of());
        when(employeeManagementInputPort.getEmployee(employeeId)).thenReturn(Optional.of(employee));

        var result = adapter.findEmployee(employeeId);

        assertThat(result).isPresent();
        assertThat(result.get().limitsByResourceType()).isEmpty();
    }
}
\end{DeskFlowCode}

\phantomsection
\addcontentsline{toc}{subsubsection}{AllocationRequestRepositoryDateEdgeCasesTest}
\DeskFlowListingCaption{allocation/infrastructure/out/persistence/AllocationRequestRepositoryDateEdgeCasesTest.java - AllocationRequestRepositoryDateEdgeCasesTest}{lst:allocation-infrastructure-out-persistence-AllocationRequestRepositoryDateEdgeCasesTest-java}
\begin{DeskFlowCode}
package com.deskflow.allocation.infrastructure.out.persistence;

import com.deskflow.common.application.domain.model.vo.EmployeeId;
import com.deskflow.allocation.application.domain.model.AllocationRequest;
import com.deskflow.allocation.application.domain.service.AllocationRequestFactory;
import com.deskflow.allocation.application.port.out.AllocationRequestRepositoryPort;
import com.deskflow.allocation.application.domain.model.vo.DateRange;
import com.deskflow.common.application.domain.model.vo.ResourceId;
import org.junit.jupiter.api.Test;
import org.springframework.beans.factory.annotation.Autowired;
import org.springframework.boot.jdbc.test.autoconfigure.AutoConfigureTestDatabase;
import org.springframework.boot.data.jpa.test.autoconfigure.DataJpaTest;
import org.springframework.context.annotation.Import;
import org.springframework.test.context.ActiveProfiles;

import java.time.LocalDate;
import java.util.List;
import java.util.UUID;

import static org.assertj.core.api.Assertions.assertThat;

@DataJpaTest
@ActiveProfiles("test")
@AutoConfigureTestDatabase(replace = AutoConfigureTestDatabase.Replace.NONE)
@Import({
        com.deskflow.allocation.infrastructure.out.persistence.HibernateAllocationRequestRepositoryAdapter.class,
        com.deskflow.allocation.infrastructure.out.persistence.AllocationRequestMapper.class
})
class AllocationRequestRepositoryDateEdgeCasesIntegrationTest {

    @Autowired
    private AllocationRequestRepositoryPort requestRepository;

    private final AllocationRequestFactory factory = new AllocationRequestFactory();
    private final EmployeeId employeeId = EmployeeId.of(UUID.randomUUID());
    private final ResourceId resourceId = ResourceId.of(UUID.randomUUID());

    @Test
    void overlapping_when_search_starts_on_request_end_is_included() {
        DateRange requestRange = new DateRange(LocalDate.of(2025, 8, 1), LocalDate.of(2025, 8, 10));
        AllocationRequest request = factory.create(employeeId, resourceId, requestRange, "test");
        requestRepository.save(request);

        DateRange searchRange = new DateRange(LocalDate.of(2025, 8, 10), LocalDate.of(2025, 8, 12));
        List<AllocationRequest> results = requestRepository.findActiveForResource(resourceId, searchRange);

        assertThat(results).anyMatch(r -> r.getId().equals(request.getId()));
    }

    @Test
    void no_overlap_when_request_ends_before_search_starts_is_not_included() {
        DateRange requestRange = new DateRange(LocalDate.of(2025, 8, 1), LocalDate.of(2025, 8, 10));
        AllocationRequest request = factory.create(employeeId, resourceId, requestRange, "test");
        requestRepository.save(request);

        DateRange searchRange = new DateRange(LocalDate.of(2025, 8, 11), LocalDate.of(2025, 8, 12));
        List<AllocationRequest> results = requestRepository.findActiveForResource(resourceId, searchRange);

        assertThat(results).noneMatch(r -> r.getId().equals(request.getId()));
    }

    @Test
    void overlapping_when_request_inside_search_is_included() {
        DateRange requestRange = new DateRange(LocalDate.of(2025, 8, 5), LocalDate.of(2025, 8, 6));
        AllocationRequest request = factory.create(employeeId, resourceId, requestRange, "test");
        requestRepository.save(request);

        DateRange searchRange = new DateRange(LocalDate.of(2025, 8, 1), LocalDate.of(2025, 8, 10));
        List<AllocationRequest> results = requestRepository.findActiveForResource(resourceId, searchRange);

        assertThat(results).anyMatch(r -> r.getId().equals(request.getId()));
    }
}
\end{DeskFlowCode}

\phantomsection
\addcontentsline{toc}{subsubsection}{AllocationRequestRepositoryIntegrationTest}
\DeskFlowListingCaption{allocation/infrastructure/out/persistence/AllocationRequestRepositoryIntegrationTest.java - AllocationRequestRepositoryIntegrationTest}{lst:allocation-infrastructure-out-persistence-AllocationRequestRepositoryIntegrationTest-java}
\begin{DeskFlowCode}
package com.deskflow.allocation.infrastructure.out.persistence;

import com.deskflow.common.application.domain.model.vo.EmployeeId;
import com.deskflow.allocation.application.domain.model.AllocationRequest;
import com.deskflow.allocation.application.domain.service.AllocationRequestFactory;
import com.deskflow.allocation.application.port.out.AllocationRequestRepositoryPort;
import com.deskflow.allocation.application.domain.model.vo.AllocationRequestId;
import com.deskflow.allocation.application.domain.model.vo.DateRange;
import com.deskflow.allocation.application.domain.model.vo.RequestStatus;
import com.deskflow.common.application.domain.model.vo.ResourceId;
import org.junit.jupiter.api.Test;
import org.springframework.beans.factory.annotation.Autowired;
import org.springframework.boot.jdbc.test.autoconfigure.AutoConfigureTestDatabase;
import org.springframework.boot.data.jpa.test.autoconfigure.DataJpaTest;
import org.springframework.context.annotation.Import;
import org.springframework.test.context.ActiveProfiles;

import java.time.LocalDate;
import java.time.LocalDateTime;
import java.util.List;
import java.util.Optional;
import java.util.UUID;

import static org.assertj.core.api.Assertions.assertThat;

@DataJpaTest
@ActiveProfiles("test")
@AutoConfigureTestDatabase(replace = AutoConfigureTestDatabase.Replace.NONE)
@Import({
        com.deskflow.allocation.infrastructure.out.persistence.HibernateAllocationRequestRepositoryAdapter.class,
        com.deskflow.allocation.infrastructure.out.persistence.AllocationRequestMapper.class
})
class AllocationRequestRepositoryIntegrationTest {

    @Autowired
    private AllocationRequestRepositoryPort requestRepository;

    private final AllocationRequestFactory factory = new AllocationRequestFactory();
    private final EmployeeId employeeId = EmployeeId.of(UUID.randomUUID());
    private final ResourceId resourceId = ResourceId.of(UUID.randomUUID());

    @Test
    void save_and_find_request() {
        DateRange dateRange = new DateRange(LocalDate.of(2025, 7, 1), LocalDate.of(2025, 7, 5));
        AllocationRequest request = factory.create(employeeId, resourceId, dateRange, "test");

        requestRepository.save(request);

        Optional<AllocationRequest> found = requestRepository.find(request.getId());
        assertThat(found).isPresent();
        assertThat(found.get().getStatus()).isEqualTo(RequestStatus.SUBMITTED);
        assertThat(found.get().getEmployeeId()).isEqualTo(employeeId);
    }

    @Test
    void find_overdue_returns_active_past_due_requests() {
        DateRange pastRange = new DateRange(LocalDate.of(2024, 1, 1), LocalDate.of(2024, 1, 10));
        AllocationRequest request = AllocationRequest.reconstitute(
                AllocationRequestId.newId(),
                employeeId, resourceId, pastRange, "test",
                LocalDateTime.of(2024, 1, 1, 9, 0),
                RequestStatus.ACTIVE, LocalDateTime.of(2024, 1, 1, 10, 0), null, null
        );
        requestRepository.save(request);

        List<AllocationRequest> overdue = requestRepository.findOverdue(LocalDateTime.now());

        assertThat(overdue).anyMatch(r -> r.getId().equals(request.getId()));
    }

    @Test
    void find_active_for_resource_with_overlapping_date_range() {
        DateRange dateRange = new DateRange(LocalDate.of(2025, 8, 1), LocalDate.of(2025, 8, 10));
        AllocationRequest request = factory.create(employeeId, resourceId, dateRange, "test");
        requestRepository.save(request);

        DateRange searchRange = new DateRange(LocalDate.of(2025, 8, 5), LocalDate.of(2025, 8, 15));
        List<AllocationRequest> active = requestRepository.findActiveForResource(resourceId, searchRange);

        assertThat(active).anyMatch(r -> r.getId().equals(request.getId()));
    }
}
\end{DeskFlowCode}

\phantomsection
\addcontentsline{toc}{subsubsection}{ResourceCatalogAdapterTest}
\DeskFlowListingCaption{allocation/infrastructure/out/resources/ResourceCatalogAdapterTest.java - ResourceCatalogAdapterTest}{lst:allocation-infrastructure-out-resources-ResourceCatalogAdapterTest-java}
\begin{DeskFlowCode}
package com.deskflow.allocation.infrastructure.out.resources;

import com.deskflow.allocation.application.domain.model.ResourceAllocationProfile;
import com.deskflow.common.application.domain.model.vo.ResourceId;
import com.deskflow.resources.application.model.ResourceView;
import com.deskflow.resources.application.port.in.ResourceManagementInputPort;
import org.junit.jupiter.api.BeforeEach;
import org.junit.jupiter.api.Test;
import org.junit.jupiter.api.extension.ExtendWith;
import org.mockito.Mock;
import org.mockito.junit.jupiter.MockitoExtension;

import java.util.List;
import java.util.Optional;
import java.util.UUID;

import static org.assertj.core.api.Assertions.*;
import static org.mockito.Mockito.*;

@ExtendWith(MockitoExtension.class)
class ResourceCatalogAdapterTest {

    @Mock private ResourceManagementInputPort resourceManagementInputPort;

    private ResourceCatalogAdapter adapter;
    private ResourceId resourceId;

    @BeforeEach
    void setUp() {
        adapter = new ResourceCatalogAdapter(resourceManagementInputPort);
        resourceId = ResourceId.of(UUID.randomUUID());
    }

    @Test
    void findResource_returns_profile_when_found() {
        var view = new ResourceView(resourceId.getValue(), "DESK", "Biurko A", "A", 1, "101", true, true);
        when(resourceManagementInputPort.getResource(resourceId)).thenReturn(Optional.of(view));

        var result = adapter.findResource(resourceId);

        assertThat(result).isPresent();
        assertThat(result.get().id()).isEqualTo(resourceId);
        assertThat(result.get().type()).isEqualTo("DESK");
        assertThat(result.get().name()).isEqualTo("Biurko A");
        assertThat(result.get().building()).isEqualTo("A");
        assertThat(result.get().floor()).isEqualTo(1);
        assertThat(result.get().number()).isEqualTo("101");
        assertThat(result.get().active()).isTrue();
        assertThat(result.get().available()).isTrue();
    }

    @Test
    void findResource_returns_empty_when_not_found() {
        when(resourceManagementInputPort.getResource(resourceId)).thenReturn(Optional.empty());

        var result = adapter.findResource(resourceId);

        assertThat(result).isEmpty();
    }

    @Test
    void findActiveResources_returns_all_active() {
        var view1 = new ResourceView(UUID.randomUUID(), "DESK", "Biurko A", "A", 1, "101", true, true);
        var view2 = new ResourceView(UUID.randomUUID(), "LAPTOP", "Laptop X", "B", 2, "201", true, false);
        when(resourceManagementInputPort.getActiveResources()).thenReturn(List.of(view1, view2));

        var profiles = adapter.findActiveResources();

        assertThat(profiles).hasSize(2);
        assertThat(profiles).extracting(ResourceAllocationProfile::active).containsOnly(true);
    }

    @Test
    void findActiveResources_returns_empty_when_none() {
        when(resourceManagementInputPort.getActiveResources()).thenReturn(List.of());

        var profiles = adapter.findActiveResources();

        assertThat(profiles).isEmpty();
    }
}
\end{DeskFlowCode}

\subsection{Testy - Resources}

\phantomsection
\addcontentsline{toc}{subsubsection}{LocationTest}
\DeskFlowListingCaption{resources/application/domain/model/LocationTest.java - LocationTest}{lst:resources-application-domain-model-LocationTest-java}
\begin{DeskFlowCode}
package com.deskflow.resources.application.domain.model;

import com.deskflow.resources.application.domain.model.vo.Location;
import org.junit.jupiter.api.Test;

import static org.assertj.core.api.Assertions.*;

class LocationTest {

    @Test
    void creates_location() {
        Location l = new Location("Building A", 2, "101");
        assertThat(l.getBuilding()).isEqualTo("Building A");
        assertThat(l.getFloor()).isEqualTo(2);
        assertThat(l.getNumber()).isEqualTo("101");
    }

    @Test
    void throws_when_building_blank() {
        assertThatThrownBy(() -> new Location("", 1, "1"))
                .isInstanceOf(IllegalArgumentException.class);
    }

    @Test
    void throws_when_number_blank() {
        assertThatThrownBy(() -> new Location("Building A", 1, ""))
                .isInstanceOf(IllegalArgumentException.class);
    }

    @Test
    void value_object_equality() {
        Location l1 = new Location("A", 1, "101");
        Location l2 = new Location("A", 1, "101");
        assertThat(l1).isEqualTo(l2);
        assertThat(l1.hashCode()).isEqualTo(l2.hashCode());
    }

    @Test
    void not_equal_when_different_fields() {
        assertThat(new Location("A", 1, "101")).isNotEqualTo(new Location("B", 1, "101"));
        assertThat(new Location("A", 1, "101")).isNotEqualTo(new Location("A", 2, "101"));
        assertThat(new Location("A", 1, "101")).isNotEqualTo(new Location("A", 1, "102"));
    }
}
\end{DeskFlowCode}

\phantomsection
\addcontentsline{toc}{subsubsection}{OfficeResourceTest}
\DeskFlowListingCaption{resources/application/domain/model/OfficeResourceTest.java - OfficeResourceTest}{lst:resources-application-domain-model-OfficeResourceTest-java}
\begin{DeskFlowCode}
package com.deskflow.resources.application.domain.model;

import com.deskflow.resources.application.domain.service.OfficeResourceFactory;
import com.deskflow.resources.application.domain.model.vo.Location;
import com.deskflow.common.application.domain.model.vo.ResourceId;
import com.deskflow.common.application.domain.model.vo.ResourceType;
import org.junit.jupiter.api.BeforeEach;
import org.junit.jupiter.api.Test;

import java.util.UUID;

import static org.assertj.core.api.Assertions.*;

class OfficeResourceTest {

    private OfficeResourceFactory factory;
    private static final Location LOC = new Location("A", 1, "101");

    @BeforeEach
    void setUp() {
        factory = new OfficeResourceFactory();
    }

    @Test
    void new_resource_is_active_and_available() {
        OfficeResource r = factory.create(ResourceType.DESK, "Desk 1", LOC);
        assertThat(r.isActive()).isTrue();
        assertThat(r.isAvailable()).isTrue();
    }

    @Test
    void throws_when_name_blank() {
        assertThatThrownBy(() -> factory.create(ResourceType.DESK, "", LOC))
                .isInstanceOf(IllegalArgumentException.class);
    }

    @Test
    void mark_as_unavailable() {
        OfficeResource r = factory.create(ResourceType.DESK, "Desk 1", LOC);
        r.markAsUnavailable();
        assertThat(r.isAvailable()).isFalse();
        assertThat(r.isActive()).isTrue();
    }

    @Test
    void mark_as_available_after_unavailable() {
        OfficeResource r = factory.create(ResourceType.DESK, "Desk 1", LOC);
        r.markAsUnavailable();
        r.markAsAvailable();
        assertThat(r.isAvailable()).isTrue();
    }

    @Test
    void deactivate_resource() {
        OfficeResource r = factory.create(ResourceType.DESK, "Desk 1", LOC);
        r.deactivate();
        assertThat(r.isActive()).isFalse();
        assertThat(r.isAvailable()).isFalse();
    }

    @Test
    void deactivate_twice_throws() {
        OfficeResource r = factory.create(ResourceType.DESK, "Desk 1", LOC);
        r.deactivate();
        assertThatThrownBy(r::deactivate).isInstanceOf(IllegalStateException.class);
    }

    @Test
    void update_changes_name_and_location() {
        OfficeResource r = factory.create(ResourceType.DESK, "Old name", LOC);
        Location newLoc = new Location("B", 3, "305");
        r.update("New name", newLoc);
        assertThat(r.getName()).isEqualTo("New name");
        assertThat(r.getLocation()).isEqualTo(newLoc);
    }

    @Test
    void cannot_mark_unavailable_on_deactivated_resource() {
        OfficeResource r = factory.create(ResourceType.DESK, "Desk", LOC);
        r.deactivate();
        assertThatThrownBy(r::markAsUnavailable).isInstanceOf(IllegalStateException.class);
    }

    @Test
    void resource_has_correct_type_and_location() {
        OfficeResource r = factory.create(ResourceType.CONFERENCE_ROOM, "Room A", LOC);
        assertThat(r.getType()).isEqualTo(ResourceType.CONFERENCE_ROOM);
        assertThat(r.getLocation()).isEqualTo(LOC);
    }

    @Test
    void reconstituted_resource_preserves_state() {
        OfficeResource r = OfficeResource.reconstitute(
                ResourceId.of(UUID.randomUUID()), ResourceType.LAPTOP, "Laptop X",
                LOC, false, false
        );
        assertThat(r.isActive()).isFalse();
        assertThat(r.isAvailable()).isFalse();
    }
}
\end{DeskFlowCode}

\phantomsection
\addcontentsline{toc}{subsubsection}{ResourceDeactivationDomainServicesTest}
\DeskFlowListingCaption{resources/application/domain/model/ResourceDeactivationDomainServicesTest.java - ResourceDeactivationDomainServicesTest}{lst:resources-application-domain-model-ResourceDeactivationDomainServicesTest-java}
\begin{DeskFlowCode}
package com.deskflow.resources.application.domain.model;

import com.deskflow.common.application.domain.model.vo.ResourceType;
import com.deskflow.resources.application.domain.model.vo.Location;
import com.deskflow.resources.application.domain.model.vo.AllocationUsageStatus;
import com.deskflow.resources.application.domain.service.DeactivateOfficeResourceDomainService;
import com.deskflow.resources.application.domain.service.OfficeResourceFactory;
import org.junit.jupiter.api.Test;
import org.junit.jupiter.params.ParameterizedTest;
import org.junit.jupiter.params.provider.EnumSource;

import java.util.List;

import static org.assertj.core.api.Assertions.assertThat;
import static org.assertj.core.api.Assertions.assertThatThrownBy;

class ResourceDeactivationDomainServicesTest {

    private final DeactivateOfficeResourceDomainService deactivateOfficeResourceDomainService =
            new DeactivateOfficeResourceDomainService();

    @ParameterizedTest
    @EnumSource(value = AllocationUsageStatus.class, names = {"IN_USE", "APPROVED", "REQUESTED", "OVERDUE"})
    void deactivate_throws_when_resource_has_blocking_allocation_status(AllocationUsageStatus status) {
        OfficeResource resource = new OfficeResourceFactory().create(ResourceType.DESK, "Desk", new Location("A", 1, "101"));

        assertThatThrownBy(() -> deactivateOfficeResourceDomainService.deactivateOfficeResource(resource, List.of(status)))
                .isInstanceOf(IllegalStateException.class)
                .hasMessageContaining("active allocation requests");
    }

    @Test
    void deactivate_resource_when_no_related_allocations() {
        OfficeResource resource = new OfficeResourceFactory().create(ResourceType.DESK, "Desk", new Location("A", 1, "101"));

        deactivateOfficeResourceDomainService.deactivateOfficeResource(resource, List.of());

        assertThat(resource.isActive()).isFalse();
        assertThat(resource.isAvailable()).isFalse();
    }
}
\end{DeskFlowCode}

\phantomsection
\addcontentsline{toc}{subsubsection}{ResourceLifecycleDomainServicesTest}
\DeskFlowListingCaption{resources/application/domain/model/ResourceLifecycleDomainServicesTest.java - ResourceLifecycleDomainServicesTest}{lst:resources-application-domain-model-ResourceLifecycleDomainServicesTest-java}
\begin{DeskFlowCode}
package com.deskflow.resources.application.domain.model;

import com.deskflow.common.application.domain.model.vo.ResourceType;
import com.deskflow.resources.application.domain.model.events.ResourceMadeAvailable;
import com.deskflow.resources.application.domain.model.events.ResourceMadeUnavailable;
import com.deskflow.resources.application.domain.model.events.ResourceUpdated;
import com.deskflow.resources.application.domain.model.vo.Location;
import com.deskflow.resources.application.domain.service.MarkOfficeResourceAvailableDomainService;
import com.deskflow.resources.application.domain.service.MarkOfficeResourceUnavailableDomainService;
import com.deskflow.resources.application.domain.service.OfficeResourceFactory;
import com.deskflow.resources.application.domain.service.UpdateOfficeResourceDomainService;
import org.junit.jupiter.api.Test;

import static org.assertj.core.api.Assertions.assertThat;

class ResourceLifecycleDomainServicesTest {

    private final UpdateOfficeResourceDomainService updateService = new UpdateOfficeResourceDomainService();
    private final MarkOfficeResourceUnavailableDomainService unavailableService = new MarkOfficeResourceUnavailableDomainService();
    private final MarkOfficeResourceAvailableDomainService availableService = new MarkOfficeResourceAvailableDomainService();

    @Test
    void update_service_updates_name_and_location_and_records_event() {
        OfficeResource resource = resource();
        Location newLocation = new Location("B", 2, "202");

        updateService.updateOfficeResource(resource, "New desk", newLocation);

        assertThat(resource.getName()).isEqualTo("New desk");
        assertThat(resource.getLocation()).isEqualTo(newLocation);
        assertThat(resource.flushEvents()).hasAtLeastOneElementOfType(ResourceUpdated.class);
    }

    @Test
    void mark_unavailable_service_marks_resource_unavailable_and_records_event() {
        OfficeResource resource = resource();

        unavailableService.markOfficeResourceUnavailable(resource);

        assertThat(resource.isAvailable()).isFalse();
        assertThat(resource.flushEvents()).hasAtLeastOneElementOfType(ResourceMadeUnavailable.class);
    }

    @Test
    void mark_available_service_marks_resource_available_and_records_event() {
        OfficeResource resource = resource();
        resource.markAsUnavailable();
        resource.flushEvents();

        availableService.markOfficeResourceAvailable(resource);

        assertThat(resource.isAvailable()).isTrue();
        assertThat(resource.flushEvents()).hasAtLeastOneElementOfType(ResourceMadeAvailable.class);
    }

    private OfficeResource resource() {
        return new OfficeResourceFactory().create(ResourceType.DESK, "Desk", new Location("A", 1, "101"));
    }
}
\end{DeskFlowCode}

\phantomsection
\addcontentsline{toc}{subsubsection}{ResourceAppServiceMissingResourceTest}
\DeskFlowListingCaption{resources/application/service/ResourceAppServiceMissingResourceTest.java - ResourceAppServiceMissingResourceTest}{lst:resources-application-service-ResourceAppServiceMissingResourceTest-java}
\begin{DeskFlowCode}
package com.deskflow.resources.application.service;

import com.deskflow.common.application.port.out.DomainEventPublisherPort;
import com.deskflow.resources.application.port.out.OfficeResourceRepositoryPort;
import com.deskflow.resources.application.port.out.ResourceAllocationRequestPort;
import com.deskflow.resources.application.domain.service.OfficeResourceFactory;
import com.deskflow.resources.application.domain.service.DeactivateOfficeResourceDomainService;
import com.deskflow.resources.application.domain.service.MarkOfficeResourceAvailableDomainService;
import com.deskflow.resources.application.domain.service.MarkOfficeResourceUnavailableDomainService;
import com.deskflow.resources.application.domain.service.UpdateOfficeResourceDomainService;
import com.deskflow.common.application.domain.model.vo.ResourceId;
import org.junit.jupiter.api.BeforeEach;
import org.junit.jupiter.api.Test;
import org.junit.jupiter.api.extension.ExtendWith;
import org.mockito.Mock;
import org.mockito.junit.jupiter.MockitoExtension;

import java.util.NoSuchElementException;
import java.util.Optional;

import static org.assertj.core.api.Assertions.assertThatThrownBy;
import static org.mockito.Mockito.when;

@ExtendWith(MockitoExtension.class)
class ResourceAppServiceMissingResourceTest {

    @Mock private OfficeResourceRepositoryPort resourceRepository;
    @Mock private ResourceAllocationRequestPort allocationRequestPort;
    @Mock private DomainEventPublisherPort eventPublisher;

    private ResourceAppService appService;

    @BeforeEach
    void setUp() {
        appService = new ResourceAppService(resourceRepository, allocationRequestPort, new OfficeResourceFactory(),
                new UpdateOfficeResourceDomainService(), new MarkOfficeResourceUnavailableDomainService(),
                new MarkOfficeResourceAvailableDomainService(),
                new DeactivateOfficeResourceDomainService(),
                eventPublisher);
    }

    @Test
    void update_throws_when_resource_missing() {
        ResourceId resourceId = ResourceId.of(java.util.UUID.randomUUID());
        when(resourceRepository.find(resourceId)).thenReturn(Optional.empty());

        assertThatThrownBy(() -> appService.updateResource(resourceId, "Desk", "A", 1, "101"))
                .isInstanceOf(NoSuchElementException.class)
                .hasMessageContaining("Resource not found");
    }
}
\end{DeskFlowCode}

\phantomsection
\addcontentsline{toc}{subsubsection}{ResourceAppServiceTest}
\DeskFlowListingCaption{resources/application/service/ResourceAppServiceTest.java - ResourceAppServiceTest}{lst:resources-application-service-ResourceAppServiceTest-java}
\begin{DeskFlowCode}
package com.deskflow.resources.application.service;

import com.deskflow.common.application.port.out.DomainEventPublisherPort;
import com.deskflow.common.application.domain.model.vo.ResourceId;
import com.deskflow.common.application.domain.model.vo.ResourceType;
import com.deskflow.resources.application.port.out.OfficeResourceRepositoryPort;
import com.deskflow.resources.application.port.out.ResourceAllocationRequestPort;
import com.deskflow.resources.application.domain.model.OfficeResource;
import com.deskflow.resources.application.domain.model.events.ResourceDeactivated;
import com.deskflow.resources.application.domain.model.events.ResourceMadeUnavailable;
import com.deskflow.resources.application.domain.model.events.ResourceUpdated;
import com.deskflow.resources.application.domain.model.vo.Location;
import com.deskflow.resources.application.domain.service.DeactivateOfficeResourceDomainService;
import com.deskflow.resources.application.domain.service.MarkOfficeResourceAvailableDomainService;
import com.deskflow.resources.application.domain.service.MarkOfficeResourceUnavailableDomainService;
import com.deskflow.resources.application.domain.service.OfficeResourceFactory;
import com.deskflow.resources.application.domain.service.UpdateOfficeResourceDomainService;
import org.junit.jupiter.api.BeforeEach;
import org.junit.jupiter.api.Test;
import org.junit.jupiter.api.extension.ExtendWith;
import org.mockito.ArgumentCaptor;
import org.mockito.Mock;
import org.mockito.junit.jupiter.MockitoExtension;

import java.util.List;
import java.util.NoSuchElementException;
import java.util.Optional;
import java.util.UUID;

import static org.assertj.core.api.Assertions.assertThat;
import static org.assertj.core.api.Assertions.assertThatThrownBy;
import static org.mockito.Mockito.any;
import static org.mockito.Mockito.verify;
import static org.mockito.Mockito.when;

@ExtendWith(MockitoExtension.class)
class ResourceAppServiceTest {

    @Mock private OfficeResourceRepositoryPort resourceRepository;
    @Mock private ResourceAllocationRequestPort allocationRequestPort;
    @Mock private DomainEventPublisherPort eventPublisher;

    private ResourceAppService appService;

    @BeforeEach
    void setUp() {
        appService = new ResourceAppService(resourceRepository, allocationRequestPort, new OfficeResourceFactory(),
                new UpdateOfficeResourceDomainService(), new MarkOfficeResourceUnavailableDomainService(),
                new MarkOfficeResourceAvailableDomainService(),
                new DeactivateOfficeResourceDomainService(),
                eventPublisher);
    }

    @Test
    void add_resource_creates_active_available_resource_and_saves() {
        ArgumentCaptor<OfficeResource> resourceCaptor = ArgumentCaptor.forClass(OfficeResource.class);

        ResourceId id = appService.addResource(ResourceType.DESK, "Desk 1", "A", 1, "101");

        verify(resourceRepository).save(resourceCaptor.capture());
        OfficeResource saved = resourceCaptor.getValue();
        assertThat(id).isEqualTo(saved.getId());
        assertThat(saved.getType()).isEqualTo(ResourceType.DESK);
        assertThat(saved.getName()).isEqualTo("Desk 1");
        assertThat(saved.isActive()).isTrue();
        assertThat(saved.isAvailable()).isTrue();
    }

    @Test
    void update_resource_saves_and_publishes_update_event() {
        OfficeResource resource = resource();
        when(resourceRepository.find(resource.getId())).thenReturn(Optional.of(resource));

        appService.updateResource(resource.getId(), "New name", "B", 2, "202");

        verify(resourceRepository).save(resource);
        verify(eventPublisher).publish(any(ResourceUpdated.class));
        assertThat(resource.getName()).isEqualTo("New name");
        assertThat(resource.getLocation()).isEqualTo(new Location("B", 2, "202"));
    }

    @Test
    void mark_unavailable_saves_and_publishes_event() {
        OfficeResource resource = resource();
        when(resourceRepository.find(resource.getId())).thenReturn(Optional.of(resource));

        appService.markUnavailable(resource.getId());

        verify(resourceRepository).save(resource);
        verify(eventPublisher).publish(any(ResourceMadeUnavailable.class));
        assertThat(resource.isAvailable()).isFalse();
    }

    @Test
    void deactivate_resource_checks_related_allocations_saves_and_publishes_event() {
        OfficeResource resource = resource();
        when(resourceRepository.find(resource.getId())).thenReturn(Optional.of(resource));
        when(allocationRequestPort.findStatusesForResource(resource.getId())).thenReturn(List.of());

        appService.deactivateResource(resource.getId());

        verify(allocationRequestPort).findStatusesForResource(resource.getId());
        verify(resourceRepository).save(resource);
        verify(eventPublisher).publish(any(ResourceDeactivated.class));
        assertThat(resource.isActive()).isFalse();
    }

    @Test
    void mark_unavailable_throws_when_resource_missing() {
        ResourceId resourceId = ResourceId.of(UUID.randomUUID());
        when(resourceRepository.find(resourceId)).thenReturn(Optional.empty());

        assertThatThrownBy(() -> appService.markUnavailable(resourceId))
                .isInstanceOf(NoSuchElementException.class)
                .hasMessageContaining("Resource not found");
    }

    private OfficeResource resource() {
        return new OfficeResourceFactory().create(ResourceType.DESK, "Desk", new Location("A", 1, "101"));
    }
}
\end{DeskFlowCode}

\phantomsection
\addcontentsline{toc}{subsubsection}{ResourceAvailabilityEventListenerTest}
\DeskFlowListingCaption{resources/infrastructure/in/events/ResourceAvailabilityEventListenerTest.java - ResourceAvailabilityEventListenerTest}{lst:resources-infrastructure-in-events-ResourceAvailabilityEventListenerTest-java}
\begin{DeskFlowCode}
package com.deskflow.resources.infrastructure.in.events;

import com.deskflow.allocation.application.domain.model.events.ResourceIssued;
import com.deskflow.allocation.application.domain.model.events.ResourceReturned;
import com.deskflow.allocation.application.domain.model.vo.AllocationRequestId;
import com.deskflow.common.application.domain.model.vo.EmployeeId;
import com.deskflow.common.application.domain.model.vo.ResourceId;
import com.deskflow.resources.application.port.in.ResourceManagementInputPort;
import org.junit.jupiter.api.BeforeEach;
import org.junit.jupiter.api.Test;
import org.junit.jupiter.api.extension.ExtendWith;
import org.mockito.Mock;
import org.mockito.junit.jupiter.MockitoExtension;

import java.time.LocalDateTime;
import java.util.UUID;

import static org.mockito.Mockito.*;

@ExtendWith(MockitoExtension.class)
class ResourceAvailabilityEventListenerTest {

    @Mock private ResourceManagementInputPort resourceManagementInputPort;

    private ResourceAvailabilityEventListener listener;
    private ResourceId resourceId;

    @BeforeEach
    void setUp() {
        listener = new ResourceAvailabilityEventListener(resourceManagementInputPort);
        resourceId = ResourceId.of(UUID.randomUUID());
    }

    @Test
    void onResourceIssued_marks_resource_unavailable() {
        var event = new ResourceIssued(AllocationRequestId.newId(), EmployeeId.of(UUID.randomUUID()), resourceId, LocalDateTime.now());

        listener.onResourceIssued(event);

        verify(resourceManagementInputPort).markUnavailable(resourceId);
    }

    @Test
    void onResourceReturned_marks_resource_available() {
        var event = new ResourceReturned(AllocationRequestId.newId(), EmployeeId.of(UUID.randomUUID()), resourceId, LocalDateTime.now());

        listener.onResourceReturned(event);

        verify(resourceManagementInputPort).markAvailable(resourceId);
    }
}
\end{DeskFlowCode}

\phantomsection
\addcontentsline{toc}{subsubsection}{AllocationRequestAdapterTest}
\DeskFlowListingCaption{resources/infrastructure/out/allocation/AllocationRequestAdapterTest.java - AllocationRequestAdapterTest}{lst:resources-infrastructure-out-allocation-AllocationRequestAdapterTest-java}
\begin{DeskFlowCode}
package com.deskflow.resources.infrastructure.out.allocation;

import com.deskflow.allocation.application.model.AllocationRequestView;
import com.deskflow.allocation.application.port.in.AllocationRequestInputPort;
import com.deskflow.common.application.domain.model.vo.ResourceId;
import com.deskflow.resources.application.domain.model.vo.AllocationUsageStatus;
import org.junit.jupiter.api.BeforeEach;
import org.junit.jupiter.api.Test;
import org.junit.jupiter.api.extension.ExtendWith;
import org.mockito.Mock;
import org.mockito.junit.jupiter.MockitoExtension;

import java.time.LocalDate;
import java.time.LocalDateTime;
import java.util.List;
import java.util.UUID;

import static org.assertj.core.api.Assertions.*;
import static org.mockito.Mockito.*;

@ExtendWith(MockitoExtension.class)
class AllocationRequestAdapterTest {

    @Mock private AllocationRequestInputPort allocationRequestInputPort;

    private AllocationRequestAdapter adapter;
    private ResourceId resourceId;

    @BeforeEach
    void setUp() {
        adapter = new AllocationRequestAdapter(allocationRequestInputPort);
        resourceId = ResourceId.of(UUID.randomUUID());
    }

    private AllocationRequestView viewWithStatus(String status) {
        return new AllocationRequestView(UUID.randomUUID(), UUID.randomUUID(), resourceId.getValue(),
                status, LocalDate.now(), LocalDate.now().plusDays(1), "test",
                LocalDateTime.now(), null, null, null);
    }

    @Test
    void submitted_maps_to_requested() {
        when(allocationRequestInputPort.findActiveForResource(resourceId))
                .thenReturn(List.of(viewWithStatus("SUBMITTED")));

        var statuses = adapter.findStatusesForResource(resourceId);

        assertThat(statuses).containsExactly(AllocationUsageStatus.REQUESTED);
    }

    @Test
    void approved_maps_to_approved() {
        when(allocationRequestInputPort.findActiveForResource(resourceId))
                .thenReturn(List.of(viewWithStatus("APPROVED")));

        var statuses = adapter.findStatusesForResource(resourceId);

        assertThat(statuses).containsExactly(AllocationUsageStatus.APPROVED);
    }

    @Test
    void active_maps_to_in_use() {
        when(allocationRequestInputPort.findActiveForResource(resourceId))
                .thenReturn(List.of(viewWithStatus("ACTIVE")));

        var statuses = adapter.findStatusesForResource(resourceId);

        assertThat(statuses).containsExactly(AllocationUsageStatus.IN_USE);
    }

    @Test
    void overdue_maps_to_overdue() {
        when(allocationRequestInputPort.findActiveForResource(resourceId))
                .thenReturn(List.of(viewWithStatus("OVERDUE")));

        var statuses = adapter.findStatusesForResource(resourceId);

        assertThat(statuses).containsExactly(AllocationUsageStatus.OVERDUE);
    }

    @Test
    void rejected_cancelled_completed_map_to_closed() {
        when(allocationRequestInputPort.findActiveForResource(resourceId))
                .thenReturn(List.of(
                        viewWithStatus("REJECTED"),
                        viewWithStatus("CANCELLED"),
                        viewWithStatus("COMPLETED")
                ));

        var statuses = adapter.findStatusesForResource(resourceId);

        assertThat(statuses).allMatch(s -> s == AllocationUsageStatus.CLOSED);
    }

    @Test
    void unknown_status_throws() {
        when(allocationRequestInputPort.findActiveForResource(resourceId))
                .thenReturn(List.of(viewWithStatus("UNKNOWN")));

        assertThatThrownBy(() -> adapter.findStatusesForResource(resourceId))
                .isInstanceOf(IllegalArgumentException.class)
                .hasMessageContaining("Unknown allocation status");
    }

    @Test
    void returns_multiple_statuses() {
        when(allocationRequestInputPort.findActiveForResource(resourceId))
                .thenReturn(List.of(
                        viewWithStatus("APPROVED"),
                        viewWithStatus("ACTIVE")
                ));

        var statuses = adapter.findStatusesForResource(resourceId);

        assertThat(statuses).containsExactly(AllocationUsageStatus.APPROVED, AllocationUsageStatus.IN_USE);
    }
}
\end{DeskFlowCode}

\phantomsection
\addcontentsline{toc}{subsubsection}{OfficeResourceRepositoryIntegrationTest}
\DeskFlowListingCaption{resources/infrastructure/out/persistence/OfficeResourceRepositoryIntegrationTest.java - OfficeResourceRepositoryIntegrationTest}{lst:resources-infrastructure-out-persistence-OfficeResourceRepositoryIntegrationTest-java}
\begin{DeskFlowCode}
package com.deskflow.resources.infrastructure.out.persistence;

import com.deskflow.resources.application.domain.model.OfficeResource;
import com.deskflow.resources.application.domain.service.OfficeResourceFactory;
import com.deskflow.resources.application.port.out.OfficeResourceRepositoryPort;
import com.deskflow.resources.application.domain.model.vo.Location;
import com.deskflow.common.application.domain.model.vo.ResourceId;
import com.deskflow.common.application.domain.model.vo.ResourceType;
import org.junit.jupiter.api.Test;
import org.springframework.beans.factory.annotation.Autowired;
import org.springframework.boot.jdbc.test.autoconfigure.AutoConfigureTestDatabase;
import org.springframework.boot.data.jpa.test.autoconfigure.DataJpaTest;
import org.springframework.context.annotation.Import;
import org.springframework.test.context.ActiveProfiles;

import java.util.List;
import java.util.UUID;

import static org.assertj.core.api.Assertions.assertThat;

@DataJpaTest
@ActiveProfiles("test")
@AutoConfigureTestDatabase(replace = AutoConfigureTestDatabase.Replace.NONE)
@Import({
        com.deskflow.resources.infrastructure.out.persistence.HibernateOfficeResourceRepositoryAdapter.class,
        com.deskflow.resources.infrastructure.out.persistence.OfficeResourceMapper.class
})
class OfficeResourceRepositoryIntegrationTest {

    @Autowired
    private OfficeResourceRepositoryPort resourceRepository;

    private final OfficeResourceFactory factory = new OfficeResourceFactory();

    @Test
    void save_and_find_resource() {
        OfficeResource resource = factory.create(ResourceType.DESK, "Desk A1", new Location("A", 1, "101"));

        resourceRepository.save(resource);

        ResourceId id = resource.getId();
        OfficeResource found = resourceRepository.find(id).orElseThrow();
        assertThat(found.getId()).isEqualTo(id);
        assertThat(found.getType()).isEqualTo(ResourceType.DESK);
        assertThat(found.getLocation()).isEqualTo(resource.getLocation());
    }

    @Test
    void find_active_returns_only_active_resources() {
        OfficeResource active = factory.create(ResourceType.DESK, "Desk A2", new Location("A", 1, "102"));
        OfficeResource inactive = factory.create(ResourceType.DESK, "Desk A3", new Location("A", 1, "103"));
        inactive.deactivate();

        resourceRepository.save(active);
        resourceRepository.save(inactive);

        List<OfficeResource> activeResources = resourceRepository.findActive();

        assertThat(activeResources).anyMatch(r -> r.getId().equals(active.getId()));
        assertThat(activeResources).noneMatch(r -> r.getId().equals(inactive.getId()));
    }

    @Test
    void find_active_returns_empty_when_no_active_resources() {
        OfficeResource inactive = factory.create(ResourceType.DESK, "Desk A4", new Location("A", 1, "104"));
        inactive.deactivate();
        resourceRepository.save(inactive);

        List<OfficeResource> activeResources = resourceRepository.findActive();

        assertThat(activeResources).isEmpty();
    }

    @Test
    void find_returns_empty_when_resource_missing() {
        var missingId = ResourceId.of(UUID.randomUUID());

        assertThat(resourceRepository.find(missingId)).isEmpty();
    }
}
\end{DeskFlowCode}

\subsection{Testy - Administration}

\phantomsection
\addcontentsline{toc}{subsubsection}{AdministeredAllocationRequestIdTest}
\DeskFlowListingCaption{administration/application/domain/model/AdministeredAllocationRequestIdTest.java - AdministeredAllocationRequestIdTest}{lst:administration-application-domain-model-AdministeredAllocationRequestIdTest-java}
\begin{DeskFlowCode}
package com.deskflow.administration.application.domain.model;

import org.junit.jupiter.api.Test;

import java.util.UUID;

import static org.assertj.core.api.Assertions.*;

class AdministeredAllocationRequestIdTest {

    @Test
    void creates_from_uuid() {
        UUID uuid = UUID.randomUUID();
        var id = AdministeredAllocationRequestId.of(uuid);
        assertThat(id.getValue()).isEqualTo(uuid);
    }

    @Test
    void rejects_null() {
        assertThatThrownBy(() -> AdministeredAllocationRequestId.of(null))
                .isInstanceOf(NullPointerException.class);
    }

    @Test
    void equality() {
        UUID uuid = UUID.randomUUID();
        assertThat(AdministeredAllocationRequestId.of(uuid)).isEqualTo(AdministeredAllocationRequestId.of(uuid));
    }

    @Test
    void inequality() {
        assertThat(AdministeredAllocationRequestId.of(UUID.randomUUID()))
                .isNotEqualTo(AdministeredAllocationRequestId.of(UUID.randomUUID()));
    }

    @Test
    void toString_returns_uuid_string() {
        UUID uuid = UUID.randomUUID();
        assertThat(AdministeredAllocationRequestId.of(uuid)).hasToString(uuid.toString());
    }
}
\end{DeskFlowCode}

\phantomsection
\addcontentsline{toc}{subsubsection}{AdministeredAllocationRequestTest}
\DeskFlowListingCaption{administration/application/domain/model/AdministeredAllocationRequestTest.java - AdministeredAllocationRequestTest}{lst:administration-application-domain-model-AdministeredAllocationRequestTest-java}
\begin{DeskFlowCode}
package com.deskflow.administration.application.domain.model;

import org.junit.jupiter.api.Test;

import java.util.UUID;

import static org.assertj.core.api.Assertions.*;

class AdministeredAllocationRequestTest {

    private AdministeredAllocationRequestId id = AdministeredAllocationRequestId.of(UUID.randomUUID());

    @Test
    void approve_from_submitted_passes() {
        var req = new AdministeredAllocationRequest(id, AdministeredAllocationStatus.SUBMITTED);

        req.ensureCanApprove();
    }

    @Test
    void approve_from_non_submitted_throws() {
        assertThatThrownBy(() -> new AdministeredAllocationRequest(id, AdministeredAllocationStatus.APPROVED).ensureCanApprove())
                .isInstanceOf(IllegalStateException.class)
                .hasMessageContaining("approve");
        assertThatThrownBy(() -> new AdministeredAllocationRequest(id, AdministeredAllocationStatus.ACTIVE).ensureCanApprove())
                .isInstanceOf(IllegalStateException.class);
        assertThatThrownBy(() -> new AdministeredAllocationRequest(id, AdministeredAllocationStatus.REJECTED).ensureCanApprove())
                .isInstanceOf(IllegalStateException.class);
        assertThatThrownBy(() -> new AdministeredAllocationRequest(id, AdministeredAllocationStatus.CANCELLED).ensureCanApprove())
                .isInstanceOf(IllegalStateException.class);
        assertThatThrownBy(() -> new AdministeredAllocationRequest(id, AdministeredAllocationStatus.COMPLETED).ensureCanApprove())
                .isInstanceOf(IllegalStateException.class);
        assertThatThrownBy(() -> new AdministeredAllocationRequest(id, AdministeredAllocationStatus.OVERDUE).ensureCanApprove())
                .isInstanceOf(IllegalStateException.class);
    }

    @Test
    void reject_from_submitted_passes() {
        var req = new AdministeredAllocationRequest(id, AdministeredAllocationStatus.SUBMITTED);

        req.ensureCanReject();
    }

    @Test
    void reject_from_non_submitted_throws() {
        assertThatThrownBy(() -> new AdministeredAllocationRequest(id, AdministeredAllocationStatus.APPROVED).ensureCanReject())
                .isInstanceOf(IllegalStateException.class)
                .hasMessageContaining("reject");
    }

    @Test
    void issue_from_approved_passes() {
        var req = new AdministeredAllocationRequest(id, AdministeredAllocationStatus.APPROVED);

        req.ensureCanIssue();
    }

    @Test
    void issue_from_non_approved_throws() {
        assertThatThrownBy(() -> new AdministeredAllocationRequest(id, AdministeredAllocationStatus.SUBMITTED).ensureCanIssue())
                .isInstanceOf(IllegalStateException.class)
                .hasMessageContaining("issue resource");
    }

    @Test
    void return_from_active_passes() {
        var req = new AdministeredAllocationRequest(id, AdministeredAllocationStatus.ACTIVE);

        req.ensureCanReturn();
    }

    @Test
    void return_from_overdue_passes() {
        var req = new AdministeredAllocationRequest(id, AdministeredAllocationStatus.OVERDUE);

        req.ensureCanReturn();
    }

    @Test
    void return_from_non_active_or_overdue_throws() {
        assertThatThrownBy(() -> new AdministeredAllocationRequest(id, AdministeredAllocationStatus.SUBMITTED).ensureCanReturn())
                .isInstanceOf(IllegalStateException.class)
                .hasMessageContaining("return resource");
        assertThatThrownBy(() -> new AdministeredAllocationRequest(id, AdministeredAllocationStatus.APPROVED).ensureCanReturn())
                .isInstanceOf(IllegalStateException.class);
        assertThatThrownBy(() -> new AdministeredAllocationRequest(id, AdministeredAllocationStatus.COMPLETED).ensureCanReturn())
                .isInstanceOf(IllegalStateException.class);
        assertThatThrownBy(() -> new AdministeredAllocationRequest(id, AdministeredAllocationStatus.REJECTED).ensureCanReturn())
                .isInstanceOf(IllegalStateException.class);
        assertThatThrownBy(() -> new AdministeredAllocationRequest(id, AdministeredAllocationStatus.CANCELLED).ensureCanReturn())
                .isInstanceOf(IllegalStateException.class);
    }
}
\end{DeskFlowCode}

\phantomsection
\addcontentsline{toc}{subsubsection}{AdministeredAllocationStatusTest}
\DeskFlowListingCaption{administration/application/domain/model/AdministeredAllocationStatusTest.java - AdministeredAllocationStatusTest}{lst:administration-application-domain-model-AdministeredAllocationStatusTest-java}
\begin{DeskFlowCode}
package com.deskflow.administration.application.domain.model;

import org.junit.jupiter.api.Test;

import static org.assertj.core.api.Assertions.*;

class AdministeredAllocationStatusTest {

    @Test
    void from_returns_correct_enum() {
        assertThat(AdministeredAllocationStatus.from("SUBMITTED")).isEqualTo(AdministeredAllocationStatus.SUBMITTED);
        assertThat(AdministeredAllocationStatus.from("APPROVED")).isEqualTo(AdministeredAllocationStatus.APPROVED);
        assertThat(AdministeredAllocationStatus.from("ACTIVE")).isEqualTo(AdministeredAllocationStatus.ACTIVE);
        assertThat(AdministeredAllocationStatus.from("OVERDUE")).isEqualTo(AdministeredAllocationStatus.OVERDUE);
        assertThat(AdministeredAllocationStatus.from("REJECTED")).isEqualTo(AdministeredAllocationStatus.REJECTED);
        assertThat(AdministeredAllocationStatus.from("CANCELLED")).isEqualTo(AdministeredAllocationStatus.CANCELLED);
        assertThat(AdministeredAllocationStatus.from("COMPLETED")).isEqualTo(AdministeredAllocationStatus.COMPLETED);
    }

    @Test
    void from_invalid_throws() {
        assertThatThrownBy(() -> AdministeredAllocationStatus.from("INVALID"))
                .isInstanceOf(IllegalArgumentException.class);
    }
}
\end{DeskFlowCode}

\phantomsection
\addcontentsline{toc}{subsubsection}{AdministrationAppServiceTest}
\DeskFlowListingCaption{administration/application/service/AdministrationAppServiceTest.java - AdministrationAppServiceTest}{lst:administration-application-service-AdministrationAppServiceTest-java}
\begin{DeskFlowCode}
package com.deskflow.administration.application.service;

import com.deskflow.administration.application.port.out.AdministrationAllocationPort;
import com.deskflow.administration.application.model.AdministeredAllocationView;
import com.deskflow.administration.application.domain.model.AdministeredAllocationRequestId;
import com.deskflow.administration.application.domain.model.AdministeredAllocationStatus;
import org.junit.jupiter.api.BeforeEach;
import org.junit.jupiter.api.Test;
import org.junit.jupiter.api.extension.ExtendWith;
import org.mockito.Mock;
import org.mockito.junit.jupiter.MockitoExtension;

import java.time.LocalDate;
import java.time.LocalDateTime;
import java.util.Optional;
import java.util.UUID;

import static org.assertj.core.api.Assertions.assertThatThrownBy;
import static org.mockito.Mockito.verify;
import static org.mockito.Mockito.when;

@ExtendWith(MockitoExtension.class)
class AdministrationAppServiceTest {

    @Mock private AdministrationAllocationPort allocationPort;

    private AdministrationAppService appService;

    @BeforeEach
    void setUp() {
        appService = new AdministrationAppService(allocationPort);
    }

    private AdministeredAllocationView requestView(AdministeredAllocationRequestId requestId, AdministeredAllocationStatus status) {
        return new AdministeredAllocationView(
                requestId.getValue(),
                UUID.randomUUID(),
                UUID.randomUUID(),
                status.name(),
                LocalDate.of(2025, 7, 1),
                LocalDate.of(2025, 7, 5),
                "test",
                LocalDateTime.now(),
                null,
                null,
                null
        );
    }

    @Test
    void approve_request_delegates_to_allocation_port() {
        AdministeredAllocationRequestId requestId = AdministeredAllocationRequestId.of(UUID.randomUUID());
        when(allocationPort.getRequest(requestId)).thenReturn(Optional.of(requestView(requestId, AdministeredAllocationStatus.SUBMITTED)));

        appService.approveRequest(requestId);

        verify(allocationPort).approveRequest(requestId);
    }

    @Test
    void list_requests_delegates_to_allocation_port() {
        appService.listRequests(AdministeredAllocationStatus.SUBMITTED);

        verify(allocationPort).listRequests(AdministeredAllocationStatus.SUBMITTED);
    }

    @Test
    void detect_overdue_delegates_to_allocation_port() {
        appService.detectOverdueAllocations();

        verify(allocationPort).detectOverdueAllocations();
    }

    @Test
    void approve_rejects_non_submitted_request_before_calling_allocation() {
        AdministeredAllocationRequestId requestId = AdministeredAllocationRequestId.of(UUID.randomUUID());
        when(allocationPort.getRequest(requestId)).thenReturn(Optional.of(requestView(requestId, AdministeredAllocationStatus.ACTIVE)));

        assertThatThrownBy(() -> appService.approveRequest(requestId))
                .isInstanceOf(IllegalStateException.class)
                .hasMessageContaining("requires status SUBMITTED");
    }
}
\end{DeskFlowCode}

\phantomsection
\addcontentsline{toc}{subsubsection}{OverdueAllocationSchedulerTest}
\DeskFlowListingCaption{administration/application/service/OverdueAllocationSchedulerTest.java - OverdueAllocationSchedulerTest}{lst:administration-application-service-OverdueAllocationSchedulerTest-java}
\begin{DeskFlowCode}
package com.deskflow.administration.application.service;

import com.deskflow.administration.application.port.out.AdministrationAllocationPort;
import org.junit.jupiter.api.Test;
import org.junit.jupiter.api.extension.ExtendWith;
import org.mockito.Mock;
import org.mockito.junit.jupiter.MockitoExtension;

import static org.mockito.Mockito.*;

@ExtendWith(MockitoExtension.class)
class OverdueAllocationSchedulerTest {

    @Mock private AdministrationAllocationPort allocationPort;

    @Test
    void detectOverdueAllocations_delegates_to_administration_service() {
        var adminService = new AdministrationAppService(allocationPort);
        var scheduler = new OverdueAllocationScheduler(adminService);

        scheduler.detectOverdueAllocations();

        verify(allocationPort).detectOverdueAllocations();
    }
}
\end{DeskFlowCode}

\phantomsection
\addcontentsline{toc}{subsubsection}{AllocationAdministrationAdapterTest}
\DeskFlowListingCaption{administration/infrastructure/out/allocation/AllocationAdministrationAdapterTest.java - AllocationAdministrationAdapterTest}{lst:administration-infrastructure-out-allocation-AllocationAdministrationAdapterTest-java}
\begin{DeskFlowCode}
package com.deskflow.administration.infrastructure.out.allocation;

import com.deskflow.administration.application.model.AdministeredAllocationView;
import com.deskflow.administration.application.domain.model.AdministeredAllocationRequestId;
import com.deskflow.administration.application.domain.model.AdministeredAllocationStatus;
import com.deskflow.allocation.application.model.AllocationRequestView;
import com.deskflow.allocation.application.port.in.AllocationAdministrationInputPort;
import com.deskflow.allocation.application.domain.model.vo.AllocationRequestId;
import org.junit.jupiter.api.BeforeEach;
import org.junit.jupiter.api.Test;
import org.junit.jupiter.api.extension.ExtendWith;
import org.mockito.Mock;
import org.mockito.junit.jupiter.MockitoExtension;

import java.time.LocalDate;
import java.time.LocalDateTime;
import java.util.List;
import java.util.Optional;
import java.util.UUID;

import static org.assertj.core.api.Assertions.*;
import static org.mockito.Mockito.*;

@ExtendWith(MockitoExtension.class)
class AllocationAdministrationAdapterTest {

    @Mock private AllocationAdministrationInputPort allocationAdministrationInputPort;

    private AllocationAdministrationAdapter adapter;

    @BeforeEach
    void setUp() {
        adapter = new AllocationAdministrationAdapter(allocationAdministrationInputPort);
    }

    @Test
    void getRequest_returns_admin_view_when_found() {
        var adminId = AdministeredAllocationRequestId.of(UUID.randomUUID());
        var allocView = new AllocationRequestView(adminId.getValue(), UUID.randomUUID(), UUID.randomUUID(),
                "SUBMITTED", LocalDate.now(), LocalDate.now().plusDays(1), "test",
                LocalDateTime.now(), null, null, null);
        when(allocationAdministrationInputPort.getRequest(AllocationRequestId.of(adminId.getValue())))
                .thenReturn(Optional.of(allocView));

        var result = adapter.getRequest(adminId);

        assertThat(result).isPresent();
        assertThat(result.get().id()).isEqualTo(adminId.getValue());
        assertThat(result.get().status()).isEqualTo("SUBMITTED");
    }

    @Test
    void getRequest_returns_empty_when_not_found() {
        var adminId = AdministeredAllocationRequestId.of(UUID.randomUUID());
        when(allocationAdministrationInputPort.getRequest(any())).thenReturn(Optional.empty());

        var result = adapter.getRequest(adminId);

        assertThat(result).isEmpty();
    }

    @Test
    void listRequests_converts_status_and_maps() {
        when(allocationAdministrationInputPort.listRequests(any()))
                .thenReturn(List.of(
                        new AllocationRequestView(UUID.randomUUID(), UUID.randomUUID(), UUID.randomUUID(),
                                "SUBMITTED", LocalDate.now(), LocalDate.now().plusDays(1), "test",
                                LocalDateTime.now(), null, null, null)
                ));

        var result = adapter.listRequests(AdministeredAllocationStatus.SUBMITTED);

        assertThat(result).hasSize(1);
        verify(allocationAdministrationInputPort).listRequests(com.deskflow.allocation.application.domain.model.vo.RequestStatus.SUBMITTED);
    }

    @Test
    void listRequests_passes_null_status() {
        when(allocationAdministrationInputPort.listRequests(null)).thenReturn(List.of());

        var result = adapter.listRequests(null);

        assertThat(result).isEmpty();
    }

    @Test
    void approveRequest_delegates() {
        var adminId = AdministeredAllocationRequestId.of(UUID.randomUUID());

        adapter.approveRequest(adminId);

        verify(allocationAdministrationInputPort).approveRequest(AllocationRequestId.of(adminId.getValue()));
    }

    @Test
    void rejectRequest_delegates() {
        var adminId = AdministeredAllocationRequestId.of(UUID.randomUUID());

        adapter.rejectRequest(adminId, "No space");

        verify(allocationAdministrationInputPort).rejectRequest(AllocationRequestId.of(adminId.getValue()), "No space");
    }

    @Test
    void issueResource_delegates() {
        var adminId = AdministeredAllocationRequestId.of(UUID.randomUUID());

        adapter.issueResource(adminId);

        verify(allocationAdministrationInputPort).issueResource(AllocationRequestId.of(adminId.getValue()));
    }

    @Test
    void returnResource_delegates() {
        var adminId = AdministeredAllocationRequestId.of(UUID.randomUUID());

        adapter.returnResource(adminId);

        verify(allocationAdministrationInputPort).returnResource(AllocationRequestId.of(adminId.getValue()));
    }

    @Test
    void monitorActiveAllocations_returns_mapped_views() {
        when(allocationAdministrationInputPort.monitorActiveAllocations())
                .thenReturn(List.of(
                        new AllocationRequestView(UUID.randomUUID(), UUID.randomUUID(), UUID.randomUUID(),
                                "ACTIVE", LocalDate.now(), LocalDate.now().plusDays(1), "test",
                                LocalDateTime.now(), LocalDateTime.now(), null, null)
                ));

        var result = adapter.monitorActiveAllocations();

        assertThat(result).hasSize(1);
        assertThat(result.getFirst().status()).isEqualTo("ACTIVE");
    }

    @Test
    void detectOverdueAllocations_delegates() {
        adapter.detectOverdueAllocations();

        verify(allocationAdministrationInputPort).detectOverdueAllocations();
    }
}
\end{DeskFlowCode}

\subsection{Testy - Employees}

\phantomsection
\addcontentsline{toc}{subsubsection}{EmployeeTest}
\DeskFlowListingCaption{employees/application/domain/model/EmployeeTest.java - EmployeeTest}{lst:employees-application-domain-model-EmployeeTest-java}
\begin{DeskFlowCode}
package com.deskflow.employees.application.domain.model;

import com.deskflow.common.application.domain.model.vo.EmployeeId;
import com.deskflow.common.application.domain.model.vo.ResourceType;
import com.deskflow.employees.application.domain.model.vo.Role;
import org.junit.jupiter.api.Test;

import java.util.Map;
import java.util.UUID;

import static org.assertj.core.api.Assertions.assertThat;
import static org.assertj.core.api.Assertions.assertThatThrownBy;

class EmployeeTest {

    @Test
    void limit_for_missing_resource_type_defaults_to_one() {
        Employee employee = new Employee(EmployeeId.of(UUID.randomUUID()), "a@test.pl", "Anna", "Test",
                Role.EMPLOYEE, Map.of(ResourceType.DESK, 3));

        assertThat(employee.limitForType(ResourceType.LAPTOP)).isEqualTo(1);
    }

    @Test
    void change_limit_updates_limit_for_resource_type() {
        Employee employee = new Employee(EmployeeId.of(UUID.randomUUID()), "a@test.pl", "Anna", "Test",
                Role.EMPLOYEE, Map.of());

        employee.changeLimit(ResourceType.DESK, 0);

        assertThat(employee.limitForType(ResourceType.DESK)).isZero();
    }

    @Test
    void change_limit_rejects_negative_amount() {
        Employee employee = new Employee(EmployeeId.of(UUID.randomUUID()), "a@test.pl", "Anna", "Test",
                Role.EMPLOYEE, Map.of());

        assertThatThrownBy(() -> employee.changeLimit(ResourceType.DESK, -1))
                .isInstanceOf(IllegalArgumentException.class)
                .hasMessageContaining("cannot be negative");
    }

    @Test
    void getters_return_expected_values() {
        var id = EmployeeId.of(UUID.randomUUID());
        var employee = new Employee(id, "test@test.pl", "Jan", "Kowalski", Role.MANAGER,
                Map.of(ResourceType.DESK, 2));

        assertThat(employee.getId()).isEqualTo(id);
        assertThat(employee.getEmail()).isEqualTo("test@test.pl");
        assertThat(employee.getFirstName()).isEqualTo("Jan");
        assertThat(employee.getLastName()).isEqualTo("Kowalski");
        assertThat(employee.getRole()).isEqualTo(Role.MANAGER);
    }

    @Test
    void change_limit_overwrites_existing_limit() {
        Employee employee = new Employee(EmployeeId.of(UUID.randomUUID()), "a@test.pl", "Anna", "Test",
                Role.EMPLOYEE, Map.of(ResourceType.DESK, 3));

        employee.changeLimit(ResourceType.DESK, 5);

        assertThat(employee.limitForType(ResourceType.DESK)).isEqualTo(5);
    }

    @Test
    void constructor_makes_defensive_copy_of_limits() {
        var mutableLimits = new java.util.HashMap<ResourceType, Integer>();
        mutableLimits.put(ResourceType.DESK, 3);
        Employee employee = new Employee(EmployeeId.of(UUID.randomUUID()), "a@test.pl", "Anna", "Test",
                Role.EMPLOYEE, mutableLimits);

        mutableLimits.put(ResourceType.LAPTOP, 2);

        assertThat(employee.limitForType(ResourceType.LAPTOP)).isEqualTo(1);
    }

    @Test
    void getLimits_returns_unmodifiable_map() {
        Employee employee = new Employee(EmployeeId.of(UUID.randomUUID()), "a@test.pl", "Anna", "Test",
                Role.EMPLOYEE, Map.of(ResourceType.DESK, 3));

        assertThatThrownBy(() -> employee.getLimits().put(ResourceType.LAPTOP, 2))
                .isInstanceOf(UnsupportedOperationException.class);
    }

    @Test
    void limit_for_zero_is_valid() {
        Employee employee = new Employee(EmployeeId.of(UUID.randomUUID()), "a@test.pl", "Anna", "Test",
                Role.EMPLOYEE, Map.of(ResourceType.DESK, 0));

        assertThat(employee.limitForType(ResourceType.DESK)).isZero();
    }
}
\end{DeskFlowCode}

\phantomsection
\addcontentsline{toc}{subsubsection}{EmployeeAppServiceTest}
\DeskFlowListingCaption{employees/application/service/EmployeeAppServiceTest.java - EmployeeAppServiceTest}{lst:employees-application-service-EmployeeAppServiceTest-java}
\begin{DeskFlowCode}
package com.deskflow.employees.application.service;

import com.deskflow.common.application.domain.model.vo.EmployeeId;
import com.deskflow.common.application.domain.model.vo.ResourceType;
import com.deskflow.employees.application.port.out.EmployeeRepositoryPort;
import com.deskflow.employees.application.domain.model.Employee;
import com.deskflow.employees.application.domain.model.vo.Role;
import org.junit.jupiter.api.BeforeEach;
import org.junit.jupiter.api.Test;
import org.junit.jupiter.api.extension.ExtendWith;
import org.mockito.Mock;
import org.mockito.junit.jupiter.MockitoExtension;

import java.util.Map;
import java.util.NoSuchElementException;
import java.util.Optional;
import java.util.UUID;

import static org.assertj.core.api.Assertions.assertThat;
import static org.assertj.core.api.Assertions.assertThatThrownBy;
import static org.mockito.Mockito.verify;
import static org.mockito.Mockito.when;

@ExtendWith(MockitoExtension.class)
class EmployeeAppServiceTest {

    @Mock private EmployeeRepositoryPort employeeRepository;

    private EmployeeAppService appService;

    @BeforeEach
    void setUp() {
        appService = new EmployeeAppService(employeeRepository);
    }

    @Test
    void upsert_limit_changes_employee_limit_and_saves() {
        EmployeeId employeeId = EmployeeId.of(UUID.randomUUID());
        Employee employee = new Employee(employeeId, "a@test.pl", "Anna", "Test", Role.EMPLOYEE, Map.of());
        when(employeeRepository.find(employeeId)).thenReturn(Optional.of(employee));

        appService.upsertLimit(employeeId, "DESK", 4);

        assertThat(employee.limitForType(ResourceType.DESK)).isEqualTo(4);
        verify(employeeRepository).save(employee);
    }

    @Test
    void upsert_limit_throws_when_employee_missing() {
        EmployeeId employeeId = EmployeeId.of(UUID.randomUUID());
        when(employeeRepository.find(employeeId)).thenReturn(Optional.empty());

        assertThatThrownBy(() -> appService.upsertLimit(employeeId, "DESK", 4))
                .isInstanceOf(NoSuchElementException.class)
                .hasMessageContaining("Employee not found");
    }

    @Test
    void upsert_limit_rejects_invalid_resource_type() {
        EmployeeId employeeId = EmployeeId.of(UUID.randomUUID());
        Employee employee = new Employee(employeeId, "a@test.pl", "Anna", "Test", Role.EMPLOYEE, Map.of());
        when(employeeRepository.find(employeeId)).thenReturn(Optional.of(employee));

        assertThatThrownBy(() -> appService.upsertLimit(employeeId, "UNKNOWN", 4))
                .isInstanceOf(IllegalArgumentException.class);
    }
}
\end{DeskFlowCode}

\phantomsection
\addcontentsline{toc}{subsubsection}{EmployeeRepositoryIntegrationTest}
\DeskFlowListingCaption{employees/infrastructure/out/persistence/EmployeeRepositoryIntegrationTest.java - EmployeeRepositoryIntegrationTest}{lst:employees-infrastructure-out-persistence-EmployeeRepositoryIntegrationTest-java}
\begin{DeskFlowCode}
package com.deskflow.employees.infrastructure.out.persistence;

import org.junit.jupiter.api.Test;
import org.springframework.beans.factory.annotation.Autowired;
import org.springframework.boot.jdbc.test.autoconfigure.AutoConfigureTestDatabase;
import org.springframework.boot.data.jpa.test.autoconfigure.DataJpaTest;
import org.springframework.test.context.ActiveProfiles;

import java.util.Map;
import java.util.UUID;

import static org.assertj.core.api.Assertions.*;

@DataJpaTest
@AutoConfigureTestDatabase(replace = AutoConfigureTestDatabase.Replace.NONE)
@ActiveProfiles("test")
class EmployeeRepositoryIntegrationTest {

    @Autowired private SpringDataEmployeeRepository jpa;

    @Test
    void save_and_find_employee() {
        EmployeeEntity entity = new EmployeeEntity();
        entity.setId(UUID.randomUUID());
        entity.setEmail("test@company.test");
        entity.setFirstName("Jan");
        entity.setLastName("Kowalski");
        entity.setRole("EMPLOYEE");
        entity.setPasswordHash("hash");
        jpa.save(entity);

        var found = jpa.findById(entity.getId());

        assertThat(found).isPresent();
        assertThat(found.get().getEmail()).isEqualTo("test@company.test");
        assertThat(found.get().getFirstName()).isEqualTo("Jan");
        assertThat(found.get().getLastName()).isEqualTo("Kowalski");
        assertThat(found.get().getRole()).isEqualTo("EMPLOYEE");
    }

    @Test
    void find_by_email() {
        EmployeeEntity entity = new EmployeeEntity();
        entity.setId(UUID.randomUUID());
        entity.setEmail("unique@company.test");
        entity.setFirstName("A");
        entity.setLastName("B");
        entity.setRole("MANAGER");
        entity.setPasswordHash("hash");
        jpa.save(entity);

        var found = jpa.findByEmail("unique@company.test");

        assertThat(found).isPresent();
        assertThat(found.get().getId()).isEqualTo(entity.getId());
    }

    @Test
    void find_by_email_returns_empty_when_missing() {
        var found = jpa.findByEmail("nonexistent@company.test");

        assertThat(found).isEmpty();
    }

    @Test
    void find_returns_empty_when_missing() {
        var found = jpa.findById(UUID.randomUUID());

        assertThat(found).isEmpty();
    }

    @Test
    void saves_and_loads_allocation_limits() {
        EmployeeEntity entity = new EmployeeEntity();
        entity.setId(UUID.randomUUID());
        entity.setEmail("limits@company.test");
        entity.setFirstName("A");
        entity.setLastName("B");
        entity.setRole("EMPLOYEE");
        entity.setPasswordHash("hash");

        AllocationLimitEntity limit = new AllocationLimitEntity();
        limit.setId(UUID.randomUUID());
        limit.setResourceType("DESK");
        limit.setLimitAmount(3);
        entity.addAllocationLimit(limit);

        jpa.save(entity);

        var found = jpa.findById(entity.getId()).orElseThrow();
        assertThat(found.getAllocationLimits()).hasSize(1);
        assertThat(found.getAllocationLimits().getFirst().getResourceType()).isEqualTo("DESK");
        assertThat(found.getAllocationLimits().getFirst().getLimitAmount()).isEqualTo(3);
    }

    @Test
    void find_all_returns_all_employees() {
        EmployeeEntity e1 = new EmployeeEntity();
        e1.setId(UUID.randomUUID());
        e1.setEmail("e1@company.test");
        e1.setFirstName("E1");
        e1.setLastName("N1");
        e1.setRole("EMPLOYEE");
        e1.setPasswordHash("hash");

        EmployeeEntity e2 = new EmployeeEntity();
        e2.setId(UUID.randomUUID());
        e2.setEmail("e2@company.test");
        e2.setFirstName("E2");
        e2.setLastName("N2");
        e2.setRole("MANAGER");
        e2.setPasswordHash("hash");

        jpa.save(e1);
        jpa.save(e2);

        var all = jpa.findAll();

        assertThat(all).hasSize(2);
    }
}
\end{DeskFlowCode}

\subsection{Testy - Identity}

\phantomsection
\addcontentsline{toc}{subsubsection}{IdentityAccessAppServiceTest}
\DeskFlowListingCaption{identity/application/service/IdentityAccessAppServiceTest.java - IdentityAccessAppServiceTest}{lst:identity-application-service-IdentityAccessAppServiceTest-java}
\begin{DeskFlowCode}
package com.deskflow.identity.application.service;

import com.deskflow.identity.application.model.IdentityProfile;
import com.deskflow.identity.application.port.out.AccessTokenIssuerPort;
import com.deskflow.identity.application.port.out.CredentialsAuthenticationPort;
import com.deskflow.identity.application.port.out.IdentityEmployeeProfilePort;
import org.junit.jupiter.api.BeforeEach;
import org.junit.jupiter.api.Test;
import org.junit.jupiter.api.extension.ExtendWith;
import org.mockito.Mock;
import org.mockito.junit.jupiter.MockitoExtension;

import java.util.NoSuchElementException;
import java.util.Optional;

import static org.assertj.core.api.Assertions.assertThat;
import static org.assertj.core.api.Assertions.assertThatThrownBy;
import static org.mockito.Mockito.verify;
import static org.mockito.Mockito.when;

@ExtendWith(MockitoExtension.class)
class IdentityAccessAppServiceTest {

    @Mock private CredentialsAuthenticationPort credentialsAuthenticationPort;
    @Mock private AccessTokenIssuerPort accessTokenIssuerPort;
    @Mock private IdentityEmployeeProfilePort employeeProfilePort;

    private IdentityAccessAppService appService;

    @BeforeEach
    void setUp() {
        appService = new IdentityAccessAppService(credentialsAuthenticationPort, accessTokenIssuerPort, employeeProfilePort);
    }

    @Test
    void login_authenticates_user_and_returns_profile_with_token() {
        IdentityProfile profile = new IdentityProfile("user@test.pl", "MANAGER", "Jan", "Testowy", "emp-1");
        when(employeeProfilePort.findByEmail("user@test.pl")).thenReturn(Optional.of(profile));
        when(accessTokenIssuerPort.issueAccessToken("user@test.pl", "MANAGER")).thenReturn("token-123");

        var identity = appService.login("user@test.pl", "password");

        verify(credentialsAuthenticationPort).authenticate("user@test.pl", "password");
        assertThat(identity.token()).isEqualTo("token-123");
        assertThat(identity.email()).isEqualTo("user@test.pl");
        assertThat(identity.role()).isEqualTo("MANAGER");
        assertThat(identity.firstName()).isEqualTo("Jan");
        assertThat(identity.lastName()).isEqualTo("Testowy");
        assertThat(identity.employeeId()).isEqualTo("emp-1");
    }

    @Test
    void login_throws_when_authenticated_user_has_no_employee_profile() {
        when(employeeProfilePort.findByEmail("missing@test.pl")).thenReturn(Optional.empty());

        assertThatThrownBy(() -> appService.login("missing@test.pl", "password"))
                .isInstanceOf(NoSuchElementException.class);

        verify(credentialsAuthenticationPort).authenticate("missing@test.pl", "password");
    }

    @Test
    void current_identity_returns_profile_by_email() {
        IdentityProfile profile = new IdentityProfile("user@test.pl", "EMPLOYEE", "Anna", "Nowak", "emp-2");
        when(employeeProfilePort.findByEmail("user@test.pl")).thenReturn(Optional.of(profile));

        assertThat(appService.currentIdentity("user@test.pl")).contains(profile);
    }
}
\end{DeskFlowCode}

\phantomsection
\addcontentsline{toc}{subsubsection}{EmployeeIdentityProfileAdapterTest}
\DeskFlowListingCaption{identity/infrastructure/out/employees/EmployeeIdentityProfileAdapterTest.java - EmployeeIdentityProfileAdapterTest}{lst:identity-infrastructure-out-employees-EmployeeIdentityProfileAdapterTest-java}
\begin{DeskFlowCode}
package com.deskflow.identity.infrastructure.out.employees;

import com.deskflow.common.application.domain.model.vo.EmployeeId;
import com.deskflow.common.application.domain.model.vo.ResourceType;
import com.deskflow.employees.application.port.in.EmployeeManagementInputPort;
import com.deskflow.employees.application.domain.model.Employee;
import com.deskflow.employees.application.domain.model.vo.Role;
import org.junit.jupiter.api.BeforeEach;
import org.junit.jupiter.api.Test;
import org.junit.jupiter.api.extension.ExtendWith;
import org.mockito.Mock;
import org.mockito.junit.jupiter.MockitoExtension;

import java.util.Map;
import java.util.Optional;
import java.util.UUID;

import static org.assertj.core.api.Assertions.*;
import static org.mockito.Mockito.*;

@ExtendWith(MockitoExtension.class)
class EmployeeIdentityProfileAdapterTest {

    @Mock private EmployeeManagementInputPort employeeManagementInputPort;

    private EmployeeIdentityProfileAdapter adapter;

    @BeforeEach
    void setUp() {
        adapter = new EmployeeIdentityProfileAdapter(employeeManagementInputPort);
    }

    @Test
    void findByEmail_returns_mapped_profile_when_found() {
        var employeeId = EmployeeId.of(UUID.randomUUID());
        var employee = new Employee(employeeId, "test@company.test", "Jan", "Kowalski", Role.MANAGER, Map.of());
        when(employeeManagementInputPort.getEmployeeByEmail("test@company.test"))
                .thenReturn(Optional.of(employee));

        var result = adapter.findByEmail("test@company.test");

        assertThat(result).isPresent();
        assertThat(result.get().email()).isEqualTo("test@company.test");
        assertThat(result.get().role()).isEqualTo("MANAGER");
        assertThat(result.get().firstName()).isEqualTo("Jan");
        assertThat(result.get().lastName()).isEqualTo("Kowalski");
        assertThat(result.get().employeeId()).isEqualTo(employeeId.toString());
    }

    @Test
    void findByEmail_returns_empty_when_not_found() {
        when(employeeManagementInputPort.getEmployeeByEmail("missing@company.test"))
                .thenReturn(Optional.empty());

        var result = adapter.findByEmail("missing@company.test");

        assertThat(result).isEmpty();
    }
}
\end{DeskFlowCode}

\phantomsection
\addcontentsline{toc}{subsubsection}{EmployeeUserCredentialsAdapterTest}
\DeskFlowListingCaption{identity/infrastructure/out/employees/EmployeeUserCredentialsAdapterTest.java - EmployeeUserCredentialsAdapterTest}{lst:identity-infrastructure-out-employees-EmployeeUserCredentialsAdapterTest-java}
\begin{DeskFlowCode}
package com.deskflow.identity.infrastructure.out.employees;

import com.deskflow.employees.application.model.EmployeeCredentials;
import com.deskflow.employees.application.port.in.EmployeeManagementInputPort;
import org.junit.jupiter.api.BeforeEach;
import org.junit.jupiter.api.Test;
import org.junit.jupiter.api.extension.ExtendWith;
import org.mockito.Mock;
import org.mockito.junit.jupiter.MockitoExtension;

import java.util.Optional;

import static org.assertj.core.api.Assertions.*;
import static org.mockito.Mockito.*;

@ExtendWith(MockitoExtension.class)
class EmployeeUserCredentialsAdapterTest {

    @Mock private EmployeeManagementInputPort employeeManagementInputPort;

    private EmployeeUserCredentialsAdapter adapter;

    @BeforeEach
    void setUp() {
        adapter = new EmployeeUserCredentialsAdapter(employeeManagementInputPort);
    }

    @Test
    void findByEmail_returns_mapped_credentials_when_found() {
        var employeeCredentials = new EmployeeCredentials("test@company.test", "hashed_pass", "EMPLOYEE");
        when(employeeManagementInputPort.getEmployeeCredentialsByEmail("test@company.test"))
                .thenReturn(Optional.of(employeeCredentials));

        var result = adapter.findByEmail("test@company.test");

        assertThat(result).isPresent();
        assertThat(result.get().email()).isEqualTo("test@company.test");
        assertThat(result.get().passwordHash()).isEqualTo("hashed_pass");
        assertThat(result.get().role()).isEqualTo("EMPLOYEE");
    }

    @Test
    void findByEmail_returns_empty_when_not_found() {
        when(employeeManagementInputPort.getEmployeeCredentialsByEmail("missing@company.test"))
                .thenReturn(Optional.empty());

        var result = adapter.findByEmail("missing@company.test");

        assertThat(result).isEmpty();
    }
}
\end{DeskFlowCode}

\phantomsection
\addcontentsline{toc}{subsubsection}{DeskflowUserDetailsServiceTest}
\DeskFlowListingCaption{identity/infrastructure/security/DeskflowUserDetailsServiceTest.java - DeskflowUserDetailsServiceTest}{lst:identity-infrastructure-security-DeskflowUserDetailsServiceTest-java}
\begin{DeskFlowCode}
package com.deskflow.identity.infrastructure.security;

import com.deskflow.identity.application.model.UserCredentials;
import com.deskflow.identity.application.port.out.UserCredentialsPort;
import org.junit.jupiter.api.BeforeEach;
import org.junit.jupiter.api.Test;
import org.junit.jupiter.api.extension.ExtendWith;
import org.mockito.Mock;
import org.mockito.junit.jupiter.MockitoExtension;
import org.springframework.security.core.userdetails.UsernameNotFoundException;

import java.util.Optional;

import static org.assertj.core.api.Assertions.*;
import static org.mockito.Mockito.*;

@ExtendWith(MockitoExtension.class)
class DeskflowUserDetailsServiceTest {

    @Mock private UserCredentialsPort userCredentialsPort;

    private DeskflowUserDetailsService service;

    @BeforeEach
    void setUp() {
        service = new DeskflowUserDetailsService(userCredentialsPort);
    }

    @Test
    void loadUserByUsername_returns_user_details() {
        when(userCredentialsPort.findByEmail("test@company.test"))
                .thenReturn(Optional.of(new UserCredentials("test@company.test", "hashed_pass", "EMPLOYEE")));

        var userDetails = service.loadUserByUsername("test@company.test");

        assertThat(userDetails.getUsername()).isEqualTo("test@company.test");
        assertThat(userDetails.getPassword()).isEqualTo("hashed_pass");
        assertThat(userDetails.getAuthorities()).hasSize(1);
        assertThat(userDetails.getAuthorities().iterator().next().getAuthority()).isEqualTo("ROLE_EMPLOYEE");
    }

    @Test
    void loadUserByUsername_throws_when_not_found() {
        when(userCredentialsPort.findByEmail("missing@company.test")).thenReturn(Optional.empty());

        assertThatThrownBy(() -> service.loadUserByUsername("missing@company.test"))
                .isInstanceOf(UsernameNotFoundException.class)
                .hasMessageContaining("missing@company.test");
    }
}
\end{DeskFlowCode}

\phantomsection
\addcontentsline{toc}{subsubsection}{JwtTokenProviderTest}
\DeskFlowListingCaption{identity/infrastructure/security/JwtTokenProviderTest.java - JwtTokenProviderTest}{lst:identity-infrastructure-security-JwtTokenProviderTest-java}
\begin{DeskFlowCode}
package com.deskflow.identity.infrastructure.security;

import org.junit.jupiter.api.Test;

import static org.assertj.core.api.Assertions.assertThat;

class JwtTokenProviderTest {

    private static final String SECRET = "test-secret-key-at-least-32-characters-long-for-hs256-tests";

    @Test
    void generated_token_contains_email_and_role() {
        JwtTokenProvider provider = new JwtTokenProvider(SECRET, 60_000);

        String token = provider.issueAccessToken("user@test.pl", "IT_ADMIN");

        assertThat(provider.validateToken(token)).isTrue();
        assertThat(provider.getEmailFromToken(token)).isEqualTo("user@test.pl");
        assertThat(provider.getRoleFromToken(token)).isEqualTo("IT_ADMIN");
    }

    @Test
    void invalid_token_is_rejected() {
        JwtTokenProvider provider = new JwtTokenProvider(SECRET, 60_000);

        assertThat(provider.validateToken("not-a-jwt")).isFalse();
    }

    @Test
    void expired_token_is_rejected() throws InterruptedException {
        JwtTokenProvider provider = new JwtTokenProvider(SECRET, 1);
        String token = provider.issueAccessToken("user@test.pl", "EMPLOYEE");

        Thread.sleep(5);

        assertThat(provider.validateToken(token)).isFalse();
    }
}
\end{DeskFlowCode}

\subsection{Testy - Notifications}

\phantomsection
\addcontentsline{toc}{subsubsection}{NotificationAppServiceTest}
\DeskFlowListingCaption{notifications/application/service/NotificationAppServiceTest.java - NotificationAppServiceTest}{lst:notifications-application-service-NotificationAppServiceTest-java}
\begin{DeskFlowCode}
package com.deskflow.notifications.application.service;

import com.deskflow.allocation.application.domain.model.vo.AllocationRequestId;
import com.deskflow.common.application.domain.model.vo.EmployeeId;
import com.deskflow.employees.application.port.in.EmployeeManagementInputPort;
import com.deskflow.employees.application.domain.model.Employee;
import com.deskflow.employees.application.domain.model.vo.Role;
import com.deskflow.notifications.application.port.out.NotificationPort;
import org.junit.jupiter.api.BeforeEach;
import org.junit.jupiter.api.Test;
import org.junit.jupiter.api.extension.ExtendWith;
import org.mockito.Mock;
import org.mockito.junit.jupiter.MockitoExtension;

import java.util.Optional;
import java.util.UUID;

import static org.mockito.Mockito.*;

@ExtendWith(MockitoExtension.class)
class NotificationAppServiceTest {

    @Mock private NotificationPort notifications;
    @Mock private EmployeeManagementInputPort employeeManagement;

    private NotificationAppService appService;
    private AllocationRequestId requestId;
    private EmployeeId employeeId;
    private Employee employee;

    @BeforeEach
    void setUp() {
        requestId = AllocationRequestId.newId();
        employeeId = EmployeeId.of(UUID.randomUUID());
        employee = new Employee(employeeId, "test@company.test", "Jan", "Kowalski", Role.EMPLOYEE, java.util.Map.of());
        appService = new NotificationAppService(notifications, employeeManagement);
    }

    @Test
    void submitted_sends_email() {
        when(employeeManagement.getEmployee(employeeId)).thenReturn(Optional.of(employee));

        appService.allocationRequestSubmitted(requestId, employeeId);

        verify(notifications).send("test@company.test", "Request submitted",
                "Your allocation request has been submitted (ID: " + requestId + ").");
    }

    @Test
    void approved_sends_email() {
        when(employeeManagement.getEmployee(employeeId)).thenReturn(Optional.of(employee));

        appService.allocationRequestApproved(requestId, employeeId);

        verify(notifications).send("test@company.test", "Request approved",
                "Your request has been approved (ID: " + requestId + "). Please collect the resource.");
    }

    @Test
    void rejected_sends_email_with_reason() {
        when(employeeManagement.getEmployee(employeeId)).thenReturn(Optional.of(employee));

        appService.allocationRequestRejected(requestId, employeeId, "Brak miejsca");

        verify(notifications).send("test@company.test", "Request rejected",
                "Your request (ID: " + requestId + ") was rejected. Reason: Brak miejsca");
    }

    @Test
    void cancelled_sends_email() {
        when(employeeManagement.getEmployee(employeeId)).thenReturn(Optional.of(employee));

        appService.allocationRequestCancelled(requestId, employeeId);

        verify(notifications).send("test@company.test", "Request cancelled",
                "Your request (ID: " + requestId + ") has been cancelled.");
    }

    @Test
    void resourceIssued_sends_email() {
        when(employeeManagement.getEmployee(employeeId)).thenReturn(Optional.of(employee));

        appService.resourceIssued(requestId, employeeId);

        verify(notifications).send("test@company.test", "Resource issued",
                "The resource from request (ID: " + requestId + ") has been issued to you.");
    }

    @Test
    void resourceReturned_sends_email() {
        when(employeeManagement.getEmployee(employeeId)).thenReturn(Optional.of(employee));

        appService.resourceReturned(requestId, employeeId);

        verify(notifications).send("test@company.test", "Resource returned",
                "We confirm receipt of the resource from request (ID: " + requestId + ").");
    }

    @Test
    void overdue_sends_email() {
        when(employeeManagement.getEmployee(employeeId)).thenReturn(Optional.of(employee));

        appService.allocationOverdue(requestId, employeeId);

        verify(notifications).send("test@company.test", "Allocation overdue",
                "Your allocation (ID: " + requestId + ") is past its return date. Please return the resource immediately.");
    }

    @Test
    void employee_not_found_skips_email() {
        when(employeeManagement.getEmployee(employeeId)).thenReturn(Optional.empty());

        appService.allocationRequestSubmitted(requestId, employeeId);

        verify(notifications, never()).send(any(), any(), any());
    }

    @Test
    void employee_lookup_failure_does_not_propagate() {
        when(employeeManagement.getEmployee(employeeId)).thenThrow(new RuntimeException("DB down"));

        appService.allocationRequestSubmitted(requestId, employeeId);

        verify(notifications, never()).send(any(), any(), any());
    }

    @Test
    void notification_failure_does_not_propagate() {
        when(employeeManagement.getEmployee(employeeId)).thenReturn(Optional.of(employee));
        doThrow(new RuntimeException("SMTP down")).when(notifications).send(any(), any(), any());

        appService.allocationRequestSubmitted(requestId, employeeId);
    }
}
\end{DeskFlowCode}

\phantomsection
\addcontentsline{toc}{subsubsection}{NotificationEventListenerTest}
\DeskFlowListingCaption{notifications/infrastructure/in/events/NotificationEventListenerTest.java - NotificationEventListenerTest}{lst:notifications-infrastructure-in-events-NotificationEventListenerTest-java}
\begin{DeskFlowCode}
package com.deskflow.notifications.infrastructure.in.events;

import com.deskflow.allocation.application.domain.model.events.*;
import com.deskflow.allocation.application.domain.model.vo.AllocationRequestId;
import com.deskflow.allocation.application.domain.model.vo.DateRange;
import com.deskflow.common.application.domain.model.vo.EmployeeId;
import com.deskflow.common.application.domain.model.vo.ResourceId;
import com.deskflow.notifications.application.port.in.NotificationEventInputPort;
import org.junit.jupiter.api.BeforeEach;
import org.junit.jupiter.api.Test;
import org.junit.jupiter.api.extension.ExtendWith;
import org.mockito.Mock;
import org.mockito.junit.jupiter.MockitoExtension;

import java.time.LocalDate;
import java.time.LocalDateTime;
import java.util.UUID;

import static org.mockito.Mockito.*;

@ExtendWith(MockitoExtension.class)
class NotificationEventListenerTest {

    @Mock private NotificationEventInputPort notificationInputPort;

    private NotificationEventListener listener;
    private AllocationRequestId requestId;
    private EmployeeId employeeId;

    private ResourceId resourceId;
    private DateRange dateRange;
    private LocalDateTime now;

    @BeforeEach
    void setUp() {
        listener = new NotificationEventListener(notificationInputPort);
        requestId = AllocationRequestId.newId();
        employeeId = EmployeeId.of(UUID.randomUUID());
        resourceId = ResourceId.of(UUID.randomUUID());
        dateRange = new DateRange(LocalDate.now(), LocalDate.now().plusDays(1));
        now = LocalDateTime.now();
    }

    @Test
    void onRequestSubmitted_delegates() {
        var event = new AllocationRequestSubmitted(requestId, employeeId, resourceId, dateRange, now);

        listener.onRequestSubmitted(event);

        verify(notificationInputPort).allocationRequestSubmitted(requestId, employeeId);
    }

    @Test
    void onRequestApproved_delegates() {
        var event = new AllocationRequestApproved(requestId, employeeId, now);

        listener.onRequestApproved(event);

        verify(notificationInputPort).allocationRequestApproved(requestId, employeeId);
    }

    @Test
    void onRequestRejected_delegates() {
        var event = new AllocationRequestRejected(requestId, employeeId, "No space", now);

        listener.onRequestRejected(event);

        verify(notificationInputPort).allocationRequestRejected(requestId, employeeId, "No space");
    }

    @Test
    void onRequestCancelled_delegates() {
        var event = new AllocationRequestCancelled(requestId, employeeId, now);

        listener.onRequestCancelled(event);

        verify(notificationInputPort).allocationRequestCancelled(requestId, employeeId);
    }

    @Test
    void onResourceIssued_delegates() {
        var event = new ResourceIssued(requestId, employeeId, resourceId, now);

        listener.onResourceIssued(event);

        verify(notificationInputPort).resourceIssued(requestId, employeeId);
    }

    @Test
    void onResourceReturned_delegates() {
        var event = new ResourceReturned(requestId, employeeId, resourceId, now);

        listener.onResourceReturned(event);

        verify(notificationInputPort).resourceReturned(requestId, employeeId);
    }

    @Test
    void onAllocationOverdue_delegates() {
        var event = new AllocationOverdue(requestId, employeeId, resourceId, now);

        listener.onAllocationOverdue(event);

        verify(notificationInputPort).allocationOverdue(requestId, employeeId);
    }
}
\end{DeskFlowCode}

\phantomsection
\addcontentsline{toc}{subsubsection}{SmtpNotificationAdapterTest}
\DeskFlowListingCaption{notifications/infrastructure/out/smtp/SmtpNotificationAdapterTest.java - SmtpNotificationAdapterTest}{lst:notifications-infrastructure-out-smtp-SmtpNotificationAdapterTest-java}
\begin{DeskFlowCode}
package com.deskflow.notifications.infrastructure.out.smtp;

import org.junit.jupiter.api.Test;
import org.junit.jupiter.api.extension.ExtendWith;
import org.mockito.ArgumentCaptor;
import org.mockito.Captor;
import org.mockito.Mock;
import org.mockito.junit.jupiter.MockitoExtension;
import org.springframework.mail.MailException;
import org.springframework.mail.SimpleMailMessage;
import org.springframework.mail.javamail.JavaMailSender;

import static org.assertj.core.api.Assertions.assertThat;
import static org.mockito.Mockito.*;

@ExtendWith(MockitoExtension.class)
class SmtpNotificationAdapterTest {

    @Mock private JavaMailSender mailSender;
    @Captor private ArgumentCaptor<SimpleMailMessage> messageCaptor;

    @Test
    void sends_email_with_correct_fields() {
        SmtpNotificationAdapter adapter = new SmtpNotificationAdapter(mailSender);

        adapter.send("user@test.com", "Subject", "Body text");

        verify(mailSender).send(messageCaptor.capture());
        SimpleMailMessage msg = messageCaptor.getValue();
        assertThat(msg.getTo()).containsExactly("user@test.com");
        assertThat(msg.getSubject()).isEqualTo("Subject");
        assertThat(msg.getText()).isEqualTo("Body text");
        assertThat(msg.getFrom()).isEqualTo("noreply@deskflow.local");
    }

    @Test
    void mail_exception_is_caught_and_logged() {
        SmtpNotificationAdapter adapter = new SmtpNotificationAdapter(mailSender);
        doThrow(new MailException("SMTP error") {}).when(mailSender).send(any(SimpleMailMessage.class));

        adapter.send("user@test.com", "Subject", "Body");
    }

    @Test
    void sends_to_different_recipients() {
        SmtpNotificationAdapter adapter = new SmtpNotificationAdapter(mailSender);

        adapter.send("alice@test.com", "A", "msg1");
        adapter.send("bob@test.com", "B", "msg2");

        verify(mailSender, times(2)).send(any(SimpleMailMessage.class));
    }
}
\end{DeskFlowCode}

\subsection{Testy - Common}

\phantomsection
\addcontentsline{toc}{subsubsection}{EmployeeIdTest}
\DeskFlowListingCaption{common/application/domain/model/vo/EmployeeIdTest.java - EmployeeIdTest}{lst:common-application-domain-model-vo-EmployeeIdTest-java}
\begin{DeskFlowCode}
package com.deskflow.common.application.domain.model.vo;

import org.junit.jupiter.api.Test;

import java.util.UUID;

import static org.assertj.core.api.Assertions.*;

class EmployeeIdTest {

    @Test
    void creates_from_uuid() {
        UUID uuid = UUID.randomUUID();
        EmployeeId id = EmployeeId.of(uuid);
        assertThat(id.getValue()).isEqualTo(uuid);
    }

    @Test
    void creates_from_string() {
        UUID uuid = UUID.randomUUID();
        EmployeeId id = EmployeeId.of(uuid.toString());
        assertThat(id.getValue()).isEqualTo(uuid);
    }

    @Test
    void rejects_null() {
        assertThatThrownBy(() -> new EmployeeId(null))
                .isInstanceOf(NullPointerException.class);
    }

    @Test
    void rejects_null_via_factory() {
        assertThatThrownBy(() -> EmployeeId.of((UUID) null))
                .isInstanceOf(NullPointerException.class);
    }

    @Test
    void equality() {
        UUID uuid = UUID.randomUUID();
        assertThat(EmployeeId.of(uuid)).isEqualTo(EmployeeId.of(uuid));
    }

    @Test
    void inequality() {
        assertThat(EmployeeId.of(UUID.randomUUID())).isNotEqualTo(EmployeeId.of(UUID.randomUUID()));
    }

    @Test
    void hashCode_is_consistent() {
        UUID uuid = UUID.randomUUID();
        assertThat(EmployeeId.of(uuid)).hasSameHashCodeAs(EmployeeId.of(uuid));
    }

    @Test
    void toString_returns_uuid_string() {
        UUID uuid = UUID.randomUUID();
        assertThat(EmployeeId.of(uuid)).hasToString(uuid.toString());
    }
}
\end{DeskFlowCode}

\phantomsection
\addcontentsline{toc}{subsubsection}{ResourceIdTest}
\DeskFlowListingCaption{common/application/domain/model/vo/ResourceIdTest.java - ResourceIdTest}{lst:common-application-domain-model-vo-ResourceIdTest-java}
\begin{DeskFlowCode}
package com.deskflow.common.application.domain.model.vo;

import org.junit.jupiter.api.Test;

import java.util.UUID;

import static org.assertj.core.api.Assertions.*;

class ResourceIdTest {

    @Test
    void creates_from_uuid() {
        UUID uuid = UUID.randomUUID();
        ResourceId id = ResourceId.of(uuid);
        assertThat(id.getValue()).isEqualTo(uuid);
    }

    @Test
    void creates_from_string() {
        UUID uuid = UUID.randomUUID();
        ResourceId id = ResourceId.of(uuid.toString());
        assertThat(id.getValue()).isEqualTo(uuid);
    }

    @Test
    void rejects_null() {
        assertThatThrownBy(() -> new ResourceId(null))
                .isInstanceOf(NullPointerException.class);
    }

    @Test
    void equality() {
        UUID uuid = UUID.randomUUID();
        assertThat(ResourceId.of(uuid)).isEqualTo(ResourceId.of(uuid));
    }

    @Test
    void inequality() {
        assertThat(ResourceId.of(UUID.randomUUID())).isNotEqualTo(ResourceId.of(UUID.randomUUID()));
    }

    @Test
    void toString_returns_uuid_string() {
        UUID uuid = UUID.randomUUID();
        assertThat(ResourceId.of(uuid)).hasToString(uuid.toString());
    }
}
\end{DeskFlowCode}

\phantomsection
\addcontentsline{toc}{subsubsection}{ResourceTypeTest}
\DeskFlowListingCaption{common/application/domain/model/vo/ResourceTypeTest.java - ResourceTypeTest}{lst:common-application-domain-model-vo-ResourceTypeTest-java}
\begin{DeskFlowCode}
package com.deskflow.common.application.domain.model.vo;

import org.junit.jupiter.api.Test;

import static org.assertj.core.api.Assertions.*;

class ResourceTypeTest {

    @Test
    void has_expected_values() {
        assertThat(ResourceType.values()).containsExactly(
                ResourceType.DESK,
                ResourceType.CONFERENCE_ROOM,
                ResourceType.LAPTOP,
                ResourceType.MONITOR,
                ResourceType.PROJECTOR,
                ResourceType.VEHICLE
        );
    }

    @Test
    void valueOf_returns_correct_enum() {
        assertThat(ResourceType.valueOf("DESK")).isEqualTo(ResourceType.DESK);
        assertThat(ResourceType.valueOf("LAPTOP")).isEqualTo(ResourceType.LAPTOP);
    }

    @Test
    void valueOf_invalid_throws() {
        assertThatThrownBy(() -> ResourceType.valueOf("INVALID"))
                .isInstanceOf(IllegalArgumentException.class);
    }
}
\end{DeskFlowCode}

\phantomsection
\addcontentsline{toc}{subsubsection}{AuthorizationSystemTest}
\DeskFlowListingCaption{common/infrastructure/in/web/AuthorizationSystemTest.java - AuthorizationSystemTest}{lst:common-infrastructure-in-web-AuthorizationSystemTest-java}
\begin{DeskFlowCode}
package com.deskflow.common.infrastructure.in.web;

import com.deskflow.identity.infrastructure.security.JwtTokenProvider;
import org.junit.jupiter.api.BeforeEach;
import org.junit.jupiter.api.Test;
import org.springframework.beans.factory.annotation.Autowired;
import org.springframework.boot.webmvc.test.autoconfigure.AutoConfigureMockMvc;
import org.springframework.boot.test.context.SpringBootTest;
import org.springframework.test.context.ActiveProfiles;
import org.springframework.test.web.servlet.MockMvc;

import java.util.UUID;

import static org.springframework.test.web.servlet.request.MockMvcRequestBuilders.delete;
import static org.springframework.test.web.servlet.request.MockMvcRequestBuilders.get;
import static org.springframework.test.web.servlet.request.MockMvcRequestBuilders.post;
import static org.springframework.test.web.servlet.result.MockMvcResultMatchers.status;

@SpringBootTest
@AutoConfigureMockMvc
@ActiveProfiles("test")
class AuthorizationSystemTest {

    @Autowired
    private MockMvc mockMvc;

    @Autowired
    private JwtTokenProvider tokenProvider;

    private String employeeToken;
    private String managerToken;
    private String adminToken;

    @BeforeEach
    void setUp() {
        employeeToken = tokenProvider.issueAccessToken("test.employee@company.test", "EMPLOYEE");
        managerToken = tokenProvider.issueAccessToken("test.manager@company.test", "MANAGER");
        adminToken = tokenProvider.issueAccessToken("test.admin@company.test", "IT_ADMIN");
    }

    @Test
    void admin_endpoint_requires_authentication() throws Exception {
        mockMvc.perform(get("/api/administration/requests"))
                .andExpect(status().isUnauthorized());
    }

    @Test
    void invalid_token_is_unauthorized() throws Exception {
        mockMvc.perform(get("/api/administration/requests")
                        .header("Authorization", "Bearer invalid-token"))
                .andExpect(status().isUnauthorized());
    }

    @Test
    void employee_role_cannot_access_admin_endpoints() throws Exception {
        mockMvc.perform(get("/api/administration/requests")
                        .header("Authorization", "Bearer " + employeeToken))
                .andExpect(status().isForbidden());
    }

    @Test
    void manager_role_can_access_admin_endpoints() throws Exception {
        mockMvc.perform(get("/api/administration/requests")
                        .header("Authorization", "Bearer " + managerToken))
                .andExpect(status().isOk());
    }

    @Test
    void resources_browse_requires_authentication() throws Exception {
        mockMvc.perform(get("/api/resources"))
                .andExpect(status().isUnauthorized());
    }

    @Test
    void auth_endpoints_are_public_but_still_validate_input() throws Exception {
        mockMvc.perform(post("/api/auth/login")
                        .contentType("application/json")
                        .content("{}"))
                .andExpect(status().isBadRequest());
    }

    @Test
    void manager_role_cannot_delete_resources() throws Exception {
        mockMvc.perform(delete("/api/resources/" + UUID.randomUUID())
                        .header("Authorization", "Bearer " + managerToken))
                .andExpect(status().isForbidden());
    }

    @Test
    void admin_role_is_authorized_to_delete_resources() throws Exception {
        mockMvc.perform(delete("/api/resources/" + UUID.randomUUID())
                        .header("Authorization", "Bearer " + adminToken))
                .andExpect(status().isNotFound());
    }
}
\end{DeskFlowCode}

\phantomsection
\addcontentsline{toc}{subsubsection}{FullAllocationCycleSystemTest}
\DeskFlowListingCaption{common/infrastructure/in/web/FullAllocationCycleSystemTest.java - FullAllocationCycleSystemTest}{lst:common-infrastructure-in-web-FullAllocationCycleSystemTest-java}
\begin{DeskFlowCode}
package com.deskflow.common.infrastructure.in.web;

import com.deskflow.common.application.domain.model.vo.EmployeeId;
import com.deskflow.allocation.application.domain.model.AllocationRequest;
import com.deskflow.allocation.application.domain.service.AllocationRequestFactory;
import com.deskflow.allocation.application.port.out.AllocationRequestRepositoryPort;
import com.deskflow.allocation.application.domain.model.vo.AllocationRequestId;
import com.deskflow.allocation.application.domain.model.vo.DateRange;
import com.deskflow.allocation.application.domain.model.vo.RequestStatus;
import com.deskflow.resources.application.domain.model.OfficeResource;
import com.deskflow.resources.application.domain.service.OfficeResourceFactory;
import com.deskflow.resources.application.port.out.OfficeResourceRepositoryPort;
import com.deskflow.resources.application.domain.model.vo.Location;
import com.deskflow.common.application.domain.model.vo.ResourceType;
import com.deskflow.employees.infrastructure.out.persistence.AllocationLimitEntity;
import com.deskflow.employees.infrastructure.out.persistence.EmployeeEntity;
import com.deskflow.employees.infrastructure.out.persistence.SpringDataEmployeeRepository;
import com.deskflow.identity.infrastructure.security.JwtTokenProvider;
import com.fasterxml.jackson.databind.ObjectMapper;
import org.junit.jupiter.api.BeforeEach;
import org.junit.jupiter.api.Test;
import org.springframework.beans.factory.annotation.Autowired;
import org.springframework.boot.webmvc.test.autoconfigure.AutoConfigureMockMvc;
import org.springframework.boot.test.context.SpringBootTest;
import org.springframework.http.MediaType;
import org.springframework.security.crypto.password.PasswordEncoder;
import org.springframework.test.context.ActiveProfiles;
import org.springframework.test.web.servlet.MockMvc;
import org.springframework.transaction.annotation.Transactional;

import java.time.LocalDate;
import java.util.Map;
import java.util.UUID;

import static org.assertj.core.api.Assertions.assertThat;
import static org.springframework.test.web.servlet.request.MockMvcRequestBuilders.*;
import static org.springframework.test.web.servlet.result.MockMvcResultMatchers.*;

@SpringBootTest
@AutoConfigureMockMvc
@ActiveProfiles("test")
@Transactional
class FullAllocationCycleSystemTest {

    @Autowired MockMvc mockMvc;
    @Autowired ObjectMapper objectMapper;
    @Autowired JwtTokenProvider tokenProvider;
    @Autowired AllocationRequestRepositoryPort requestRepository;
    @Autowired OfficeResourceRepositoryPort resourceRepository;
    @Autowired SpringDataEmployeeRepository employeeRepo;
    @Autowired PasswordEncoder passwordEncoder;

    private String employeeToken;
    private String managerToken;
    private EmployeeId employeeId;
    private UUID resourceUuid;

    @BeforeEach
    void setUp() {
        EmployeeEntity employee = new EmployeeEntity();
        employee.setId(UUID.randomUUID());
        employee.setEmail("test.employee@company.test");
        employee.setFirstName("Test");
        employee.setLastName("Employee");
        employee.setRole("EMPLOYEE");
        employee.setPasswordHash(passwordEncoder.encode("pass"));
        AllocationLimitEntity limit = new AllocationLimitEntity();
        limit.setId(UUID.randomUUID());
        limit.setResourceType("DESK");
        limit.setLimitAmount(3);
        employee.addAllocationLimit(limit);
        employeeRepo.save(employee);
        employeeId = EmployeeId.of(employee.getId());

        EmployeeEntity manager = new EmployeeEntity();
        manager.setId(UUID.randomUUID());
        manager.setEmail("test.kierownik@firma.pl");
        manager.setFirstName("Test");
        manager.setLastName("Kierownik");
        manager.setRole("MANAGER");
        manager.setPasswordHash(passwordEncoder.encode("pass"));
        employeeRepo.save(manager);

        employeeToken = tokenProvider.issueAccessToken("test.employee@company.test", "EMPLOYEE");
        managerToken = tokenProvider.issueAccessToken("test.kierownik@firma.pl", "MANAGER");

        OfficeResource resource = new OfficeResourceFactory().create(
                ResourceType.DESK, "Biurko testowe", new Location("A", 1, "T01")
        );
        resourceRepository.save(resource);
        resourceUuid = resource.getId().getValue();
    }

    @Test
    void full_cycle_submit_approve_issue_return() throws Exception {
        // UC-02: Submit request
        String requestBody = objectMapper.writeValueAsString(Map.of(
                "resourceId", resourceUuid,
                "dateFrom", LocalDate.now().plusDays(1).toString(),
                "dateTo", LocalDate.now().plusDays(5).toString(),
                "justification", "Test systemowy"
        ));

        String responseBody = mockMvc.perform(post("/api/requests")
                        .header("Authorization", "Bearer " + employeeToken)
                        .contentType(MediaType.APPLICATION_JSON)
                        .content(requestBody))
                .andExpect(status().isCreated())
                .andReturn().getResponse().getContentAsString();

        UUID requestId = UUID.fromString(
                objectMapper.readTree(responseBody).get("requestId").asText()
        );

        assertThat(requestRepository.find(AllocationRequestId.of(requestId)).orElseThrow().getStatus())
                .isEqualTo(RequestStatus.SUBMITTED);

        // UC-05: Approve request
        mockMvc.perform(put("/api/administration/requests/" + requestId + "/approve")
                        .header("Authorization", "Bearer " + managerToken))
                .andExpect(status().isOk());
        assertThat(requestRepository.find(AllocationRequestId.of(requestId)).orElseThrow().getStatus())
                .isEqualTo(RequestStatus.APPROVED);

        // UC-07: Issue resource
        mockMvc.perform(put("/api/administration/requests/" + requestId + "/issue")
                        .header("Authorization", "Bearer " + managerToken))
                .andExpect(status().isOk());
        assertThat(requestRepository.find(AllocationRequestId.of(requestId)).orElseThrow().getStatus())
                .isEqualTo(RequestStatus.ACTIVE);

        // UC-08: Return resource
        mockMvc.perform(put("/api/administration/requests/" + requestId + "/return")
                        .header("Authorization", "Bearer " + managerToken))
                .andExpect(status().isOk());
        assertThat(requestRepository.find(AllocationRequestId.of(requestId)).orElseThrow().getStatus())
                .isEqualTo(RequestStatus.COMPLETED);
    }

    @Test
    void employee_can_cancel_submitted_request() throws Exception {
        DateRange dateRange = new DateRange(LocalDate.now().plusDays(10), LocalDate.now().plusDays(15));
        AllocationRequest request = new AllocationRequestFactory().create(
                employeeId, com.deskflow.common.application.domain.model.vo.ResourceId.of(resourceUuid), dateRange, "test"
        );
        requestRepository.save(request);

        mockMvc.perform(put("/api/requests/" + request.getId().getValue() + "/cancel")
                        .header("Authorization", "Bearer " + employeeToken))
                .andExpect(status().isOk());

        assertThat(requestRepository.find(request.getId()).orElseThrow().getStatus())
                .isEqualTo(RequestStatus.CANCELLED);
    }
}
\end{DeskFlowCode}

\phantomsection
\addcontentsline{toc}{subsubsection}{GlobalExceptionHandlerTest}
\DeskFlowListingCaption{common/infrastructure/in/web/GlobalExceptionHandlerTest.java - GlobalExceptionHandlerTest}{lst:common-infrastructure-in-web-GlobalExceptionHandlerTest-java}
\begin{DeskFlowCode}
package com.deskflow.common.infrastructure.in.web;

import org.junit.jupiter.api.BeforeEach;
import org.junit.jupiter.api.Test;
import org.springframework.core.MethodParameter;
import org.springframework.http.HttpStatus;
import org.springframework.http.ResponseEntity;
import org.springframework.validation.BeanPropertyBindingResult;
import org.springframework.web.bind.MethodArgumentNotValidException;

import java.util.Map;
import java.util.NoSuchElementException;

import static org.assertj.core.api.Assertions.*;

class GlobalExceptionHandlerTest {

    private GlobalExceptionHandler handler;

    @BeforeEach
    void setUp() {
        handler = new GlobalExceptionHandler();
    }

    @Test
    void illegalState_returns_409() {
        var ex = new IllegalStateException("Resource is not available");

        ResponseEntity<?> response = handler.handleBusinessRule(ex);

        assertThat(response.getStatusCode()).isEqualTo(HttpStatus.CONFLICT);
        assertThat(response.getBody()).isEqualTo(Map.of("error", "Resource is not available"));
    }

    @Test
    void illegalArgument_returns_400() {
        var ex = new IllegalArgumentException("Invalid resource type");

        ResponseEntity<?> response = handler.handleBadArgument(ex);

        assertThat(response.getStatusCode()).isEqualTo(HttpStatus.BAD_REQUEST);
        assertThat(response.getBody()).isEqualTo(Map.of("error", "Invalid resource type"));
    }

    @Test
    void noSuchElement_returns_404() {
        var ex = new NoSuchElementException("Employee not found");

        ResponseEntity<?> response = handler.handleNotFound(ex);

        assertThat(response.getStatusCode()).isEqualTo(HttpStatus.NOT_FOUND);
        assertThat(response.getBody()).isEqualTo(Map.of("error", "Employee not found"));
    }

    @Test
    void validationError_returns_400_with_details() throws Exception {
        var method = GlobalExceptionHandlerTest.class.getDeclaredMethod("validationError_returns_400_with_details");
        var param = new MethodParameter(method, -1);
        var binding = new BeanPropertyBindingResult(new Object(), "object");
        var ex = new MethodArgumentNotValidException(param, binding);

        ResponseEntity<?> response = handler.handleValidation(ex);

        assertThat(response.getStatusCode()).isEqualTo(HttpStatus.BAD_REQUEST);
        assertThat(response.getBody()).isInstanceOf(Map.class);
        @SuppressWarnings("unchecked")
        Map<String, Object> body = (Map<String, Object>) response.getBody();
        assertThat(body).containsKey("error");
        assertThat(body).containsKey("details");
    }
}
\end{DeskFlowCode}

\phantomsection
\addcontentsline{toc}{subsubsection}{InputValidationSystemTest}
\DeskFlowListingCaption{common/infrastructure/in/web/InputValidationSystemTest.java - InputValidationSystemTest}{lst:common-infrastructure-in-web-InputValidationSystemTest-java}
\begin{DeskFlowCode}
package com.deskflow.common.infrastructure.in.web;

import com.deskflow.employees.infrastructure.out.persistence.AllocationLimitEntity;
import com.deskflow.employees.infrastructure.out.persistence.EmployeeEntity;
import com.deskflow.employees.infrastructure.out.persistence.SpringDataEmployeeRepository;
import com.deskflow.identity.infrastructure.security.JwtTokenProvider;
import org.junit.jupiter.api.BeforeEach;
import org.junit.jupiter.api.Test;
import org.springframework.beans.factory.annotation.Autowired;
import org.springframework.boot.webmvc.test.autoconfigure.AutoConfigureMockMvc;
import org.springframework.boot.test.context.SpringBootTest;
import org.springframework.security.crypto.password.PasswordEncoder;
import org.springframework.test.context.ActiveProfiles;
import org.springframework.test.web.servlet.MockMvc;
import org.springframework.transaction.annotation.Transactional;

import java.util.UUID;

import static org.springframework.test.web.servlet.request.MockMvcRequestBuilders.post;
import static org.springframework.test.web.servlet.request.MockMvcRequestBuilders.put;
import static org.springframework.test.web.servlet.result.MockMvcResultMatchers.status;

@SpringBootTest
@AutoConfigureMockMvc
@ActiveProfiles("test")
@Transactional
class InputValidationSystemTest {

    @Autowired
    private MockMvc mockMvc;

    @Autowired
    private JwtTokenProvider tokenProvider;

    @Autowired
    private SpringDataEmployeeRepository employeeRepo;

    @Autowired
    private PasswordEncoder passwordEncoder;

    private String employeeToken;
    private String managerToken;

    @BeforeEach
    void setUp() {
        EmployeeEntity employee = new EmployeeEntity();
        employee.setId(UUID.randomUUID());
        employee.setEmail("test.employee@company.test");
        employee.setFirstName("Test");
        employee.setLastName("Employee");
        employee.setRole("EMPLOYEE");
        employee.setPasswordHash(passwordEncoder.encode("pass"));
        AllocationLimitEntity limit = new AllocationLimitEntity();
        limit.setId(UUID.randomUUID());
        limit.setResourceType("DESK");
        limit.setLimitAmount(3);
        employee.addAllocationLimit(limit);
        employeeRepo.save(employee);

        employeeToken = tokenProvider.issueAccessToken("test.employee@company.test", "EMPLOYEE");
        managerToken = tokenProvider.issueAccessToken("test.kierownik@firma.pl", "MANAGER");
    }

    @Test
    void submit_request_requires_resource_id() throws Exception {
        String body = "{" +
                "\"dateFrom\":\"2099-01-01\"," +
                "\"dateTo\":\"2099-01-02\"," +
                "\"justification\":\"test\"" +
                "}";

        mockMvc.perform(post("/api/requests")
                        .header("Authorization", "Bearer " + employeeToken)
                        .contentType("application/json")
                        .content(body))
                .andExpect(status().isBadRequest());
    }

    @Test
    void submit_request_requires_date_from() throws Exception {
        String body = "{" +
                "\"resourceId\":\"" + UUID.randomUUID() + "\"," +
                "\"dateTo\":\"2099-01-02\"," +
                "\"justification\":\"test\"" +
                "}";

        mockMvc.perform(post("/api/requests")
                        .header("Authorization", "Bearer " + employeeToken)
                        .contentType("application/json")
                        .content(body))
                .andExpect(status().isBadRequest());
    }

    @Test
    void submit_request_rejects_past_date_from() throws Exception {
        String body = "{" +
                "\"resourceId\":\"" + UUID.randomUUID() + "\"," +
                "\"dateFrom\":\"2000-01-01\"," +
                "\"dateTo\":\"2099-01-02\"," +
                "\"justification\":\"test\"" +
                "}";

        mockMvc.perform(post("/api/requests")
                        .header("Authorization", "Bearer " + employeeToken)
                        .contentType("application/json")
                        .content(body))
                .andExpect(status().isBadRequest());
    }

    @Test
    void submit_request_requires_date_to() throws Exception {
        String body = "{" +
                "\"resourceId\":\"" + UUID.randomUUID() + "\"," +
                "\"dateFrom\":\"2099-01-01\"," +
                "\"justification\":\"test\"" +
                "}";

        mockMvc.perform(post("/api/requests")
                        .header("Authorization", "Bearer " + employeeToken)
                        .contentType("application/json")
                        .content(body))
                .andExpect(status().isBadRequest());
    }

    @Test
    void add_resource_requires_type() throws Exception {
        String body = "{" +
                "\"name\":\"Desk 1\"," +
                "\"building\":\"A\"," +
                "\"floor\":1," +
                "\"number\":\"101\"" +
                "}";

        mockMvc.perform(post("/api/resources")
                        .header("Authorization", "Bearer " + managerToken)
                        .contentType("application/json")
                        .content(body))
                .andExpect(status().isBadRequest());
    }

    @Test
    void add_resource_requires_name() throws Exception {
        String body = "{" +
                "\"type\":\"DESK\"," +
                "\"name\":\"\"," +
                "\"building\":\"A\"," +
                "\"floor\":1," +
                "\"number\":\"101\"" +
                "}";

        mockMvc.perform(post("/api/resources")
                        .header("Authorization", "Bearer " + managerToken)
                        .contentType("application/json")
                        .content(body))
                .andExpect(status().isBadRequest());
    }

    @Test
    void reject_request_requires_reason() throws Exception {
        String body = "{" +
                "\"reason\":\"\"" +
                "}";

        mockMvc.perform(put("/api/administration/requests/" + UUID.randomUUID() + "/reject")
                        .header("Authorization", "Bearer " + managerToken)
                        .contentType("application/json")
                        .content(body))
                .andExpect(status().isBadRequest());
    }
}
\end{DeskFlowCode}

\phantomsection
\addcontentsline{toc}{subsubsection}{RestControllerWebMvcTest}
\DeskFlowListingCaption{common/infrastructure/in/web/RestControllerWebMvcTest.java - RestControllerWebMvcTest}{lst:common-infrastructure-in-web-RestControllerWebMvcTest-java}
\begin{DeskFlowCode}
package com.deskflow.common.infrastructure.in.web;

import com.deskflow.administration.application.domain.model.AdministeredAllocationRequestId;
import com.deskflow.administration.application.domain.model.AdministeredAllocationStatus;
import com.deskflow.administration.application.model.AdministeredAllocationView;
import com.deskflow.administration.application.port.in.AdministrationInputPort;
import com.deskflow.administration.infrastructure.in.web.AdministrationController;
import com.deskflow.allocation.application.model.ResourceCatalogItem;
import com.deskflow.allocation.application.port.in.AllocationRequestInputPort;
import com.deskflow.allocation.infrastructure.in.web.AllocationRequestController;
import com.deskflow.allocation.infrastructure.in.web.CurrentUserResolver;
import com.deskflow.common.application.domain.model.vo.EmployeeId;
import com.deskflow.common.application.domain.model.vo.ResourceId;
import com.deskflow.common.application.domain.model.vo.ResourceType;
import com.deskflow.employees.application.domain.model.Employee;
import com.deskflow.employees.application.domain.model.vo.Role;
import com.deskflow.employees.application.port.in.EmployeeManagementInputPort;
import com.deskflow.employees.infrastructure.in.web.EmployeeAdminController;
import com.deskflow.identity.application.model.AuthenticatedIdentity;
import com.deskflow.identity.application.model.IdentityProfile;
import com.deskflow.identity.application.port.in.IdentityAccessInputPort;
import com.deskflow.identity.infrastructure.in.web.AuthController;
import com.deskflow.identity.infrastructure.security.JwtAuthenticationFilter;
import com.deskflow.identity.infrastructure.security.JwtTokenProvider;
import com.deskflow.identity.infrastructure.security.SecurityConfig;
import com.deskflow.resources.application.port.in.ResourceManagementInputPort;
import com.deskflow.resources.infrastructure.in.web.ResourceController;
import org.junit.jupiter.api.BeforeEach;
import org.junit.jupiter.api.Test;
import org.springframework.beans.factory.annotation.Autowired;
import org.springframework.boot.webmvc.test.autoconfigure.WebMvcTest;
import org.springframework.context.annotation.Import;
import org.springframework.http.MediaType;
import org.springframework.test.context.bean.override.mockito.MockitoBean;
import org.springframework.test.web.servlet.MockMvc;
import org.springframework.test.web.servlet.request.RequestPostProcessor;

import java.time.LocalDate;
import java.time.LocalDateTime;
import java.util.List;
import java.util.Map;
import java.util.NoSuchElementException;
import java.util.Optional;
import java.util.UUID;

import static org.mockito.ArgumentMatchers.any;
import static org.mockito.ArgumentMatchers.eq;
import static org.mockito.Mockito.doThrow;
import static org.mockito.Mockito.verify;
import static org.mockito.Mockito.when;
import static org.springframework.test.web.servlet.request.MockMvcRequestBuilders.delete;
import static org.springframework.test.web.servlet.request.MockMvcRequestBuilders.get;
import static org.springframework.test.web.servlet.request.MockMvcRequestBuilders.post;
import static org.springframework.test.web.servlet.request.MockMvcRequestBuilders.put;
import static org.springframework.security.test.web.servlet.request.SecurityMockMvcRequestPostProcessors.user;
import static org.springframework.test.web.servlet.result.MockMvcResultMatchers.jsonPath;
import static org.springframework.test.web.servlet.result.MockMvcResultMatchers.status;

@WebMvcTest(controllers = {
        AllocationRequestController.class,
        AdministrationController.class,
        EmployeeAdminController.class,
        AuthController.class,
        ResourceController.class
})
@Import({SecurityConfig.class, JwtAuthenticationFilter.class, GlobalExceptionHandler.class})
class RestControllerWebMvcTest {

    @Autowired private MockMvc mockMvc;

    @MockitoBean private JwtTokenProvider tokenProvider;
    @MockitoBean private AllocationRequestInputPort allocationRequestInputPort;
    @MockitoBean private AdministrationInputPort administrationInputPort;
    @MockitoBean private EmployeeManagementInputPort employeeManagementInputPort;
    @MockitoBean private IdentityAccessInputPort identityAccessInputPort;
    @MockitoBean private ResourceManagementInputPort resourceManagementInputPort;
    @MockitoBean private CurrentUserResolver currentUserResolver;

    private final UUID employeeUuid = UUID.fromString("00000000-0000-0000-0000-000000000001");
    private final UUID resourceUuid = UUID.fromString("00000000-0000-0000-0000-000000000002");
    private final UUID requestUuid = UUID.fromString("00000000-0000-0000-0000-000000000003");

    @BeforeEach
    void setUpTokens() {
        token("employee-token", "employee@test.pl", "EMPLOYEE");
        token("manager-token", "manager@test.pl", "MANAGER");
        token("admin-token", "admin@test.pl", "IT_ADMIN");
        when(tokenProvider.validateToken("invalid-token")).thenReturn(false);
    }

    @Test
    void auth_login_returns_identity_with_token() throws Exception {
        when(identityAccessInputPort.login("user@test.pl", "secret"))
                .thenReturn(new AuthenticatedIdentity("jwt", "user@test.pl", "MANAGER", "Jan", "Kowalski", employeeUuid.toString()));

        mockMvc.perform(post("/api/auth/login")
                        .contentType(MediaType.APPLICATION_JSON)
                        .content("{\"email\":\"user@test.pl\",\"password\":\"secret\"}"))
                .andExpect(status().isOk())
                .andExpect(jsonPath("$.token").value("jwt"))
                .andExpect(jsonPath("$.email").value("user@test.pl"))
                .andExpect(jsonPath("$.role").value("MANAGER"));
    }

    @Test
    void auth_me_returns_current_identity_for_valid_token() throws Exception {
        when(identityAccessInputPort.currentIdentity("employee@test.pl"))
                .thenReturn(Optional.of(new IdentityProfile("employee@test.pl", "EMPLOYEE", "Anna", "Nowak", employeeUuid.toString())));

        mockMvc.perform(get("/api/auth/me").with(authenticated("employee-token")))
                .andExpect(status().isOk())
                .andExpect(jsonPath("$.email").value("employee@test.pl"))
                .andExpect(jsonPath("$.token").doesNotExist());
    }

    @Test
    void auth_me_returns_not_found_when_profile_is_missing() throws Exception {
        when(identityAccessInputPort.currentIdentity("employee@test.pl")).thenReturn(Optional.empty());

        mockMvc.perform(get("/api/auth/me").with(authenticated("employee-token")))
                .andExpect(status().isNotFound());
    }

    @Test
    void resource_browsing_requires_authentication() throws Exception {
        mockMvc.perform(get("/api/resources"))
                .andExpect(status().isForbidden());
    }

    @Test
    void resource_browsing_returns_catalog_for_authenticated_user() throws Exception {
        when(allocationRequestInputPort.browseResources()).thenReturn(List.of(
                new ResourceCatalogItem(resourceUuid, "DESK", "Desk 1", "A", 1, "101", true)
        ));

        mockMvc.perform(get("/api/resources").with(authenticated("employee-token")))
                .andExpect(status().isOk())
                .andExpect(jsonPath("$[0].id").value(resourceUuid.toString()))
                .andExpect(jsonPath("$[0].available").value(true));
    }

    @Test
    void employee_cannot_add_resource() throws Exception {
        mockMvc.perform(post("/api/resources")
                        .with(authenticated("employee-token"))
                        .contentType(MediaType.APPLICATION_JSON)
                        .content(resourceBody("DESK", "Desk 1", "A", "101")))
                .andExpect(status().isForbidden());
    }

    @Test
    void manager_can_add_resource() throws Exception {
        when(resourceManagementInputPort.addResource(ResourceType.DESK, "Desk 1", "A", 1, "101"))
                .thenReturn(ResourceId.of(resourceUuid));

        mockMvc.perform(post("/api/resources")
                        .with(authenticated("manager-token"))
                        .contentType(MediaType.APPLICATION_JSON)
                        .content(resourceBody("DESK", "Desk 1", "A", "101")))
                .andExpect(status().isCreated())
                .andExpect(jsonPath("$.resourceId").value(resourceUuid.toString()));

        verify(resourceManagementInputPort).addResource(ResourceType.DESK, "Desk 1", "A", 1, "101");
    }

    @Test
    void add_resource_rejects_blank_building_and_number() throws Exception {
        mockMvc.perform(post("/api/resources")
                        .with(authenticated("manager-token"))
                        .contentType(MediaType.APPLICATION_JSON)
                        .content(resourceBody("DESK", "Desk 1", "", "")))
                .andExpect(status().isBadRequest());
    }

    @Test
    void add_resource_rejects_invalid_type() throws Exception {
        mockMvc.perform(post("/api/resources")
                        .with(authenticated("manager-token"))
                        .contentType(MediaType.APPLICATION_JSON)
                        .content(resourceBody("INVALID", "Desk 1", "A", "101")))
                .andExpect(status().isBadRequest());
    }

    @Test
    void add_resource_rejects_malformed_json() throws Exception {
        mockMvc.perform(post("/api/resources")
                        .with(authenticated("manager-token"))
                        .contentType(MediaType.APPLICATION_JSON)
                        .content("{"))
                .andExpect(status().isBadRequest());
    }

    @Test
    void add_resource_rejects_missing_content_type() throws Exception {
        mockMvc.perform(post("/api/resources")
                        .with(authenticated("manager-token"))
                        .content(resourceBody("DESK", "Desk 1", "A", "101")))
                .andExpect(status().isUnsupportedMediaType());
    }

    @Test
    void manager_can_update_resource_availability() throws Exception {
        mockMvc.perform(put("/api/resources/" + resourceUuid + "/unavailable")
                        .with(authenticated("manager-token")))
                .andExpect(status().isOk());

        mockMvc.perform(put("/api/resources/" + resourceUuid + "/available")
                        .with(authenticated("manager-token")))
                .andExpect(status().isOk());

        verify(resourceManagementInputPort).markUnavailable(ResourceId.of(resourceUuid));
        verify(resourceManagementInputPort).markAvailable(ResourceId.of(resourceUuid));
    }

    @Test
    void resource_path_rejects_invalid_uuid() throws Exception {
        mockMvc.perform(put("/api/resources/not-a-uuid/unavailable")
                        .with(authenticated("manager-token")))
                .andExpect(status().isBadRequest());
    }

    @Test
    void only_it_admin_can_delete_resource() throws Exception {
        mockMvc.perform(delete("/api/resources/" + resourceUuid)
                        .with(authenticated("manager-token")))
                .andExpect(status().isForbidden());

        mockMvc.perform(delete("/api/resources/" + resourceUuid)
                        .with(authenticated("admin-token")))
                .andExpect(status().isOk());

        verify(resourceManagementInputPort).deactivateResource(ResourceId.of(resourceUuid));
    }

    @Test
    void employee_admin_lists_employees_with_limits() throws Exception {
        Employee employee = new Employee(EmployeeId.of(employeeUuid), "employee@test.pl", "Anna", "Nowak", Role.EMPLOYEE,
                Map.of(ResourceType.DESK, 2));
        when(employeeManagementInputPort.listEmployees()).thenReturn(List.of(employee));

        mockMvc.perform(get("/api/administration/employees").with(authenticated("manager-token")))
                .andExpect(status().isOk())
                .andExpect(jsonPath("$[0].id").value(employeeUuid.toString()))
                .andExpect(jsonPath("$[0].limits.DESK").value(2));
    }

    @Test
    void employee_admin_rejects_invalid_limit_update() throws Exception {
        mockMvc.perform(put("/api/administration/employees/" + employeeUuid + "/limits")
                        .with(authenticated("manager-token"))
                        .contentType(MediaType.APPLICATION_JSON)
                        .content("{\"resourceType\":\"DESK\",\"amount\":-1}"))
                .andExpect(status().isBadRequest());

        mockMvc.perform(put("/api/administration/employees/" + employeeUuid + "/limits")
                        .with(authenticated("manager-token"))
                        .contentType(MediaType.APPLICATION_JSON)
                        .content("{\"resourceType\":\"INVALID\",\"amount\":1}"))
                .andExpect(status().isBadRequest());
    }

    @Test
    void employee_admin_updates_limit_and_later_lists_updated_employee() throws Exception {
        mockMvc.perform(put("/api/administration/employees/" + employeeUuid + "/limits")
                        .with(authenticated("manager-token"))
                        .contentType(MediaType.APPLICATION_JSON)
                        .content("{\"resourceType\":\"DESK\",\"amount\":5}"))
                .andExpect(status().isOk());

        verify(employeeManagementInputPort).upsertLimit(EmployeeId.of(employeeUuid), "DESK", 5);

        Employee updatedEmployee = new Employee(EmployeeId.of(employeeUuid), "employee@test.pl", "Anna", "Nowak", Role.EMPLOYEE,
                Map.of(ResourceType.DESK, 5));
        when(employeeManagementInputPort.listEmployees()).thenReturn(List.of(updatedEmployee));

        mockMvc.perform(get("/api/administration/employees").with(authenticated("manager-token")))
                .andExpect(status().isOk())
                .andExpect(jsonPath("$[0].id").value(employeeUuid.toString()))
                .andExpect(jsonPath("$[0].limits.DESK").value(5));
    }

    @Test
    void administration_lists_requests_with_status_filter_and_employee_names() throws Exception {
        when(administrationInputPort.listRequests(AdministeredAllocationStatus.SUBMITTED))
                .thenReturn(List.of(adminView("SUBMITTED")));
        when(employeeManagementInputPort.getEmployee(EmployeeId.of(employeeUuid)))
                .thenReturn(Optional.of(new Employee(EmployeeId.of(employeeUuid), "employee@test.pl", "Anna", "Nowak", Role.EMPLOYEE, Map.of())));

        mockMvc.perform(get("/api/administration/requests?status=SUBMITTED")
                .with(authenticated("manager-token")))
                .andExpect(status().isOk())
                .andExpect(jsonPath("$[0].id").value(requestUuid.toString()))
                .andExpect(jsonPath("$[0].employeeFirstName").value("Anna"));

        verify(administrationInputPort).listRequests(AdministeredAllocationStatus.SUBMITTED);
    }

    @Test
    void administration_rejects_invalid_status_filter() throws Exception {
        mockMvc.perform(get("/api/administration/requests?status=BAD_STATUS")
                        .with(authenticated("manager-token")))
                .andExpect(status().isBadRequest());
    }

    @Test
    void administration_returns_conflict_for_invalid_lifecycle_operation() throws Exception {
        doThrow(new IllegalStateException("Request requires status SUBMITTED"))
                .when(administrationInputPort).approveRequest(AdministeredAllocationRequestId.of(requestUuid));

        mockMvc.perform(put("/api/administration/requests/" + requestUuid + "/approve")
                        .with(authenticated("manager-token")))
                .andExpect(status().isConflict())
                .andExpect(jsonPath("$.error").value("Request requires status SUBMITTED"));
    }

    @Test
    void submit_request_rejects_date_to_before_date_from() throws Exception {
        when(currentUserResolver.resolveEmployeeId(any())).thenReturn(EmployeeId.of(employeeUuid));
        doThrow(new IllegalArgumentException("dateTo must not be before dateFrom"))
                .when(allocationRequestInputPort).submitRequest(
                        eq(EmployeeId.of(employeeUuid)), eq(ResourceId.of(resourceUuid)),
                        eq(LocalDate.of(2099, 1, 10)), eq(LocalDate.of(2099, 1, 1)), eq("test"));

        mockMvc.perform(post("/api/requests")
                        .with(authenticated("employee-token"))
                        .contentType(MediaType.APPLICATION_JSON)
                        .content("{\"resourceId\":\"" + resourceUuid + "\",\"dateFrom\":\"2099-01-10\",\"dateTo\":\"2099-01-01\",\"justification\":\"test\"}"))
                .andExpect(status().isBadRequest())
                .andExpect(jsonPath("$.error").value("dateTo must not be before dateFrom"));
    }

    @Test
    void missing_resource_from_controller_returns_not_found() throws Exception {
        doThrow(new NoSuchElementException("Resource not found"))
                .when(resourceManagementInputPort).updateResource(ResourceId.of(resourceUuid), "Desk 1", "A", 1, "101");

        mockMvc.perform(put("/api/resources/" + resourceUuid)
                        .with(authenticated("manager-token"))
                        .contentType(MediaType.APPLICATION_JSON)
                        .content(resourceBody("DESK", "Desk 1", "A", "101")))
                .andExpect(status().isNotFound())
                .andExpect(jsonPath("$.error").value("Resource not found"));
    }

    private void token(String token, String email, String role) {
        when(tokenProvider.validateToken(token)).thenReturn(true);
        when(tokenProvider.getEmailFromToken(token)).thenReturn(email);
        when(tokenProvider.getRoleFromToken(token)).thenReturn(role);
    }

    private RequestPostProcessor authenticated(String token) {
        return switch (token) {
            case "employee-token" -> user("employee@test.pl").roles("EMPLOYEE");
            case "manager-token" -> user("manager@test.pl").roles("MANAGER");
            case "admin-token" -> user("admin@test.pl").roles("IT_ADMIN");
            default -> user("unknown@test.pl");
        };
    }

    private String resourceBody(String type, String name, String building, String number) {
        return "{\"type\":\"" + type + "\",\"name\":\"" + name + "\",\"building\":\"" + building +
                "\",\"floor\":1,\"number\":\"" + number + "\"}";
    }

    private AdministeredAllocationView adminView(String status) {
        return new AdministeredAllocationView(
                requestUuid,
                employeeUuid,
                resourceUuid,
                status,
                LocalDate.of(2099, 1, 1),
                LocalDate.of(2099, 1, 2),
                "test",
                LocalDateTime.of(2098, 12, 1, 10, 0),
                null,
                null,
                null
        );
    }
}
\end{DeskFlowCode}

\phantomsection
\addcontentsline{toc}{subsubsection}{SpringDomainEventPublisherAdapterTest}
\DeskFlowListingCaption{common/infrastructure/out/events/SpringDomainEventPublisherAdapterTest.java - SpringDomainEventPublisherAdapterTest}{lst:common-infrastructure-out-events-SpringDomainEventPublisherAdapterTest-java}
\begin{DeskFlowCode}
package com.deskflow.common.infrastructure.out.events;

import com.deskflow.common.application.domain.DomainEvent;
import org.junit.jupiter.api.Test;
import org.junit.jupiter.api.extension.ExtendWith;
import org.mockito.Mock;
import org.mockito.junit.jupiter.MockitoExtension;
import org.springframework.context.ApplicationEventPublisher;

import static org.mockito.Mockito.*;

@ExtendWith(MockitoExtension.class)
class SpringDomainEventPublisherAdapterTest {

    @Mock private ApplicationEventPublisher publisher;

    @Test
    void publish_delegates_to_spring_publisher() {
        var adapter = new SpringDomainEventPublisherAdapter(publisher);
        DomainEvent event = mock(DomainEvent.class);

        adapter.publish(event);

        verify(publisher).publishEvent(event);
    }
}
\end{DeskFlowCode}

\clearpage
\section{Pliki konfiguracyjne i migracje}

\phantomsection
\addcontentsline{toc}{subsubsection}{application-test.properties}
\DeskFlowListingCaption{resources/application-test.properties - application-test.properties}{lst:resources-application-test-properties}
\begin{DeskFlowCode}
spring.datasource.url=jdbc:tc:postgresql:17-alpine:///deskflow
spring.datasource.driver-class-name=org.testcontainers.jdbc.ContainerDatabaseDriver
spring.flyway.enabled=true
spring.jpa.hibernate.ddl-auto=validate
spring.mail.host=localhost
spring.mail.port=3025
app.jwt.secret=test-secret-key-at-least-32-characters-long-for-hs256-tests
\end{DeskFlowCode}

\phantomsection
\addcontentsline{toc}{subsubsection}{application.properties}
\DeskFlowListingCaption{resources/application.properties - application.properties}{lst:resources-application-properties}
\begin{DeskFlowCode}
spring.application.name=deskflow-backend

# DataSource
spring.datasource.url=${SPRING_DATASOURCE_URL:jdbc:postgresql://localhost:5432/deskflow}
spring.datasource.username=${SPRING_DATASOURCE_USERNAME:deskflow}
spring.datasource.password=${SPRING_DATASOURCE_PASSWORD:deskflow}
spring.datasource.driver-class-name=org.postgresql.Driver

# JPA
spring.jpa.hibernate.ddl-auto=validate
spring.jpa.show-sql=false
spring.jpa.properties.hibernate.dialect=org.hibernate.dialect.PostgreSQLDialect
spring.jpa.properties.hibernate.format_sql=true

# Flyway
spring.flyway.enabled=true
spring.flyway.locations=classpath:db/migration

# Mail
spring.mail.host=${SPRING_MAIL_HOST:localhost}
spring.mail.port=${SPRING_MAIL_PORT:1025}
spring.mail.username=${SPRING_MAIL_USERNAME:}
spring.mail.password=${SPRING_MAIL_PASSWORD:}
spring.mail.properties.mail.smtp.auth=false
spring.mail.properties.mail.smtp.starttls.enable=false

# JWT
app.jwt.secret=${JWT_SECRET:deskflow-secret-key-at-least-32-characters-long-for-hs256}
app.jwt.expiration-ms=86400000

# CORS
app.cors.allowed-origin=${ALLOWED_ORIGIN:http://localhost:3000}

# Scheduler
spring.task.scheduling.enabled=true

# Actuator
management.endpoints.web.exposure.include=health,info
management.health.mail.enabled=false
\end{DeskFlowCode}

\phantomsection
\addcontentsline{toc}{subsubsection}{V1\_\_init\_schema.sql}
\DeskFlowListingCaption{resources/db/migration/V1\_\_init\_schema.sql - V1\_\_init\_schema.sql}{lst:resources-db-migration-V1-init-schema-sql}
\begin{DeskFlowCode}
CREATE TABLE employees (
    id UUID PRIMARY KEY,
    email VARCHAR(255) NOT NULL UNIQUE,
    first_name VARCHAR(100) NOT NULL,
    last_name VARCHAR(100) NOT NULL,
    role VARCHAR(50) NOT NULL,
    password_hash VARCHAR(255) NOT NULL
);

CREATE TABLE allocation_limits (
    id UUID PRIMARY KEY,
    employee_id UUID NOT NULL REFERENCES employees(id),
    resource_type VARCHAR(50) NOT NULL,
    limit_amount INT NOT NULL DEFAULT 1,
    UNIQUE (employee_id, resource_type)
);

CREATE TABLE office_resources (
    id UUID PRIMARY KEY,
    type VARCHAR(50) NOT NULL,
    name VARCHAR(255) NOT NULL,
    building VARCHAR(100) NOT NULL,
    floor INT NOT NULL,
    number VARCHAR(50) NOT NULL,
    active BOOLEAN NOT NULL DEFAULT TRUE,
    available BOOLEAN NOT NULL DEFAULT TRUE
);

CREATE TABLE allocation_requests (
    id UUID PRIMARY KEY,
    employee_id UUID NOT NULL REFERENCES employees(id),
    resource_id UUID NOT NULL REFERENCES office_resources(id),
    status VARCHAR(50) NOT NULL,
    date_from DATE NOT NULL,
    date_to DATE NOT NULL,
    justification TEXT,
    submitted_at TIMESTAMP NOT NULL,
    issued_at TIMESTAMP,
    returned_at TIMESTAMP,
    rejection_reason TEXT
);

CREATE INDEX idx_allocation_requests_status ON allocation_requests(status);
CREATE INDEX idx_allocation_requests_employee_id ON allocation_requests(employee_id);
CREATE INDEX idx_allocation_requests_resource_id ON allocation_requests(resource_id);
CREATE INDEX idx_allocation_requests_date_to ON allocation_requests(date_to);
\end{DeskFlowCode}

\clearpage
\section{Pliki konfiguracji budowania (pom.xml)}

\phantomsection
\addcontentsline{toc}{subsubsection}{deskflow-backend}
\DeskFlowListingCaption{deskflow-backend/pom.xml - deskflow-backend}{lst:pom-xml}
\begin{DeskFlowCode}
<?xml version="1.0" encoding="UTF-8"?>
<project xmlns="http://maven.apache.org/POM/4.0.0"
         xmlns:xsi="http://www.w3.org/2001/XMLSchema-instance"
         xsi:schemaLocation="http://maven.apache.org/POM/4.0.0 https://maven.apache.org/xsd/maven-4.0.0.xsd">
    <modelVersion>4.0.0</modelVersion>

    <parent>
        <groupId>org.springframework.boot</groupId>
        <artifactId>spring-boot-starter-parent</artifactId>
        <version>4.1.0</version>
        <relativePath/>
    </parent>

    <groupId>com.deskflow</groupId>
    <artifactId>deskflow-backend</artifactId>
    <version>1.0.0</version>
    <name>deskflow-backend</name>
    <description>DeskFlow — zarządzanie zasobami biurowymi (DDD + hexagonal architecture)</description>

    <properties>
        <java.version>21</java.version>
        <archunit.version>1.3.0</archunit.version>
        <testcontainers.version>1.20.3</testcontainers.version>
        <jjwt.version>0.12.6</jjwt.version>
    </properties>

    <dependencies>
        <!- Spring Boot ->
        <dependency>
            <groupId>org.springframework.boot</groupId>
            <artifactId>spring-boot-starter-web</artifactId>
        </dependency>
        <dependency>
            <groupId>org.springframework.boot</groupId>
            <artifactId>spring-boot-starter-data-jpa</artifactId>
        </dependency>
        <dependency>
            <groupId>org.springframework.boot</groupId>
            <artifactId>spring-boot-starter-security</artifactId>
        </dependency>
        <dependency>
            <groupId>org.springframework.boot</groupId>
            <artifactId>spring-boot-starter-mail</artifactId>
        </dependency>
        <dependency>
            <groupId>org.springframework.boot</groupId>
            <artifactId>spring-boot-starter-validation</artifactId>
        </dependency>
        <dependency>
            <groupId>org.springframework.boot</groupId>
            <artifactId>spring-boot-starter-actuator</artifactId>
        </dependency>

        <!- PostgreSQL ->
        <dependency>
            <groupId>org.postgresql</groupId>
            <artifactId>postgresql</artifactId>
            <scope>runtime</scope>
        </dependency>

        <!- JWT ->
        <dependency>
            <groupId>io.jsonwebtoken</groupId>
            <artifactId>jjwt-api</artifactId>
            <version>${jjwt.version}</version>
        </dependency>
        <dependency>
            <groupId>io.jsonwebtoken</groupId>
            <artifactId>jjwt-impl</artifactId>
            <version>${jjwt.version}</version>
            <scope>runtime</scope>
        </dependency>
        <dependency>
            <groupId>io.jsonwebtoken</groupId>
            <artifactId>jjwt-jackson</artifactId>
            <version>${jjwt.version}</version>
            <scope>runtime</scope>
        </dependency>

        <!- Flyway ->
        <dependency>
            <groupId>org.flywaydb</groupId>
            <artifactId>flyway-core</artifactId>
        </dependency>
        <dependency>
            <groupId>org.flywaydb</groupId>
            <artifactId>flyway-database-postgresql</artifactId>
        </dependency>

        <!- Testing ->
        <dependency>
            <groupId>org.springframework.boot</groupId>
            <artifactId>spring-boot-starter-test</artifactId>
            <scope>test</scope>
        </dependency>
        <dependency>
            <groupId>org.springframework.boot</groupId>
            <artifactId>spring-boot-starter-webmvc-test</artifactId>
            <scope>test</scope>
        </dependency>
        <dependency>
            <groupId>org.springframework.boot</groupId>
            <artifactId>spring-boot-starter-data-jpa-test</artifactId>
            <scope>test</scope>
        </dependency>
        <dependency>
            <groupId>org.springframework.security</groupId>
            <artifactId>spring-security-test</artifactId>
            <scope>test</scope>
        </dependency>

        <!- ArchUnit ->
        <dependency>
            <groupId>com.tngtech.archunit</groupId>
            <artifactId>archunit-junit5</artifactId>
            <version>${archunit.version}</version>
            <scope>test</scope>
        </dependency>

        <!- TestContainers ->
        <dependency>
            <groupId>org.testcontainers</groupId>
            <artifactId>junit-jupiter</artifactId>
            <version>${testcontainers.version}</version>
            <scope>test</scope>
        </dependency>
        <dependency>
            <groupId>org.testcontainers</groupId>
            <artifactId>postgresql</artifactId>
            <version>${testcontainers.version}</version>
            <scope>test</scope>
        </dependency>
    </dependencies>

    <build>
        <plugins>
            <plugin>
                <groupId>org.springframework.boot</groupId>
                <artifactId>spring-boot-maven-plugin</artifactId>
                <configuration />
            </plugin>
            <plugin>
                <groupId>org.apache.maven.plugins</groupId>
                <artifactId>maven-surefire-plugin</artifactId>
                <configuration>
                    <argLine>-Dnet.bytebuddy.experimental=true</argLine>
                </configuration>
            </plugin>
        </plugins>
    </build>
</project>
\end{DeskFlowCode}

\section{Skrypty wdrożeniowe i konfiguracja Docker}

\phantomsection
\addcontentsline{toc}{subsubsection}{.dockerignore}
\DeskFlowListingCaption{.dockerignore}{lst:dockerignore}
\begin{DeskFlowCode}
.git
.gitignore
.DS_Store
**/.next
**/node_modules
**/target
**/*.log
**/.env*
**/.idea
**/.vscode
\end{DeskFlowCode}

\phantomsection
\addcontentsline{toc}{subsubsection}{Dockerfile.prod}
\DeskFlowListingCaption{Dockerfile.prod}{lst:Dockerfile-prod}
\begin{DeskFlowCode}
FROM maven:3.9-eclipse-temurin-21 AS backend-build
WORKDIR /app
COPY deskflow-backend/pom.xml .
RUN mvn dependency:go-offline -q
COPY deskflow-backend/src ./src
RUN mvn package -Dmaven.test.skip=true

FROM eclipse-temurin:21-jre-alpine AS backend
RUN apk add -no-cache wget \
    && addgroup -S deskflow \
    && adduser -S deskflow -G deskflow
WORKDIR /app
COPY -from=backend-build /app/target/deskflow-backend-*.jar app.jar
EXPOSE 8080
USER deskflow
ENTRYPOINT ["java", "-jar", "app.jar"]

FROM node:22-alpine AS frontend-build
WORKDIR /app
COPY deskflow-frontend/package.json deskflow-frontend/package-lock.json ./
RUN npm ci
COPY deskflow-frontend .
ARG NEXT_PUBLIC_API_URL=http://localhost:8080
ENV NEXT_PUBLIC_API_URL=$NEXT_PUBLIC_API_URL
ENV NEXT_TELEMETRY_DISABLED=1
RUN npm run build

FROM node:22-alpine AS frontend
WORKDIR /app
ENV NODE_ENV=production
ENV NEXT_TELEMETRY_DISABLED=1
COPY -from=frontend-build /app/.next/standalone ./
COPY -from=frontend-build /app/.next/static ./.next/static
COPY -from=frontend-build /app/public ./public
EXPOSE 3000
USER node
CMD ["node", "server.js"]
\end{DeskFlowCode}

\phantomsection
\addcontentsline{toc}{subsubsection}{deskflow-backend/Dockerfile.dev}
\DeskFlowListingCaption{deskflow-backend/Dockerfile.dev}{lst:deskflow-backend-Dockerfile-dev}
\begin{DeskFlowCode}
FROM maven:3.9-eclipse-temurin-21
WORKDIR /app
COPY pom.xml .
RUN mvn dependency:go-offline -q
COPY src ./src
EXPOSE 8080 5005
CMD ["mvn", "spring-boot:run", \
  "-Dspring-boot.run.jvmArguments=-agentlib:jdwp=transport=dt_socket,server=y,suspend=n,address=*:5005"]
\end{DeskFlowCode}

\phantomsection
\addcontentsline{toc}{subsubsection}{docker-compose.dev.yml}
\DeskFlowListingCaption{docker-compose.dev.yml}{lst:docker-compose-dev-yml}
\begin{DeskFlowCode}
# DeskFlow — lokalne środowisko deweloperskie
# Uruchom: docker compose -f docker-compose.dev.yml up -build

services:
  postgres:
    image: postgres:17-alpine
    container_name: deskflow-postgres
    environment:
      POSTGRES_DB: deskflow
      POSTGRES_USER: deskflow
      POSTGRES_PASSWORD: deskflow
    ports:
      - "5432:5432"
    volumes:
      - postgres_data:/var/lib/postgresql/data
    healthcheck:
      test: ["CMD-SHELL", "pg_isready -U deskflow -d deskflow"]
      interval: 5s
      timeout: 5s
      retries: 10

  mailhog:
    image: mailhog/mailhog:latest
    container_name: deskflow-mailhog
    ports:
      - "1025:1025"
      - "8025:8025"

  backend-dev:
    build:
      context: ./deskflow-backend
      dockerfile: Dockerfile.dev
    container_name: deskflow-backend-dev
    ports:
      - "8080:8080"
      - "5005:5005"
    environment:
      SPRING_DATASOURCE_URL: jdbc:postgresql://postgres:5432/deskflow
      SPRING_DATASOURCE_USERNAME: deskflow
      SPRING_DATASOURCE_PASSWORD: deskflow
      SPRING_MAIL_HOST: mailhog
      SPRING_MAIL_PORT: 1025
      JWT_SECRET: deskflow-dev-secret-key-change-in-production
      ALLOWED_ORIGIN: http://localhost:3000
      SPRING_JPA_SHOW_SQL: "true"
    volumes:
      - ./deskflow-backend:/app
      - /app/target
    depends_on:
      postgres:
        condition: service_healthy

  frontend-dev:
    build:
      context: ./deskflow-frontend
      dockerfile: Dockerfile.dev
    container_name: deskflow-frontend-dev
    ports:
      - "3000:3000"
    environment:
      NEXT_PUBLIC_API_URL: http://localhost:8080
    volumes:
      - ./deskflow-frontend:/app
      - /app/node_modules
      - /app/.next
    depends_on:
      - backend-dev

volumes:
  postgres_data:
\end{DeskFlowCode}

\phantomsection
\addcontentsline{toc}{subsubsection}{docker-compose.prod.yml}
\DeskFlowListingCaption{docker-compose.prod.yml}{lst:docker-compose-prod-yml}
\begin{DeskFlowCode}
# DeskFlow — deployment (Coolify / production)
# Użyj w Coolify jako serwis typu "Docker Compose".
# Zmienne środowiskowe ustaw w panelu Coolify (Stacks → Environment).

services:
  postgres:
    image: postgres:17-alpine
    container_name: deskflow-postgres
    environment:
      POSTGRES_DB: deskflow
      POSTGRES_USER: deskflow
      POSTGRES_PASSWORD: ${POSTGRES_PASSWORD}
    volumes:
      - postgres_data_v2:/var/lib/postgresql/data
    healthcheck:
      test: ["CMD-SHELL", "pg_isready -U deskflow -d deskflow"]
      interval: 5s
      timeout: 5s
      retries: 10
    restart: unless-stopped

  backend:
    build:
      context: .
      dockerfile: Dockerfile.prod
      target: backend
    container_name: deskflow-backend
    environment:
      SPRING_DATASOURCE_URL: jdbc:postgresql://postgres:5432/deskflow
      SPRING_DATASOURCE_USERNAME: deskflow
      SPRING_DATASOURCE_PASSWORD: ${POSTGRES_PASSWORD}
      SPRING_MAIL_HOST: ${SMTP_HOST:-}
      SPRING_MAIL_PORT: ${SMTP_PORT:-}
      SPRING_MAIL_USERNAME: ${SMTP_USERNAME:-}
      SPRING_MAIL_PASSWORD: ${SMTP_PASSWORD:-}
      JWT_SECRET: ${JWT_SECRET}
      ALLOWED_ORIGIN: ${ALLOWED_ORIGIN}
    depends_on:
      postgres:
        condition: service_healthy
    healthcheck:
      test: ["CMD", "wget", "-q", "-spider", "http://localhost:8080/actuator/health"]
      interval: 15s
      timeout: 10s
      retries: 5
      start_period: 60s
    restart: unless-stopped

  frontend:
    build:
      context: .
      dockerfile: Dockerfile.prod
      target: frontend
      args:
        NEXT_PUBLIC_API_URL: ${NEXT_PUBLIC_API_URL}
    container_name: deskflow-frontend
    environment:
      NEXT_PUBLIC_API_URL: ${NEXT_PUBLIC_API_URL}
    depends_on:
      backend:
        condition: service_healthy
    restart: unless-stopped

volumes:
  postgres_data_v2:
\end{DeskFlowCode}

\phantomsection
\addcontentsline{toc}{subsubsection}{.env.example}
\DeskFlowListingCaption{.env.example}{lst:env-example}
\begin{DeskFlowCode}
# ─────────────────────────────────────────────
# DeskFlow — zmienne środowiskowe
# ─────────────────────────────────────────────
# Dla dev: cp .env.example .env.dev
# Dla prod: ustaw w panelu Coolify

# ── Postgres ─────────────────────────────────
POSTGRES_PASSWORD=deskflow

# ── Backend JWT ──────────────────────────────
# Minimum 32 znaki dla HS256
JWT_SECRET=zmien-mnie-na-cos-bardziej-skomplikowanego

# ── Backend CORS ─────────────────────────────
ALLOWED_ORIGIN=http://localhost:3000

# ── SMTP (opcjonalny, MailHog domyślnie) ─────
SMTP_HOST=mailhog
SMTP_PORT=1025
SMTP_USERNAME=
SMTP_PASSWORD=

# ── Frontend ─────────────────────────────────
# URL backendu osiągalny z przeglądarki
NEXT_PUBLIC_API_URL=http://localhost:8080
\end{DeskFlowCode}


% ─────────────────────────────────────────────
%  Kod źródłowy frontendu
% ─────────────────────────────────────────────
\section{Kod źródłowy frontendu}

% ── FRONTEND CONFIG ──


\phantomsection
\addcontentsline{toc}{subsubsection}{deskflow-frontend/Dockerfile.dev}
\DeskFlowListingCaption{deskflow-frontend/Dockerfile.dev}{lst:deskflow-frontend-Dockerfile-dev}
\begin{DeskFlowCode}
FROM node:22-alpine
WORKDIR /app
COPY package.json package-lock.json ./
RUN npm ci
COPY . .
ENV NEXT_TELEMETRY_DISABLED=1
EXPOSE 3000
CMD ["npm", "run", "dev"]
\end{DeskFlowCode}

\phantomsection
\addcontentsline{toc}{subsubsection}{deskflow-frontend/Dockerfile.frontend}
\DeskFlowListingCaption{deskflow-frontend/Dockerfile.frontend}{lst:deskflow-frontend-Dockerfile-frontend}
\begin{DeskFlowCode}
FROM node:22-alpine AS build
WORKDIR /app
COPY package.json package-lock.json ./
RUN npm ci
COPY . .
ARG NEXT_PUBLIC_API_URL=http://localhost:8080
ENV NEXT_PUBLIC_API_URL=$NEXT_PUBLIC_API_URL
ENV NEXT_TELEMETRY_DISABLED=1
RUN npm run build

FROM node:22-alpine AS runner
WORKDIR /app
ENV NODE_ENV=production
ENV NEXT_TELEMETRY_DISABLED=1
COPY -from=build /app/.next/standalone ./
COPY -from=build /app/.next/static ./.next/static
COPY -from=build /app/public ./public
EXPOSE 3000
USER node
CMD ["node", "server.js"]
\end{DeskFlowCode}

\phantomsection
\addcontentsline{toc}{subsubsection}{deskflow-frontend/docker-compose.yml}
\DeskFlowListingCaption{deskflow-frontend/docker-compose.yml}{lst:deskflow-frontend-docker-compose-yml}
\begin{DeskFlowCode}
services:
  db:
    image: postgres:17-alpine
    container_name: deskflow-db
    restart: unless-stopped
    environment:
      POSTGRES_USER: postgres
      POSTGRES_PASSWORD: postgres
      POSTGRES_DB: deskflow
    ports:
      - '5432:5432'
    volumes:
      - deskflow-db-data:/var/lib/postgresql/data
    healthcheck:
      test: ['CMD-SHELL', 'pg_isready -U postgres -d deskflow']
      interval: 5s
      timeout: 5s
      retries: 10

  app:
    build:
      context: .
      dockerfile: Dockerfile.frontend
    container_name: deskflow-app
    restart: unless-stopped
    depends_on:
      db:
        condition: service_healthy
    env_file:
      - .env.local
    environment:
      NODE_ENV: development
      DATABASE_URL: postgres://postgres:postgres@db:5432/deskflow
      NEXT_TELEMETRY_DISABLED: '1'
    ports:
      - '3000:3000'

volumes:
  deskflow-db-data:
\end{DeskFlowCode}

\phantomsection
\addcontentsline{toc}{subsubsection}{deskflow-frontend/.gitignore}
\DeskFlowListingCaption{deskflow-frontend/.gitignore}{lst:deskflow-frontend-gitignore}
\begin{DeskFlowCode}
# v0 runtime files
__v0_runtime_loader.js
__v0_devtools.tsx
__v0_jsx-dev-runtime.ts

# Common ignores
node_modules/
.next/
.env*.local
.DS_Store

tsconfig.tsbuildinfo
# clerk configuration (can include secrets)
/.clerk/
\end{DeskFlowCode}

\phantomsection
\addcontentsline{toc}{subsubsection}{deskflow-frontend/.dockerignore}
\DeskFlowListingCaption{deskflow-frontend/.dockerignore}{lst:deskflow-frontend-dockerignore}
\begin{DeskFlowCode}
node_modules
.next
.git
.gitignore
pnpm-debug.log*
*.log
.vscode
\end{DeskFlowCode}

\phantomsection
\addcontentsline{toc}{subsubsection}{deskflow-frontend/package.json}
\DeskFlowListingCaption{deskflow-frontend/package.json}{lst:deskflow-frontend-package-json}
\begin{DeskFlowCode}
{
  "name": "deskflow-frontend",
  "version": "1.0.0",
  "private": true,
  "scripts": {
    "dev": "next dev",
    "build": "next build",
    "start": "next start",
    "lint": "eslint ."
  },
  "dependencies": {
    "@hookform/resolvers": "^3.9.1",
    "@radix-ui/react-accordion": "1.2.12",
    "@radix-ui/react-alert-dialog": "1.1.15",
    "@radix-ui/react-avatar": "1.1.11",
    "@radix-ui/react-checkbox": "1.3.3",
    "@radix-ui/react-dialog": "1.1.15",
    "@radix-ui/react-dropdown-menu": "2.1.16",
    "@radix-ui/react-label": "2.1.8",
    "@radix-ui/react-popover": "1.1.15",
    "@radix-ui/react-progress": "1.1.8",
    "@radix-ui/react-radio-group": "1.3.8",
    "@radix-ui/react-scroll-area": "1.2.10",
    "@radix-ui/react-select": "2.2.6",
    "@radix-ui/react-separator": "1.1.8",
    "@radix-ui/react-slot": "1.2.4",
    "@radix-ui/react-tabs": "1.1.13",
    "@radix-ui/react-toast": "1.2.15",
    "@radix-ui/react-tooltip": "1.2.8",
    "class-variance-authority": "^0.7.1",
    "clsx": "^2.1.1",
    "cmdk": "1.1.1",
    "date-fns": "4.1.0",
    "lucide-react": "^0.564.0",
    "next": "16.2.9",
    "next-themes": "^0.4.6",
    "react": "^19.0.0",
    "react-dom": "^19.0.0",
    "react-hook-form": "^7.54.1",
    "recharts": "2.15.0",
    "sonner": "^1.7.1",
    "tailwind-merge": "^3.3.1",
    "zod": "^3.24.1"
  },
  "devDependencies": {
    "@tailwindcss/postcss": "^4.2.0",
    "@types/node": "^22",
    "@types/react": "^19",
    "@types/react-dom": "^19",
    "eslint": "^9.39.2",
    "eslint-config-next": "16.2.9",
    "postcss": "^8.5",
    "tailwindcss": "^4.2.0",
    "tw-animate-css": "1.3.3",
    "typescript": "5.7.3"
  }
}
\end{DeskFlowCode}

\phantomsection
\addcontentsline{toc}{subsubsection}{deskflow-frontend/tsconfig.json}
\DeskFlowListingCaption{deskflow-frontend/tsconfig.json}{lst:deskflow-frontend-tsconfig-json}
\begin{DeskFlowCode}
{
  "compilerOptions": {
    "lib": [
      "dom",
      "dom.iterable",
      "esnext"
    ],
    "allowJs": true,
    "target": "ES6",
    "skipLibCheck": true,
    "strict": true,
    "noEmit": true,
    "esModuleInterop": true,
    "module": "esnext",
    "moduleResolution": "bundler",
    "resolveJsonModule": true,
    "isolatedModules": true,
    "jsx": "react-jsx",
    "incremental": true,
    "plugins": [
      {
        "name": "next"
      }
    ],
    "paths": {
      "@/*": [
        "./*"
      ]
    }
  },
  "include": [
    "next-env.d.ts",
    "**/*.ts",
    "**/*.tsx",
    ".next/types/**/*.ts",
    ".next/dev/types/**/*.ts"
  ],
  "exclude": [
    "node_modules"
  ]
}
\end{DeskFlowCode}

\phantomsection
\addcontentsline{toc}{subsubsection}{deskflow-frontend/next.config.mjs}
\DeskFlowListingCaption{deskflow-frontend/next.config.mjs}{lst:deskflow-frontend-next-config-mjs}
\begin{DeskFlowCode}
/** @type {import('next').NextConfig} */
const nextConfig = {
  output: 'standalone',
  typescript: {
    ignoreBuildErrors: true,
  },
  images: {
    unoptimized: true,
  },
}

export default nextConfig
\end{DeskFlowCode}

\phantomsection
\addcontentsline{toc}{subsubsection}{deskflow-frontend/eslint.config.mjs}
\DeskFlowListingCaption{deskflow-frontend/eslint.config.mjs}{lst:deskflow-frontend-eslint-config-mjs}
\begin{DeskFlowCode}
import nextVitals from 'eslint-config-next/core-web-vitals'

const eslintConfig = [
  ...nextVitals,
  {
    rules: {
      'react-hooks/immutability': 'off',
      'react-hooks/purity': 'off',
      'react-hooks/set-state-in-effect': 'off',
      'react/no-unescaped-entities': 'off',
    },
  },
]

export default eslintConfig
\end{DeskFlowCode}

\phantomsection
\addcontentsline{toc}{subsubsection}{deskflow-frontend/postcss.config.mjs}
\DeskFlowListingCaption{deskflow-frontend/postcss.config.mjs}{lst:deskflow-frontend-postcss-config-mjs}
\begin{DeskFlowCode}
/** @type {import('postcss-load-config').Config} */
const config = {
  plugins: {
    '@tailwindcss/postcss': {},
  },
}

export default config
\end{DeskFlowCode}

\phantomsection
\addcontentsline{toc}{subsubsection}{deskflow-frontend/components.json}
\DeskFlowListingCaption{deskflow-frontend/components.json}{lst:deskflow-frontend-components-json}
\begin{DeskFlowCode}
{
  "$schema": "https://ui.shadcn.com/schema.json",
  "style": "new-york",
  "rsc": true,
  "tsx": true,
  "tailwind": {
    "config": "",
    "css": "app/globals.css",
    "baseColor": "neutral",
    "cssVariables": true,
    "prefix": ""
  },
  "aliases": {
    "components": "@/components",
    "utils": "@/lib/utils",
    "ui": "@/components/ui",
    "lib": "@/lib",
    "hooks": "@/hooks"
  },
  "iconLibrary": "lucide"
}
\end{DeskFlowCode}

\phantomsection
\addcontentsline{toc}{subsubsection}{deskflow-frontend/next-env.d.ts}
\DeskFlowListingCaption{deskflow-frontend/next-env.d.ts}{lst:deskflow-frontend-next-env-dts}
\begin{DeskFlowCode}
/// <reference types="next" />
/// <reference types="next/image-types/global" />

// NOTE: This file should not be edited
// see https://nextjs.org/docs/app/api-reference/config/typescript for more information.
\end{DeskFlowCode}

% ── APP ──

\phantomsection
\addcontentsline{toc}{subsubsection}{deskflow\-frontend/app/layout.tsx}
\DeskFlowListingCaption{deskflow-frontend/app/layout.tsx}{lst:fd-app-layout}
\begin{DeskFlowCode}
import type { Metadata } from 'next'
import { Toaster } from '@/components/ui/toaster'
import { AuthProvider } from '@/lib/auth-context'
import './globals.css'

export const metadata: Metadata = {
  title: 'DeskFlow - Zarządzanie Zasobami Biurowymi',
  description: 'System udostępniania zasobów biurowych — biurka, sale, sprzęt',
  icons: {
    icon: '/favicon.ico',
  },
}

export default function RootLayout({
  children,
}: Readonly<{
  children: React.ReactNode
}>) {
  return (
    <html lang="pl">
      <body className="font-sans antialiased">
        <AuthProvider>
          {children}
          <Toaster />
        </AuthProvider>
      </body>
    </html>
  )
}
\end{DeskFlowCode}

\phantomsection
\addcontentsline{toc}{subsubsection}{deskflow\-frontend/app/globals.css}
\DeskFlowListingCaption{deskflow-frontend/app/globals.css}{lst:fd-app-globals.css}
\begin{DeskFlowCode}
@import 'tailwindcss';
@import 'tw-animate-css';

@custom-variant dark (&:is(.dark *));

:root {
  -background: oklch(0.985 0 0);
  -foreground: oklch(0.145 0 0);
  -card: oklch(1 0 0);
  -card-foreground: oklch(0.145 0 0);
  -popover: oklch(1 0 0);
  -popover-foreground: oklch(0.145 0 0);
  -primary: oklch(0.55 0.2 250);
  -primary-foreground: oklch(1 0 0);
  -secondary: oklch(0.96 0.01 250);
  -secondary-foreground: oklch(0.25 0.05 250);
  -muted: oklch(0.96 0 0);
  -muted-foreground: oklch(0.5 0 0);
  -accent: oklch(0.7 0.15 165);
  -accent-foreground: oklch(0.2 0.05 165);
  -destructive: oklch(0.55 0.22 27);
  -destructive-foreground: oklch(1 0 0);
  -border: oklch(0.9 0 0);
  -input: oklch(0.9 0 0);
  -ring: oklch(0.55 0.2 250);
  -chart-1: oklch(0.55 0.2 250);
  -chart-2: oklch(0.7 0.15 165);
  -chart-3: oklch(0.65 0.18 45);
  -chart-4: oklch(0.6 0.15 300);
  -chart-5: oklch(0.55 0.2 15);
  -radius: 0.625rem;
  -sidebar: oklch(0.15 0.03 250);
  -sidebar-foreground: oklch(0.95 0 0);
  -sidebar-primary: oklch(0.55 0.2 250);
  -sidebar-primary-foreground: oklch(1 0 0);
  -sidebar-accent: oklch(0.25 0.04 250);
  -sidebar-accent-foreground: oklch(0.95 0 0);
  -sidebar-border: oklch(0.25 0.03 250);
  -sidebar-ring: oklch(0.55 0.2 250);
}

.dark {
  -background: oklch(0.12 0.02 250);
  -foreground: oklch(0.95 0 0);
  -card: oklch(0.15 0.02 250);
  -card-foreground: oklch(0.95 0 0);
  -popover: oklch(0.15 0.02 250);
  -popover-foreground: oklch(0.95 0 0);
  -primary: oklch(0.6 0.2 250);
  -primary-foreground: oklch(1 0 0);
  -secondary: oklch(0.25 0.03 250);
  -secondary-foreground: oklch(0.95 0 0);
  -muted: oklch(0.2 0.02 250);
  -muted-foreground: oklch(0.65 0 0);
  -accent: oklch(0.65 0.15 165);
  -accent-foreground: oklch(0.1 0.02 165);
  -destructive: oklch(0.55 0.22 27);
  -destructive-foreground: oklch(1 0 0);
  -border: oklch(0.25 0.02 250);
  -input: oklch(0.25 0.02 250);
  -ring: oklch(0.6 0.2 250);
  -chart-1: oklch(0.6 0.2 250);
  -chart-2: oklch(0.65 0.15 165);
  -chart-3: oklch(0.7 0.18 45);
  -chart-4: oklch(0.65 0.15 300);
  -chart-5: oklch(0.6 0.2 15);
  -sidebar: oklch(0.12 0.02 250);
  -sidebar-foreground: oklch(0.95 0 0);
  -sidebar-primary: oklch(0.6 0.2 250);
  -sidebar-primary-foreground: oklch(1 0 0);
  -sidebar-accent: oklch(0.2 0.03 250);
  -sidebar-accent-foreground: oklch(0.95 0 0);
  -sidebar-border: oklch(0.2 0.02 250);
  -sidebar-ring: oklch(0.6 0.2 250);
}

@theme inline {
  -font-sans: 'Geist', 'Geist Fallback';
  -font-mono: 'Geist Mono', 'Geist Mono Fallback';
  -color-background: var(-background);
  -color-foreground: var(-foreground);
  -color-card: var(-card);
  -color-card-foreground: var(-card-foreground);
  -color-popover: var(-popover);
  -color-popover-foreground: var(-popover-foreground);
  -color-primary: var(-primary);
  -color-primary-foreground: var(-primary-foreground);
  -color-secondary: var(-secondary);
  -color-secondary-foreground: var(-secondary-foreground);
  -color-muted: var(-muted);
  -color-muted-foreground: var(-muted-foreground);
  -color-accent: var(-accent);
  -color-accent-foreground: var(-accent-foreground);
  -color-destructive: var(-destructive);
  -color-destructive-foreground: var(-destructive-foreground);
  -color-border: var(-border);
  -color-input: var(-input);
  -color-ring: var(-ring);
  -color-chart-1: var(-chart-1);
  -color-chart-2: var(-chart-2);
  -color-chart-3: var(-chart-3);
  -color-chart-4: var(-chart-4);
  -color-chart-5: var(-chart-5);
  -radius-sm: calc(var(-radius) - 4px);
  -radius-md: calc(var(-radius) - 2px);
  -radius-lg: var(-radius);
  -radius-xl: calc(var(-radius) + 4px);
  -color-sidebar: var(-sidebar);
  -color-sidebar-foreground: var(-sidebar-foreground);
  -color-sidebar-primary: var(-sidebar-primary);
  -color-sidebar-primary-foreground: var(-sidebar-primary-foreground);
  -color-sidebar-accent: var(-sidebar-accent);
  -color-sidebar-accent-foreground: var(-sidebar-accent-foreground);
  -color-sidebar-border: var(-sidebar-border);
  -color-sidebar-ring: var(-sidebar-ring);
}

/* Dynamic Company Branding Variables */
:root {
  /* Floor plan element colors - these will be updated by JavaScript */
  -desk-available: #10b981;
  -desk-reserved: #3b82f6;  
  -desk-occupied: #ef4444;
  -desk-maintenance: #f97316;
  
  -room-available: #10b981;
  -room-reserved: #3b82f6;
  -room-occupied: #ef4444;
  -room-maintenance: #f97316;
  
  /* Company colors - these will be updated by the BrandingContext */
  -company-primary: #3b82f6;
  -company-secondary: #10b981;
  -company-text: #111827;
  -company-active-button-text: #ffffff;
  -company-primary-rgb: 59, 130, 246;
  -company-secondary-rgb: 16, 185, 129;
  -company-text-rgb: 17, 24, 39;
  -company-active-button-text-rgb: 255, 255, 255;
}

/* Utility classes for dynamic theming */
.company-primary {
  background-color: var(-company-primary);
}

.company-secondary {
  background-color: var(-company-secondary);
}

.text-company-primary {
  color: var(-company-primary);
}

.text-company-secondary {
  color: var(-company-secondary);
}

.text-company-text {
  color: var(-company-text);
}

.text-company-active-button {
  color: var(-company-active-button-text);
}

.border-company-primary {
  border-color: var(-company-primary);
}

.border-company-secondary {
  border-color: var(-company-secondary);
}

/* Floor plan specific styles */
.floor-plan-desk-available {
  background-color: var(-desk-available);
}

.floor-plan-desk-reserved {
  background-color: var(-desk-reserved);
}

.floor-plan-desk-occupied {
  background-color: var(-desk-occupied);
}

.floor-plan-room-available {
  background-color: var(-room-available);
}

.floor-plan-room-reserved {
  background-color: var(-room-reserved);
}

.floor-plan-room-occupied {
  background-color: var(-room-occupied);
}

@layer base {
  * {
    @apply border-border outline-ring/50;
  }
  body {
    @apply bg-background text-foreground;
  }
}
\end{DeskFlowCode}

\phantomsection
\addcontentsline{toc}{subsubsection}{deskflow\-frontend/app/(auth)/login/page.tsx}
\DeskFlowListingCaption{deskflow-frontend/app/(auth)/login/page.tsx}{lst:fd-app-auth-login-page}
\begin{DeskFlowCode}
'use client';

import { useState } from 'react';
import { useRouter } from 'next/navigation';
import { Building2 } from 'lucide-react';
import { useAuth } from '@/lib/auth-context';
import { authApi } from '@/lib/api';
import { Button } from '@/components/ui/button';
import { Input } from '@/components/ui/input';
import { Card, CardContent, CardDescription, CardHeader, CardTitle } from '@/components/ui/card';
import { Label } from '@/components/ui/label';
import { useToast } from '@/hooks/use-toast';

export default function LoginPage() {
  const [email, setEmail] = useState('');
  const [password, setPassword] = useState('');
  const [loading, setLoading] = useState(false);
  const { login } = useAuth();
  const router = useRouter();
  const { toast } = useToast();

  const handleSubmit = async (e: React.FormEvent) => {
    e.preventDefault();
    setLoading(true);
    try {
      const res = await authApi.login(email, password);
      login(res.token, {
        email: res.email,
        role: res.role,
        firstName: res.firstName,
        lastName: res.lastName,
        employeeId: res.employeeId,
      });
      router.push('/');
    } catch (err: unknown) {
      toast({
        title: 'Błąd logowania',
        description: err instanceof Error ? err.message : 'Nieprawidłowe dane logowania',
        variant: 'destructive',
      });
    } finally {
      setLoading(false);
    }
  };

  return (
    <div className="min-h-screen flex items-center justify-center bg-gray-50 px-4">
      <Card className="w-full max-w-md">
        <CardHeader className="text-center space-y-2">
          <div className="flex justify-center">
            <Building2 className="h-10 w-10 text-blue-600" />
          </div>
          <CardTitle className="text-2xl font-bold">DeskFlow</CardTitle>
          <CardDescription>Zarządzanie zasobami biurowymi</CardDescription>
        </CardHeader>
        <CardContent>
          <form onSubmit={handleSubmit} className="space-y-4">
            <div className="space-y-2">
              <Label htmlFor="email">Email</Label>
              <Input
                id="email"
                type="email"
                value={email}
                onChange={e => setEmail(e.target.value)}
                placeholder="jan.kowalski@firma.pl"
                required
                autoFocus
              />
            </div>
            <div className="space-y-2">
              <Label htmlFor="password">Hasło</Label>
              <Input
                id="password"
                type="password"
                value={password}
                onChange={e => setPassword(e.target.value)}
                required
              />
            </div>
            <Button type="submit" className="w-full" disabled={loading}>
              {loading ? 'Logowanie...' : 'Zaloguj się'}
            </Button>
          </form>
          <div className="mt-6 p-4 bg-slate-50 border rounded-lg text-sm text-slate-600">
            <p className="font-semibold mb-2">Dane testowe:</p>
            <div className="space-y-1">
              <p>�� Pracownik: anna.kowalska@firma.pl / password123</p>
              <p>�� Kierownik: marek.kierownik@firma.pl / password123</p>
              <p>⚙️ Admin IT: admin@firma.pl / admin123</p>
            </div>
          </div>
        </CardContent>
      </Card>
    </div>
  );
}
\end{DeskFlowCode}

\phantomsection
\addcontentsline{toc}{subsubsection}{deskflow\-frontend/app/(auth)/register/page.tsx}
\DeskFlowListingCaption{deskflow-frontend/app/(auth)/register/page.tsx}{lst:fd-app-auth-register-page}
\begin{DeskFlowCode}
"use client"

import { useState } from "react"
import Link from "next/link"
import { Building2, Eye, EyeOff } from "lucide-react"
import { Button } from "@/components/ui/button"
import { Input } from "@/components/ui/input"
import { Card, CardContent, CardDescription, CardHeader, CardTitle } from "@/components/ui/card"
import { Checkbox } from "@/components/ui/checkbox"
import { Field, FieldGroup, FieldLabel } from "@/components/ui/field"
import { Select, SelectContent, SelectItem, SelectTrigger, SelectValue } from "@/components/ui/select"

export default function RegisterPage() {
  const [showPassword, setShowPassword] = useState(false)

  return (
    <div className="min-h-screen flex items-center justify-center bg-background p-4">
      <div className="w-full max-w-md">
        <div className="flex flex-col items-center mb-8">
          <div className="flex h-14 w-14 items-center justify-center rounded-xl bg-primary mb-4">
            <Building2 className="h-8 w-8 text-primary-foreground" />
          </div>
          <h1 className="text-2xl font-bold">DeskFlow</h1>
          <p className="text-muted-foreground">System rezerwacji biura</p>
        </div>

        <Card>
          <CardHeader className="text-center">
            <CardTitle>Utworz konto</CardTitle>
            <CardDescription>
              Wypelnij formularz, aby zarejestrowac sie w systemie
            </CardDescription>
          </CardHeader>
          <CardContent>
            <form className="space-y-4">
              <div className="grid grid-cols-2 gap-4">
                <FieldGroup>
                  <Field>
                    <FieldLabel>Imie</FieldLabel>
                    <Input placeholder="Jan" />
                  </Field>
                </FieldGroup>
                <FieldGroup>
                  <Field>
                    <FieldLabel>Nazwisko</FieldLabel>
                    <Input placeholder="Kowalski" />
                  </Field>
                </FieldGroup>
              </div>

              <FieldGroup>
                <Field>
                  <FieldLabel>Email sluzbowy</FieldLabel>
                  <Input
                    type="email"
                    placeholder="jan.kowalski@firma.pl"
                  />
                </Field>
              </FieldGroup>

              <FieldGroup>
                <Field>
                  <FieldLabel>Dzial</FieldLabel>
                  <Select>
                    <SelectTrigger>
                      <SelectValue placeholder="Wybierz dzial" />
                    </SelectTrigger>
                    <SelectContent>
                      <SelectItem value="it">IT</SelectItem>
                      <SelectItem value="hr">HR</SelectItem>
                      <SelectItem value="sales">Sprzedaz</SelectItem>
                      <SelectItem value="marketing">Marketing</SelectItem>
                      <SelectItem value="finance">Finanse</SelectItem>
                      <SelectItem value="operations">Operacje</SelectItem>
                    </SelectContent>
                  </Select>
                </Field>
              </FieldGroup>

              <FieldGroup>
                <Field>
                  <FieldLabel>Haslo</FieldLabel>
                  <div className="relative">
                    <Input
                      type={showPassword ? "text" : "password"}
                      placeholder="Minimum 8 znakow"
                      className="pr-10"
                    />
                    <button
                      type="button"
                      onClick={() => setShowPassword(!showPassword)}
                      className="absolute right-3 top-1/2 -translate-y-1/2 text-muted-foreground hover:text-foreground"
                    >
                      {showPassword ? (
                        <EyeOff className="h-4 w-4" />
                      ) : (
                        <Eye className="h-4 w-4" />
                      )}
                    </button>
                  </div>
                </Field>
              </FieldGroup>

              <FieldGroup>
                <Field>
                  <FieldLabel>Powtorz haslo</FieldLabel>
                  <Input type="password" placeholder="Powtorz haslo" />
                </Field>
              </FieldGroup>

              <div className="flex items-start gap-2">
                <Checkbox id="terms" className="mt-1" />
                <label
                  htmlFor="terms"
                  className="text-sm text-muted-foreground cursor-pointer"
                >
                  Akceptuje{" "}
                  <Link href="#" className="text-primary hover:underline">
                    Regulamin
                  </Link>{" "}
                  oraz{" "}
                  <Link href="#" className="text-primary hover:underline">
                    Polityke prywatnosci
                  </Link>
                </label>
              </div>

              <Link href="/">
                <Button type="button" className="w-full">
                  Zarejestruj sie
                </Button>
              </Link>
            </form>

            <div className="mt-6 text-center text-sm">
              <span className="text-muted-foreground">Masz juz konto? </span>
              <Link href="/login" className="text-primary hover:underline font-medium">
                Zaloguj sie
              </Link>
            </div>
          </CardContent>
        </Card>
      </div>
    </div>
  )
}
\end{DeskFlowCode}

\phantomsection
\addcontentsline{toc}{subsubsection}{deskflow\-frontend/app/(dashboard)/layout.tsx}
\DeskFlowListingCaption{deskflow-frontend/app/(dashboard)/layout.tsx}{lst:fd-app-dashboard-layout}
\begin{DeskFlowCode}
export default function DashboardLayout({ children }: { children: React.ReactNode }) {
  return <>{children}</>;
}
\end{DeskFlowCode}

\phantomsection
\addcontentsline{toc}{subsubsection}{deskflow\-frontend/app/(dashboard)/page.tsx}
\DeskFlowListingCaption{deskflow-frontend/app/(dashboard)/page.tsx}{lst:fd-app-dashboard-page}
\begin{DeskFlowCode}
'use client';

import { useEffect, useState } from 'react';
import Link from 'next/link';
import { Monitor, Users, Package, Calendar, ArrowRight, Clock, Building2, LogOut, Settings } from 'lucide-react';
import { Card, CardContent, CardDescription, CardHeader, CardTitle } from '@/components/ui/card';
import { Badge } from '@/components/ui/badge';
import { Button } from '@/components/ui/button';
import { useAuth } from '@/lib/auth-context';
import { requestsApi, resourcesApi, AllocationRequest, Resource } from '@/lib/api';
import { useRouter } from 'next/navigation';

const STATUS_LABELS: Record<string, { label: string; color: string }> = {
  SUBMITTED: { label: 'Złożony', color: 'bg-yellow-100 text-yellow-800' },
  APPROVED: { label: 'Zatwierdzony', color: 'bg-blue-100 text-blue-800' },
  ACTIVE: { label: 'Aktywny', color: 'bg-green-100 text-green-800' },
  COMPLETED: { label: 'Zakończony', color: 'bg-gray-100 text-gray-800' },
  REJECTED: { label: 'Odrzucony', color: 'bg-red-100 text-red-800' },
  CANCELLED: { label: 'Anulowany', color: 'bg-gray-100 text-gray-500' },
  OVERDUE: { label: 'Przeterminowane', color: 'bg-orange-100 text-orange-800' },
};

export default function DashboardPage() {
  const { user, loading: authLoading, logout, isManager, isItAdmin } = useAuth();
  const router = useRouter();
  const [myRequests, setMyRequests] = useState<AllocationRequest[]>([]);
  const [resources, setResources] = useState<Resource[]>([]);
  const [loading, setLoading] = useState(true);

  useEffect(() => {
    if (authLoading) return;

    if (!user) {
      router.push('/login');
      return;
    }
    Promise.all([
      requestsApi.my().then(setMyRequests).catch(() => {}),
      resourcesApi.list().then(setResources).catch(() => {}),
    ]).finally(() => setLoading(false));
  }, [user, authLoading, router]);

  if (authLoading || !user) return null;

  const activeRequests = myRequests.filter(r => ['SUBMITTED', 'APPROVED', 'ACTIVE'].includes(r.status));
  const availableResources = resources.filter(r => r.available);

  return (
    <div className="min-h-screen bg-gray-50">
      {/* Navbar */}
      <nav className="bg-white border-b px-6 py-3 flex items-center justify-between">
        <div className="flex items-center gap-2">
          <Building2 className="h-6 w-6 text-blue-600" />
          <span className="font-bold text-lg">DeskFlow</span>
        </div>
        <div className="flex items-center gap-4">
          <span className="text-sm text-gray-600">
            {user.firstName} {user.lastName} ({user.role})
          </span>
          <div className="flex gap-2">
            <Link href="/wnioski">
              <Button variant="ghost" size="sm">Moje wnioski</Button>
            </Link>
            <Link href="/zasoby">
              <Button variant="ghost" size="sm">Zasoby</Button>
            </Link>
            {isManager() && (
              <Link href="/administracja">
                <Button variant="ghost" size="sm">Wnioski</Button>
              </Link>
            )}
            {isItAdmin() && (
              <Link href="/admin/it">
                <Button variant="ghost" size="sm">IT Admin</Button>
              </Link>
            )}
            <Button variant="ghost" size="sm" onClick={logout}>
              <LogOut className="h-4 w-4" />
            </Button>
          </div>
        </div>
      </nav>

      <div className="p-8 max-w-7xl mx-auto">
        <div className="mb-8">
          <h1 className="text-2xl font-bold">Witaj, {user.firstName}!</h1>
          <p className="text-gray-600 mt-1">System zarządzania zasobami biurowymi</p>
        </div>

        {/* Stats */}
        <div className="grid gap-4 md:grid-cols-3 mb-8">
          <Card>
            <CardContent className="p-6">
              <div className="flex items-center justify-between">
                <div>
                  <p className="text-sm text-gray-600">Moje aktywne wnioski</p>
                  <p className="text-3xl font-bold mt-1">{activeRequests.length}</p>
                </div>
                <Calendar className="h-8 w-8 text-blue-500" />
              </div>
            </CardContent>
          </Card>
          <Card>
            <CardContent className="p-6">
              <div className="flex items-center justify-between">
                <div>
                  <p className="text-sm text-gray-600">Dostępne zasoby</p>
                  <p className="text-3xl font-bold mt-1">{availableResources.length}</p>
                </div>
                <Monitor className="h-8 w-8 text-green-500" />
              </div>
            </CardContent>
          </Card>
          <Card>
            <CardContent className="p-6">
              <div className="flex items-center justify-between">
                <div>
                  <p className="text-sm text-gray-600">Wszystkie zasoby</p>
                  <p className="text-3xl font-bold mt-1">{resources.length}</p>
                </div>
                <Package className="h-8 w-8 text-purple-500" />
              </div>
            </CardContent>
          </Card>
        </div>

        {/* Quick Actions */}
        <Card className="mb-6">
          <CardHeader>
            <CardTitle>Szybkie akcje</CardTitle>
          </CardHeader>
          <CardContent>
            <div className="flex flex-wrap gap-3">
              <Link href="/zasoby">
                <Button className="gap-2">
                  <Monitor className="h-4 w-4" />
                  Przeglądaj zasoby i złóż wniosek
                </Button>
              </Link>
              <Link href="/wnioski">
                <Button variant="outline" className="gap-2">
                  <Clock className="h-4 w-4" />
                  Historia wniosków
                </Button>
              </Link>
              {isManager() && (
                <Link href="/administracja">
                  <Button variant="outline" className="gap-2">
                    <Users className="h-4 w-4" />
                    Panel administracji
                  </Button>
                </Link>
              )}
              {isItAdmin() && (
                <Link href="/admin/it">
                  <Button variant="outline" className="gap-2">
                    <Settings className="h-4 w-4" />
                    IT Admin
                  </Button>
                </Link>
              )}
            </div>
          </CardContent>
        </Card>

        {/* Active requests */}
        <Card>
          <CardHeader className="flex flex-row items-center justify-between">
            <div>
              <CardTitle>Moje aktywne wnioski</CardTitle>
              <CardDescription>Wnioski oczekujące i aktywne udostępnienia</CardDescription>
            </div>
            <Link href="/wnioski">
              <Button variant="outline" size="sm" className="gap-2">
                Wszystkie <ArrowRight className="h-4 w-4" />
              </Button>
            </Link>
          </CardHeader>
          <CardContent>
            {loading ? (
              <p className="text-gray-500 text-sm">Ładowanie...</p>
            ) : activeRequests.length === 0 ? (
              <p className="text-gray-500 text-sm">Brak aktywnych wniosków</p>
            ) : (
              <div className="space-y-3">
                {activeRequests.slice(0, 5).map(r => (
                  <div key={r.id} className="flex items-center justify-between p-3 bg-gray-50 rounded-lg">
                    <div>
                      <div className="flex items-center gap-2 mb-1">
                        <span className={`text-xs px-2 py-0.5 rounded-full font-medium ${STATUS_LABELS[r.status]?.color}`}>
                          {STATUS_LABELS[r.status]?.label}
                        </span>
                      </div>
                      <p className="text-sm font-medium">{r.dateFrom} → {r.dateTo}</p>
                      {r.justification && (
                        <p className="text-xs text-gray-500">{r.justification}</p>
                      )}
                    </div>
                  </div>
                ))}
              </div>
            )}
          </CardContent>
        </Card>
      </div>
    </div>
  );
}
\end{DeskFlowCode}

\phantomsection
\addcontentsline{toc}{subsubsection}{deskflow\-frontend/app/(dashboard)/biurka/page.tsx}
\DeskFlowListingCaption{deskflow-frontend/app/(dashboard)/biurka/page.tsx}{lst:fd-app-dashboard-biurka-page}
\begin{DeskFlowCode}
'use client';

import { useEffect, useState } from 'react';
import { useRouter } from 'next/navigation';
import Link from 'next/link';
import { ArrowLeft, Plus, MapPin, CheckCircle, XCircle } from 'lucide-react';
import { Card, CardContent, CardHeader, CardTitle } from '@/components/ui/card';
import { Button } from '@/components/ui/button';
import { Input } from '@/components/ui/input';
import { Label } from '@/components/ui/label';
import { Dialog, DialogContent, DialogHeader, DialogTitle, DialogTrigger } from '@/components/ui/dialog';
import { useAuth } from '@/lib/auth-context';
import { resourcesApi, requestsApi, Resource } from '@/lib/api';
import { useToast } from '@/hooks/use-toast';

const TYPE_LABELS: Record<string, string> = {
  DESK: 'Biurko', CONFERENCE_ROOM: 'Sala', LAPTOP: 'Laptop',
  MONITOR: 'Monitor', PROJECTOR: 'Projektor', VEHICLE: 'Pojazd', ACCESSORY: 'Akcesoria',
};

const today = new Date().toISOString().split('T')[0];

function nextDay(date: string) {
  const match = /^(\d{4})-(\d{2})-(\d{2})$/.exec(date);
  if (!match) return today;

  const [, year, month, day] = match.map(Number);
  const value = new Date(Date.UTC(year, month - 1, day));
  if (
    Number.isNaN(value.getTime()) ||
    value.getUTCFullYear() !== year ||
    value.getUTCMonth() !== month - 1 ||
    value.getUTCDate() !== day
  ) return today;

  value.setUTCDate(value.getUTCDate() + 1);
  return value.toISOString().split('T')[0];
}

export default function ResourcesPage() {
  const { user, loading: authLoading } = useAuth();
  const router = useRouter();
  const { toast } = useToast();
  const [resources, setResources] = useState<Resource[]>([]);
  const [loading, setLoading] = useState(true);
  const [open, setOpen] = useState(false);
  const [selected, setSelected] = useState<Resource | null>(null);
  const [form, setForm] = useState({ dateFrom: '', dateTo: '', justification: '' });

  const load = () => resourcesApi.list().then(setResources).catch(() => {}).finally(() => setLoading(false));

  useEffect(() => {
    if (authLoading) return;

    if (!user) { router.push('/login'); return; }
    load();
  }, [user, authLoading, router]);

  const handleSubmit = async () => {
    if (!selected) return;

    if (form.dateFrom < today) {
      toast({ title: 'Błąd', description: 'Data od nie może być wcześniejsza niż dzisiaj.', variant: 'destructive' });
      return;
    }

    if (form.dateTo <= form.dateFrom) {
      toast({ title: 'Błąd', description: 'Data do musi być późniejsza niż Data od.', variant: 'destructive' });
      return;
    }

    try {
      await requestsApi.submit({ resourceId: selected.id, ...form });
      toast({ title: 'Wniosek złożony pomyślnie' });
      setOpen(false);
      setForm({ dateFrom: '', dateTo: '', justification: '' });
    } catch (e: unknown) {
      toast({ title: 'Błąd', description: e instanceof Error ? e.message : 'Błąd', variant: 'destructive' });
    }
  };

  return (
    <div className="min-h-screen bg-gray-50 p-8">
      <div className="max-w-6xl mx-auto">
        <div className="flex items-center gap-4 mb-6">
          <Link href="/"><Button variant="ghost" size="sm"><ArrowLeft className="h-4 w-4 mr-1" />Powrót</Button></Link>
          <div>
            <h1 className="text-2xl font-bold">Oferta zasobów</h1>
          </div>
        </div>

        {loading ? (
          <p className="text-gray-500">Ładowanie zasobów...</p>
        ) : (
          <div className="grid gap-4 md:grid-cols-2 lg:grid-cols-3">
            {resources.map(r => (
              <Card key={r.id} className={`${!r.available ? 'opacity-60' : 'hover:shadow-md transition-shadow'}`}>
                <CardHeader className="pb-2">
                  <div className="flex items-start justify-between">
                    <CardTitle className="text-base">{r.name}</CardTitle>
                    {r.available ? (
                      <span className="flex items-center gap-1 text-xs text-green-600">
                        <CheckCircle className="h-3 w-3" />Dostępny
                      </span>
                    ) : (
                      <span className="flex items-center gap-1 text-xs text-red-500">
                        <XCircle className="h-3 w-3" />Niedostępny
                      </span>
                    )}
                  </div>
                </CardHeader>
                <CardContent>
                  <div className="space-y-1 mb-3">
                    <span className="text-xs bg-blue-100 text-blue-700 px-2 py-0.5 rounded-full">
                      {TYPE_LABELS[r.type] || r.type}
                    </span>
                    <p className="text-sm text-gray-600 flex items-center gap-1 mt-2">
                      <MapPin className="h-3 w-3" />
                      {r.building}, piętro {r.floor}, nr {r.number}
                    </p>
                  </div>
                  <Dialog open={open && selected?.id === r.id} onOpenChange={v => { setOpen(v); if (!v) setSelected(null); }}>
                    <DialogTrigger asChild>
                      <Button
                        size="sm" className="w-full gap-2"
                        disabled={!r.available}
                        onClick={() => { setSelected(r); setOpen(true); }}
                      >
                        <Plus className="h-4 w-4" />Złóż wniosek
                      </Button>
                    </DialogTrigger>
                    <DialogContent>
                      <DialogHeader>
                        <DialogTitle>Złóż wniosek: {r.name}</DialogTitle>
                      </DialogHeader>
                      <div className="space-y-4 pt-2">
                        <div className="grid grid-cols-2 gap-3">
                          <div>
                            <Label>Data od</Label>
                            <Input
                              type="date"
                              min={today}
                              value={form.dateFrom}
                              onChange={e => setForm(f => ({ ...f, dateFrom: e.target.value }))}
                            />
                          </div>
                          <div>
                            <Label>Data do</Label>
                            <Input
                              type="date"
                              min={form.dateFrom ? nextDay(form.dateFrom) : today}
                              value={form.dateTo}
                              onChange={e => setForm(f => ({ ...f, dateTo: e.target.value }))}
                            />
                          </div>
                        </div>
                        <div>
                          <Label>Uzasadnienie</Label>
                          <Input
                            value={form.justification}
                            onChange={e => setForm(f => ({ ...f, justification: e.target.value }))}
                            placeholder="Cel wypożyczenia..."
                          />
                        </div>
                        <Button className="w-full" onClick={handleSubmit} disabled={!form.dateFrom || !form.dateTo}>
                          Złóż wniosek
                        </Button>
                      </div>
                    </DialogContent>
                  </Dialog>
                </CardContent>
              </Card>
            ))}
          </div>
        )}
      </div>
    </div>
  );
}
\end{DeskFlowCode}

\phantomsection
\addcontentsline{toc}{subsubsection}{deskflow\-frontend/app/(dashboard)/rezerwacje/page.tsx}
\DeskFlowListingCaption{deskflow-frontend/app/(dashboard)/rezerwacje/page.tsx}{lst:fd-app-dashboard-rezerwacje-page}
\begin{DeskFlowCode}
'use client';

import { useEffect, useState } from 'react';
import { useRouter } from 'next/navigation';
import Link from 'next/link';
import { ArrowLeft, XCircle, Box, Calendar, MapPin } from 'lucide-react';
import { Card, CardContent } from '@/components/ui/card';
import { Button } from '@/components/ui/button';
import { useAuth } from '@/lib/auth-context';
import { requestsApi, resourcesApi, AllocationRequest, Resource } from '@/lib/api';
import { useToast } from '@/hooks/use-toast';

const STATUS_LABELS: Record<string, { label: string; color: string }> = {
  SUBMITTED: { label: 'Złożony', color: 'bg-yellow-100 text-yellow-800' },
  APPROVED: { label: 'Zatwierdzony', color: 'bg-blue-100 text-blue-800' },
  ACTIVE: { label: 'Aktywny', color: 'bg-green-100 text-green-800' },
  COMPLETED: { label: 'Zakończony', color: 'bg-gray-100 text-gray-700' },
  REJECTED: { label: 'Odrzucony', color: 'bg-red-100 text-red-800' },
  CANCELLED: { label: 'Anulowany', color: 'bg-gray-100 text-gray-500' },
  OVERDUE: { label: 'Przeterminowane', color: 'bg-orange-100 text-orange-800' },
};

const TYPE_LABELS: Record<string, string> = {
  DESK: 'Biurko', CONFERENCE_ROOM: 'Sala konf.', LAPTOP: 'Laptop',
  MONITOR: 'Monitor', PROJECTOR: 'Projektor', VEHICLE: 'Pojazd', ACCESSORY: 'Akcesoria',
};

export default function MyRequestsPage() {
  const { user, loading: authLoading } = useAuth();
  const router = useRouter();
  const { toast } = useToast();
  const [requests, setRequests] = useState<AllocationRequest[]>([]);
  const [resourceMap, setResourceMap] = useState<Record<string, Resource>>({});
  const [loading, setLoading] = useState(true);

  const load = async () => {
    try {
      const [reqs, resources] = await Promise.all([requestsApi.my(), resourcesApi.list()]);
      setRequests(reqs);
      setResourceMap(Object.fromEntries(resources.map(r => [r.id, r])));
    } catch (e: unknown) {
      toast({ title: 'Błąd', description: e instanceof Error ? e.message : 'Błąd', variant: 'destructive' });
    } finally {
      setLoading(false);
    }
  };

  useEffect(() => {
    if (authLoading) return;

    if (!user) { router.push('/login'); return; }
    load();
  }, [user, authLoading, router]);

  const handleCancel = async (id: string) => {
    try {
      await requestsApi.cancel(id);
      toast({ title: 'Wniosek anulowany' });
      load();
    } catch (e: unknown) {
      toast({ title: 'Błąd', description: e instanceof Error ? e.message : 'Błąd', variant: 'destructive' });
    }
  };

  return (
    <div className="min-h-screen bg-gray-50 p-8">
      <div className="max-w-4xl mx-auto">
        <div className="flex items-center gap-4 mb-6">
          <Link href="/"><Button variant="ghost" size="sm"><ArrowLeft className="h-4 w-4 mr-1" />Powrót</Button></Link>
          <div>
            <h1 className="text-2xl font-bold">Moje wnioski</h1>
            <p className="text-gray-600 text-sm">Historia rezerwacji zasobów</p>
          </div>
        </div>

        {loading ? (
          <p className="text-gray-500">Ładowanie...</p>
        ) : requests.length === 0 ? (
          <Card>
            <CardContent className="p-8 text-center text-gray-500">
              Brak wniosków. <Link href="/zasoby" className="text-blue-600 hover:underline">Złóż pierwszy wniosek</Link>.
            </CardContent>
          </Card>
        ) : (
          <div className="space-y-3">
            {requests.map(r => {
              const resource = resourceMap[r.resourceId];
              return (
                <Card key={r.id} className="hover:shadow-sm transition-shadow">
                  <CardContent className="p-4">
                    <div className="flex items-start justify-between gap-4">
                      <div className="flex-1 min-w-0">
                        <div className="flex items-center gap-2 mb-2 flex-wrap">
                          <span className={`text-xs px-2 py-0.5 rounded-full font-medium ${STATUS_LABELS[r.status]?.color}`}>
                            {STATUS_LABELS[r.status]?.label}
                          </span>
                          <span className="text-xs text-gray-400">
                            złożony {new Date(r.submittedAt).toLocaleDateString('pl-PL')}
                          </span>
                        </div>

                        <div className="flex items-center gap-1.5 mb-1">
                          <Box className="h-4 w-4 text-gray-500 flex-shrink-0" />
                          <span className="font-semibold text-sm">
                            {resource ? resource.name : `Zasób #${r.resourceId.slice(0, 8)}`}
                          </span>
                          {resource && (
                            <span className="text-xs bg-blue-50 text-blue-700 px-1.5 py-0.5 rounded">
                              {TYPE_LABELS[resource.type] || resource.type}
                            </span>
                          )}
                        </div>

                        {resource && (
                          <div className="flex items-center gap-1 text-xs text-gray-500 mb-1">
                            <MapPin className="h-3 w-3" />
                            {resource.building}, piętro {resource.floor}, nr {resource.number}
                          </div>
                        )}

                        <div className="flex items-center gap-1 text-sm text-gray-700">
                          <Calendar className="h-3.5 w-3.5 text-gray-400" />
                          {r.dateFrom} — {r.dateTo}
                        </div>

                        {r.justification && (
                          <p className="text-xs text-gray-500 mt-1 italic">„{r.justification}"</p>
                        )}
                        {r.rejectionReason && (
                          <p className="text-xs text-red-600 mt-1">Powód odrzucenia: {r.rejectionReason}</p>
                        )}
                      </div>

                      {(r.status === 'SUBMITTED' || r.status === 'APPROVED') && (
                        <Button
                          variant="outline"
                          size="sm"
                          className="text-red-600 border-red-200 hover:bg-red-50 flex-shrink-0"
                          onClick={() => handleCancel(r.id)}
                        >
                          <XCircle className="h-4 w-4 mr-1" />
                          Anuluj
                        </Button>
                      )}
                    </div>
                  </CardContent>
                </Card>
              );
            })}
          </div>
        )}
      </div>
    </div>
  );
}
\end{DeskFlowCode}

\phantomsection
\addcontentsline{toc}{subsubsection}{deskflow\-frontend/app/(dashboard)/sale/page.tsx}
\DeskFlowListingCaption{deskflow-frontend/app/(dashboard)/sale/page.tsx}{lst:fd-app-dashboard-sale-page}
\begin{DeskFlowCode}
import { redirect } from 'next/navigation';
export default function RoomsRedirect() { redirect('/biurka'); }
\end{DeskFlowCode}

\phantomsection
\addcontentsline{toc}{subsubsection}{deskflow\-frontend/app/(dashboard)/sprzet/page.tsx}
\DeskFlowListingCaption{deskflow-frontend/app/(dashboard)/sprzet/page.tsx}{lst:fd-app-dashboard-sprzet-page}
\begin{DeskFlowCode}
'use client';
import { useEffect } from 'react';
import { useRouter } from 'next/navigation';

export default function EquipmentRedirect() {
  const router = useRouter();
  useEffect(() => { router.replace('/biurka'); }, [router]);
  return null;
}
\end{DeskFlowCode}

\phantomsection
\addcontentsline{toc}{subsubsection}{deskflow\-frontend/app/(dashboard)/zasoby/page.tsx}
\DeskFlowListingCaption{deskflow-frontend/app/(dashboard)/zasoby/page.tsx}{lst:fd-app-dashboard-zasoby-page}
\begin{DeskFlowCode}
import { redirect } from 'next/navigation';
export default function ResourcesRedirect() { redirect('/biurka'); }
\end{DeskFlowCode}

\phantomsection
\addcontentsline{toc}{subsubsection}{deskflow\-frontend/app/(dashboard)/zasoby\-admin/page.tsx}
\DeskFlowListingCaption{deskflow-frontend/app/(dashboard)/zasoby-admin/page.tsx}{lst:fd-app-dashboard-zasoby-admin-page}
\begin{DeskFlowCode}
import { redirect } from 'next/navigation';
export default function AdminITRedirect() { redirect('/admin/it'); }
\end{DeskFlowCode}

\phantomsection
\addcontentsline{toc}{subsubsection}{deskflow\-frontend/app/(dashboard)/wnioski/page.tsx}
\DeskFlowListingCaption{deskflow-frontend/app/(dashboard)/wnioski/page.tsx}{lst:fd-app-dashboard-wnioski-page}
\begin{DeskFlowCode}
import { redirect } from 'next/navigation';
export default function RequestsRedirect() { redirect('/rezerwacje'); }
\end{DeskFlowCode}

\phantomsection
\addcontentsline{toc}{subsubsection}{deskflow\-frontend/app/(dashboard)/powiadomienia/page.tsx}
\DeskFlowListingCaption{deskflow-frontend/app/(dashboard)/powiadomienia/page.tsx}{lst:fd-app-dashboard-powiadomienia-page}
\begin{DeskFlowCode}
'use client';
import { useEffect } from 'react';
import { useRouter } from 'next/navigation';

export default function NotificationsRedirect() {
  const router = useRouter();
  useEffect(() => { router.replace('/'); }, [router]);
  return null;
}
\end{DeskFlowCode}

\phantomsection
\addcontentsline{toc}{subsubsection}{deskflow\-frontend/app/(dashboard)/administracja/page.tsx}
\DeskFlowListingCaption{deskflow-frontend/app/(dashboard)/administracja/page.tsx}{lst:fd-app-dashboard-administracja-page}
\begin{DeskFlowCode}
import { redirect } from 'next/navigation';
export default function AdminRedirect() { redirect('/admin'); }
\end{DeskFlowCode}

\phantomsection
\addcontentsline{toc}{subsubsection}{deskflow\-frontend/app/(dashboard)/admin/page.tsx}
\DeskFlowListingCaption{deskflow-frontend/app/(dashboard)/admin/page.tsx}{lst:fd-app-dashboard-admin-page}
\begin{DeskFlowCode}
'use client';

import { useEffect, useState } from 'react';
import { useRouter } from 'next/navigation';
import Link from 'next/link';
import { ArrowLeft, Check, X, Package, RotateCcw, Box, Calendar, MapPin, User } from 'lucide-react';
import { Card, CardContent } from '@/components/ui/card';
import { Button } from '@/components/ui/button';
import { Input } from '@/components/ui/input';
import { Dialog, DialogContent, DialogHeader, DialogTitle } from '@/components/ui/dialog';
import { useAuth } from '@/lib/auth-context';
import { administrationApi, resourcesApi, AllocationRequest, Resource } from '@/lib/api';
import { useToast } from '@/hooks/use-toast';

const STATUS_LABELS: Record<string, { label: string; color: string }> = {
  SUBMITTED: { label: 'Złożony', color: 'bg-yellow-100 text-yellow-800' },
  APPROVED: { label: 'Zatwierdzony', color: 'bg-blue-100 text-blue-800' },
  ACTIVE: { label: 'Aktywny', color: 'bg-green-100 text-green-800' },
  COMPLETED: { label: 'Zakończony', color: 'bg-gray-100 text-gray-700' },
  REJECTED: { label: 'Odrzucony', color: 'bg-red-100 text-red-800' },
  CANCELLED: { label: 'Anulowany', color: 'bg-gray-100 text-gray-500' },
  OVERDUE: { label: 'Przeterminowane', color: 'bg-orange-100 text-orange-800' },
};

const TYPE_LABELS: Record<string, string> = {
  DESK: 'Biurko', CONFERENCE_ROOM: 'Sala konf.', LAPTOP: 'Laptop',
  MONITOR: 'Monitor', PROJECTOR: 'Projektor', VEHICLE: 'Pojazd', ACCESSORY: 'Akcesoria',
};

export default function AdminPanelPage() {
  const { user, loading: authLoading, isManager } = useAuth();
  const router = useRouter();
  const { toast } = useToast();
  const [requests, setRequests] = useState<AllocationRequest[]>([]);
  const [resourceMap, setResourceMap] = useState<Record<string, Resource>>({});
  const [filterStatus, setFilterStatus] = useState('');
  const [loading, setLoading] = useState(true);
  const [rejectDialog, setRejectDialog] = useState<string | null>(null);
  const [reason, setReason] = useState('');

  const load = async () => {
    try {
      const [reqs, resources] = await Promise.all([
        administrationApi.requests(filterStatus || undefined),
        resourcesApi.list(),
      ]);
      setRequests(reqs);
      setResourceMap(Object.fromEntries(resources.map(r => [r.id, r])));
    } catch (e: unknown) {
      toast({ title: 'Błąd', description: e instanceof Error ? e.message : 'Błąd', variant: 'destructive' });
    } finally {
      setLoading(false);
    }
  };

  useEffect(() => {
    if (authLoading) return;

    if (!user) { router.push('/login'); return; }
    if (!isManager()) { router.push('/'); return; }
    load();
  }, [user, authLoading, filterStatus, router]);

  const action = async (fn: () => Promise<unknown>, label: string) => {
    try {
      await fn();
      toast({ title: `${label} — sukces` });
      load();
    } catch (e: unknown) {
      toast({ title: 'Błąd', description: e instanceof Error ? e.message : 'Błąd', variant: 'destructive' });
    }
  };

  const handleReject = async () => {
    if (!rejectDialog) return;
    await action(() => administrationApi.reject(rejectDialog, reason), 'Wniosek odrzucony');
    setRejectDialog(null);
    setReason('');
  };

  return (
    <div className="min-h-screen bg-gray-50 p-8">
      <div className="max-w-5xl mx-auto">
        <div className="flex items-center gap-4 mb-6">
          <Link href="/"><Button variant="ghost" size="sm"><ArrowLeft className="h-4 w-4 mr-1" />Powrót</Button></Link>
          <div>
            <h1 className="text-2xl font-bold">Panel administracji</h1>
            <p className="text-gray-600 text-sm">Zarządzanie wnioskami o zasoby</p>
          </div>
        </div>

        <div className="flex gap-2 mb-6 flex-wrap">
          {['', 'SUBMITTED', 'APPROVED', 'ACTIVE', 'COMPLETED', 'OVERDUE'].map(s => (
            <Button
              key={s}
              variant={filterStatus === s ? 'default' : 'outline'}
              size="sm"
              onClick={() => setFilterStatus(s)}
            >
              {s === '' ? 'Wszystkie' : STATUS_LABELS[s]?.label}
            </Button>
          ))}
        </div>

        {loading ? (
          <p className="text-gray-500">Ładowanie...</p>
        ) : requests.length === 0 ? (
          <Card><CardContent className="p-8 text-center text-gray-500">Brak wniosków.</CardContent></Card>
        ) : (
          <div className="space-y-3">
            {requests.map(r => {
              const resource = resourceMap[r.resourceId];
              return (
                <Card key={r.id} className="hover:shadow-sm transition-shadow">
                  <CardContent className="p-4">
                    <div className="flex items-start justify-between gap-4">
                      <div className="flex-1 min-w-0">
                        <div className="flex items-center gap-2 mb-2 flex-wrap">
                          <span className={`text-xs px-2 py-0.5 rounded-full font-medium ${STATUS_LABELS[r.status]?.color}`}>
                            {STATUS_LABELS[r.status]?.label}
                          </span>
                          <span className="text-xs text-gray-400">
                            złożony {new Date(r.submittedAt).toLocaleDateString('pl-PL')}
                          </span>
                        </div>

                        <div className="flex items-center gap-1.5 mb-1">
                          <User className="h-3.5 w-3.5 text-gray-400 flex-shrink-0" />
                          <span className="text-sm text-gray-700">
                            {r.employeeFirstName && r.employeeLastName
                              ? <span className="font-medium">{r.employeeFirstName} {r.employeeLastName}</span>
                              : <span className="font-mono text-xs text-gray-400">{r.employeeId}</span>
                            }
                          </span>
                        </div>

                        <div className="flex items-center gap-1.5 mb-1">
                          <Box className="h-3.5 w-3.5 text-gray-500 flex-shrink-0" />
                          <span className="font-semibold text-sm">
                            {resource ? resource.name : `Zasób #${r.resourceId.slice(0, 8)}`}
                          </span>
                          {resource && (
                            <span className="text-xs bg-blue-50 text-blue-700 px-1.5 py-0.5 rounded">
                              {TYPE_LABELS[resource.type] || resource.type}
                            </span>
                          )}
                        </div>

                        {resource && (
                          <div className="flex items-center gap-1 text-xs text-gray-500 mb-1">
                            <MapPin className="h-3 w-3" />
                            {resource.building}, piętro {resource.floor}, nr {resource.number}
                          </div>
                        )}

                        <div className="flex items-center gap-1 text-sm text-gray-700">
                          <Calendar className="h-3.5 w-3.5 text-gray-400" />
                          {r.dateFrom} — {r.dateTo}
                        </div>

                        {r.justification && <p className="text-xs text-gray-500 mt-1 italic">„{r.justification}"</p>}
                        {r.rejectionReason && <p className="text-xs text-red-500 mt-1">Powód: {r.rejectionReason}</p>}
                      </div>

                      <div className="flex gap-2 flex-shrink-0 flex-wrap justify-end">
                        {r.status === 'SUBMITTED' && (
                          <>
                            <Button size="sm" className="gap-1 bg-green-600 hover:bg-green-700"
                              onClick={() => action(() => administrationApi.approve(r.id), 'Zatwierdzono')}>
                              <Check className="h-3 w-3" />Zatwierdź
                            </Button>
                            <Button size="sm" variant="destructive" className="gap-1"
                              onClick={() => setRejectDialog(r.id)}>
                              <X className="h-3 w-3" />Odrzuć
                            </Button>
                          </>
                        )}
                        {r.status === 'APPROVED' && (
                          <Button size="sm" className="gap-1"
                            onClick={() => action(() => administrationApi.issue(r.id), 'Zasób wydany')}>
                            <Package className="h-3 w-3" />Wydaj zasób
                          </Button>
                        )}
                        {(r.status === 'ACTIVE' || r.status === 'OVERDUE') && (
                          <Button size="sm" variant="outline" className="gap-1"
                            onClick={() => action(() => administrationApi.return(r.id), 'Zasób zwrócony')}>
                            <RotateCcw className="h-3 w-3" />Zwróć zasób
                          </Button>
                        )}
                      </div>
                    </div>
                  </CardContent>
                </Card>
              );
            })}
          </div>
        )}
      </div>

      <Dialog open={!!rejectDialog} onOpenChange={v => !v && setRejectDialog(null)}>
        <DialogContent>
          <DialogHeader>
            <DialogTitle>Odrzuć wniosek</DialogTitle>
          </DialogHeader>
          <div className="space-y-4">
            <Input
              placeholder="Podaj powód odrzucenia..."
              value={reason}
              onChange={e => setReason(e.target.value)}
            />
            <div className="flex gap-2 justify-end">
              <Button variant="outline" onClick={() => setRejectDialog(null)}>Anuluj</Button>
              <Button variant="destructive" onClick={handleReject} disabled={!reason.trim()}>
                Odrzuć
              </Button>
            </div>
          </div>
        </DialogContent>
      </Dialog>
    </div>
  );
}
\end{DeskFlowCode}

\phantomsection
\addcontentsline{toc}{subsubsection}{deskflow\-frontend/app/(dashboard)/admin/layout.tsx}
\DeskFlowListingCaption{deskflow-frontend/app/(dashboard)/admin/layout.tsx}{lst:fd-app-dashboard-admin-layout}
\begin{DeskFlowCode}
export default function AdminLayout({ children }: { children: React.ReactNode }) {
  return <>{children}</>;
}
\end{DeskFlowCode}

\phantomsection
\addcontentsline{toc}{subsubsection}{deskflow\-frontend/app/(dashboard)/admin/it/page.tsx}
\DeskFlowListingCaption{deskflow-frontend/app/(dashboard)/admin/it/page.tsx}{lst:fd-app-dashboard-admin-it-page}
\begin{DeskFlowCode}
'use client';

import { useEffect, useState } from 'react';
import { useRouter } from 'next/navigation';
import Link from 'next/link';
import {
  ArrowLeft, Plus, Pencil, Trash2, CheckCircle, XCircle,
  Box, Users, History, Save, X,
} from 'lucide-react';
import { Card, CardContent, CardHeader, CardTitle } from '@/components/ui/card';
import { Button } from '@/components/ui/button';
import { Input } from '@/components/ui/input';
import { Label } from '@/components/ui/label';
import { Select, SelectContent, SelectItem, SelectTrigger, SelectValue } from '@/components/ui/select';
import { Dialog, DialogContent, DialogHeader, DialogTitle } from '@/components/ui/dialog';
import { Tabs, TabsContent, TabsList, TabsTrigger } from '@/components/ui/tabs';
import { useAuth } from '@/lib/auth-context';
import {
  resourcesApi, employeeAdminApi, administrationApi,
  Resource, ResourceCreate, EmployeeWithLimits, AllocationRequest,
} from '@/lib/api';
import { useToast } from '@/hooks/use-toast';

const TYPE_LABELS: Record<string, string> = {
  DESK: 'Biurko', CONFERENCE_ROOM: 'Sala konf.', LAPTOP: 'Laptop',
  MONITOR: 'Monitor', PROJECTOR: 'Projektor', VEHICLE: 'Pojazd', ACCESSORY: 'Akcesoria',
};
const RESOURCE_TYPES = Object.keys(TYPE_LABELS);

const STATUS_LABELS: Record<string, { label: string; color: string }> = {
  SUBMITTED: { label: 'Złożony', color: 'bg-yellow-100 text-yellow-800' },
  APPROVED: { label: 'Zatwierdzony', color: 'bg-blue-100 text-blue-800' },
  ACTIVE: { label: 'Aktywny', color: 'bg-green-100 text-green-800' },
  COMPLETED: { label: 'Zakończony', color: 'bg-gray-100 text-gray-700' },
  REJECTED: { label: 'Odrzucony', color: 'bg-red-100 text-red-800' },
  CANCELLED: { label: 'Anulowany', color: 'bg-gray-100 text-gray-500' },
  OVERDUE: { label: 'Przeterminowane', color: 'bg-orange-100 text-orange-800' },
};

const EMPTY_FORM: ResourceCreate = { type: 'DESK', name: '', building: '', floor: 1, number: '' };

export default function ItAdminPage() {
  const { user, loading: authLoading, isItAdmin } = useAuth();
  const router = useRouter();
  const { toast } = useToast();

  // - Resources -
  const [resources, setResources] = useState<Resource[]>([]);
  const [resLoading, setResLoading] = useState(true);
  const [resDialog, setResDialog] = useState<'add' | 'edit' | null>(null);
  const [editTarget, setEditTarget] = useState<Resource | null>(null);
  const [form, setForm] = useState<ResourceCreate>(EMPTY_FORM);

  // - Limits -
  const [employees, setEmployees] = useState<EmployeeWithLimits[]>([]);
  const [empLoading, setEmpLoading] = useState(true);
  const [editingLimit, setEditingLimit] = useState<{ empId: string; type: string; value: number } | null>(null);

  // - History -
  const [history, setHistory] = useState<AllocationRequest[]>([]);
  const [histLoading, setHistLoading] = useState(true);
  const [histFilter, setHistFilter] = useState('');
  const [resourceMap, setResourceMap] = useState<Record<string, Resource>>({});

  useEffect(() => {
    if (authLoading) return;

    if (!user) { router.push('/login'); return; }
    if (!isItAdmin()) { router.push('/'); return; }
    loadResources();
    loadEmployees();
    loadHistory();
  }, [user, authLoading, router]);

  const loadResources = () =>
    resourcesApi.list().then(setResources).catch(() => {}).finally(() => setResLoading(false));

  const loadEmployees = () =>
    employeeAdminApi.list().then(setEmployees).catch(() => {}).finally(() => setEmpLoading(false));

  const loadHistory = () =>
    Promise.all([administrationApi.requests(), resourcesApi.list()])
      .then(([reqs, res]) => {
        setHistory(reqs);
        setResourceMap(Object.fromEntries(res.map(r => [r.id, r])));
      })
      .catch(() => {})
      .finally(() => setHistLoading(false));

  // - Resources actions -
  const openAdd = () => { setForm(EMPTY_FORM); setEditTarget(null); setResDialog('add'); };
  const openEdit = (r: Resource) => {
    setForm({ type: r.type, name: r.name, building: r.building, floor: r.floor, number: r.number });
    setEditTarget(r);
    setResDialog('edit');
  };

  const handleSaveResource = async () => {
    try {
      if (resDialog === 'add') {
        await resourcesApi.add(form);
        toast({ title: 'Zasób dodany' });
      } else if (editTarget) {
        await resourcesApi.update(editTarget.id, form);
        toast({ title: 'Zasób zaktualizowany' });
      }
      setResDialog(null);
      loadResources();
    } catch (e: unknown) {
      toast({ title: 'Błąd', description: e instanceof Error ? e.message : 'Błąd', variant: 'destructive' });
    }
  };

  const handleToggleAvailability = async (r: Resource) => {
    try {
      if (r.available) await resourcesApi.markUnavailable(r.id);
      else await resourcesApi.markAvailable(r.id);
      loadResources();
    } catch (e: unknown) {
      toast({ title: 'Błąd', description: e instanceof Error ? e.message : 'Błąd', variant: 'destructive' });
    }
  };

  const handleDeactivate = async (id: string) => {
    try {
      await resourcesApi.deactivate(id);
      toast({ title: 'Zasób dezaktywowany' });
      loadResources();
    } catch (e: unknown) {
      toast({ title: 'Błąd', description: e instanceof Error ? e.message : 'Błąd', variant: 'destructive' });
    }
  };

  // - Limits actions -
  const handleSaveLimit = async () => {
    if (!editingLimit) return;
    try {
      await employeeAdminApi.upsertLimit(editingLimit.empId, editingLimit.type, editingLimit.value);
      toast({ title: 'Limit zaktualizowany' });
      setEditingLimit(null);
      loadEmployees();
    } catch (e: unknown) {
      toast({ title: 'Błąd', description: e instanceof Error ? e.message : 'Błąd', variant: 'destructive' });
    }
  };

  const filteredHistory = histFilter
    ? history.filter(r => r.status === histFilter)
    : history;

  return (
    <div className="min-h-screen bg-gray-50 p-8">
      <div className="max-w-6xl mx-auto">
        <div className="flex items-center gap-4 mb-6">
          <Link href="/"><Button variant="ghost" size="sm"><ArrowLeft className="h-4 w-4 mr-1" />Powrót</Button></Link>
          <div>
            <h1 className="text-2xl font-bold">Panel IT Admin</h1>
            <p className="text-gray-600 text-sm">Zarządzanie zasobami, limitami i historia zdarzeń</p>
          </div>
        </div>

        <Tabs defaultValue="resources">
          <TabsList className="mb-6">
            <TabsTrigger value="resources" className="gap-2"><Box className="h-4 w-4" />Zasoby</TabsTrigger>
            <TabsTrigger value="limits" className="gap-2"><Users className="h-4 w-4" />Limity pracowników</TabsTrigger>
            <TabsTrigger value="history" className="gap-2"><History className="h-4 w-4" />Historia zdarzeń</TabsTrigger>
          </TabsList>

          {/* ===== RESOURCES ===== */}
          <TabsContent value="resources">
            <div className="flex justify-between items-center mb-4">
              <p className="text-sm text-gray-500">{resources.length} zasobów w systemie</p>
              <Button size="sm" className="gap-2" onClick={openAdd}>
                <Plus className="h-4 w-4" />Dodaj zasób
              </Button>
            </div>
            {resLoading ? (
              <p className="text-gray-500">Ładowanie...</p>
            ) : (
              <div className="space-y-2">
                {resources.map(r => (
                  <Card key={r.id} className={`${!r.available ? 'opacity-60' : ''}`}>
                    <CardContent className="p-4 flex items-center justify-between gap-4">
                      <div className="flex-1 min-w-0">
                        <div className="flex items-center gap-2 flex-wrap">
                          <span className="font-medium text-sm">{r.name}</span>
                          <span className="text-xs bg-blue-50 text-blue-700 px-1.5 py-0.5 rounded">
                            {TYPE_LABELS[r.type] || r.type}
                          </span>
                          {r.available
                            ? <span className="text-xs text-green-600 flex items-center gap-1"><CheckCircle className="h-3 w-3" />Dostępny</span>
                            : <span className="text-xs text-red-500 flex items-center gap-1"><XCircle className="h-3 w-3" />Niedostępny</span>
                          }
                        </div>
                        <p className="text-xs text-gray-500 mt-0.5">{r.building}, piętro {r.floor}, nr {r.number}</p>
                      </div>
                      <div className="flex gap-2 flex-shrink-0">
                        <Button size="sm" variant="outline" onClick={() => handleToggleAvailability(r)}>
                          {r.available ? 'Zablokuj' : 'Odblokuj'}
                        </Button>
                        <Button size="sm" variant="outline" onClick={() => openEdit(r)}>
                          <Pencil className="h-3.5 w-3.5" />
                        </Button>
                        <Button size="sm" variant="outline" className="text-red-600 border-red-200 hover:bg-red-50"
                          onClick={() => handleDeactivate(r.id)}>
                          <Trash2 className="h-3.5 w-3.5" />
                        </Button>
                      </div>
                    </CardContent>
                  </Card>
                ))}
              </div>
            )}
          </TabsContent>

          {/* ===== LIMITS ===== */}
          <TabsContent value="limits">
            {empLoading ? (
              <p className="text-gray-500">Ładowanie...</p>
            ) : (
              <div className="space-y-4">
                {employees.map(emp => (
                  <Card key={emp.id}>
                    <CardHeader className="pb-2">
                      <CardTitle className="text-base flex items-center gap-2">
                        {emp.firstName} {emp.lastName}
                        <span className="text-xs font-normal text-gray-500">{emp.email}</span>
                        <span className="text-xs bg-gray-100 text-gray-600 px-1.5 py-0.5 rounded">{emp.role}</span>
                      </CardTitle>
                    </CardHeader>
                    <CardContent className="pt-0">
                      <div className="flex flex-wrap gap-3">
                        {RESOURCE_TYPES.map(type => {
                          const current = emp.limits[type] ?? 0;
                          const isEditing = editingLimit?.empId === emp.id && editingLimit?.type === type;
                          return (
                            <div key={type} className="flex items-center gap-1.5 bg-gray-50 border rounded-lg px-3 py-2 text-sm">
                              <span className="text-gray-600 w-24 text-xs">{TYPE_LABELS[type]}</span>
                              {isEditing ? (
                                <>
                                  <Input
                                    type="number"
                                    min={0}
                                    className="w-16 h-7 text-center text-sm"
                                    value={editingLimit.value}
                                    onChange={e => setEditingLimit({ ...editingLimit, value: parseInt(e.target.value) || 0 })}
                                    autoFocus
                                  />
                                  <Button size="sm" variant="ghost" className="h-7 w-7 p-0 text-green-600" onClick={handleSaveLimit}>
                                    <Save className="h-3.5 w-3.5" />
                                  </Button>
                                  <Button size="sm" variant="ghost" className="h-7 w-7 p-0" onClick={() => setEditingLimit(null)}>
                                    <X className="h-3.5 w-3.5" />
                                  </Button>
                                </>
                              ) : (
                                <>
                                  <span className="font-semibold w-6 text-center">{current}</span>
                                  <Button
                                    size="sm" variant="ghost" className="h-7 w-7 p-0 text-gray-400 hover:text-gray-700"
                                    onClick={() => setEditingLimit({ empId: emp.id, type, value: current })}
                                  >
                                    <Pencil className="h-3 w-3" />
                                  </Button>
                                </>
                              )}
                            </div>
                          );
                        })}
                      </div>
                    </CardContent>
                  </Card>
                ))}
              </div>
            )}
          </TabsContent>

          {/* ===== HISTORY ===== */}
          <TabsContent value="history">
            <div className="flex gap-2 mb-4 flex-wrap">
              {['', ...Object.keys(STATUS_LABELS)].map(s => (
                <Button key={s} variant={histFilter === s ? 'default' : 'outline'} size="sm"
                  onClick={() => setHistFilter(s)}>
                  {s === '' ? 'Wszystkie' : STATUS_LABELS[s]?.label}
                </Button>
              ))}
            </div>
            {histLoading ? (
              <p className="text-gray-500">Ładowanie...</p>
            ) : filteredHistory.length === 0 ? (
              <Card><CardContent className="p-8 text-center text-gray-500">Brak zdarzeń.</CardContent></Card>
            ) : (
              <div className="space-y-2">
                {filteredHistory.map(r => {
                  const res = resourceMap[r.resourceId];
                  return (
                    <Card key={r.id}>
                      <CardContent className="p-3 flex items-center gap-4">
                        <div className="flex-1 min-w-0">
                          <div className="flex items-center gap-2 flex-wrap mb-1">
                            <span className={`text-xs px-2 py-0.5 rounded-full font-medium ${STATUS_LABELS[r.status]?.color}`}>
                              {STATUS_LABELS[r.status]?.label}
                            </span>
                            <span className="text-xs text-gray-500">
                              {r.employeeFirstName && r.employeeLastName
                                ? `${r.employeeFirstName} ${r.employeeLastName}`
                                : r.employeeId.slice(0, 8)}
                            </span>
                            <span className="text-xs text-gray-400">→</span>
                            <span className="text-xs font-medium">{res ? res.name : r.resourceId.slice(0, 8)}</span>
                            {res && <span className="text-xs text-gray-400">{TYPE_LABELS[res.type] || res.type}</span>}
                          </div>
                          <div className="flex items-center gap-3 text-xs text-gray-500">
                            <span>{r.dateFrom} — {r.dateTo}</span>
                            <span>złożony: {new Date(r.submittedAt).toLocaleDateString('pl-PL')}</span>
                            {r.issuedAt && <span>wydany: {new Date(r.issuedAt).toLocaleDateString('pl-PL')}</span>}
                            {r.returnedAt && <span>zwrócony: {new Date(r.returnedAt).toLocaleDateString('pl-PL')}</span>}
                          </div>
                          {r.rejectionReason && <p className="text-xs text-red-500 mt-0.5">Powód: {r.rejectionReason}</p>}
                        </div>
                      </CardContent>
                    </Card>
                  );
                })}
              </div>
            )}
          </TabsContent>
        </Tabs>
      </div>

      {/* Resource add/edit dialog */}
      <Dialog open={!!resDialog} onOpenChange={v => !v && setResDialog(null)}>
        <DialogContent>
          <DialogHeader>
            <DialogTitle>{resDialog === 'add' ? 'Dodaj zasób' : 'Edytuj zasób'}</DialogTitle>
          </DialogHeader>
          <div className="space-y-3">
            {resDialog === 'add' && (
              <div>
                <Label>Typ</Label>
                <Select value={form.type} onValueChange={v => setForm(f => ({ ...f, type: v }))}>
                  <SelectTrigger><SelectValue /></SelectTrigger>
                  <SelectContent>
                    {RESOURCE_TYPES.map(t => <SelectItem key={t} value={t}>{TYPE_LABELS[t]}</SelectItem>)}
                  </SelectContent>
                </Select>
              </div>
            )}
            <div>
              <Label>Nazwa</Label>
              <Input value={form.name} onChange={e => setForm(f => ({ ...f, name: e.target.value }))} placeholder="np. Biurko A-105" />
            </div>
            <div className="grid grid-cols-3 gap-2">
              <div>
                <Label>Budynek</Label>
                <Input value={form.building} onChange={e => setForm(f => ({ ...f, building: e.target.value }))} placeholder="A" />
              </div>
              <div>
                <Label>Piętro</Label>
                <Input type="number" value={form.floor} onChange={e => setForm(f => ({ ...f, floor: parseInt(e.target.value) || 0 }))} />
              </div>
              <div>
                <Label>Numer</Label>
                <Input value={form.number} onChange={e => setForm(f => ({ ...f, number: e.target.value }))} placeholder="105" />
              </div>
            </div>
            <div className="flex gap-2 justify-end pt-2">
              <Button variant="outline" onClick={() => setResDialog(null)}>Anuluj</Button>
              <Button onClick={handleSaveResource} disabled={!form.name || !form.building || !form.number}>
                {resDialog === 'add' ? 'Dodaj' : 'Zapisz'}
              </Button>
            </div>
          </div>
        </DialogContent>
      </Dialog>
    </div>
  );
}
\end{DeskFlowCode}

% ── LIB ──

\phantomsection
\addcontentsline{toc}{subsubsection}{deskflow\-frontend/lib/api.ts}
\DeskFlowListingCaption{deskflow-frontend/lib/api.ts}{lst:fd-lib-api}
\begin{DeskFlowCode}
const API_URL = process.env.NEXT_PUBLIC_API_URL || 'http://localhost:8080';

function getToken(): string | null {
  if (typeof window === 'undefined') return null;
  return localStorage.getItem('deskflow_token');
}

async function request<T>(path: string, options: RequestInit = {}): Promise<T> {
  const token = getToken();
  const headers: Record<string, string> = {
    'Content-Type': 'application/json',
    ...(options.headers as Record<string, string>),
  };
  if (token) headers['Authorization'] = `Bearer ${token}`;

  const res = await fetch(`${API_URL}${path}`, { ...options, headers });

  if (res.status === 401) {
    localStorage.removeItem('deskflow_token');
    localStorage.removeItem('deskflow_user');
    window.location.href = '/login';
    throw new Error('Session expired');
  }

  if (!res.ok) {
    const err = await res.json().catch(() => ({ error: res.statusText }));
    throw new Error(err.error || 'Server error');
  }

  const text = await res.text();
  return text ? JSON.parse(text) : null;
}

export const api = {
  get: <T>(path: string) => request<T>(path, { method: 'GET' }),
  post: <T>(path: string, body: unknown) => request<T>(path, { method: 'POST', body: JSON.stringify(body) }),
  put: <T>(path: string, body?: unknown) => request<T>(path, { method: 'PUT', body: body ? JSON.stringify(body) : undefined }),
  delete: <T>(path: string) => request<T>(path, { method: 'DELETE' }),
};

// Auth
export const authApi = {
  login: (email: string, password: string) =>
    api.post<{ token: string; email: string; role: string; firstName: string; lastName: string; employeeId: string }>(
      '/api/auth/login', { email, password }
    ),
  me: () => api.get<{ email: string; role: string; firstName: string; lastName: string; employeeId: string }>('/api/auth/me'),
};

// Resources (UC-01)
export const resourcesApi = {
  list: () => api.get<Resource[]>('/api/resources'),
  add: (data: ResourceCreate) => api.post<{ resourceId: string }>('/api/resources', data),
  update: (id: string, data: ResourceCreate) => api.put(`/api/resources/${id}`, data),
  markUnavailable: (id: string) => api.put(`/api/resources/${id}/unavailable`),
  markAvailable: (id: string) => api.put(`/api/resources/${id}/available`),
  deactivate: (id: string) => api.delete(`/api/resources/${id}`),
};

// Allocation requests (UC-02, 03, 14)
export const requestsApi = {
  submit: (data: SubmitRequest) => api.post<{ requestId: string }>('/api/requests', data),
  cancel: (id: string) => api.put(`/api/requests/${id}/cancel`),
  my: () => api.get<AllocationRequest[]>('/api/requests/my'),
};

// Administration (UC-04-09, 15)
export const administrationApi = {
  requests: (status?: string) => api.get<AllocationRequest[]>(`/api/administration/requests${status ? `?status=${status}` : ''}`),
  approve: (id: string) => api.put(`/api/administration/requests/${id}/approve`),
  reject: (id: string, reason: string) => api.put(`/api/administration/requests/${id}/reject`, { reason }),
  issue: (id: string) => api.put(`/api/administration/requests/${id}/issue`),
  return: (id: string) => api.put(`/api/administration/requests/${id}/return`),
  allocations: () => api.get<AllocationRequest[]>('/api/administration/allocations'),
};

// Employee admin (IT_ADMIN — UC-10, UC-11, limits)
export const employeeAdminApi = {
  list: () => api.get<EmployeeWithLimits[]>('/api/administration/employees'),
  upsertLimit: (employeeId: string, resourceType: string, amount: number) =>
    api.put(`/api/administration/employees/${employeeId}/limits`, { resourceType, amount }),
};

// Types
export interface Resource {
  id: string;
  type: string;
  name: string;
  building: string;
  floor: number;
  number: string;
  available: boolean;
}

export interface ResourceCreate {
  type: string;
  name: string;
  building: string;
  floor: number;
  number: string;
}

export interface AllocationRequest {
  id: string;
  employeeId: string;
  employeeFirstName: string | null;
  employeeLastName: string | null;
  resourceId: string;
  status: string;
  dateFrom: string;
  dateTo: string;
  justification: string;
  submittedAt: string;
  issuedAt: string | null;
  returnedAt: string | null;
  rejectionReason: string | null;
}

export interface SubmitRequest {
  resourceId: string;
  dateFrom: string;
  dateTo: string;
  justification: string;
}

export interface EmployeeWithLimits {
  id: string;
  email: string;
  firstName: string;
  lastName: string;
  role: string;
  limits: Record<string, number>;
}

export interface CurrentUser {
  email: string;
  role: string;
  firstName: string;
  lastName: string;
  employeeId: string;
}
\end{DeskFlowCode}

\phantomsection
\addcontentsline{toc}{subsubsection}{deskflow\-frontend/lib/utils.ts}
\DeskFlowListingCaption{deskflow-frontend/lib/utils.ts}{lst:fd-lib-utils}
\begin{DeskFlowCode}
import { clsx, type ClassValue } from 'clsx'
import { twMerge } from 'tailwind-merge'

export function cn(...inputs: ClassValue[]) {
  return twMerge(clsx(inputs))
}
\end{DeskFlowCode}

\phantomsection
\addcontentsline{toc}{subsubsection}{deskflow\-frontend/lib/auth\-context.tsx}
\DeskFlowListingCaption{deskflow-frontend/lib/auth-context.tsx}{lst:fd-lib-auth-context}
\begin{DeskFlowCode}
'use client';

import { createContext, useContext, useEffect, useState, ReactNode } from 'react';
import { CurrentUser, authApi } from './api';

interface AuthContextType {
  user: CurrentUser | null;
  loading: boolean;
  login: (token: string, user: CurrentUser) => void;
  logout: () => void;
  isManager: () => boolean;
  isItAdmin: () => boolean;
}

const AuthContext = createContext<AuthContextType | null>(null);

export function AuthProvider({ children }: { children: ReactNode }) {
  const [user, setUser] = useState<CurrentUser | null>(null);
  const [loading, setLoading] = useState(true);

  useEffect(() => {
    const token = localStorage.getItem('deskflow_token');
    const stored = localStorage.getItem('deskflow_user');
    if (token && stored) {
      setUser(JSON.parse(stored));
    }
    setLoading(false);
  }, []);

  const login = (token: string, userData: CurrentUser) => {
    localStorage.setItem('deskflow_token', token);
    localStorage.setItem('deskflow_user', JSON.stringify(userData));
    setUser(userData);
  };

  const logout = () => {
    localStorage.removeItem('deskflow_token');
    localStorage.removeItem('deskflow_user');
    setUser(null);
    window.location.href = '/login';
  };

  const isManager = () => user?.role === 'MANAGER' || user?.role === 'IT_ADMIN';
  const isItAdmin = () => user?.role === 'IT_ADMIN';

  return (
    <AuthContext.Provider value={{ user, loading, login, logout, isManager, isItAdmin }}>
      {children}
    </AuthContext.Provider>
  );
}

export function useAuth() {
  const ctx = useContext(AuthContext);
  if (!ctx) throw new Error('useAuth must be used within AuthProvider');
  return ctx;
}
\end{DeskFlowCode}

% ── HOOKS ──

\phantomsection
\addcontentsline{toc}{subsubsection}{deskflow\-frontend/hooks/use\-toast.ts}
\DeskFlowListingCaption{deskflow-frontend/hooks/use-toast.ts}{lst:fd-hooks-use-toast}
\begin{DeskFlowCode}
'use client'

// Inspired by react-hot-toast library
import * as React from 'react'

import type { ToastActionElement, ToastProps } from '@/components/ui/toast'

const TOAST_LIMIT = 1
const TOAST_REMOVE_DELAY = 1000000

type ToasterToast = ToastProps & {
  id: string
  title?: React.ReactNode
  description?: React.ReactNode
  action?: ToastActionElement
}

const actionTypes = {
  ADD_TOAST: 'ADD_TOAST',
  UPDATE_TOAST: 'UPDATE_TOAST',
  DISMISS_TOAST: 'DISMISS_TOAST',
  REMOVE_TOAST: 'REMOVE_TOAST',
} as const

let count = 0

function genId() {
  count = (count + 1) % Number.MAX_SAFE_INTEGER
  return count.toString()
}

type ActionType = typeof actionTypes

type Action =
  | {
      type: ActionType['ADD_TOAST']
      toast: ToasterToast
    }
  | {
      type: ActionType['UPDATE_TOAST']
      toast: Partial<ToasterToast>
    }
  | {
      type: ActionType['DISMISS_TOAST']
      toastId?: ToasterToast['id']
    }
  | {
      type: ActionType['REMOVE_TOAST']
      toastId?: ToasterToast['id']
    }

interface State {
  toasts: ToasterToast[]
}

const toastTimeouts = new Map<string, ReturnType<typeof setTimeout>>()

const addToRemoveQueue = (toastId: string) => {
  if (toastTimeouts.has(toastId)) {
    return
  }

  const timeout = setTimeout(() => {
    toastTimeouts.delete(toastId)
    dispatch({
      type: 'REMOVE_TOAST',
      toastId: toastId,
    })
  }, TOAST_REMOVE_DELAY)

  toastTimeouts.set(toastId, timeout)
}

export const reducer = (state: State, action: Action): State => {
  switch (action.type) {
    case 'ADD_TOAST':
      return {
        ...state,
        toasts: [action.toast, ...state.toasts].slice(0, TOAST_LIMIT),
      }

    case 'UPDATE_TOAST':
      return {
        ...state,
        toasts: state.toasts.map((t) =>
          t.id === action.toast.id ? { ...t, ...action.toast } : t,
        ),
      }

    case 'DISMISS_TOAST': {
      const { toastId } = action

      // ! Side effects ! - This could be extracted into a dismissToast() action,
      // but I'll keep it here for simplicity
      if (toastId) {
        addToRemoveQueue(toastId)
      } else {
        state.toasts.forEach((toast) => {
          addToRemoveQueue(toast.id)
        })
      }

      return {
        ...state,
        toasts: state.toasts.map((t) =>
          t.id === toastId || toastId === undefined
            ? {
                ...t,
                open: false,
              }
            : t,
        ),
      }
    }
    case 'REMOVE_TOAST':
      if (action.toastId === undefined) {
        return {
          ...state,
          toasts: [],
        }
      }
      return {
        ...state,
        toasts: state.toasts.filter((t) => t.id !== action.toastId),
      }
  }
}

const listeners: Array<(state: State) => void> = []

let memoryState: State = { toasts: [] }

function dispatch(action: Action) {
  memoryState = reducer(memoryState, action)
  listeners.forEach((listener) => {
    listener(memoryState)
  })
}

type Toast = Omit<ToasterToast, 'id'>

function toast({ ...props }: Toast) {
  const id = genId()

  const update = (props: ToasterToast) =>
    dispatch({
      type: 'UPDATE_TOAST',
      toast: { ...props, id },
    })
  const dismiss = () => dispatch({ type: 'DISMISS_TOAST', toastId: id })

  dispatch({
    type: 'ADD_TOAST',
    toast: {
      ...props,
      id,
      open: true,
      onOpenChange: (open) => {
        if (!open) dismiss()
      },
    },
  })

  return {
    id: id,
    dismiss,
    update,
  }
}

function useToast() {
  const [state, setState] = React.useState<State>(memoryState)

  React.useEffect(() => {
    listeners.push(setState)
    return () => {
      const index = listeners.indexOf(setState)
      if (index > -1) {
        listeners.splice(index, 1)
      }
    }
  }, [state])

  return {
    ...state,
    toast,
    dismiss: (toastId?: string) => dispatch({ type: 'DISMISS_TOAST', toastId }),
  }
}

export { useToast, toast }
\end{DeskFlowCode}

\phantomsection
\addcontentsline{toc}{subsubsection}{deskflow\-frontend/hooks/use\-mobile.ts}
\DeskFlowListingCaption{deskflow-frontend/hooks/use-mobile.ts}{lst:fd-hooks-use-mobile}
\begin{DeskFlowCode}
import * as React from 'react'

const MOBILE_BREAKPOINT = 768

export function useIsMobile() {
  const [isMobile, setIsMobile] = React.useState<boolean | undefined>(undefined)

  React.useEffect(() => {
    const mql = window.matchMedia(`(max-width: ${MOBILE_BREAKPOINT - 1}px)`)
    const onChange = () => {
      setIsMobile(window.innerWidth < MOBILE_BREAKPOINT)
    }
    mql.addEventListener('change', onChange)
    setIsMobile(window.innerWidth < MOBILE_BREAKPOINT)
    return () => mql.removeEventListener('change', onChange)
  }, [])

  return !!isMobile
}
\end{DeskFlowCode}

% ── STYLES ──

\phantomsection
\addcontentsline{toc}{subsubsection}{deskflow\-frontend/styles/globals.css}
\DeskFlowListingCaption{deskflow-frontend/styles/globals.css}{lst:fd-styles-globals.css}
\begin{DeskFlowCode}
@import 'tailwindcss';
@import 'tw-animate-css';

@custom-variant dark (&:is(.dark *));

:root {
  -background: oklch(1 0 0);
  -foreground: oklch(0.145 0 0);
  -card: oklch(1 0 0);
  -card-foreground: oklch(0.145 0 0);
  -popover: oklch(1 0 0);
  -popover-foreground: oklch(0.145 0 0);
  -primary: oklch(0.205 0 0);
  -primary-foreground: oklch(0.985 0 0);
  -secondary: oklch(0.97 0 0);
  -secondary-foreground: oklch(0.205 0 0);
  -muted: oklch(0.97 0 0);
  -muted-foreground: oklch(0.556 0 0);
  -accent: oklch(0.97 0 0);
  -accent-foreground: oklch(0.205 0 0);
  -destructive: oklch(0.577 0.245 27.325);
  -destructive-foreground: oklch(0.577 0.245 27.325);
  -border: oklch(0.922 0 0);
  -input: oklch(0.922 0 0);
  -ring: oklch(0.708 0 0);
  -chart-1: oklch(0.646 0.222 41.116);
  -chart-2: oklch(0.6 0.118 184.704);
  -chart-3: oklch(0.398 0.07 227.392);
  -chart-4: oklch(0.828 0.189 84.429);
  -chart-5: oklch(0.769 0.188 70.08);
  -radius: 0.625rem;
  -sidebar: oklch(0.985 0 0);
  -sidebar-foreground: oklch(0.145 0 0);
  -sidebar-primary: oklch(0.205 0 0);
  -sidebar-primary-foreground: oklch(0.985 0 0);
  -sidebar-accent: oklch(0.97 0 0);
  -sidebar-accent-foreground: oklch(0.205 0 0);
  -sidebar-border: oklch(0.922 0 0);
  -sidebar-ring: oklch(0.708 0 0);
}

.dark {
  -background: oklch(0.145 0 0);
  -foreground: oklch(0.985 0 0);
  -card: oklch(0.145 0 0);
  -card-foreground: oklch(0.985 0 0);
  -popover: oklch(0.145 0 0);
  -popover-foreground: oklch(0.985 0 0);
  -primary: oklch(0.985 0 0);
  -primary-foreground: oklch(0.205 0 0);
  -secondary: oklch(0.269 0 0);
  -secondary-foreground: oklch(0.985 0 0);
  -muted: oklch(0.269 0 0);
  -muted-foreground: oklch(0.708 0 0);
  -accent: oklch(0.269 0 0);
  -accent-foreground: oklch(0.985 0 0);
  -destructive: oklch(0.396 0.141 25.723);
  -destructive-foreground: oklch(0.637 0.237 25.331);
  -border: oklch(0.269 0 0);
  -input: oklch(0.269 0 0);
  -ring: oklch(0.439 0 0);
  -chart-1: oklch(0.488 0.243 264.376);
  -chart-2: oklch(0.696 0.17 162.48);
  -chart-3: oklch(0.769 0.188 70.08);
  -chart-4: oklch(0.627 0.265 303.9);
  -chart-5: oklch(0.645 0.246 16.439);
  -sidebar: oklch(0.205 0 0);
  -sidebar-foreground: oklch(0.985 0 0);
  -sidebar-primary: oklch(0.488 0.243 264.376);
  -sidebar-primary-foreground: oklch(0.985 0 0);
  -sidebar-accent: oklch(0.269 0 0);
  -sidebar-accent-foreground: oklch(0.985 0 0);
  -sidebar-border: oklch(0.269 0 0);
  -sidebar-ring: oklch(0.439 0 0);
}

@theme inline {
  -font-sans: 'Geist', 'Geist Fallback';
  -font-mono: 'Geist Mono', 'Geist Mono Fallback';
  -color-background: var(-background);
  -color-foreground: var(-foreground);
  -color-card: var(-card);
  -color-card-foreground: var(-card-foreground);
  -color-popover: var(-popover);
  -color-popover-foreground: var(-popover-foreground);
  -color-primary: var(-primary);
  -color-primary-foreground: var(-primary-foreground);
  -color-secondary: var(-secondary);
  -color-secondary-foreground: var(-secondary-foreground);
  -color-muted: var(-muted);
  -color-muted-foreground: var(-muted-foreground);
  -color-accent: var(-accent);
  -color-accent-foreground: var(-accent-foreground);
  -color-destructive: var(-destructive);
  -color-destructive-foreground: var(-destructive-foreground);
  -color-border: var(-border);
  -color-input: var(-input);
  -color-ring: var(-ring);
  -color-chart-1: var(-chart-1);
  -color-chart-2: var(-chart-2);
  -color-chart-3: var(-chart-3);
  -color-chart-4: var(-chart-4);
  -color-chart-5: var(-chart-5);
  -radius-sm: calc(var(-radius) - 4px);
  -radius-md: calc(var(-radius) - 2px);
  -radius-lg: var(-radius);
  -radius-xl: calc(var(-radius) + 4px);
  -color-sidebar: var(-sidebar);
  -color-sidebar-foreground: var(-sidebar-foreground);
  -color-sidebar-primary: var(-sidebar-primary);
  -color-sidebar-primary-foreground: var(-sidebar-primary-foreground);
  -color-sidebar-accent: var(-sidebar-accent);
  -color-sidebar-accent-foreground: var(-sidebar-accent-foreground);
  -color-sidebar-border: var(-sidebar-border);
  -color-sidebar-ring: var(-sidebar-ring);
}

@layer base {
  * {
    @apply border-border outline-ring/50;
  }
  body {
    @apply bg-background text-foreground;
  }
}
\end{DeskFlowCode}

% ── COMPONENTS ──

\phantomsection
\addcontentsline{toc}{subsubsection}{deskflow\-frontend/components/auth\-landing.tsx}
\DeskFlowListingCaption{deskflow-frontend/components/auth-landing.tsx}{lst:fd-components-auth-landing}
\begin{DeskFlowCode}
"use client"

import { SignInButton, SignOutButton, SignUpButton } from "@clerk/nextjs"
import { Button } from "@/components/ui/button"

interface AuthLandingProps {
  title: string
  subtitle: string
  showAuthActions?: boolean
  showSignOutAction?: boolean
  statusMessage?: string
}

export function AuthLanding({
  title,
  subtitle,
  showAuthActions = true,
  showSignOutAction = false,
  statusMessage = "Brak przypisanej organizacji. Skontaktuj sie z administratorem firmy.",
}: AuthLandingProps) {
  return (
    <div className="relative flex min-h-screen items-center justify-center overflow-hidden bg-slate-950 px-6 py-16 text-slate-100">
      <div className="pointer-events-none absolute inset-0">
        <div className="absolute -left-28 -top-28 h-80 w-80 rounded-full bg-cyan-500/20 blur-3xl" />
        <div className="absolute -right-24 bottom-0 h-96 w-96 rounded-full bg-emerald-500/20 blur-3xl" />
        <div className="absolute inset-0 bg-[radial-gradient(circle_at_20%_20%,rgba(34,211,238,0.12),transparent_35%),radial-gradient(circle_at_80%_80%,rgba(16,185,129,0.12),transparent_40%)]" />
      </div>

      <div className="relative w-full max-w-3xl rounded-2xl border border-slate-700/60 bg-slate-900/70 p-8 shadow-2xl backdrop-blur md:p-12">
        <p className="inline-flex items-center rounded-full border border-cyan-400/30 bg-cyan-500/10 px-3 py-1 text-xs font-semibold tracking-wide text-cyan-300">
          DESKFLOW
        </p>

        <h1 className="mt-6 text-3xl font-bold leading-tight md:text-5xl">{title}</h1>
        <p className="mt-4 max-w-2xl text-sm text-slate-300 md:text-base">{subtitle}</p>

        {showAuthActions ? (
          <div className="mt-8 flex flex-wrap items-center gap-3">
            <SignInButton>
              <Button className="h-11 bg-cyan-500 px-6 text-slate-950 hover:bg-cyan-400">Zaloguj sie</Button>
            </SignInButton>
            <SignUpButton>
              <Button variant="outline" className="h-11 border-slate-300 bg-white px-6 text-black hover:bg-slate-100 hover:text-black">
                Zarejestruj sie
              </Button>
            </SignUpButton>
          </div>
        ) : (
          <div className="mt-8 space-y-3">
            <div className="rounded-lg border border-amber-300/40 bg-amber-500/10 p-3 text-sm text-amber-100">
              {statusMessage}
            </div>
            {showSignOutAction && (
              <SignOutButton>
                <Button variant="outline" className="h-11 border-slate-300 bg-white px-6 text-black hover:bg-slate-100 hover:text-black">
                  Wyloguj sie
                </Button>
              </SignOutButton>
            )}
          </div>
        )}

        <div className="mt-10 grid gap-3 text-xs text-slate-400 md:grid-cols-3">
          <div className="rounded-lg border border-slate-700/70 bg-slate-900/70 p-3">Rezerwacje biurek i sal</div>
          <div className="rounded-lg border border-slate-700/70 bg-slate-900/70 p-3">Wypozyczanie zasobow</div>
          <div className="rounded-lg border border-slate-700/70 bg-slate-900/70 p-3">Panel administracyjny firmy</div>
        </div>
      </div>
    </div>
  )
}
\end{DeskFlowCode}

\phantomsection
\addcontentsline{toc}{subsubsection}{deskflow\-frontend/components/theme\-provider.tsx}
\DeskFlowListingCaption{deskflow-frontend/components/theme-provider.tsx}{lst:fd-components-theme-provider}
\begin{DeskFlowCode}
'use client'

import * as React from 'react'
import {
  ThemeProvider as NextThemesProvider,
  type ThemeProviderProps,
} from 'next-themes'

export function ThemeProvider({ children, ...props }: ThemeProviderProps) {
  return <NextThemesProvider {...props}>{children}</NextThemesProvider>
}
\end{DeskFlowCode}

% ── UI ──

\phantomsection
\addcontentsline{toc}{subsubsection}{deskflow\-frontend/components/ui/accordion.tsx}
\DeskFlowListingCaption{deskflow-frontend/components/ui/accordion.tsx}{lst:fd-ui-accordion}
\begin{DeskFlowCode}
'use client'

import * as React from 'react'
import * as AccordionPrimitive from '@radix-ui/react-accordion'
import { ChevronDownIcon } from 'lucide-react'

import { cn } from '@/lib/utils'

function Accordion({
  ...props
}: React.ComponentProps<typeof AccordionPrimitive.Root>) {
  return <AccordionPrimitive.Root data-slot="accordion" {...props} />
}

function AccordionItem({
  className,
  ...props
}: React.ComponentProps<typeof AccordionPrimitive.Item>) {
  return (
    <AccordionPrimitive.Item
      data-slot="accordion-item"
      className={cn('border-b last:border-b-0', className)}
      {...props}
    />
  )
}

function AccordionTrigger({
  className,
  children,
  ...props
}: React.ComponentProps<typeof AccordionPrimitive.Trigger>) {
  return (
    <AccordionPrimitive.Header className="flex">
      <AccordionPrimitive.Trigger
        data-slot="accordion-trigger"
        className={cn(
          'focus-visible:border-ring focus-visible:ring-ring/50 flex flex-1 items-start justify-between gap-4 rounded-md py-4 text-left text-sm font-medium transition-all outline-none hover:underline focus-visible:ring-[3px] disabled:pointer-events-none disabled:opacity-50 [&[data-state=open]>svg]:rotate-180',
          className,
        )}
        {...props}
      >
        {children}
        <ChevronDownIcon className="text-muted-foreground pointer-events-none size-4 shrink-0 translate-y-0.5 transition-transform duration-200" />
      </AccordionPrimitive.Trigger>
    </AccordionPrimitive.Header>
  )
}

function AccordionContent({
  className,
  children,
  ...props
}: React.ComponentProps<typeof AccordionPrimitive.Content>) {
  return (
    <AccordionPrimitive.Content
      data-slot="accordion-content"
      className="data-[state=closed]:animate-accordion-up data-[state=open]:animate-accordion-down overflow-hidden text-sm"
      {...props}
    >
      <div className={cn('pt-0 pb-4', className)}>{children}</div>
    </AccordionPrimitive.Content>
  )
}

export { Accordion, AccordionItem, AccordionTrigger, AccordionContent }
\end{DeskFlowCode}

\phantomsection
\addcontentsline{toc}{subsubsection}{deskflow\-frontend/components/ui/alert\-dialog.tsx}
\DeskFlowListingCaption{deskflow-frontend/components/ui/alert-dialog.tsx}{lst:fd-ui-alert-dialog}
\begin{DeskFlowCode}
'use client'

import * as React from 'react'
import * as AlertDialogPrimitive from '@radix-ui/react-alert-dialog'

import { cn } from '@/lib/utils'
import { buttonVariants } from '@/components/ui/button'

function AlertDialog({
  ...props
}: React.ComponentProps<typeof AlertDialogPrimitive.Root>) {
  return <AlertDialogPrimitive.Root data-slot="alert-dialog" {...props} />
}

function AlertDialogTrigger({
  ...props
}: React.ComponentProps<typeof AlertDialogPrimitive.Trigger>) {
  return (
    <AlertDialogPrimitive.Trigger data-slot="alert-dialog-trigger" {...props} />
  )
}

function AlertDialogPortal({
  ...props
}: React.ComponentProps<typeof AlertDialogPrimitive.Portal>) {
  return (
    <AlertDialogPrimitive.Portal data-slot="alert-dialog-portal" {...props} />
  )
}

function AlertDialogOverlay({
  className,
  ...props
}: React.ComponentProps<typeof AlertDialogPrimitive.Overlay>) {
  return (
    <AlertDialogPrimitive.Overlay
      data-slot="alert-dialog-overlay"
      className={cn(
        'data-[state=open]:animate-in data-[state=closed]:animate-out data-[state=closed]:fade-out-0 data-[state=open]:fade-in-0 fixed inset-0 z-50 bg-black/50',
        className,
      )}
      {...props}
    />
  )
}

function AlertDialogContent({
  className,
  ...props
}: React.ComponentProps<typeof AlertDialogPrimitive.Content>) {
  return (
    <AlertDialogPortal>
      <AlertDialogOverlay />
      <AlertDialogPrimitive.Content
        data-slot="alert-dialog-content"
        className={cn(
          'bg-background data-[state=open]:animate-in data-[state=closed]:animate-out data-[state=closed]:fade-out-0 data-[state=open]:fade-in-0 data-[state=closed]:zoom-out-95 data-[state=open]:zoom-in-95 fixed top-[50%] left-[50%] z-50 grid w-full max-w-[calc(100%-2rem)] translate-x-[-50%] translate-y-[-50%] gap-4 rounded-lg border p-6 shadow-lg duration-200 sm:max-w-lg',
          className,
        )}
        {...props}
      />
    </AlertDialogPortal>
  )
}

function AlertDialogHeader({
  className,
  ...props
}: React.ComponentProps<'div'>) {
  return (
    <div
      data-slot="alert-dialog-header"
      className={cn('flex flex-col gap-2 text-center sm:text-left', className)}
      {...props}
    />
  )
}

function AlertDialogFooter({
  className,
  ...props
}: React.ComponentProps<'div'>) {
  return (
    <div
      data-slot="alert-dialog-footer"
      className={cn(
        'flex flex-col-reverse gap-2 sm:flex-row sm:justify-end',
        className,
      )}
      {...props}
    />
  )
}

function AlertDialogTitle({
  className,
  ...props
}: React.ComponentProps<typeof AlertDialogPrimitive.Title>) {
  return (
    <AlertDialogPrimitive.Title
      data-slot="alert-dialog-title"
      className={cn('text-lg font-semibold', className)}
      {...props}
    />
  )
}

function AlertDialogDescription({
  className,
  ...props
}: React.ComponentProps<typeof AlertDialogPrimitive.Description>) {
  return (
    <AlertDialogPrimitive.Description
      data-slot="alert-dialog-description"
      className={cn('text-muted-foreground text-sm', className)}
      {...props}
    />
  )
}

function AlertDialogAction({
  className,
  ...props
}: React.ComponentProps<typeof AlertDialogPrimitive.Action>) {
  return (
    <AlertDialogPrimitive.Action
      className={cn(buttonVariants(), className)}
      {...props}
    />
  )
}

function AlertDialogCancel({
  className,
  ...props
}: React.ComponentProps<typeof AlertDialogPrimitive.Cancel>) {
  return (
    <AlertDialogPrimitive.Cancel
      className={cn(buttonVariants({ variant: 'outline' }), className)}
      {...props}
    />
  )
}

export {
  AlertDialog,
  AlertDialogPortal,
  AlertDialogOverlay,
  AlertDialogTrigger,
  AlertDialogContent,
  AlertDialogHeader,
  AlertDialogFooter,
  AlertDialogTitle,
  AlertDialogDescription,
  AlertDialogAction,
  AlertDialogCancel,
}
\end{DeskFlowCode}

\phantomsection
\addcontentsline{toc}{subsubsection}{deskflow\-frontend/components/ui/alert.tsx}
\DeskFlowListingCaption{deskflow-frontend/components/ui/alert.tsx}{lst:fd-ui-alert}
\begin{DeskFlowCode}
import * as React from 'react'
import { cva, type VariantProps } from 'class-variance-authority'

import { cn } from '@/lib/utils'

const alertVariants = cva(
  'relative w-full rounded-lg border px-4 py-3 text-sm grid has-[>svg]:grid-cols-[calc(var(-spacing)*4)_1fr] grid-cols-[0_1fr] has-[>svg]:gap-x-3 gap-y-0.5 items-start [&>svg]:size-4 [&>svg]:translate-y-0.5 [&>svg]:text-current',
  {
    variants: {
      variant: {
        default: 'bg-card text-card-foreground',
        destructive:
          'text-destructive bg-card [&>svg]:text-current *:data-[slot=alert-description]:text-destructive/90',
      },
    },
    defaultVariants: {
      variant: 'default',
    },
  },
)

function Alert({
  className,
  variant,
  ...props
}: React.ComponentProps<'div'> & VariantProps<typeof alertVariants>) {
  return (
    <div
      data-slot="alert"
      role="alert"
      className={cn(alertVariants({ variant }), className)}
      {...props}
    />
  )
}

function AlertTitle({ className, ...props }: React.ComponentProps<'div'>) {
  return (
    <div
      data-slot="alert-title"
      className={cn(
        'col-start-2 line-clamp-1 min-h-4 font-medium tracking-tight',
        className,
      )}
      {...props}
    />
  )
}

function AlertDescription({
  className,
  ...props
}: React.ComponentProps<'div'>) {
  return (
    <div
      data-slot="alert-description"
      className={cn(
        'text-muted-foreground col-start-2 grid justify-items-start gap-1 text-sm [&_p]:leading-relaxed',
        className,
      )}
      {...props}
    />
  )
}

export { Alert, AlertTitle, AlertDescription }
\end{DeskFlowCode}

\phantomsection
\addcontentsline{toc}{subsubsection}{deskflow\-frontend/components/ui/aspect\-ratio.tsx}
\DeskFlowListingCaption{deskflow-frontend/components/ui/aspect-ratio.tsx}{lst:fd-ui-aspect-ratio}
\begin{DeskFlowCode}
'use client'

import * as AspectRatioPrimitive from '@radix-ui/react-aspect-ratio'

function AspectRatio({
  ...props
}: React.ComponentProps<typeof AspectRatioPrimitive.Root>) {
  return <AspectRatioPrimitive.Root data-slot="aspect-ratio" {...props} />
}

export { AspectRatio }
\end{DeskFlowCode}

\phantomsection
\addcontentsline{toc}{subsubsection}{deskflow\-frontend/components/ui/avatar.tsx}
\DeskFlowListingCaption{deskflow-frontend/components/ui/avatar.tsx}{lst:fd-ui-avatar}
\begin{DeskFlowCode}
'use client'

import * as React from 'react'
import * as AvatarPrimitive from '@radix-ui/react-avatar'

import { cn } from '@/lib/utils'

function Avatar({
  className,
  ...props
}: React.ComponentProps<typeof AvatarPrimitive.Root>) {
  return (
    <AvatarPrimitive.Root
      data-slot="avatar"
      className={cn(
        'relative flex size-8 shrink-0 overflow-hidden rounded-full',
        className,
      )}
      {...props}
    />
  )
}

function AvatarImage({
  className,
  ...props
}: React.ComponentProps<typeof AvatarPrimitive.Image>) {
  return (
    <AvatarPrimitive.Image
      data-slot="avatar-image"
      className={cn('aspect-square size-full', className)}
      {...props}
    />
  )
}

function AvatarFallback({
  className,
  ...props
}: React.ComponentProps<typeof AvatarPrimitive.Fallback>) {
  return (
    <AvatarPrimitive.Fallback
      data-slot="avatar-fallback"
      className={cn(
        'bg-muted flex size-full items-center justify-center rounded-full',
        className,
      )}
      {...props}
    />
  )
}

export { Avatar, AvatarImage, AvatarFallback }
\end{DeskFlowCode}

\phantomsection
\addcontentsline{toc}{subsubsection}{deskflow\-frontend/components/ui/badge.tsx}
\DeskFlowListingCaption{deskflow-frontend/components/ui/badge.tsx}{lst:fd-ui-badge}
\begin{DeskFlowCode}
import * as React from 'react'
import { Slot } from '@radix-ui/react-slot'
import { cva, type VariantProps } from 'class-variance-authority'

import { cn } from '@/lib/utils'

const badgeVariants = cva(
  'inline-flex items-center justify-center rounded-md border px-2 py-0.5 text-xs font-medium w-fit whitespace-nowrap shrink-0 [&>svg]:size-3 gap-1 [&>svg]:pointer-events-none focus-visible:border-ring focus-visible:ring-ring/50 focus-visible:ring-[3px] aria-invalid:ring-destructive/20 dark:aria-invalid:ring-destructive/40 aria-invalid:border-destructive transition-[color,box-shadow] overflow-hidden',
  {
    variants: {
      variant: {
        default:
          'border-transparent bg-primary text-primary-foreground [a&]:hover:bg-primary/90',
        secondary:
          'border-transparent bg-secondary text-secondary-foreground [a&]:hover:bg-secondary/90',
        destructive:
          'border-transparent bg-destructive text-white [a&]:hover:bg-destructive/90 focus-visible:ring-destructive/20 dark:focus-visible:ring-destructive/40 dark:bg-destructive/60',
        outline:
          'text-foreground [a&]:hover:bg-accent [a&]:hover:text-accent-foreground',
      },
    },
    defaultVariants: {
      variant: 'default',
    },
  },
)

function Badge({
  className,
  variant,
  asChild = false,
  ...props
}: React.ComponentProps<'span'> &
  VariantProps<typeof badgeVariants> & { asChild?: boolean }) {
  const Comp = asChild ? Slot : 'span'

  return (
    <Comp
      data-slot="badge"
      className={cn(badgeVariants({ variant }), className)}
      {...props}
    />
  )
}

export { Badge, badgeVariants }
\end{DeskFlowCode}

\phantomsection
\addcontentsline{toc}{subsubsection}{deskflow\-frontend/components/ui/breadcrumb.tsx}
\DeskFlowListingCaption{deskflow-frontend/components/ui/breadcrumb.tsx}{lst:fd-ui-breadcrumb}
\begin{DeskFlowCode}
import * as React from 'react'
import { Slot } from '@radix-ui/react-slot'
import { ChevronRight, MoreHorizontal } from 'lucide-react'

import { cn } from '@/lib/utils'

function Breadcrumb({ ...props }: React.ComponentProps<'nav'>) {
  return <nav aria-label="breadcrumb" data-slot="breadcrumb" {...props} />
}

function BreadcrumbList({ className, ...props }: React.ComponentProps<'ol'>) {
  return (
    <ol
      data-slot="breadcrumb-list"
      className={cn(
        'text-muted-foreground flex flex-wrap items-center gap-1.5 text-sm break-words sm:gap-2.5',
        className,
      )}
      {...props}
    />
  )
}

function BreadcrumbItem({ className, ...props }: React.ComponentProps<'li'>) {
  return (
    <li
      data-slot="breadcrumb-item"
      className={cn('inline-flex items-center gap-1.5', className)}
      {...props}
    />
  )
}

function BreadcrumbLink({
  asChild,
  className,
  ...props
}: React.ComponentProps<'a'> & {
  asChild?: boolean
}) {
  const Comp = asChild ? Slot : 'a'

  return (
    <Comp
      data-slot="breadcrumb-link"
      className={cn('hover:text-foreground transition-colors', className)}
      {...props}
    />
  )
}

function BreadcrumbPage({ className, ...props }: React.ComponentProps<'span'>) {
  return (
    <span
      data-slot="breadcrumb-page"
      role="link"
      aria-disabled="true"
      aria-current="page"
      className={cn('text-foreground font-normal', className)}
      {...props}
    />
  )
}

function BreadcrumbSeparator({
  children,
  className,
  ...props
}: React.ComponentProps<'li'>) {
  return (
    <li
      data-slot="breadcrumb-separator"
      role="presentation"
      aria-hidden="true"
      className={cn('[&>svg]:size-3.5', className)}
      {...props}
    >
      {children ?? <ChevronRight />}
    </li>
  )
}

function BreadcrumbEllipsis({
  className,
  ...props
}: React.ComponentProps<'span'>) {
  return (
    <span
      data-slot="breadcrumb-ellipsis"
      role="presentation"
      aria-hidden="true"
      className={cn('flex size-9 items-center justify-center', className)}
      {...props}
    >
      <MoreHorizontal className="size-4" />
      <span className="sr-only">More</span>
    </span>
  )
}

export {
  Breadcrumb,
  BreadcrumbList,
  BreadcrumbItem,
  BreadcrumbLink,
  BreadcrumbPage,
  BreadcrumbSeparator,
  BreadcrumbEllipsis,
}
\end{DeskFlowCode}

\phantomsection
\addcontentsline{toc}{subsubsection}{deskflow\-frontend/components/ui/button\-group.tsx}
\DeskFlowListingCaption{deskflow-frontend/components/ui/button-group.tsx}{lst:fd-ui-button-group}
\begin{DeskFlowCode}
import { Slot } from '@radix-ui/react-slot'
import { cva, type VariantProps } from 'class-variance-authority'

import { cn } from '@/lib/utils'
import { Separator } from '@/components/ui/separator'

const buttonGroupVariants = cva(
  "flex w-fit items-stretch [&>*]:focus-visible:z-10 [&>*]:focus-visible:relative [&>[data-slot=select-trigger]:not([class*='w-'])]:w-fit [&>input]:flex-1 has-[select[aria-hidden=true]:last-child]:[&>[data-slot=select-trigger]:last-of-type]:rounded-r-md has-[>[data-slot=button-group]]:gap-2",
  {
    variants: {
      orientation: {
        horizontal:
          '[&>*:not(:first-child)]:rounded-l-none [&>*:not(:first-child)]:border-l-0 [&>*:not(:last-child)]:rounded-r-none',
        vertical:
          'flex-col [&>*:not(:first-child)]:rounded-t-none [&>*:not(:first-child)]:border-t-0 [&>*:not(:last-child)]:rounded-b-none',
      },
    },
    defaultVariants: {
      orientation: 'horizontal',
    },
  },
)

function ButtonGroup({
  className,
  orientation,
  ...props
}: React.ComponentProps<'div'> & VariantProps<typeof buttonGroupVariants>) {
  return (
    <div
      role="group"
      data-slot="button-group"
      data-orientation={orientation}
      className={cn(buttonGroupVariants({ orientation }), className)}
      {...props}
    />
  )
}

function ButtonGroupText({
  className,
  asChild = false,
  ...props
}: React.ComponentProps<'div'> & {
  asChild?: boolean
}) {
  const Comp = asChild ? Slot : 'div'

  return (
    <Comp
      className={cn(
        "bg-muted flex items-center gap-2 rounded-md border px-4 text-sm font-medium shadow-xs [&_svg]:pointer-events-none [&_svg:not([class*='size-'])]:size-4",
        className,
      )}
      {...props}
    />
  )
}

function ButtonGroupSeparator({
  className,
  orientation = 'vertical',
  ...props
}: React.ComponentProps<typeof Separator>) {
  return (
    <Separator
      data-slot="button-group-separator"
      orientation={orientation}
      className={cn(
        'bg-input relative !m-0 self-stretch data-[orientation=vertical]:h-auto',
        className,
      )}
      {...props}
    />
  )
}

export {
  ButtonGroup,
  ButtonGroupSeparator,
  ButtonGroupText,
  buttonGroupVariants,
}
\end{DeskFlowCode}

\phantomsection
\addcontentsline{toc}{subsubsection}{deskflow\-frontend/components/ui/button.tsx}
\DeskFlowListingCaption{deskflow-frontend/components/ui/button.tsx}{lst:fd-ui-button}
\begin{DeskFlowCode}
import * as React from 'react'
import { Slot } from '@radix-ui/react-slot'
import { cva, type VariantProps } from 'class-variance-authority'

import { cn } from '@/lib/utils'

const buttonVariants = cva(
  "inline-flex items-center justify-center gap-2 whitespace-nowrap rounded-md text-sm font-medium transition-all disabled:pointer-events-none disabled:opacity-50 [&_svg]:pointer-events-none [&_svg:not([class*='size-'])]:size-4 shrink-0 [&_svg]:shrink-0 outline-none focus-visible:border-ring focus-visible:ring-ring/50 focus-visible:ring-[3px] aria-invalid:ring-destructive/20 dark:aria-invalid:ring-destructive/40 aria-invalid:border-destructive",
  {
    variants: {
      variant: {
        default: 'bg-primary text-primary-foreground hover:bg-primary/90',
        destructive:
          'bg-destructive text-white hover:bg-destructive/90 focus-visible:ring-destructive/20 dark:focus-visible:ring-destructive/40 dark:bg-destructive/60',
        outline:
          'border bg-background shadow-xs hover:bg-accent hover:text-accent-foreground dark:bg-input/30 dark:border-input dark:hover:bg-input/50',
        secondary:
          'bg-secondary text-secondary-foreground hover:bg-secondary/80',
        ghost:
          'hover:bg-accent hover:text-accent-foreground dark:hover:bg-accent/50',
        link: 'text-primary underline-offset-4 hover:underline',
      },
      size: {
        default: 'h-9 px-4 py-2 has-[>svg]:px-3',
        sm: 'h-8 rounded-md gap-1.5 px-3 has-[>svg]:px-2.5',
        lg: 'h-10 rounded-md px-6 has-[>svg]:px-4',
        icon: 'size-9',
        'icon-sm': 'size-8',
        'icon-lg': 'size-10',
      },
    },
    defaultVariants: {
      variant: 'default',
      size: 'default',
    },
  },
)

function Button({
  className,
  variant,
  size,
  asChild = false,
  ...props
}: React.ComponentProps<'button'> &
  VariantProps<typeof buttonVariants> & {
    asChild?: boolean
  }) {
  const Comp = asChild ? Slot : 'button'

  return (
    <Comp
      data-slot="button"
      className={cn(buttonVariants({ variant, size, className }))}
      {...props}
    />
  )
}

export { Button, buttonVariants }
\end{DeskFlowCode}

\phantomsection
\addcontentsline{toc}{subsubsection}{deskflow\-frontend/components/ui/calendar.tsx}
\DeskFlowListingCaption{deskflow-frontend/components/ui/calendar.tsx}{lst:fd-ui-calendar}
\begin{DeskFlowCode}
'use client'

import * as React from 'react'
import {
  ChevronDownIcon,
  ChevronLeftIcon,
  ChevronRightIcon,
} from 'lucide-react'
import { DayButton, DayPicker, getDefaultClassNames } from 'react-day-picker'

import { cn } from '@/lib/utils'
import { Button, buttonVariants } from '@/components/ui/button'

function Calendar({
  className,
  classNames,
  showOutsideDays = true,
  captionLayout = 'label',
  buttonVariant = 'ghost',
  formatters,
  components,
  ...props
}: React.ComponentProps<typeof DayPicker> & {
  buttonVariant?: React.ComponentProps<typeof Button>['variant']
}) {
  const defaultClassNames = getDefaultClassNames()

  return (
    <DayPicker
      showOutsideDays={showOutsideDays}
      className={cn(
        'bg-background group/calendar p-3 [-cell-size:-spacing(8)] [[data-slot=card-content]_&]:bg-transparent [[data-slot=popover-content]_&]:bg-transparent',
        String.raw`rtl:**:[.rdp-button\_next>svg]:rotate-180`,
        String.raw`rtl:**:[.rdp-button\_previous>svg]:rotate-180`,
        className,
      )}
      captionLayout={captionLayout}
      formatters={{
        formatMonthDropdown: (date) =>
          date.toLocaleString('default', { month: 'short' }),
        ...formatters,
      }}
      classNames={{
        root: cn('w-fit', defaultClassNames.root),
        months: cn(
          'flex gap-4 flex-col md:flex-row relative',
          defaultClassNames.months,
        ),
        month: cn('flex flex-col w-full gap-4', defaultClassNames.month),
        nav: cn(
          'flex items-center gap-1 w-full absolute top-0 inset-x-0 justify-between',
          defaultClassNames.nav,
        ),
        button_previous: cn(
          buttonVariants({ variant: buttonVariant }),
          'size-(-cell-size) aria-disabled:opacity-50 p-0 select-none',
          defaultClassNames.button_previous,
        ),
        button_next: cn(
          buttonVariants({ variant: buttonVariant }),
          'size-(-cell-size) aria-disabled:opacity-50 p-0 select-none',
          defaultClassNames.button_next,
        ),
        month_caption: cn(
          'flex items-center justify-center h-(-cell-size) w-full px-(-cell-size)',
          defaultClassNames.month_caption,
        ),
        dropdowns: cn(
          'w-full flex items-center text-sm font-medium justify-center h-(-cell-size) gap-1.5',
          defaultClassNames.dropdowns,
        ),
        dropdown_root: cn(
          'relative has-focus:border-ring border border-input shadow-xs has-focus:ring-ring/50 has-focus:ring-[3px] rounded-md',
          defaultClassNames.dropdown_root,
        ),
        dropdown: cn(
          'absolute bg-popover inset-0 opacity-0',
          defaultClassNames.dropdown,
        ),
        caption_label: cn(
          'select-none font-medium',
          captionLayout === 'label'
            ? 'text-sm'
            : 'rounded-md pl-2 pr-1 flex items-center gap-1 text-sm h-8 [&>svg]:text-muted-foreground [&>svg]:size-3.5',
          defaultClassNames.caption_label,
        ),
        table: 'w-full border-collapse',
        weekdays: cn('flex', defaultClassNames.weekdays),
        weekday: cn(
          'text-muted-foreground rounded-md flex-1 font-normal text-[0.8rem] select-none',
          defaultClassNames.weekday,
        ),
        week: cn('flex w-full mt-2', defaultClassNames.week),
        week_number_header: cn(
          'select-none w-(-cell-size)',
          defaultClassNames.week_number_header,
        ),
        week_number: cn(
          'text-[0.8rem] select-none text-muted-foreground',
          defaultClassNames.week_number,
        ),
        day: cn(
          'relative w-full h-full p-0 text-center [&:first-child[data-selected=true]_button]:rounded-l-md [&:last-child[data-selected=true]_button]:rounded-r-md group/day aspect-square select-none',
          defaultClassNames.day,
        ),
        range_start: cn(
          'rounded-l-md bg-accent',
          defaultClassNames.range_start,
        ),
        range_middle: cn('rounded-none', defaultClassNames.range_middle),
        range_end: cn('rounded-r-md bg-accent', defaultClassNames.range_end),
        today: cn(
          'bg-accent text-accent-foreground rounded-md data-[selected=true]:rounded-none',
          defaultClassNames.today,
        ),
        outside: cn(
          'text-muted-foreground aria-selected:text-muted-foreground',
          defaultClassNames.outside,
        ),
        disabled: cn(
          'text-muted-foreground opacity-50',
          defaultClassNames.disabled,
        ),
        hidden: cn('invisible', defaultClassNames.hidden),
        ...classNames,
      }}
      components={{
        Root: ({ className, rootRef, ...props }) => {
          return (
            <div
              data-slot="calendar"
              ref={rootRef}
              className={cn(className)}
              {...props}
            />
          )
        },
        Chevron: ({ className, orientation, ...props }) => {
          if (orientation === 'left') {
            return (
              <ChevronLeftIcon className={cn('size-4', className)} {...props} />
            )
          }

          if (orientation === 'right') {
            return (
              <ChevronRightIcon
                className={cn('size-4', className)}
                {...props}
              />
            )
          }

          return (
            <ChevronDownIcon className={cn('size-4', className)} {...props} />
          )
        },
        DayButton: CalendarDayButton,
        WeekNumber: ({ children, ...props }) => {
          return (
            <td {...props}>
              <div className="flex size-(-cell-size) items-center justify-center text-center">
                {children}
              </div>
            </td>
          )
        },
        ...components,
      }}
      {...props}
    />
  )
}

function CalendarDayButton({
  className,
  day,
  modifiers,
  ...props
}: React.ComponentProps<typeof DayButton>) {
  const defaultClassNames = getDefaultClassNames()

  const ref = React.useRef<HTMLButtonElement>(null)
  React.useEffect(() => {
    if (modifiers.focused) ref.current?.focus()
  }, [modifiers.focused])

  return (
    <Button
      ref={ref}
      variant="ghost"
      size="icon"
      data-day={day.date.toLocaleDateString()}
      data-selected-single={
        modifiers.selected &&
        !modifiers.range_start &&
        !modifiers.range_end &&
        !modifiers.range_middle
      }
      data-range-start={modifiers.range_start}
      data-range-end={modifiers.range_end}
      data-range-middle={modifiers.range_middle}
      className={cn(
        'data-[selected-single=true]:bg-primary data-[selected-single=true]:text-primary-foreground data-[range-middle=true]:bg-accent data-[range-middle=true]:text-accent-foreground data-[range-start=true]:bg-primary data-[range-start=true]:text-primary-foreground data-[range-end=true]:bg-primary data-[range-end=true]:text-primary-foreground group-data-[focused=true]/day:border-ring group-data-[focused=true]/day:ring-ring/50 dark:hover:text-accent-foreground flex aspect-square size-auto w-full min-w-(-cell-size) flex-col gap-1 leading-none font-normal group-data-[focused=true]/day:relative group-data-[focused=true]/day:z-10 group-data-[focused=true]/day:ring-[3px] data-[range-end=true]:rounded-md data-[range-end=true]:rounded-r-md data-[range-middle=true]:rounded-none data-[range-start=true]:rounded-md data-[range-start=true]:rounded-l-md [&>span]:text-xs [&>span]:opacity-70',
        defaultClassNames.day,
        className,
      )}
      {...props}
    />
  )
}

export { Calendar, CalendarDayButton }
\end{DeskFlowCode}

\phantomsection
\addcontentsline{toc}{subsubsection}{deskflow\-frontend/components/ui/card.tsx}
\DeskFlowListingCaption{deskflow-frontend/components/ui/card.tsx}{lst:fd-ui-card}
\begin{DeskFlowCode}
import * as React from 'react'

import { cn } from '@/lib/utils'

function Card({ className, ...props }: React.ComponentProps<'div'>) {
  return (
    <div
      data-slot="card"
      className={cn(
        'bg-card text-card-foreground flex flex-col gap-6 rounded-xl border py-6 shadow-sm',
        className,
      )}
      {...props}
    />
  )
}

function CardHeader({ className, ...props }: React.ComponentProps<'div'>) {
  return (
    <div
      data-slot="card-header"
      className={cn(
        '@container/card-header grid auto-rows-min grid-rows-[auto_auto] items-start gap-2 px-6 has-data-[slot=card-action]:grid-cols-[1fr_auto] [.border-b]:pb-6',
        className,
      )}
      {...props}
    />
  )
}

function CardTitle({ className, ...props }: React.ComponentProps<'div'>) {
  return (
    <div
      data-slot="card-title"
      className={cn('leading-none font-semibold', className)}
      {...props}
    />
  )
}

function CardDescription({ className, ...props }: React.ComponentProps<'div'>) {
  return (
    <div
      data-slot="card-description"
      className={cn('text-muted-foreground text-sm', className)}
      {...props}
    />
  )
}

function CardAction({ className, ...props }: React.ComponentProps<'div'>) {
  return (
    <div
      data-slot="card-action"
      className={cn(
        'col-start-2 row-span-2 row-start-1 self-start justify-self-end',
        className,
      )}
      {...props}
    />
  )
}

function CardContent({ className, ...props }: React.ComponentProps<'div'>) {
  return (
    <div
      data-slot="card-content"
      className={cn('px-6', className)}
      {...props}
    />
  )
}

function CardFooter({ className, ...props }: React.ComponentProps<'div'>) {
  return (
    <div
      data-slot="card-footer"
      className={cn('flex items-center px-6 [.border-t]:pt-6', className)}
      {...props}
    />
  )
}

export {
  Card,
  CardHeader,
  CardFooter,
  CardTitle,
  CardAction,
  CardDescription,
  CardContent,
}
\end{DeskFlowCode}

\phantomsection
\addcontentsline{toc}{subsubsection}{deskflow\-frontend/components/ui/carousel.tsx}
\DeskFlowListingCaption{deskflow-frontend/components/ui/carousel.tsx}{lst:fd-ui-carousel}
\begin{DeskFlowCode}
'use client'

import * as React from 'react'
import useEmblaCarousel, {
  type UseEmblaCarouselType,
} from 'embla-carousel-react'
import { ArrowLeft, ArrowRight } from 'lucide-react'

import { cn } from '@/lib/utils'
import { Button } from '@/components/ui/button'

type CarouselApi = UseEmblaCarouselType[1]
type UseCarouselParameters = Parameters<typeof useEmblaCarousel>
type CarouselOptions = UseCarouselParameters[0]
type CarouselPlugin = UseCarouselParameters[1]

type CarouselProps = {
  opts?: CarouselOptions
  plugins?: CarouselPlugin
  orientation?: 'horizontal' | 'vertical'
  setApi?: (api: CarouselApi) => void
}

type CarouselContextProps = {
  carouselRef: ReturnType<typeof useEmblaCarousel>[0]
  api: ReturnType<typeof useEmblaCarousel>[1]
  scrollPrev: () => void
  scrollNext: () => void
  canScrollPrev: boolean
  canScrollNext: boolean
} & CarouselProps

const CarouselContext = React.createContext<CarouselContextProps | null>(null)

function useCarousel() {
  const context = React.useContext(CarouselContext)

  if (!context) {
    throw new Error('useCarousel must be used within a <Carousel />')
  }

  return context
}

function Carousel({
  orientation = 'horizontal',
  opts,
  setApi,
  plugins,
  className,
  children,
  ...props
}: React.ComponentProps<'div'> & CarouselProps) {
  const [carouselRef, api] = useEmblaCarousel(
    {
      ...opts,
      axis: orientation === 'horizontal' ? 'x' : 'y',
    },
    plugins,
  )
  const [canScrollPrev, setCanScrollPrev] = React.useState(false)
  const [canScrollNext, setCanScrollNext] = React.useState(false)

  const onSelect = React.useCallback((api: CarouselApi) => {
    if (!api) return
    setCanScrollPrev(api.canScrollPrev())
    setCanScrollNext(api.canScrollNext())
  }, [])

  const scrollPrev = React.useCallback(() => {
    api?.scrollPrev()
  }, [api])

  const scrollNext = React.useCallback(() => {
    api?.scrollNext()
  }, [api])

  const handleKeyDown = React.useCallback(
    (event: React.KeyboardEvent<HTMLDivElement>) => {
      if (event.key === 'ArrowLeft') {
        event.preventDefault()
        scrollPrev()
      } else if (event.key === 'ArrowRight') {
        event.preventDefault()
        scrollNext()
      }
    },
    [scrollPrev, scrollNext],
  )

  React.useEffect(() => {
    if (!api || !setApi) return
    setApi(api)
  }, [api, setApi])

  React.useEffect(() => {
    if (!api) return
    onSelect(api)
    api.on('reInit', onSelect)
    api.on('select', onSelect)

    return () => {
      api?.off('select', onSelect)
    }
  }, [api, onSelect])

  return (
    <CarouselContext.Provider
      value={{
        carouselRef,
        api: api,
        opts,
        orientation:
          orientation || (opts?.axis === 'y' ? 'vertical' : 'horizontal'),
        scrollPrev,
        scrollNext,
        canScrollPrev,
        canScrollNext,
      }}
    >
      <div
        onKeyDownCapture={handleKeyDown}
        className={cn('relative', className)}
        role="region"
        aria-roledescription="carousel"
        data-slot="carousel"
        {...props}
      >
        {children}
      </div>
    </CarouselContext.Provider>
  )
}

function CarouselContent({ className, ...props }: React.ComponentProps<'div'>) {
  const { carouselRef, orientation } = useCarousel()

  return (
    <div
      ref={carouselRef}
      className="overflow-hidden"
      data-slot="carousel-content"
    >
      <div
        className={cn(
          'flex',
          orientation === 'horizontal' ? '-ml-4' : '-mt-4 flex-col',
          className,
        )}
        {...props}
      />
    </div>
  )
}

function CarouselItem({ className, ...props }: React.ComponentProps<'div'>) {
  const { orientation } = useCarousel()

  return (
    <div
      role="group"
      aria-roledescription="slide"
      data-slot="carousel-item"
      className={cn(
        'min-w-0 shrink-0 grow-0 basis-full',
        orientation === 'horizontal' ? 'pl-4' : 'pt-4',
        className,
      )}
      {...props}
    />
  )
}

function CarouselPrevious({
  className,
  variant = 'outline',
  size = 'icon',
  ...props
}: React.ComponentProps<typeof Button>) {
  const { orientation, scrollPrev, canScrollPrev } = useCarousel()

  return (
    <Button
      data-slot="carousel-previous"
      variant={variant}
      size={size}
      className={cn(
        'absolute size-8 rounded-full',
        orientation === 'horizontal'
          ? 'top-1/2 -left-12 -translate-y-1/2'
          : '-top-12 left-1/2 -translate-x-1/2 rotate-90',
        className,
      )}
      disabled={!canScrollPrev}
      onClick={scrollPrev}
      {...props}
    >
      <ArrowLeft />
      <span className="sr-only">Previous slide</span>
    </Button>
  )
}

function CarouselNext({
  className,
  variant = 'outline',
  size = 'icon',
  ...props
}: React.ComponentProps<typeof Button>) {
  const { orientation, scrollNext, canScrollNext } = useCarousel()

  return (
    <Button
      data-slot="carousel-next"
      variant={variant}
      size={size}
      className={cn(
        'absolute size-8 rounded-full',
        orientation === 'horizontal'
          ? 'top-1/2 -right-12 -translate-y-1/2'
          : '-bottom-12 left-1/2 -translate-x-1/2 rotate-90',
        className,
      )}
      disabled={!canScrollNext}
      onClick={scrollNext}
      {...props}
    >
      <ArrowRight />
      <span className="sr-only">Next slide</span>
    </Button>
  )
}

export {
  type CarouselApi,
  Carousel,
  CarouselContent,
  CarouselItem,
  CarouselPrevious,
  CarouselNext,
}
\end{DeskFlowCode}

\phantomsection
\addcontentsline{toc}{subsubsection}{deskflow\-frontend/components/ui/chart.tsx}
\DeskFlowListingCaption{deskflow-frontend/components/ui/chart.tsx}{lst:fd-ui-chart}
\begin{DeskFlowCode}
'use client'

import * as React from 'react'
import * as RechartsPrimitive from 'recharts'

import { cn } from '@/lib/utils'

// Format: { THEME_NAME: CSS_SELECTOR }
const THEMES = { light: '', dark: '.dark' } as const

export type ChartConfig = {
  [k in string]: {
    label?: React.ReactNode
    icon?: React.ComponentType
  } & (
    | { color?: string; theme?: never }
    | { color?: never; theme: Record<keyof typeof THEMES, string> }
  )
}

type ChartContextProps = {
  config: ChartConfig
}

const ChartContext = React.createContext<ChartContextProps | null>(null)

function useChart() {
  const context = React.useContext(ChartContext)

  if (!context) {
    throw new Error('useChart must be used within a <ChartContainer />')
  }

  return context
}

function ChartContainer({
  id,
  className,
  children,
  config,
  ...props
}: React.ComponentProps<'div'> & {
  config: ChartConfig
  children: React.ComponentProps<
    typeof RechartsPrimitive.ResponsiveContainer
  >['children']
}) {
  const uniqueId = React.useId()
  const chartId = `chart-${id || uniqueId.replace(/:/g, '')}`

  return (
    <ChartContext.Provider value={{ config }}>
      <div
        data-slot="chart"
        data-chart={chartId}
        className={cn(
          "[&_.recharts-cartesian-axis-tick_text]:fill-muted-foreground [&_.recharts-cartesian-grid_line[stroke='#ccc']]:stroke-border/50 [&_.recharts-curve.recharts-tooltip-cursor]:stroke-border [&_.recharts-polar-grid_[stroke='#ccc']]:stroke-border [&_.recharts-radial-bar-background-sector]:fill-muted [&_.recharts-rectangle.recharts-tooltip-cursor]:fill-muted [&_.recharts-reference-line_[stroke='#ccc']]:stroke-border flex aspect-video justify-center text-xs [&_.recharts-dot[stroke='#fff']]:stroke-transparent [&_.recharts-layer]:outline-hidden [&_.recharts-sector]:outline-hidden [&_.recharts-sector[stroke='#fff']]:stroke-transparent [&_.recharts-surface]:outline-hidden",
          className,
        )}
        {...props}
      >
        <ChartStyle id={chartId} config={config} />
        <RechartsPrimitive.ResponsiveContainer>
          {children}
        </RechartsPrimitive.ResponsiveContainer>
      </div>
    </ChartContext.Provider>
  )
}

const ChartStyle = ({ id, config }: { id: string; config: ChartConfig }) => {
  const colorConfig = Object.entries(config).filter(
    ([, config]) => config.theme || config.color,
  )

  if (!colorConfig.length) {
    return null
  }

  return (
    <style
      dangerouslySetInnerHTML={{
        __html: Object.entries(THEMES)
          .map(
            ([theme, prefix]) => `
${prefix} [data-chart=${id}] {
${colorConfig
  .map(([key, itemConfig]) => {
    const color =
      itemConfig.theme?.[theme as keyof typeof itemConfig.theme] ||
      itemConfig.color
    return color ? `  -color-${key}: ${color};` : null
  })
  .join('\n')}
}
`,
          )
          .join('\n'),
      }}
    />
  )
}

const ChartTooltip = RechartsPrimitive.Tooltip

function ChartTooltipContent({
  active,
  payload,
  className,
  indicator = 'dot',
  hideLabel = false,
  hideIndicator = false,
  label,
  labelFormatter,
  labelClassName,
  formatter,
  color,
  nameKey,
  labelKey,
}: React.ComponentProps<typeof RechartsPrimitive.Tooltip> &
  React.ComponentProps<'div'> & {
    hideLabel?: boolean
    hideIndicator?: boolean
    indicator?: 'line' | 'dot' | 'dashed'
    nameKey?: string
    labelKey?: string
  }) {
  const { config } = useChart()

  const tooltipLabel = React.useMemo(() => {
    if (hideLabel || !payload?.length) {
      return null
    }

    const [item] = payload
    const key = `${labelKey || item?.dataKey || item?.name || 'value'}`
    const itemConfig = getPayloadConfigFromPayload(config, item, key)
    const value =
      !labelKey && typeof label === 'string'
        ? config[label as keyof typeof config]?.label || label
        : itemConfig?.label

    if (labelFormatter) {
      return (
        <div className={cn('font-medium', labelClassName)}>
          {labelFormatter(value, payload)}
        </div>
      )
    }

    if (!value) {
      return null
    }

    return <div className={cn('font-medium', labelClassName)}>{value}</div>
  }, [
    label,
    labelFormatter,
    payload,
    hideLabel,
    labelClassName,
    config,
    labelKey,
  ])

  if (!active || !payload?.length) {
    return null
  }

  const nestLabel = payload.length === 1 && indicator !== 'dot'

  return (
    <div
      className={cn(
        'border-border/50 bg-background grid min-w-[8rem] items-start gap-1.5 rounded-lg border px-2.5 py-1.5 text-xs shadow-xl',
        className,
      )}
    >
      {!nestLabel ? tooltipLabel : null}
      <div className="grid gap-1.5">
        {payload.map((item, index) => {
          const key = `${nameKey || item.name || item.dataKey || 'value'}`
          const itemConfig = getPayloadConfigFromPayload(config, item, key)
          const indicatorColor = color || item.payload.fill || item.color

          return (
            <div
              key={item.dataKey}
              className={cn(
                '[&>svg]:text-muted-foreground flex w-full flex-wrap items-stretch gap-2 [&>svg]:h-2.5 [&>svg]:w-2.5',
                indicator === 'dot' && 'items-center',
              )}
            >
              {formatter && item?.value !== undefined && item.name ? (
                formatter(item.value, item.name, item, index, item.payload)
              ) : (
                <>
                  {itemConfig?.icon ? (
                    <itemConfig.icon />
                  ) : (
                    !hideIndicator && (
                      <div
                        className={cn(
                          'shrink-0 rounded-[2px] border-(-color-border) bg-(-color-bg)',
                          {
                            'h-2.5 w-2.5': indicator === 'dot',
                            'w-1': indicator === 'line',
                            'w-0 border-[1.5px] border-dashed bg-transparent':
                              indicator === 'dashed',
                            'my-0.5': nestLabel && indicator === 'dashed',
                          },
                        )}
                        style={
                          {
                            '-color-bg': indicatorColor,
                            '-color-border': indicatorColor,
                          } as React.CSSProperties
                        }
                      />
                    )
                  )}
                  <div
                    className={cn(
                      'flex flex-1 justify-between leading-none',
                      nestLabel ? 'items-end' : 'items-center',
                    )}
                  >
                    <div className="grid gap-1.5">
                      {nestLabel ? tooltipLabel : null}
                      <span className="text-muted-foreground">
                        {itemConfig?.label || item.name}
                      </span>
                    </div>
                    {item.value && (
                      <span className="text-foreground font-mono font-medium tabular-nums">
                        {item.value.toLocaleString()}
                      </span>
                    )}
                  </div>
                </>
              )}
            </div>
          )
        })}
      </div>
    </div>
  )
}

const ChartLegend = RechartsPrimitive.Legend

function ChartLegendContent({
  className,
  hideIcon = false,
  payload,
  verticalAlign = 'bottom',
  nameKey,
}: React.ComponentProps<'div'> &
  Pick<RechartsPrimitive.LegendProps, 'payload' | 'verticalAlign'> & {
    hideIcon?: boolean
    nameKey?: string
  }) {
  const { config } = useChart()

  if (!payload?.length) {
    return null
  }

  return (
    <div
      className={cn(
        'flex items-center justify-center gap-4',
        verticalAlign === 'top' ? 'pb-3' : 'pt-3',
        className,
      )}
    >
      {payload.map((item) => {
        const key = `${nameKey || item.dataKey || 'value'}`
        const itemConfig = getPayloadConfigFromPayload(config, item, key)

        return (
          <div
            key={item.value}
            className={
              '[&>svg]:text-muted-foreground flex items-center gap-1.5 [&>svg]:h-3 [&>svg]:w-3'
            }
          >
            {itemConfig?.icon && !hideIcon ? (
              <itemConfig.icon />
            ) : (
              <div
                className="h-2 w-2 shrink-0 rounded-[2px]"
                style={{
                  backgroundColor: item.color,
                }}
              />
            )}
            {itemConfig?.label}
          </div>
        )
      })}
    </div>
  )
}

// Helper to extract item config from a payload.
function getPayloadConfigFromPayload(
  config: ChartConfig,
  payload: unknown,
  key: string,
) {
  if (typeof payload !== 'object' || payload === null) {
    return undefined
  }

  const payloadPayload =
    'payload' in payload &&
    typeof payload.payload === 'object' &&
    payload.payload !== null
      ? payload.payload
      : undefined

  let configLabelKey: string = key

  if (
    key in payload &&
    typeof payload[key as keyof typeof payload] === 'string'
  ) {
    configLabelKey = payload[key as keyof typeof payload] as string
  } else if (
    payloadPayload &&
    key in payloadPayload &&
    typeof payloadPayload[key as keyof typeof payloadPayload] === 'string'
  ) {
    configLabelKey = payloadPayload[
      key as keyof typeof payloadPayload
    ] as string
  }

  return configLabelKey in config
    ? config[configLabelKey]
    : config[key as keyof typeof config]
}

export {
  ChartContainer,
  ChartTooltip,
  ChartTooltipContent,
  ChartLegend,
  ChartLegendContent,
  ChartStyle,
}
\end{DeskFlowCode}

\phantomsection
\addcontentsline{toc}{subsubsection}{deskflow\-frontend/components/ui/checkbox.tsx}
\DeskFlowListingCaption{deskflow-frontend/components/ui/checkbox.tsx}{lst:fd-ui-checkbox}
\begin{DeskFlowCode}
'use client'

import * as React from 'react'
import * as CheckboxPrimitive from '@radix-ui/react-checkbox'
import { CheckIcon } from 'lucide-react'

import { cn } from '@/lib/utils'

function Checkbox({
  className,
  ...props
}: React.ComponentProps<typeof CheckboxPrimitive.Root>) {
  return (
    <CheckboxPrimitive.Root
      data-slot="checkbox"
      className={cn(
        'peer border-input dark:bg-input/30 data-[state=checked]:bg-primary data-[state=checked]:text-primary-foreground dark:data-[state=checked]:bg-primary data-[state=checked]:border-primary focus-visible:border-ring focus-visible:ring-ring/50 aria-invalid:ring-destructive/20 dark:aria-invalid:ring-destructive/40 aria-invalid:border-destructive size-4 shrink-0 rounded-[4px] border shadow-xs transition-shadow outline-none focus-visible:ring-[3px] disabled:cursor-not-allowed disabled:opacity-50',
        className,
      )}
      {...props}
    >
      <CheckboxPrimitive.Indicator
        data-slot="checkbox-indicator"
        className="flex items-center justify-center text-current transition-none"
      >
        <CheckIcon className="size-3.5" />
      </CheckboxPrimitive.Indicator>
    </CheckboxPrimitive.Root>
  )
}

export { Checkbox }
\end{DeskFlowCode}

\phantomsection
\addcontentsline{toc}{subsubsection}{deskflow\-frontend/components/ui/collapsible.tsx}
\DeskFlowListingCaption{deskflow-frontend/components/ui/collapsible.tsx}{lst:fd-ui-collapsible}
\begin{DeskFlowCode}
'use client'

import * as CollapsiblePrimitive from '@radix-ui/react-collapsible'

function Collapsible({
  ...props
}: React.ComponentProps<typeof CollapsiblePrimitive.Root>) {
  return <CollapsiblePrimitive.Root data-slot="collapsible" {...props} />
}

function CollapsibleTrigger({
  ...props
}: React.ComponentProps<typeof CollapsiblePrimitive.CollapsibleTrigger>) {
  return (
    <CollapsiblePrimitive.CollapsibleTrigger
      data-slot="collapsible-trigger"
      {...props}
    />
  )
}

function CollapsibleContent({
  ...props
}: React.ComponentProps<typeof CollapsiblePrimitive.CollapsibleContent>) {
  return (
    <CollapsiblePrimitive.CollapsibleContent
      data-slot="collapsible-content"
      {...props}
    />
  )
}

export { Collapsible, CollapsibleTrigger, CollapsibleContent }
\end{DeskFlowCode}

\phantomsection
\addcontentsline{toc}{subsubsection}{deskflow\-frontend/components/ui/command.tsx}
\DeskFlowListingCaption{deskflow-frontend/components/ui/command.tsx}{lst:fd-ui-command}
\begin{DeskFlowCode}
'use client'

import * as React from 'react'
import { Command as CommandPrimitive } from 'cmdk'
import { SearchIcon } from 'lucide-react'

import { cn } from '@/lib/utils'
import {
  Dialog,
  DialogContent,
  DialogDescription,
  DialogHeader,
  DialogTitle,
} from '@/components/ui/dialog'

function Command({
  className,
  ...props
}: React.ComponentProps<typeof CommandPrimitive>) {
  return (
    <CommandPrimitive
      data-slot="command"
      className={cn(
        'bg-popover text-popover-foreground flex h-full w-full flex-col overflow-hidden rounded-md',
        className,
      )}
      {...props}
    />
  )
}

function CommandDialog({
  title = 'Command Palette',
  description = 'Search for a command to run...',
  children,
  className,
  showCloseButton = true,
  ...props
}: React.ComponentProps<typeof Dialog> & {
  title?: string
  description?: string
  className?: string
  showCloseButton?: boolean
}) {
  return (
    <Dialog {...props}>
      <DialogHeader className="sr-only">
        <DialogTitle>{title}</DialogTitle>
        <DialogDescription>{description}</DialogDescription>
      </DialogHeader>
      <DialogContent
        className={cn('overflow-hidden p-0', className)}
        showCloseButton={showCloseButton}
      >
        <Command className="[&_[cmdk-group-heading]]:text-muted-foreground **:data-[slot=command-input-wrapper]:h-12 [&_[cmdk-group-heading]]:px-2 [&_[cmdk-group-heading]]:font-medium [&_[cmdk-group]]:px-2 [&_[cmdk-group]:not([hidden])_~[cmdk-group]]:pt-0 [&_[cmdk-input-wrapper]_svg]:h-5 [&_[cmdk-input-wrapper]_svg]:w-5 [&_[cmdk-input]]:h-12 [&_[cmdk-item]]:px-2 [&_[cmdk-item]]:py-3 [&_[cmdk-item]_svg]:h-5 [&_[cmdk-item]_svg]:w-5">
          {children}
        </Command>
      </DialogContent>
    </Dialog>
  )
}

function CommandInput({
  className,
  ...props
}: React.ComponentProps<typeof CommandPrimitive.Input>) {
  return (
    <div
      data-slot="command-input-wrapper"
      className="flex h-9 items-center gap-2 border-b px-3"
    >
      <SearchIcon className="size-4 shrink-0 opacity-50" />
      <CommandPrimitive.Input
        data-slot="command-input"
        className={cn(
          'placeholder:text-muted-foreground flex h-10 w-full rounded-md bg-transparent py-3 text-sm outline-hidden disabled:cursor-not-allowed disabled:opacity-50',
          className,
        )}
        {...props}
      />
    </div>
  )
}

function CommandList({
  className,
  ...props
}: React.ComponentProps<typeof CommandPrimitive.List>) {
  return (
    <CommandPrimitive.List
      data-slot="command-list"
      className={cn(
        'max-h-[300px] scroll-py-1 overflow-x-hidden overflow-y-auto',
        className,
      )}
      {...props}
    />
  )
}

function CommandEmpty({
  ...props
}: React.ComponentProps<typeof CommandPrimitive.Empty>) {
  return (
    <CommandPrimitive.Empty
      data-slot="command-empty"
      className="py-6 text-center text-sm"
      {...props}
    />
  )
}

function CommandGroup({
  className,
  ...props
}: React.ComponentProps<typeof CommandPrimitive.Group>) {
  return (
    <CommandPrimitive.Group
      data-slot="command-group"
      className={cn(
        'text-foreground [&_[cmdk-group-heading]]:text-muted-foreground overflow-hidden p-1 [&_[cmdk-group-heading]]:px-2 [&_[cmdk-group-heading]]:py-1.5 [&_[cmdk-group-heading]]:text-xs [&_[cmdk-group-heading]]:font-medium',
        className,
      )}
      {...props}
    />
  )
}

function CommandSeparator({
  className,
  ...props
}: React.ComponentProps<typeof CommandPrimitive.Separator>) {
  return (
    <CommandPrimitive.Separator
      data-slot="command-separator"
      className={cn('bg-border -mx-1 h-px', className)}
      {...props}
    />
  )
}

function CommandItem({
  className,
  ...props
}: React.ComponentProps<typeof CommandPrimitive.Item>) {
  return (
    <CommandPrimitive.Item
      data-slot="command-item"
      className={cn(
        "data-[selected=true]:bg-accent data-[selected=true]:text-accent-foreground [&_svg:not([class*='text-'])]:text-muted-foreground relative flex cursor-default items-center gap-2 rounded-sm px-2 py-1.5 text-sm outline-hidden select-none data-[disabled=true]:pointer-events-none data-[disabled=true]:opacity-50 [&_svg]:pointer-events-none [&_svg]:shrink-0 [&_svg:not([class*='size-'])]:size-4",
        className,
      )}
      {...props}
    />
  )
}

function CommandShortcut({
  className,
  ...props
}: React.ComponentProps<'span'>) {
  return (
    <span
      data-slot="command-shortcut"
      className={cn(
        'text-muted-foreground ml-auto text-xs tracking-widest',
        className,
      )}
      {...props}
    />
  )
}

export {
  Command,
  CommandDialog,
  CommandInput,
  CommandList,
  CommandEmpty,
  CommandGroup,
  CommandItem,
  CommandShortcut,
  CommandSeparator,
}
\end{DeskFlowCode}

\phantomsection
\addcontentsline{toc}{subsubsection}{deskflow\-frontend/components/ui/context\-menu.tsx}
\DeskFlowListingCaption{deskflow-frontend/components/ui/context-menu.tsx}{lst:fd-ui-context-menu}
\begin{DeskFlowCode}
'use client'

import * as React from 'react'
import * as ContextMenuPrimitive from '@radix-ui/react-context-menu'
import { CheckIcon, ChevronRightIcon, CircleIcon } from 'lucide-react'

import { cn } from '@/lib/utils'

function ContextMenu({
  ...props
}: React.ComponentProps<typeof ContextMenuPrimitive.Root>) {
  return <ContextMenuPrimitive.Root data-slot="context-menu" {...props} />
}

function ContextMenuTrigger({
  ...props
}: React.ComponentProps<typeof ContextMenuPrimitive.Trigger>) {
  return (
    <ContextMenuPrimitive.Trigger data-slot="context-menu-trigger" {...props} />
  )
}

function ContextMenuGroup({
  ...props
}: React.ComponentProps<typeof ContextMenuPrimitive.Group>) {
  return (
    <ContextMenuPrimitive.Group data-slot="context-menu-group" {...props} />
  )
}

function ContextMenuPortal({
  ...props
}: React.ComponentProps<typeof ContextMenuPrimitive.Portal>) {
  return (
    <ContextMenuPrimitive.Portal data-slot="context-menu-portal" {...props} />
  )
}

function ContextMenuSub({
  ...props
}: React.ComponentProps<typeof ContextMenuPrimitive.Sub>) {
  return <ContextMenuPrimitive.Sub data-slot="context-menu-sub" {...props} />
}

function ContextMenuRadioGroup({
  ...props
}: React.ComponentProps<typeof ContextMenuPrimitive.RadioGroup>) {
  return (
    <ContextMenuPrimitive.RadioGroup
      data-slot="context-menu-radio-group"
      {...props}
    />
  )
}

function ContextMenuSubTrigger({
  className,
  inset,
  children,
  ...props
}: React.ComponentProps<typeof ContextMenuPrimitive.SubTrigger> & {
  inset?: boolean
}) {
  return (
    <ContextMenuPrimitive.SubTrigger
      data-slot="context-menu-sub-trigger"
      data-inset={inset}
      className={cn(
        "focus:bg-accent focus:text-accent-foreground data-[state=open]:bg-accent data-[state=open]:text-accent-foreground [&_svg:not([class*='text-'])]:text-muted-foreground flex cursor-default items-center rounded-sm px-2 py-1.5 text-sm outline-hidden select-none data-[inset]:pl-8 [&_svg]:pointer-events-none [&_svg]:shrink-0 [&_svg:not([class*='size-'])]:size-4",
        className,
      )}
      {...props}
    >
      {children}
      <ChevronRightIcon className="ml-auto" />
    </ContextMenuPrimitive.SubTrigger>
  )
}

function ContextMenuSubContent({
  className,
  ...props
}: React.ComponentProps<typeof ContextMenuPrimitive.SubContent>) {
  return (
    <ContextMenuPrimitive.SubContent
      data-slot="context-menu-sub-content"
      className={cn(
        'bg-popover text-popover-foreground data-[state=open]:animate-in data-[state=closed]:animate-out data-[state=closed]:fade-out-0 data-[state=open]:fade-in-0 data-[state=closed]:zoom-out-95 data-[state=open]:zoom-in-95 data-[side=bottom]:slide-in-from-top-2 data-[side=left]:slide-in-from-right-2 data-[side=right]:slide-in-from-left-2 data-[side=top]:slide-in-from-bottom-2 z-50 min-w-[8rem] origin-(-radix-context-menu-content-transform-origin) overflow-hidden rounded-md border p-1 shadow-lg',
        className,
      )}
      {...props}
    />
  )
}

function ContextMenuContent({
  className,
  ...props
}: React.ComponentProps<typeof ContextMenuPrimitive.Content>) {
  return (
    <ContextMenuPrimitive.Portal>
      <ContextMenuPrimitive.Content
        data-slot="context-menu-content"
        className={cn(
          'bg-popover text-popover-foreground data-[state=open]:animate-in data-[state=closed]:animate-out data-[state=closed]:fade-out-0 data-[state=open]:fade-in-0 data-[state=closed]:zoom-out-95 data-[state=open]:zoom-in-95 data-[side=bottom]:slide-in-from-top-2 data-[side=left]:slide-in-from-right-2 data-[side=right]:slide-in-from-left-2 data-[side=top]:slide-in-from-bottom-2 z-50 max-h-(-radix-context-menu-content-available-height) min-w-[8rem] origin-(-radix-context-menu-content-transform-origin) overflow-x-hidden overflow-y-auto rounded-md border p-1 shadow-md',
          className,
        )}
        {...props}
      />
    </ContextMenuPrimitive.Portal>
  )
}

function ContextMenuItem({
  className,
  inset,
  variant = 'default',
  ...props
}: React.ComponentProps<typeof ContextMenuPrimitive.Item> & {
  inset?: boolean
  variant?: 'default' | 'destructive'
}) {
  return (
    <ContextMenuPrimitive.Item
      data-slot="context-menu-item"
      data-inset={inset}
      data-variant={variant}
      className={cn(
        "focus:bg-accent focus:text-accent-foreground data-[variant=destructive]:text-destructive data-[variant=destructive]:focus:bg-destructive/10 dark:data-[variant=destructive]:focus:bg-destructive/20 data-[variant=destructive]:focus:text-destructive data-[variant=destructive]:*:[svg]:!text-destructive [&_svg:not([class*='text-'])]:text-muted-foreground relative flex cursor-default items-center gap-2 rounded-sm px-2 py-1.5 text-sm outline-hidden select-none data-[disabled]:pointer-events-none data-[disabled]:opacity-50 data-[inset]:pl-8 [&_svg]:pointer-events-none [&_svg]:shrink-0 [&_svg:not([class*='size-'])]:size-4",
        className,
      )}
      {...props}
    />
  )
}

function ContextMenuCheckboxItem({
  className,
  children,
  checked,
  ...props
}: React.ComponentProps<typeof ContextMenuPrimitive.CheckboxItem>) {
  return (
    <ContextMenuPrimitive.CheckboxItem
      data-slot="context-menu-checkbox-item"
      className={cn(
        "focus:bg-accent focus:text-accent-foreground relative flex cursor-default items-center gap-2 rounded-sm py-1.5 pr-2 pl-8 text-sm outline-hidden select-none data-[disabled]:pointer-events-none data-[disabled]:opacity-50 [&_svg]:pointer-events-none [&_svg]:shrink-0 [&_svg:not([class*='size-'])]:size-4",
        className,
      )}
      checked={checked}
      {...props}
    >
      <span className="pointer-events-none absolute left-2 flex size-3.5 items-center justify-center">
        <ContextMenuPrimitive.ItemIndicator>
          <CheckIcon className="size-4" />
        </ContextMenuPrimitive.ItemIndicator>
      </span>
      {children}
    </ContextMenuPrimitive.CheckboxItem>
  )
}

function ContextMenuRadioItem({
  className,
  children,
  ...props
}: React.ComponentProps<typeof ContextMenuPrimitive.RadioItem>) {
  return (
    <ContextMenuPrimitive.RadioItem
      data-slot="context-menu-radio-item"
      className={cn(
        "focus:bg-accent focus:text-accent-foreground relative flex cursor-default items-center gap-2 rounded-sm py-1.5 pr-2 pl-8 text-sm outline-hidden select-none data-[disabled]:pointer-events-none data-[disabled]:opacity-50 [&_svg]:pointer-events-none [&_svg]:shrink-0 [&_svg:not([class*='size-'])]:size-4",
        className,
      )}
      {...props}
    >
      <span className="pointer-events-none absolute left-2 flex size-3.5 items-center justify-center">
        <ContextMenuPrimitive.ItemIndicator>
          <CircleIcon className="size-2 fill-current" />
        </ContextMenuPrimitive.ItemIndicator>
      </span>
      {children}
    </ContextMenuPrimitive.RadioItem>
  )
}

function ContextMenuLabel({
  className,
  inset,
  ...props
}: React.ComponentProps<typeof ContextMenuPrimitive.Label> & {
  inset?: boolean
}) {
  return (
    <ContextMenuPrimitive.Label
      data-slot="context-menu-label"
      data-inset={inset}
      className={cn(
        'text-foreground px-2 py-1.5 text-sm font-medium data-[inset]:pl-8',
        className,
      )}
      {...props}
    />
  )
}

function ContextMenuSeparator({
  className,
  ...props
}: React.ComponentProps<typeof ContextMenuPrimitive.Separator>) {
  return (
    <ContextMenuPrimitive.Separator
      data-slot="context-menu-separator"
      className={cn('bg-border -mx-1 my-1 h-px', className)}
      {...props}
    />
  )
}

function ContextMenuShortcut({
  className,
  ...props
}: React.ComponentProps<'span'>) {
  return (
    <span
      data-slot="context-menu-shortcut"
      className={cn(
        'text-muted-foreground ml-auto text-xs tracking-widest',
        className,
      )}
      {...props}
    />
  )
}

export {
  ContextMenu,
  ContextMenuTrigger,
  ContextMenuContent,
  ContextMenuItem,
  ContextMenuCheckboxItem,
  ContextMenuRadioItem,
  ContextMenuLabel,
  ContextMenuSeparator,
  ContextMenuShortcut,
  ContextMenuGroup,
  ContextMenuPortal,
  ContextMenuSub,
  ContextMenuSubContent,
  ContextMenuSubTrigger,
  ContextMenuRadioGroup,
}
\end{DeskFlowCode}

\phantomsection
\addcontentsline{toc}{subsubsection}{deskflow\-frontend/components/ui/dialog.tsx}
\DeskFlowListingCaption{deskflow-frontend/components/ui/dialog.tsx}{lst:fd-ui-dialog}
\begin{DeskFlowCode}
'use client'

import * as React from 'react'
import * as DialogPrimitive from '@radix-ui/react-dialog'
import { XIcon } from 'lucide-react'

import { cn } from '@/lib/utils'

function Dialog({
  ...props
}: React.ComponentProps<typeof DialogPrimitive.Root>) {
  return <DialogPrimitive.Root data-slot="dialog" {...props} />
}

function DialogTrigger({
  ...props
}: React.ComponentProps<typeof DialogPrimitive.Trigger>) {
  return <DialogPrimitive.Trigger data-slot="dialog-trigger" {...props} />
}

function DialogPortal({
  ...props
}: React.ComponentProps<typeof DialogPrimitive.Portal>) {
  return <DialogPrimitive.Portal data-slot="dialog-portal" {...props} />
}

function DialogClose({
  ...props
}: React.ComponentProps<typeof DialogPrimitive.Close>) {
  return <DialogPrimitive.Close data-slot="dialog-close" {...props} />
}

function DialogOverlay({
  className,
  ...props
}: React.ComponentProps<typeof DialogPrimitive.Overlay>) {
  return (
    <DialogPrimitive.Overlay
      data-slot="dialog-overlay"
      className={cn(
        'data-[state=open]:animate-in data-[state=closed]:animate-out data-[state=closed]:fade-out-0 data-[state=open]:fade-in-0 fixed inset-0 z-50 bg-black/50',
        className,
      )}
      {...props}
    />
  )
}

function DialogContent({
  className,
  children,
  showCloseButton = true,
  ...props
}: React.ComponentProps<typeof DialogPrimitive.Content> & {
  showCloseButton?: boolean
}) {
  return (
    <DialogPortal data-slot="dialog-portal">
      <DialogOverlay />
      <DialogPrimitive.Content
        data-slot="dialog-content"
        className={cn(
          'bg-background data-[state=open]:animate-in data-[state=closed]:animate-out data-[state=closed]:fade-out-0 data-[state=open]:fade-in-0 data-[state=closed]:zoom-out-95 data-[state=open]:zoom-in-95 fixed top-[50%] left-[50%] z-50 grid w-full max-w-[calc(100%-2rem)] translate-x-[-50%] translate-y-[-50%] gap-4 rounded-lg border p-6 shadow-lg duration-200 sm:max-w-lg',
          className,
        )}
        {...props}
      >
        {children}
        {showCloseButton && (
          <DialogPrimitive.Close
            data-slot="dialog-close"
            className="ring-offset-background focus:ring-ring data-[state=open]:bg-accent data-[state=open]:text-muted-foreground absolute top-4 right-4 rounded-xs opacity-70 transition-opacity hover:opacity-100 focus:ring-2 focus:ring-offset-2 focus:outline-hidden disabled:pointer-events-none [&_svg]:pointer-events-none [&_svg]:shrink-0 [&_svg:not([class*='size-'])]:size-4"
          >
            <XIcon />
            <span className="sr-only">Close</span>
          </DialogPrimitive.Close>
        )}
      </DialogPrimitive.Content>
    </DialogPortal>
  )
}

function DialogHeader({ className, ...props }: React.ComponentProps<'div'>) {
  return (
    <div
      data-slot="dialog-header"
      className={cn('flex flex-col gap-2 text-center sm:text-left', className)}
      {...props}
    />
  )
}

function DialogFooter({ className, ...props }: React.ComponentProps<'div'>) {
  return (
    <div
      data-slot="dialog-footer"
      className={cn(
        'flex flex-col-reverse gap-2 sm:flex-row sm:justify-end',
        className,
      )}
      {...props}
    />
  )
}

function DialogTitle({
  className,
  ...props
}: React.ComponentProps<typeof DialogPrimitive.Title>) {
  return (
    <DialogPrimitive.Title
      data-slot="dialog-title"
      className={cn('text-lg leading-none font-semibold', className)}
      {...props}
    />
  )
}

function DialogDescription({
  className,
  ...props
}: React.ComponentProps<typeof DialogPrimitive.Description>) {
  return (
    <DialogPrimitive.Description
      data-slot="dialog-description"
      className={cn('text-muted-foreground text-sm', className)}
      {...props}
    />
  )
}

export {
  Dialog,
  DialogClose,
  DialogContent,
  DialogDescription,
  DialogFooter,
  DialogHeader,
  DialogOverlay,
  DialogPortal,
  DialogTitle,
  DialogTrigger,
}
\end{DeskFlowCode}

\phantomsection
\addcontentsline{toc}{subsubsection}{deskflow\-frontend/components/ui/drawer.tsx}
\DeskFlowListingCaption{deskflow-frontend/components/ui/drawer.tsx}{lst:fd-ui-drawer}
\begin{DeskFlowCode}
'use client'

import * as React from 'react'
import { Drawer as DrawerPrimitive } from 'vaul'

import { cn } from '@/lib/utils'

function Drawer({
  ...props
}: React.ComponentProps<typeof DrawerPrimitive.Root>) {
  return <DrawerPrimitive.Root data-slot="drawer" {...props} />
}

function DrawerTrigger({
  ...props
}: React.ComponentProps<typeof DrawerPrimitive.Trigger>) {
  return <DrawerPrimitive.Trigger data-slot="drawer-trigger" {...props} />
}

function DrawerPortal({
  ...props
}: React.ComponentProps<typeof DrawerPrimitive.Portal>) {
  return <DrawerPrimitive.Portal data-slot="drawer-portal" {...props} />
}

function DrawerClose({
  ...props
}: React.ComponentProps<typeof DrawerPrimitive.Close>) {
  return <DrawerPrimitive.Close data-slot="drawer-close" {...props} />
}

function DrawerOverlay({
  className,
  ...props
}: React.ComponentProps<typeof DrawerPrimitive.Overlay>) {
  return (
    <DrawerPrimitive.Overlay
      data-slot="drawer-overlay"
      className={cn(
        'data-[state=open]:animate-in data-[state=closed]:animate-out data-[state=closed]:fade-out-0 data-[state=open]:fade-in-0 fixed inset-0 z-50 bg-black/50',
        className,
      )}
      {...props}
    />
  )
}

function DrawerContent({
  className,
  children,
  ...props
}: React.ComponentProps<typeof DrawerPrimitive.Content>) {
  return (
    <DrawerPortal data-slot="drawer-portal">
      <DrawerOverlay />
      <DrawerPrimitive.Content
        data-slot="drawer-content"
        className={cn(
          'group/drawer-content bg-background fixed z-50 flex h-auto flex-col',
          'data-[vaul-drawer-direction=top]:inset-x-0 data-[vaul-drawer-direction=top]:top-0 data-[vaul-drawer-direction=top]:mb-24 data-[vaul-drawer-direction=top]:max-h-[80vh] data-[vaul-drawer-direction=top]:rounded-b-lg data-[vaul-drawer-direction=top]:border-b',
          'data-[vaul-drawer-direction=bottom]:inset-x-0 data-[vaul-drawer-direction=bottom]:bottom-0 data-[vaul-drawer-direction=bottom]:mt-24 data-[vaul-drawer-direction=bottom]:max-h-[80vh] data-[vaul-drawer-direction=bottom]:rounded-t-lg data-[vaul-drawer-direction=bottom]:border-t',
          'data-[vaul-drawer-direction=right]:inset-y-0 data-[vaul-drawer-direction=right]:right-0 data-[vaul-drawer-direction=right]:w-3/4 data-[vaul-drawer-direction=right]:border-l data-[vaul-drawer-direction=right]:sm:max-w-sm',
          'data-[vaul-drawer-direction=left]:inset-y-0 data-[vaul-drawer-direction=left]:left-0 data-[vaul-drawer-direction=left]:w-3/4 data-[vaul-drawer-direction=left]:border-r data-[vaul-drawer-direction=left]:sm:max-w-sm',
          className,
        )}
        {...props}
      >
        <div className="bg-muted mx-auto mt-4 hidden h-2 w-[100px] shrink-0 rounded-full group-data-[vaul-drawer-direction=bottom]/drawer-content:block" />
        {children}
      </DrawerPrimitive.Content>
    </DrawerPortal>
  )
}

function DrawerHeader({ className, ...props }: React.ComponentProps<'div'>) {
  return (
    <div
      data-slot="drawer-header"
      className={cn(
        'flex flex-col gap-0.5 p-4 group-data-[vaul-drawer-direction=bottom]/drawer-content:text-center group-data-[vaul-drawer-direction=top]/drawer-content:text-center md:gap-1.5 md:text-left',
        className,
      )}
      {...props}
    />
  )
}

function DrawerFooter({ className, ...props }: React.ComponentProps<'div'>) {
  return (
    <div
      data-slot="drawer-footer"
      className={cn('mt-auto flex flex-col gap-2 p-4', className)}
      {...props}
    />
  )
}

function DrawerTitle({
  className,
  ...props
}: React.ComponentProps<typeof DrawerPrimitive.Title>) {
  return (
    <DrawerPrimitive.Title
      data-slot="drawer-title"
      className={cn('text-foreground font-semibold', className)}
      {...props}
    />
  )
}

function DrawerDescription({
  className,
  ...props
}: React.ComponentProps<typeof DrawerPrimitive.Description>) {
  return (
    <DrawerPrimitive.Description
      data-slot="drawer-description"
      className={cn('text-muted-foreground text-sm', className)}
      {...props}
    />
  )
}

export {
  Drawer,
  DrawerPortal,
  DrawerOverlay,
  DrawerTrigger,
  DrawerClose,
  DrawerContent,
  DrawerHeader,
  DrawerFooter,
  DrawerTitle,
  DrawerDescription,
}
\end{DeskFlowCode}

\phantomsection
\addcontentsline{toc}{subsubsection}{deskflow\-frontend/components/ui/dropdown\-menu.tsx}
\DeskFlowListingCaption{deskflow-frontend/components/ui/dropdown-menu.tsx}{lst:fd-ui-dropdown-menu}
\begin{DeskFlowCode}
'use client'

import * as React from 'react'
import * as DropdownMenuPrimitive from '@radix-ui/react-dropdown-menu'
import { CheckIcon, ChevronRightIcon, CircleIcon } from 'lucide-react'

import { cn } from '@/lib/utils'

function DropdownMenu({
  ...props
}: React.ComponentProps<typeof DropdownMenuPrimitive.Root>) {
  return <DropdownMenuPrimitive.Root data-slot="dropdown-menu" {...props} />
}

function DropdownMenuPortal({
  ...props
}: React.ComponentProps<typeof DropdownMenuPrimitive.Portal>) {
  return (
    <DropdownMenuPrimitive.Portal data-slot="dropdown-menu-portal" {...props} />
  )
}

function DropdownMenuTrigger({
  ...props
}: React.ComponentProps<typeof DropdownMenuPrimitive.Trigger>) {
  return (
    <DropdownMenuPrimitive.Trigger
      data-slot="dropdown-menu-trigger"
      {...props}
    />
  )
}

function DropdownMenuContent({
  className,
  sideOffset = 4,
  ...props
}: React.ComponentProps<typeof DropdownMenuPrimitive.Content>) {
  return (
    <DropdownMenuPrimitive.Portal>
      <DropdownMenuPrimitive.Content
        data-slot="dropdown-menu-content"
        sideOffset={sideOffset}
        className={cn(
          'bg-popover text-popover-foreground data-[state=open]:animate-in data-[state=closed]:animate-out data-[state=closed]:fade-out-0 data-[state=open]:fade-in-0 data-[state=closed]:zoom-out-95 data-[state=open]:zoom-in-95 data-[side=bottom]:slide-in-from-top-2 data-[side=left]:slide-in-from-right-2 data-[side=right]:slide-in-from-left-2 data-[side=top]:slide-in-from-bottom-2 z-50 max-h-(-radix-dropdown-menu-content-available-height) min-w-[8rem] origin-(-radix-dropdown-menu-content-transform-origin) overflow-x-hidden overflow-y-auto rounded-md border p-1 shadow-md',
          className,
        )}
        {...props}
      />
    </DropdownMenuPrimitive.Portal>
  )
}

function DropdownMenuGroup({
  ...props
}: React.ComponentProps<typeof DropdownMenuPrimitive.Group>) {
  return (
    <DropdownMenuPrimitive.Group data-slot="dropdown-menu-group" {...props} />
  )
}

function DropdownMenuItem({
  className,
  inset,
  variant = 'default',
  ...props
}: React.ComponentProps<typeof DropdownMenuPrimitive.Item> & {
  inset?: boolean
  variant?: 'default' | 'destructive'
}) {
  return (
    <DropdownMenuPrimitive.Item
      data-slot="dropdown-menu-item"
      data-inset={inset}
      data-variant={variant}
      className={cn(
        "focus:bg-accent focus:text-accent-foreground data-[variant=destructive]:text-destructive data-[variant=destructive]:focus:bg-destructive/10 dark:data-[variant=destructive]:focus:bg-destructive/20 data-[variant=destructive]:focus:text-destructive data-[variant=destructive]:*:[svg]:!text-destructive [&_svg:not([class*='text-'])]:text-muted-foreground relative flex cursor-default items-center gap-2 rounded-sm px-2 py-1.5 text-sm outline-hidden select-none data-[disabled]:pointer-events-none data-[disabled]:opacity-50 data-[inset]:pl-8 [&_svg]:pointer-events-none [&_svg]:shrink-0 [&_svg:not([class*='size-'])]:size-4",
        className,
      )}
      {...props}
    />
  )
}

function DropdownMenuCheckboxItem({
  className,
  children,
  checked,
  ...props
}: React.ComponentProps<typeof DropdownMenuPrimitive.CheckboxItem>) {
  return (
    <DropdownMenuPrimitive.CheckboxItem
      data-slot="dropdown-menu-checkbox-item"
      className={cn(
        "focus:bg-accent focus:text-accent-foreground relative flex cursor-default items-center gap-2 rounded-sm py-1.5 pr-2 pl-8 text-sm outline-hidden select-none data-[disabled]:pointer-events-none data-[disabled]:opacity-50 [&_svg]:pointer-events-none [&_svg]:shrink-0 [&_svg:not([class*='size-'])]:size-4",
        className,
      )}
      checked={checked}
      {...props}
    >
      <span className="pointer-events-none absolute left-2 flex size-3.5 items-center justify-center">
        <DropdownMenuPrimitive.ItemIndicator>
          <CheckIcon className="size-4" />
        </DropdownMenuPrimitive.ItemIndicator>
      </span>
      {children}
    </DropdownMenuPrimitive.CheckboxItem>
  )
}

function DropdownMenuRadioGroup({
  ...props
}: React.ComponentProps<typeof DropdownMenuPrimitive.RadioGroup>) {
  return (
    <DropdownMenuPrimitive.RadioGroup
      data-slot="dropdown-menu-radio-group"
      {...props}
    />
  )
}

function DropdownMenuRadioItem({
  className,
  children,
  ...props
}: React.ComponentProps<typeof DropdownMenuPrimitive.RadioItem>) {
  return (
    <DropdownMenuPrimitive.RadioItem
      data-slot="dropdown-menu-radio-item"
      className={cn(
        "focus:bg-accent focus:text-accent-foreground relative flex cursor-default items-center gap-2 rounded-sm py-1.5 pr-2 pl-8 text-sm outline-hidden select-none data-[disabled]:pointer-events-none data-[disabled]:opacity-50 [&_svg]:pointer-events-none [&_svg]:shrink-0 [&_svg:not([class*='size-'])]:size-4",
        className,
      )}
      {...props}
    >
      <span className="pointer-events-none absolute left-2 flex size-3.5 items-center justify-center">
        <DropdownMenuPrimitive.ItemIndicator>
          <CircleIcon className="size-2 fill-current" />
        </DropdownMenuPrimitive.ItemIndicator>
      </span>
      {children}
    </DropdownMenuPrimitive.RadioItem>
  )
}

function DropdownMenuLabel({
  className,
  inset,
  ...props
}: React.ComponentProps<typeof DropdownMenuPrimitive.Label> & {
  inset?: boolean
}) {
  return (
    <DropdownMenuPrimitive.Label
      data-slot="dropdown-menu-label"
      data-inset={inset}
      className={cn(
        'px-2 py-1.5 text-sm font-medium data-[inset]:pl-8',
        className,
      )}
      {...props}
    />
  )
}

function DropdownMenuSeparator({
  className,
  ...props
}: React.ComponentProps<typeof DropdownMenuPrimitive.Separator>) {
  return (
    <DropdownMenuPrimitive.Separator
      data-slot="dropdown-menu-separator"
      className={cn('bg-border -mx-1 my-1 h-px', className)}
      {...props}
    />
  )
}

function DropdownMenuShortcut({
  className,
  ...props
}: React.ComponentProps<'span'>) {
  return (
    <span
      data-slot="dropdown-menu-shortcut"
      className={cn(
        'text-muted-foreground ml-auto text-xs tracking-widest',
        className,
      )}
      {...props}
    />
  )
}

function DropdownMenuSub({
  ...props
}: React.ComponentProps<typeof DropdownMenuPrimitive.Sub>) {
  return <DropdownMenuPrimitive.Sub data-slot="dropdown-menu-sub" {...props} />
}

function DropdownMenuSubTrigger({
  className,
  inset,
  children,
  ...props
}: React.ComponentProps<typeof DropdownMenuPrimitive.SubTrigger> & {
  inset?: boolean
}) {
  return (
    <DropdownMenuPrimitive.SubTrigger
      data-slot="dropdown-menu-sub-trigger"
      data-inset={inset}
      className={cn(
        "focus:bg-accent focus:text-accent-foreground data-[state=open]:bg-accent data-[state=open]:text-accent-foreground [&_svg:not([class*='text-'])]:text-muted-foreground flex cursor-default items-center gap-2 rounded-sm px-2 py-1.5 text-sm outline-hidden select-none data-[inset]:pl-8 [&_svg]:pointer-events-none [&_svg]:shrink-0 [&_svg:not([class*='size-'])]:size-4",
        className,
      )}
      {...props}
    >
      {children}
      <ChevronRightIcon className="ml-auto size-4" />
    </DropdownMenuPrimitive.SubTrigger>
  )
}

function DropdownMenuSubContent({
  className,
  ...props
}: React.ComponentProps<typeof DropdownMenuPrimitive.SubContent>) {
  return (
    <DropdownMenuPrimitive.SubContent
      data-slot="dropdown-menu-sub-content"
      className={cn(
        'bg-popover text-popover-foreground data-[state=open]:animate-in data-[state=closed]:animate-out data-[state=closed]:fade-out-0 data-[state=open]:fade-in-0 data-[state=closed]:zoom-out-95 data-[state=open]:zoom-in-95 data-[side=bottom]:slide-in-from-top-2 data-[side=left]:slide-in-from-right-2 data-[side=right]:slide-in-from-left-2 data-[side=top]:slide-in-from-bottom-2 z-50 min-w-[8rem] origin-(-radix-dropdown-menu-content-transform-origin) overflow-hidden rounded-md border p-1 shadow-lg',
        className,
      )}
      {...props}
    />
  )
}

export {
  DropdownMenu,
  DropdownMenuPortal,
  DropdownMenuTrigger,
  DropdownMenuContent,
  DropdownMenuGroup,
  DropdownMenuLabel,
  DropdownMenuItem,
  DropdownMenuCheckboxItem,
  DropdownMenuRadioGroup,
  DropdownMenuRadioItem,
  DropdownMenuSeparator,
  DropdownMenuShortcut,
  DropdownMenuSub,
  DropdownMenuSubTrigger,
  DropdownMenuSubContent,
}
\end{DeskFlowCode}

\phantomsection
\addcontentsline{toc}{subsubsection}{deskflow\-frontend/components/ui/empty.tsx}
\DeskFlowListingCaption{deskflow-frontend/components/ui/empty.tsx}{lst:fd-ui-empty}
\begin{DeskFlowCode}
import { cva, type VariantProps } from 'class-variance-authority'

import { cn } from '@/lib/utils'

function Empty({ className, ...props }: React.ComponentProps<'div'>) {
  return (
    <div
      data-slot="empty"
      className={cn(
        'flex min-w-0 flex-1 flex-col items-center justify-center gap-6 rounded-lg border-dashed p-6 text-center text-balance md:p-12',
        className,
      )}
      {...props}
    />
  )
}

function EmptyHeader({ className, ...props }: React.ComponentProps<'div'>) {
  return (
    <div
      data-slot="empty-header"
      className={cn(
        'flex max-w-sm flex-col items-center gap-2 text-center',
        className,
      )}
      {...props}
    />
  )
}

const emptyMediaVariants = cva(
  'flex shrink-0 items-center justify-center mb-2 [&_svg]:pointer-events-none [&_svg]:shrink-0',
  {
    variants: {
      variant: {
        default: 'bg-transparent',
        icon: "bg-muted text-foreground flex size-10 shrink-0 items-center justify-center rounded-lg [&_svg:not([class*='size-'])]:size-6",
      },
    },
    defaultVariants: {
      variant: 'default',
    },
  },
)

function EmptyMedia({
  className,
  variant = 'default',
  ...props
}: React.ComponentProps<'div'> & VariantProps<typeof emptyMediaVariants>) {
  return (
    <div
      data-slot="empty-icon"
      data-variant={variant}
      className={cn(emptyMediaVariants({ variant, className }))}
      {...props}
    />
  )
}

function EmptyTitle({ className, ...props }: React.ComponentProps<'div'>) {
  return (
    <div
      data-slot="empty-title"
      className={cn('text-lg font-medium tracking-tight', className)}
      {...props}
    />
  )
}

function EmptyDescription({ className, ...props }: React.ComponentProps<'p'>) {
  return (
    <div
      data-slot="empty-description"
      className={cn(
        'text-muted-foreground [&>a:hover]:text-primary text-sm/relaxed [&>a]:underline [&>a]:underline-offset-4',
        className,
      )}
      {...props}
    />
  )
}

function EmptyContent({ className, ...props }: React.ComponentProps<'div'>) {
  return (
    <div
      data-slot="empty-content"
      className={cn(
        'flex w-full max-w-sm min-w-0 flex-col items-center gap-4 text-sm text-balance',
        className,
      )}
      {...props}
    />
  )
}

export {
  Empty,
  EmptyHeader,
  EmptyTitle,
  EmptyDescription,
  EmptyContent,
  EmptyMedia,
}
\end{DeskFlowCode}

\phantomsection
\addcontentsline{toc}{subsubsection}{deskflow\-frontend/components/ui/field.tsx}
\DeskFlowListingCaption{deskflow-frontend/components/ui/field.tsx}{lst:fd-ui-field}
\begin{DeskFlowCode}
'use client'

import { useMemo } from 'react'
import { cva, type VariantProps } from 'class-variance-authority'

import { cn } from '@/lib/utils'
import { Label } from '@/components/ui/label'
import { Separator } from '@/components/ui/separator'

function FieldSet({ className, ...props }: React.ComponentProps<'fieldset'>) {
  return (
    <fieldset
      data-slot="field-set"
      className={cn(
        'flex flex-col gap-6',
        'has-[>[data-slot=checkbox-group]]:gap-3 has-[>[data-slot=radio-group]]:gap-3',
        className,
      )}
      {...props}
    />
  )
}

function FieldLegend({
  className,
  variant = 'legend',
  ...props
}: React.ComponentProps<'legend'> & { variant?: 'legend' | 'label' }) {
  return (
    <legend
      data-slot="field-legend"
      data-variant={variant}
      className={cn(
        'mb-3 font-medium',
        'data-[variant=legend]:text-base',
        'data-[variant=label]:text-sm',
        className,
      )}
      {...props}
    />
  )
}

function FieldGroup({ className, ...props }: React.ComponentProps<'div'>) {
  return (
    <div
      data-slot="field-group"
      className={cn(
        'group/field-group @container/field-group flex w-full flex-col gap-7 data-[slot=checkbox-group]:gap-3 [&>[data-slot=field-group]]:gap-4',
        className,
      )}
      {...props}
    />
  )
}

const fieldVariants = cva(
  'group/field flex w-full gap-3 data-[invalid=true]:text-destructive',
  {
    variants: {
      orientation: {
        vertical: ['flex-col [&>*]:w-full [&>.sr-only]:w-auto'],
        horizontal: [
          'flex-row items-center',
          '[&>[data-slot=field-label]]:flex-auto',
          'has-[>[data-slot=field-content]]:items-start has-[>[data-slot=field-content]]:[&>[role=checkbox],[role=radio]]:mt-px',
        ],
        responsive: [
          'flex-col [&>*]:w-full [&>.sr-only]:w-auto @md/field-group:flex-row @md/field-group:items-center @md/field-group:[&>*]:w-auto',
          '@md/field-group:[&>[data-slot=field-label]]:flex-auto',
          '@md/field-group:has-[>[data-slot=field-content]]:items-start @md/field-group:has-[>[data-slot=field-content]]:[&>[role=checkbox],[role=radio]]:mt-px',
        ],
      },
    },
    defaultVariants: {
      orientation: 'vertical',
    },
  },
)

function Field({
  className,
  orientation = 'vertical',
  ...props
}: React.ComponentProps<'div'> & VariantProps<typeof fieldVariants>) {
  return (
    <div
      role="group"
      data-slot="field"
      data-orientation={orientation}
      className={cn(fieldVariants({ orientation }), className)}
      {...props}
    />
  )
}

function FieldContent({ className, ...props }: React.ComponentProps<'div'>) {
  return (
    <div
      data-slot="field-content"
      className={cn(
        'group/field-content flex flex-1 flex-col gap-1.5 leading-snug',
        className,
      )}
      {...props}
    />
  )
}

function FieldLabel({
  className,
  ...props
}: React.ComponentProps<typeof Label>) {
  return (
    <Label
      data-slot="field-label"
      className={cn(
        'group/field-label peer/field-label flex w-fit gap-2 leading-snug group-data-[disabled=true]/field:opacity-50',
        'has-[>[data-slot=field]]:w-full has-[>[data-slot=field]]:flex-col has-[>[data-slot=field]]:rounded-md has-[>[data-slot=field]]:border [&>*]:data-[slot=field]:p-4',
        'has-data-[state=checked]:bg-primary/5 has-data-[state=checked]:border-primary dark:has-data-[state=checked]:bg-primary/10',
        className,
      )}
      {...props}
    />
  )
}

function FieldTitle({ className, ...props }: React.ComponentProps<'div'>) {
  return (
    <div
      data-slot="field-label"
      className={cn(
        'flex w-fit items-center gap-2 text-sm leading-snug font-medium group-data-[disabled=true]/field:opacity-50',
        className,
      )}
      {...props}
    />
  )
}

function FieldDescription({ className, ...props }: React.ComponentProps<'p'>) {
  return (
    <p
      data-slot="field-description"
      className={cn(
        'text-muted-foreground text-sm leading-normal font-normal group-has-[[data-orientation=horizontal]]/field:text-balance',
        'last:mt-0 nth-last-2:-mt-1 [[data-variant=legend]+&]:-mt-1.5',
        '[&>a:hover]:text-primary [&>a]:underline [&>a]:underline-offset-4',
        className,
      )}
      {...props}
    />
  )
}

function FieldSeparator({
  children,
  className,
  ...props
}: React.ComponentProps<'div'> & {
  children?: React.ReactNode
}) {
  return (
    <div
      data-slot="field-separator"
      data-content={!!children}
      className={cn(
        'relative -my-2 h-5 text-sm group-data-[variant=outline]/field-group:-mb-2',
        className,
      )}
      {...props}
    >
      <Separator className="absolute inset-0 top-1/2" />
      {children && (
        <span
          className="bg-background text-muted-foreground relative mx-auto block w-fit px-2"
          data-slot="field-separator-content"
        >
          {children}
        </span>
      )}
    </div>
  )
}

function FieldError({
  className,
  children,
  errors,
  ...props
}: React.ComponentProps<'div'> & {
  errors?: Array<{ message?: string } | undefined>
}) {
  const content = useMemo(() => {
    if (children) {
      return children
    }

    if (!errors) {
      return null
    }

    if (errors.length === 1 && errors[0]?.message) {
      return errors[0].message
    }

    return (
      <ul className="ml-4 flex list-disc flex-col gap-1">
        {errors.map(
          (error, index) =>
            error?.message && <li key={index}>{error.message}</li>,
        )}
      </ul>
    )
  }, [children, errors])

  if (!content) {
    return null
  }

  return (
    <div
      role="alert"
      data-slot="field-error"
      className={cn('text-destructive text-sm font-normal', className)}
      {...props}
    >
      {content}
    </div>
  )
}

export {
  Field,
  FieldLabel,
  FieldDescription,
  FieldError,
  FieldGroup,
  FieldLegend,
  FieldSeparator,
  FieldSet,
  FieldContent,
  FieldTitle,
}
\end{DeskFlowCode}

\phantomsection
\addcontentsline{toc}{subsubsection}{deskflow\-frontend/components/ui/form.tsx}
\DeskFlowListingCaption{deskflow-frontend/components/ui/form.tsx}{lst:fd-ui-form}
\begin{DeskFlowCode}
'use client'

import * as React from 'react'
import * as LabelPrimitive from '@radix-ui/react-label'
import { Slot } from '@radix-ui/react-slot'
import {
  Controller,
  FormProvider,
  useFormContext,
  useFormState,
  type ControllerProps,
  type FieldPath,
  type FieldValues,
} from 'react-hook-form'

import { cn } from '@/lib/utils'
import { Label } from '@/components/ui/label'

const Form = FormProvider

type FormFieldContextValue<
  TFieldValues extends FieldValues = FieldValues,
  TName extends FieldPath<TFieldValues> = FieldPath<TFieldValues>,
> = {
  name: TName
}

const FormFieldContext = React.createContext<FormFieldContextValue>(
  {} as FormFieldContextValue,
)

const FormField = <
  TFieldValues extends FieldValues = FieldValues,
  TName extends FieldPath<TFieldValues> = FieldPath<TFieldValues>,
>({
  ...props
}: ControllerProps<TFieldValues, TName>) => {
  return (
    <FormFieldContext.Provider value={{ name: props.name }}>
      <Controller {...props} />
    </FormFieldContext.Provider>
  )
}

const useFormField = () => {
  const fieldContext = React.useContext(FormFieldContext)
  const itemContext = React.useContext(FormItemContext)
  const { getFieldState } = useFormContext()
  const formState = useFormState({ name: fieldContext.name })
  const fieldState = getFieldState(fieldContext.name, formState)

  if (!fieldContext) {
    throw new Error('useFormField should be used within <FormField>')
  }

  const { id } = itemContext

  return {
    id,
    name: fieldContext.name,
    formItemId: `${id}-form-item`,
    formDescriptionId: `${id}-form-item-description`,
    formMessageId: `${id}-form-item-message`,
    ...fieldState,
  }
}

type FormItemContextValue = {
  id: string
}

const FormItemContext = React.createContext<FormItemContextValue>(
  {} as FormItemContextValue,
)

function FormItem({ className, ...props }: React.ComponentProps<'div'>) {
  const id = React.useId()

  return (
    <FormItemContext.Provider value={{ id }}>
      <div
        data-slot="form-item"
        className={cn('grid gap-2', className)}
        {...props}
      />
    </FormItemContext.Provider>
  )
}

function FormLabel({
  className,
  ...props
}: React.ComponentProps<typeof LabelPrimitive.Root>) {
  const { error, formItemId } = useFormField()

  return (
    <Label
      data-slot="form-label"
      data-error={!!error}
      className={cn('data-[error=true]:text-destructive', className)}
      htmlFor={formItemId}
      {...props}
    />
  )
}

function FormControl({ ...props }: React.ComponentProps<typeof Slot>) {
  const { error, formItemId, formDescriptionId, formMessageId } = useFormField()

  return (
    <Slot
      data-slot="form-control"
      id={formItemId}
      aria-describedby={
        !error
          ? `${formDescriptionId}`
          : `${formDescriptionId} ${formMessageId}`
      }
      aria-invalid={!!error}
      {...props}
    />
  )
}

function FormDescription({ className, ...props }: React.ComponentProps<'p'>) {
  const { formDescriptionId } = useFormField()

  return (
    <p
      data-slot="form-description"
      id={formDescriptionId}
      className={cn('text-muted-foreground text-sm', className)}
      {...props}
    />
  )
}

function FormMessage({ className, ...props }: React.ComponentProps<'p'>) {
  const { error, formMessageId } = useFormField()
  const body = error ? String(error?.message ?? '') : props.children

  if (!body) {
    return null
  }

  return (
    <p
      data-slot="form-message"
      id={formMessageId}
      className={cn('text-destructive text-sm', className)}
      {...props}
    >
      {body}
    </p>
  )
}

export {
  useFormField,
  Form,
  FormItem,
  FormLabel,
  FormControl,
  FormDescription,
  FormMessage,
  FormField,
}
\end{DeskFlowCode}

\phantomsection
\addcontentsline{toc}{subsubsection}{deskflow\-frontend/components/ui/hover\-card.tsx}
\DeskFlowListingCaption{deskflow-frontend/components/ui/hover-card.tsx}{lst:fd-ui-hover-card}
\begin{DeskFlowCode}
'use client'

import * as React from 'react'
import * as HoverCardPrimitive from '@radix-ui/react-hover-card'

import { cn } from '@/lib/utils'

function HoverCard({
  ...props
}: React.ComponentProps<typeof HoverCardPrimitive.Root>) {
  return <HoverCardPrimitive.Root data-slot="hover-card" {...props} />
}

function HoverCardTrigger({
  ...props
}: React.ComponentProps<typeof HoverCardPrimitive.Trigger>) {
  return (
    <HoverCardPrimitive.Trigger data-slot="hover-card-trigger" {...props} />
  )
}

function HoverCardContent({
  className,
  align = 'center',
  sideOffset = 4,
  ...props
}: React.ComponentProps<typeof HoverCardPrimitive.Content>) {
  return (
    <HoverCardPrimitive.Portal data-slot="hover-card-portal">
      <HoverCardPrimitive.Content
        data-slot="hover-card-content"
        align={align}
        sideOffset={sideOffset}
        className={cn(
          'bg-popover text-popover-foreground data-[state=open]:animate-in data-[state=closed]:animate-out data-[state=closed]:fade-out-0 data-[state=open]:fade-in-0 data-[state=closed]:zoom-out-95 data-[state=open]:zoom-in-95 data-[side=bottom]:slide-in-from-top-2 data-[side=left]:slide-in-from-right-2 data-[side=right]:slide-in-from-left-2 data-[side=top]:slide-in-from-bottom-2 z-50 w-64 origin-(-radix-hover-card-content-transform-origin) rounded-md border p-4 shadow-md outline-hidden',
          className,
        )}
        {...props}
      />
    </HoverCardPrimitive.Portal>
  )
}

export { HoverCard, HoverCardTrigger, HoverCardContent }
\end{DeskFlowCode}

\phantomsection
\addcontentsline{toc}{subsubsection}{deskflow\-frontend/components/ui/input\-group.tsx}
\DeskFlowListingCaption{deskflow-frontend/components/ui/input-group.tsx}{lst:fd-ui-input-group}
\begin{DeskFlowCode}
'use client'

import { cva, type VariantProps } from 'class-variance-authority'

import { cn } from '@/lib/utils'
import { Button } from '@/components/ui/button'
import { Input } from '@/components/ui/input'
import { Textarea } from '@/components/ui/textarea'

function InputGroup({ className, ...props }: React.ComponentProps<'div'>) {
  return (
    <div
      data-slot="input-group"
      role="group"
      className={cn(
        'group/input-group border-input dark:bg-input/30 relative flex w-full items-center rounded-md border shadow-xs transition-[color,box-shadow] outline-none',
        'h-9 has-[>textarea]:h-auto',

        // Variants based on alignment.
        'has-[>[data-align=inline-start]]:[&>input]:pl-2',
        'has-[>[data-align=inline-end]]:[&>input]:pr-2',
        'has-[>[data-align=block-start]]:h-auto has-[>[data-align=block-start]]:flex-col has-[>[data-align=block-start]]:[&>input]:pb-3',
        'has-[>[data-align=block-end]]:h-auto has-[>[data-align=block-end]]:flex-col has-[>[data-align=block-end]]:[&>input]:pt-3',

        // Focus state.
        'has-[[data-slot=input-group-control]:focus-visible]:border-ring has-[[data-slot=input-group-control]:focus-visible]:ring-ring/50 has-[[data-slot=input-group-control]:focus-visible]:ring-[3px]',

        // Error state.
        'has-[[data-slot][aria-invalid=true]]:ring-destructive/20 has-[[data-slot][aria-invalid=true]]:border-destructive dark:has-[[data-slot][aria-invalid=true]]:ring-destructive/40',

        className,
      )}
      {...props}
    />
  )
}

const inputGroupAddonVariants = cva(
  "text-muted-foreground flex h-auto cursor-text items-center justify-center gap-2 py-1.5 text-sm font-medium select-none [&>svg:not([class*='size-'])]:size-4 [&>kbd]:rounded-[calc(var(-radius)-5px)] group-data-[disabled=true]/input-group:opacity-50",
  {
    variants: {
      align: {
        'inline-start':
          'order-first pl-3 has-[>button]:ml-[-0.45rem] has-[>kbd]:ml-[-0.35rem]',
        'inline-end':
          'order-last pr-3 has-[>button]:mr-[-0.4rem] has-[>kbd]:mr-[-0.35rem]',
        'block-start':
          'order-first w-full justify-start px-3 pt-3 [.border-b]:pb-3 group-has-[>input]/input-group:pt-2.5',
        'block-end':
          'order-last w-full justify-start px-3 pb-3 [.border-t]:pt-3 group-has-[>input]/input-group:pb-2.5',
      },
    },
    defaultVariants: {
      align: 'inline-start',
    },
  },
)

function InputGroupAddon({
  className,
  align = 'inline-start',
  ...props
}: React.ComponentProps<'div'> & VariantProps<typeof inputGroupAddonVariants>) {
  return (
    <div
      role="group"
      data-slot="input-group-addon"
      data-align={align}
      className={cn(inputGroupAddonVariants({ align }), className)}
      onClick={(e) => {
        if ((e.target as HTMLElement).closest('button')) {
          return
        }
        e.currentTarget.parentElement?.querySelector('input')?.focus()
      }}
      {...props}
    />
  )
}

const inputGroupButtonVariants = cva(
  'text-sm shadow-none flex gap-2 items-center',
  {
    variants: {
      size: {
        xs: "h-6 gap-1 px-2 rounded-[calc(var(-radius)-5px)] [&>svg:not([class*='size-'])]:size-3.5 has-[>svg]:px-2",
        sm: 'h-8 px-2.5 gap-1.5 rounded-md has-[>svg]:px-2.5',
        'icon-xs':
          'size-6 rounded-[calc(var(-radius)-5px)] p-0 has-[>svg]:p-0',
        'icon-sm': 'size-8 p-0 has-[>svg]:p-0',
      },
    },
    defaultVariants: {
      size: 'xs',
    },
  },
)

function InputGroupButton({
  className,
  type = 'button',
  variant = 'ghost',
  size = 'xs',
  ...props
}: Omit<React.ComponentProps<typeof Button>, 'size'> &
  VariantProps<typeof inputGroupButtonVariants>) {
  return (
    <Button
      type={type}
      data-size={size}
      variant={variant}
      className={cn(inputGroupButtonVariants({ size }), className)}
      {...props}
    />
  )
}

function InputGroupText({ className, ...props }: React.ComponentProps<'span'>) {
  return (
    <span
      className={cn(
        "text-muted-foreground flex items-center gap-2 text-sm [&_svg]:pointer-events-none [&_svg:not([class*='size-'])]:size-4",
        className,
      )}
      {...props}
    />
  )
}

function InputGroupInput({
  className,
  ...props
}: React.ComponentProps<'input'>) {
  return (
    <Input
      data-slot="input-group-control"
      className={cn(
        'flex-1 rounded-none border-0 bg-transparent shadow-none focus-visible:ring-0 dark:bg-transparent',
        className,
      )}
      {...props}
    />
  )
}

function InputGroupTextarea({
  className,
  ...props
}: React.ComponentProps<'textarea'>) {
  return (
    <Textarea
      data-slot="input-group-control"
      className={cn(
        'flex-1 resize-none rounded-none border-0 bg-transparent py-3 shadow-none focus-visible:ring-0 dark:bg-transparent',
        className,
      )}
      {...props}
    />
  )
}

export {
  InputGroup,
  InputGroupAddon,
  InputGroupButton,
  InputGroupText,
  InputGroupInput,
  InputGroupTextarea,
}
\end{DeskFlowCode}

\phantomsection
\addcontentsline{toc}{subsubsection}{deskflow\-frontend/components/ui/input\-otp.tsx}
\DeskFlowListingCaption{deskflow-frontend/components/ui/input-otp.tsx}{lst:fd-ui-input-otp}
\begin{DeskFlowCode}
'use client'

import * as React from 'react'
import { OTPInput, OTPInputContext } from 'input-otp'
import { MinusIcon } from 'lucide-react'

import { cn } from '@/lib/utils'

function InputOTP({
  className,
  containerClassName,
  ...props
}: React.ComponentProps<typeof OTPInput> & {
  containerClassName?: string
}) {
  return (
    <OTPInput
      data-slot="input-otp"
      containerClassName={cn(
        'flex items-center gap-2 has-disabled:opacity-50',
        containerClassName,
      )}
      className={cn('disabled:cursor-not-allowed', className)}
      {...props}
    />
  )
}

function InputOTPGroup({ className, ...props }: React.ComponentProps<'div'>) {
  return (
    <div
      data-slot="input-otp-group"
      className={cn('flex items-center', className)}
      {...props}
    />
  )
}

function InputOTPSlot({
  index,
  className,
  ...props
}: React.ComponentProps<'div'> & {
  index: number
}) {
  const inputOTPContext = React.useContext(OTPInputContext)
  const { char, hasFakeCaret, isActive } = inputOTPContext?.slots[index] ?? {}

  return (
    <div
      data-slot="input-otp-slot"
      data-active={isActive}
      className={cn(
        'data-[active=true]:border-ring data-[active=true]:ring-ring/50 data-[active=true]:aria-invalid:ring-destructive/20 dark:data-[active=true]:aria-invalid:ring-destructive/40 aria-invalid:border-destructive data-[active=true]:aria-invalid:border-destructive dark:bg-input/30 border-input relative flex h-9 w-9 items-center justify-center border-y border-r text-sm shadow-xs transition-all outline-none first:rounded-l-md first:border-l last:rounded-r-md data-[active=true]:z-10 data-[active=true]:ring-[3px]',
        className,
      )}
      {...props}
    >
      {char}
      {hasFakeCaret && (
        <div className="pointer-events-none absolute inset-0 flex items-center justify-center">
          <div className="animate-caret-blink bg-foreground h-4 w-px duration-1000" />
        </div>
      )}
    </div>
  )
}

function InputOTPSeparator({ ...props }: React.ComponentProps<'div'>) {
  return (
    <div data-slot="input-otp-separator" role="separator" {...props}>
      <MinusIcon />
    </div>
  )
}

export { InputOTP, InputOTPGroup, InputOTPSlot, InputOTPSeparator }
\end{DeskFlowCode}

\phantomsection
\addcontentsline{toc}{subsubsection}{deskflow\-frontend/components/ui/input.tsx}
\DeskFlowListingCaption{deskflow-frontend/components/ui/input.tsx}{lst:fd-ui-input}
\begin{DeskFlowCode}
import * as React from 'react'

import { cn } from '@/lib/utils'

function Input({ className, type, ...props }: React.ComponentProps<'input'>) {
  return (
    <input
      type={type}
      data-slot="input"
      className={cn(
        'file:text-foreground placeholder:text-muted-foreground selection:bg-primary selection:text-primary-foreground dark:bg-input/30 border-input h-9 w-full min-w-0 rounded-md border bg-transparent px-3 py-1 text-base shadow-xs transition-[color,box-shadow] outline-none file:inline-flex file:h-7 file:border-0 file:bg-transparent file:text-sm file:font-medium disabled:pointer-events-none disabled:cursor-not-allowed disabled:opacity-50 md:text-sm',
        'focus-visible:border-ring focus-visible:ring-ring/50 focus-visible:ring-[3px]',
        'aria-invalid:ring-destructive/20 dark:aria-invalid:ring-destructive/40 aria-invalid:border-destructive',
        className,
      )}
      {...props}
    />
  )
}

export { Input }
\end{DeskFlowCode}

\phantomsection
\addcontentsline{toc}{subsubsection}{deskflow\-frontend/components/ui/item.tsx}
\DeskFlowListingCaption{deskflow-frontend/components/ui/item.tsx}{lst:fd-ui-item}
\begin{DeskFlowCode}
import * as React from 'react'
import { Slot } from '@radix-ui/react-slot'
import { cva, type VariantProps } from 'class-variance-authority'

import { cn } from '@/lib/utils'
import { Separator } from '@/components/ui/separator'

function ItemGroup({ className, ...props }: React.ComponentProps<'div'>) {
  return (
    <div
      role="list"
      data-slot="item-group"
      className={cn('group/item-group flex flex-col', className)}
      {...props}
    />
  )
}

function ItemSeparator({
  className,
  ...props
}: React.ComponentProps<typeof Separator>) {
  return (
    <Separator
      data-slot="item-separator"
      orientation="horizontal"
      className={cn('my-0', className)}
      {...props}
    />
  )
}

const itemVariants = cva(
  'group/item flex items-center border border-transparent text-sm rounded-md transition-colors [a&]:hover:bg-accent/50 [a&]:transition-colors duration-100 flex-wrap outline-none focus-visible:border-ring focus-visible:ring-ring/50 focus-visible:ring-[3px]',
  {
    variants: {
      variant: {
        default: 'bg-transparent',
        outline: 'border-border',
        muted: 'bg-muted/50',
      },
      size: {
        default: 'p-4 gap-4 ',
        sm: 'py-3 px-4 gap-2.5',
      },
    },
    defaultVariants: {
      variant: 'default',
      size: 'default',
    },
  },
)

function Item({
  className,
  variant = 'default',
  size = 'default',
  asChild = false,
  ...props
}: React.ComponentProps<'div'> &
  VariantProps<typeof itemVariants> & { asChild?: boolean }) {
  const Comp = asChild ? Slot : 'div'
  return (
    <Comp
      data-slot="item"
      data-variant={variant}
      data-size={size}
      className={cn(itemVariants({ variant, size, className }))}
      {...props}
    />
  )
}

const itemMediaVariants = cva(
  'flex shrink-0 items-center justify-center gap-2 group-has-[[data-slot=item-description]]/item:self-start [&_svg]:pointer-events-none group-has-[[data-slot=item-description]]/item:translate-y-0.5',
  {
    variants: {
      variant: {
        default: 'bg-transparent',
        icon: "size-8 border rounded-sm bg-muted [&_svg:not([class*='size-'])]:size-4",
        image:
          'size-10 rounded-sm overflow-hidden [&_img]:size-full [&_img]:object-cover',
      },
    },
    defaultVariants: {
      variant: 'default',
    },
  },
)

function ItemMedia({
  className,
  variant = 'default',
  ...props
}: React.ComponentProps<'div'> & VariantProps<typeof itemMediaVariants>) {
  return (
    <div
      data-slot="item-media"
      data-variant={variant}
      className={cn(itemMediaVariants({ variant, className }))}
      {...props}
    />
  )
}

function ItemContent({ className, ...props }: React.ComponentProps<'div'>) {
  return (
    <div
      data-slot="item-content"
      className={cn(
        'flex flex-1 flex-col gap-1 [&+[data-slot=item-content]]:flex-none',
        className,
      )}
      {...props}
    />
  )
}

function ItemTitle({ className, ...props }: React.ComponentProps<'div'>) {
  return (
    <div
      data-slot="item-title"
      className={cn(
        'flex w-fit items-center gap-2 text-sm leading-snug font-medium',
        className,
      )}
      {...props}
    />
  )
}

function ItemDescription({ className, ...props }: React.ComponentProps<'p'>) {
  return (
    <p
      data-slot="item-description"
      className={cn(
        'text-muted-foreground line-clamp-2 text-sm leading-normal font-normal text-balance',
        '[&>a:hover]:text-primary [&>a]:underline [&>a]:underline-offset-4',
        className,
      )}
      {...props}
    />
  )
}

function ItemActions({ className, ...props }: React.ComponentProps<'div'>) {
  return (
    <div
      data-slot="item-actions"
      className={cn('flex items-center gap-2', className)}
      {...props}
    />
  )
}

function ItemHeader({ className, ...props }: React.ComponentProps<'div'>) {
  return (
    <div
      data-slot="item-header"
      className={cn(
        'flex basis-full items-center justify-between gap-2',
        className,
      )}
      {...props}
    />
  )
}

function ItemFooter({ className, ...props }: React.ComponentProps<'div'>) {
  return (
    <div
      data-slot="item-footer"
      className={cn(
        'flex basis-full items-center justify-between gap-2',
        className,
      )}
      {...props}
    />
  )
}

export {
  Item,
  ItemMedia,
  ItemContent,
  ItemActions,
  ItemGroup,
  ItemSeparator,
  ItemTitle,
  ItemDescription,
  ItemHeader,
  ItemFooter,
}
\end{DeskFlowCode}

\phantomsection
\addcontentsline{toc}{subsubsection}{deskflow\-frontend/components/ui/kbd.tsx}
\DeskFlowListingCaption{deskflow-frontend/components/ui/kbd.tsx}{lst:fd-ui-kbd}
\begin{DeskFlowCode}
import { cn } from '@/lib/utils'

function Kbd({ className, ...props }: React.ComponentProps<'kbd'>) {
  return (
    <kbd
      data-slot="kbd"
      className={cn(
        'bg-muted w-fit text-muted-foreground pointer-events-none inline-flex h-5 min-w-5 items-center justify-center gap-1 rounded-sm px-1 font-sans text-xs font-medium select-none',
        "[&_svg:not([class*='size-'])]:size-3",
        '[[data-slot=tooltip-content]_&]:bg-background/20 [[data-slot=tooltip-content]_&]:text-background dark:[[data-slot=tooltip-content]_&]:bg-background/10',
        className,
      )}
      {...props}
    />
  )
}

function KbdGroup({ className, ...props }: React.ComponentProps<'div'>) {
  return (
    <kbd
      data-slot="kbd-group"
      className={cn('inline-flex items-center gap-1', className)}
      {...props}
    />
  )
}

export { Kbd, KbdGroup }
\end{DeskFlowCode}

\phantomsection
\addcontentsline{toc}{subsubsection}{deskflow\-frontend/components/ui/label.tsx}
\DeskFlowListingCaption{deskflow-frontend/components/ui/label.tsx}{lst:fd-ui-label}
\begin{DeskFlowCode}
'use client'

import * as React from 'react'
import * as LabelPrimitive from '@radix-ui/react-label'

import { cn } from '@/lib/utils'

function Label({
  className,
  ...props
}: React.ComponentProps<typeof LabelPrimitive.Root>) {
  return (
    <LabelPrimitive.Root
      data-slot="label"
      className={cn(
        'flex items-center gap-2 text-sm leading-none font-medium select-none group-data-[disabled=true]:pointer-events-none group-data-[disabled=true]:opacity-50 peer-disabled:cursor-not-allowed peer-disabled:opacity-50',
        className,
      )}
      {...props}
    />
  )
}

export { Label }
\end{DeskFlowCode}

\phantomsection
\addcontentsline{toc}{subsubsection}{deskflow\-frontend/components/ui/menubar.tsx}
\DeskFlowListingCaption{deskflow-frontend/components/ui/menubar.tsx}{lst:fd-ui-menubar}
\begin{DeskFlowCode}
'use client'

import * as React from 'react'
import * as MenubarPrimitive from '@radix-ui/react-menubar'
import { CheckIcon, ChevronRightIcon, CircleIcon } from 'lucide-react'

import { cn } from '@/lib/utils'

function Menubar({
  className,
  ...props
}: React.ComponentProps<typeof MenubarPrimitive.Root>) {
  return (
    <MenubarPrimitive.Root
      data-slot="menubar"
      className={cn(
        'bg-background flex h-9 items-center gap-1 rounded-md border p-1 shadow-xs',
        className,
      )}
      {...props}
    />
  )
}

function MenubarMenu({
  ...props
}: React.ComponentProps<typeof MenubarPrimitive.Menu>) {
  return <MenubarPrimitive.Menu data-slot="menubar-menu" {...props} />
}

function MenubarGroup({
  ...props
}: React.ComponentProps<typeof MenubarPrimitive.Group>) {
  return <MenubarPrimitive.Group data-slot="menubar-group" {...props} />
}

function MenubarPortal({
  ...props
}: React.ComponentProps<typeof MenubarPrimitive.Portal>) {
  return <MenubarPrimitive.Portal data-slot="menubar-portal" {...props} />
}

function MenubarRadioGroup({
  ...props
}: React.ComponentProps<typeof MenubarPrimitive.RadioGroup>) {
  return (
    <MenubarPrimitive.RadioGroup data-slot="menubar-radio-group" {...props} />
  )
}

function MenubarTrigger({
  className,
  ...props
}: React.ComponentProps<typeof MenubarPrimitive.Trigger>) {
  return (
    <MenubarPrimitive.Trigger
      data-slot="menubar-trigger"
      className={cn(
        'focus:bg-accent focus:text-accent-foreground data-[state=open]:bg-accent data-[state=open]:text-accent-foreground flex items-center rounded-sm px-2 py-1 text-sm font-medium outline-hidden select-none',
        className,
      )}
      {...props}
    />
  )
}

function MenubarContent({
  className,
  align = 'start',
  alignOffset = -4,
  sideOffset = 8,
  ...props
}: React.ComponentProps<typeof MenubarPrimitive.Content>) {
  return (
    <MenubarPortal>
      <MenubarPrimitive.Content
        data-slot="menubar-content"
        align={align}
        alignOffset={alignOffset}
        sideOffset={sideOffset}
        className={cn(
          'bg-popover text-popover-foreground data-[state=open]:animate-in data-[state=closed]:fade-out-0 data-[state=open]:fade-in-0 data-[state=closed]:zoom-out-95 data-[state=open]:zoom-in-95 data-[side=bottom]:slide-in-from-top-2 data-[side=left]:slide-in-from-right-2 data-[side=right]:slide-in-from-left-2 data-[side=top]:slide-in-from-bottom-2 z-50 min-w-[12rem] origin-(-radix-menubar-content-transform-origin) overflow-hidden rounded-md border p-1 shadow-md',
          className,
        )}
        {...props}
      />
    </MenubarPortal>
  )
}

function MenubarItem({
  className,
  inset,
  variant = 'default',
  ...props
}: React.ComponentProps<typeof MenubarPrimitive.Item> & {
  inset?: boolean
  variant?: 'default' | 'destructive'
}) {
  return (
    <MenubarPrimitive.Item
      data-slot="menubar-item"
      data-inset={inset}
      data-variant={variant}
      className={cn(
        "focus:bg-accent focus:text-accent-foreground data-[variant=destructive]:text-destructive data-[variant=destructive]:focus:bg-destructive/10 dark:data-[variant=destructive]:focus:bg-destructive/20 data-[variant=destructive]:focus:text-destructive data-[variant=destructive]:*:[svg]:!text-destructive [&_svg:not([class*='text-'])]:text-muted-foreground relative flex cursor-default items-center gap-2 rounded-sm px-2 py-1.5 text-sm outline-hidden select-none data-[disabled]:pointer-events-none data-[disabled]:opacity-50 data-[inset]:pl-8 [&_svg]:pointer-events-none [&_svg]:shrink-0 [&_svg:not([class*='size-'])]:size-4",
        className,
      )}
      {...props}
    />
  )
}

function MenubarCheckboxItem({
  className,
  children,
  checked,
  ...props
}: React.ComponentProps<typeof MenubarPrimitive.CheckboxItem>) {
  return (
    <MenubarPrimitive.CheckboxItem
      data-slot="menubar-checkbox-item"
      className={cn(
        "focus:bg-accent focus:text-accent-foreground relative flex cursor-default items-center gap-2 rounded-xs py-1.5 pr-2 pl-8 text-sm outline-hidden select-none data-[disabled]:pointer-events-none data-[disabled]:opacity-50 [&_svg]:pointer-events-none [&_svg]:shrink-0 [&_svg:not([class*='size-'])]:size-4",
        className,
      )}
      checked={checked}
      {...props}
    >
      <span className="pointer-events-none absolute left-2 flex size-3.5 items-center justify-center">
        <MenubarPrimitive.ItemIndicator>
          <CheckIcon className="size-4" />
        </MenubarPrimitive.ItemIndicator>
      </span>
      {children}
    </MenubarPrimitive.CheckboxItem>
  )
}

function MenubarRadioItem({
  className,
  children,
  ...props
}: React.ComponentProps<typeof MenubarPrimitive.RadioItem>) {
  return (
    <MenubarPrimitive.RadioItem
      data-slot="menubar-radio-item"
      className={cn(
        "focus:bg-accent focus:text-accent-foreground relative flex cursor-default items-center gap-2 rounded-xs py-1.5 pr-2 pl-8 text-sm outline-hidden select-none data-[disabled]:pointer-events-none data-[disabled]:opacity-50 [&_svg]:pointer-events-none [&_svg]:shrink-0 [&_svg:not([class*='size-'])]:size-4",
        className,
      )}
      {...props}
    >
      <span className="pointer-events-none absolute left-2 flex size-3.5 items-center justify-center">
        <MenubarPrimitive.ItemIndicator>
          <CircleIcon className="size-2 fill-current" />
        </MenubarPrimitive.ItemIndicator>
      </span>
      {children}
    </MenubarPrimitive.RadioItem>
  )
}

function MenubarLabel({
  className,
  inset,
  ...props
}: React.ComponentProps<typeof MenubarPrimitive.Label> & {
  inset?: boolean
}) {
  return (
    <MenubarPrimitive.Label
      data-slot="menubar-label"
      data-inset={inset}
      className={cn(
        'px-2 py-1.5 text-sm font-medium data-[inset]:pl-8',
        className,
      )}
      {...props}
    />
  )
}

function MenubarSeparator({
  className,
  ...props
}: React.ComponentProps<typeof MenubarPrimitive.Separator>) {
  return (
    <MenubarPrimitive.Separator
      data-slot="menubar-separator"
      className={cn('bg-border -mx-1 my-1 h-px', className)}
      {...props}
    />
  )
}

function MenubarShortcut({
  className,
  ...props
}: React.ComponentProps<'span'>) {
  return (
    <span
      data-slot="menubar-shortcut"
      className={cn(
        'text-muted-foreground ml-auto text-xs tracking-widest',
        className,
      )}
      {...props}
    />
  )
}

function MenubarSub({
  ...props
}: React.ComponentProps<typeof MenubarPrimitive.Sub>) {
  return <MenubarPrimitive.Sub data-slot="menubar-sub" {...props} />
}

function MenubarSubTrigger({
  className,
  inset,
  children,
  ...props
}: React.ComponentProps<typeof MenubarPrimitive.SubTrigger> & {
  inset?: boolean
}) {
  return (
    <MenubarPrimitive.SubTrigger
      data-slot="menubar-sub-trigger"
      data-inset={inset}
      className={cn(
        'focus:bg-accent focus:text-accent-foreground data-[state=open]:bg-accent data-[state=open]:text-accent-foreground flex cursor-default items-center rounded-sm px-2 py-1.5 text-sm outline-none select-none data-[inset]:pl-8',
        className,
      )}
      {...props}
    >
      {children}
      <ChevronRightIcon className="ml-auto h-4 w-4" />
    </MenubarPrimitive.SubTrigger>
  )
}

function MenubarSubContent({
  className,
  ...props
}: React.ComponentProps<typeof MenubarPrimitive.SubContent>) {
  return (
    <MenubarPrimitive.SubContent
      data-slot="menubar-sub-content"
      className={cn(
        'bg-popover text-popover-foreground data-[state=open]:animate-in data-[state=closed]:animate-out data-[state=closed]:fade-out-0 data-[state=open]:fade-in-0 data-[state=closed]:zoom-out-95 data-[state=open]:zoom-in-95 data-[side=bottom]:slide-in-from-top-2 data-[side=left]:slide-in-from-right-2 data-[side=right]:slide-in-from-left-2 data-[side=top]:slide-in-from-bottom-2 z-50 min-w-[8rem] origin-(-radix-menubar-content-transform-origin) overflow-hidden rounded-md border p-1 shadow-lg',
        className,
      )}
      {...props}
    />
  )
}

export {
  Menubar,
  MenubarPortal,
  MenubarMenu,
  MenubarTrigger,
  MenubarContent,
  MenubarGroup,
  MenubarSeparator,
  MenubarLabel,
  MenubarItem,
  MenubarShortcut,
  MenubarCheckboxItem,
  MenubarRadioGroup,
  MenubarRadioItem,
  MenubarSub,
  MenubarSubTrigger,
  MenubarSubContent,
}
\end{DeskFlowCode}

\phantomsection
\addcontentsline{toc}{subsubsection}{deskflow\-frontend/components/ui/navigation\-menu.tsx}
\DeskFlowListingCaption{deskflow-frontend/components/ui/navigation-menu.tsx}{lst:fd-ui-navigation-menu}
\begin{DeskFlowCode}
import * as React from 'react'
import * as NavigationMenuPrimitive from '@radix-ui/react-navigation-menu'
import { cva } from 'class-variance-authority'
import { ChevronDownIcon } from 'lucide-react'

import { cn } from '@/lib/utils'

function NavigationMenu({
  className,
  children,
  viewport = true,
  ...props
}: React.ComponentProps<typeof NavigationMenuPrimitive.Root> & {
  viewport?: boolean
}) {
  return (
    <NavigationMenuPrimitive.Root
      data-slot="navigation-menu"
      data-viewport={viewport}
      className={cn(
        'group/navigation-menu relative flex max-w-max flex-1 items-center justify-center',
        className,
      )}
      {...props}
    >
      {children}
      {viewport && <NavigationMenuViewport />}
    </NavigationMenuPrimitive.Root>
  )
}

function NavigationMenuList({
  className,
  ...props
}: React.ComponentProps<typeof NavigationMenuPrimitive.List>) {
  return (
    <NavigationMenuPrimitive.List
      data-slot="navigation-menu-list"
      className={cn(
        'group flex flex-1 list-none items-center justify-center gap-1',
        className,
      )}
      {...props}
    />
  )
}

function NavigationMenuItem({
  className,
  ...props
}: React.ComponentProps<typeof NavigationMenuPrimitive.Item>) {
  return (
    <NavigationMenuPrimitive.Item
      data-slot="navigation-menu-item"
      className={cn('relative', className)}
      {...props}
    />
  )
}

const navigationMenuTriggerStyle = cva(
  'group inline-flex h-9 w-max items-center justify-center rounded-md bg-background px-4 py-2 text-sm font-medium hover:bg-accent hover:text-accent-foreground focus:bg-accent focus:text-accent-foreground disabled:pointer-events-none disabled:opacity-50 data-[state=open]:hover:bg-accent data-[state=open]:text-accent-foreground data-[state=open]:focus:bg-accent data-[state=open]:bg-accent/50 focus-visible:ring-ring/50 outline-none transition-[color,box-shadow] focus-visible:ring-[3px] focus-visible:outline-1',
)

function NavigationMenuTrigger({
  className,
  children,
  ...props
}: React.ComponentProps<typeof NavigationMenuPrimitive.Trigger>) {
  return (
    <NavigationMenuPrimitive.Trigger
      data-slot="navigation-menu-trigger"
      className={cn(navigationMenuTriggerStyle(), 'group', className)}
      {...props}
    >
      {children}{' '}
      <ChevronDownIcon
        className="relative top-[1px] ml-1 size-3 transition duration-300 group-data-[state=open]:rotate-180"
        aria-hidden="true"
      />
    </NavigationMenuPrimitive.Trigger>
  )
}

function NavigationMenuContent({
  className,
  ...props
}: React.ComponentProps<typeof NavigationMenuPrimitive.Content>) {
  return (
    <NavigationMenuPrimitive.Content
      data-slot="navigation-menu-content"
      className={cn(
        'data-[motion^=from-]:animate-in data-[motion^=to-]:animate-out data-[motion^=from-]:fade-in data-[motion^=to-]:fade-out data-[motion=from-end]:slide-in-from-right-52 data-[motion=from-start]:slide-in-from-left-52 data-[motion=to-end]:slide-out-to-right-52 data-[motion=to-start]:slide-out-to-left-52 top-0 left-0 w-full p-2 pr-2.5 md:absolute md:w-auto',
        'group-data-[viewport=false]/navigation-menu:bg-popover group-data-[viewport=false]/navigation-menu:text-popover-foreground group-data-[viewport=false]/navigation-menu:data-[state=open]:animate-in group-data-[viewport=false]/navigation-menu:data-[state=closed]:animate-out group-data-[viewport=false]/navigation-menu:data-[state=closed]:zoom-out-95 group-data-[viewport=false]/navigation-menu:data-[state=open]:zoom-in-95 group-data-[viewport=false]/navigation-menu:data-[state=open]:fade-in-0 group-data-[viewport=false]/navigation-menu:data-[state=closed]:fade-out-0 group-data-[viewport=false]/navigation-menu:top-full group-data-[viewport=false]/navigation-menu:mt-1.5 group-data-[viewport=false]/navigation-menu:overflow-hidden group-data-[viewport=false]/navigation-menu:rounded-md group-data-[viewport=false]/navigation-menu:border group-data-[viewport=false]/navigation-menu:shadow group-data-[viewport=false]/navigation-menu:duration-200 **:data-[slot=navigation-menu-link]:focus:ring-0 **:data-[slot=navigation-menu-link]:focus:outline-none',
        className,
      )}
      {...props}
    />
  )
}

function NavigationMenuViewport({
  className,
  ...props
}: React.ComponentProps<typeof NavigationMenuPrimitive.Viewport>) {
  return (
    <div
      className="absolute top-full left-0 isolate z-50 flex justify-center"
    >
      <NavigationMenuPrimitive.Viewport
        data-slot="navigation-menu-viewport"
        className={cn(
          'origin-top-center bg-popover text-popover-foreground data-[state=open]:animate-in data-[state=closed]:animate-out data-[state=closed]:zoom-out-95 data-[state=open]:zoom-in-90 relative mt-1.5 h-[var(-radix-navigation-menu-viewport-height)] w-full overflow-hidden rounded-md border shadow md:w-[var(-radix-navigation-menu-viewport-width)]',
          className,
        )}
        {...props}
      />
    </div>
  )
}

function NavigationMenuLink({
  className,
  ...props
}: React.ComponentProps<typeof NavigationMenuPrimitive.Link>) {
  return (
    <NavigationMenuPrimitive.Link
      data-slot="navigation-menu-link"
      className={cn(
        "data-[active=true]:focus:bg-accent data-[active=true]:hover:bg-accent data-[active=true]:bg-accent/50 data-[active=true]:text-accent-foreground hover:bg-accent hover:text-accent-foreground focus:bg-accent focus:text-accent-foreground focus-visible:ring-ring/50 [&_svg:not([class*='text-'])]:text-muted-foreground flex flex-col gap-1 rounded-sm p-2 text-sm transition-all outline-none focus-visible:ring-[3px] focus-visible:outline-1 [&_svg:not([class*='size-'])]:size-4",
        className,
      )}
      {...props}
    />
  )
}

function NavigationMenuIndicator({
  className,
  ...props
}: React.ComponentProps<typeof NavigationMenuPrimitive.Indicator>) {
  return (
    <NavigationMenuPrimitive.Indicator
      data-slot="navigation-menu-indicator"
      className={cn(
        'data-[state=visible]:animate-in data-[state=hidden]:animate-out data-[state=hidden]:fade-out data-[state=visible]:fade-in top-full z-[1] flex h-1.5 items-end justify-center overflow-hidden',
        className,
      )}
      {...props}
    >
      <div className="bg-border relative top-[60%] h-2 w-2 rotate-45 rounded-tl-sm shadow-md" />
    </NavigationMenuPrimitive.Indicator>
  )
}

export {
  NavigationMenu,
  NavigationMenuList,
  NavigationMenuItem,
  NavigationMenuContent,
  NavigationMenuTrigger,
  NavigationMenuLink,
  NavigationMenuIndicator,
  NavigationMenuViewport,
  navigationMenuTriggerStyle,
}
\end{DeskFlowCode}

\phantomsection
\addcontentsline{toc}{subsubsection}{deskflow\-frontend/components/ui/pagination.tsx}
\DeskFlowListingCaption{deskflow-frontend/components/ui/pagination.tsx}{lst:fd-ui-pagination}
\begin{DeskFlowCode}
import * as React from 'react'
import {
  ChevronLeftIcon,
  ChevronRightIcon,
  MoreHorizontalIcon,
} from 'lucide-react'

import { cn } from '@/lib/utils'
import { Button, buttonVariants } from '@/components/ui/button'

function Pagination({ className, ...props }: React.ComponentProps<'nav'>) {
  return (
    <nav
      role="navigation"
      aria-label="pagination"
      data-slot="pagination"
      className={cn('mx-auto flex w-full justify-center', className)}
      {...props}
    />
  )
}

function PaginationContent({
  className,
  ...props
}: React.ComponentProps<'ul'>) {
  return (
    <ul
      data-slot="pagination-content"
      className={cn('flex flex-row items-center gap-1', className)}
      {...props}
    />
  )
}

function PaginationItem({ ...props }: React.ComponentProps<'li'>) {
  return <li data-slot="pagination-item" {...props} />
}

type PaginationLinkProps = {
  isActive?: boolean
} & Pick<React.ComponentProps<typeof Button>, 'size'> &
  React.ComponentProps<'a'>

function PaginationLink({
  className,
  isActive,
  size = 'icon',
  ...props
}: PaginationLinkProps) {
  return (
    <a
      aria-current={isActive ? 'page' : undefined}
      data-slot="pagination-link"
      data-active={isActive}
      className={cn(
        buttonVariants({
          variant: isActive ? 'outline' : 'ghost',
          size,
        }),
        className,
      )}
      {...props}
    />
  )
}

function PaginationPrevious({
  className,
  ...props
}: React.ComponentProps<typeof PaginationLink>) {
  return (
    <PaginationLink
      aria-label="Go to previous page"
      size="default"
      className={cn('gap-1 px-2.5 sm:pl-2.5', className)}
      {...props}
    >
      <ChevronLeftIcon />
      <span className="hidden sm:block">Previous</span>
    </PaginationLink>
  )
}

function PaginationNext({
  className,
  ...props
}: React.ComponentProps<typeof PaginationLink>) {
  return (
    <PaginationLink
      aria-label="Go to next page"
      size="default"
      className={cn('gap-1 px-2.5 sm:pr-2.5', className)}
      {...props}
    >
      <span className="hidden sm:block">Next</span>
      <ChevronRightIcon />
    </PaginationLink>
  )
}

function PaginationEllipsis({
  className,
  ...props
}: React.ComponentProps<'span'>) {
  return (
    <span
      aria-hidden
      data-slot="pagination-ellipsis"
      className={cn('flex size-9 items-center justify-center', className)}
      {...props}
    >
      <MoreHorizontalIcon className="size-4" />
      <span className="sr-only">More pages</span>
    </span>
  )
}

export {
  Pagination,
  PaginationContent,
  PaginationLink,
  PaginationItem,
  PaginationPrevious,
  PaginationNext,
  PaginationEllipsis,
}
\end{DeskFlowCode}

\phantomsection
\addcontentsline{toc}{subsubsection}{deskflow\-frontend/components/ui/popover.tsx}
\DeskFlowListingCaption{deskflow-frontend/components/ui/popover.tsx}{lst:fd-ui-popover}
\begin{DeskFlowCode}
'use client'

import * as React from 'react'
import * as PopoverPrimitive from '@radix-ui/react-popover'

import { cn } from '@/lib/utils'

function Popover({
  ...props
}: React.ComponentProps<typeof PopoverPrimitive.Root>) {
  return <PopoverPrimitive.Root data-slot="popover" {...props} />
}

function PopoverTrigger({
  ...props
}: React.ComponentProps<typeof PopoverPrimitive.Trigger>) {
  return <PopoverPrimitive.Trigger data-slot="popover-trigger" {...props} />
}

function PopoverContent({
  className,
  align = 'center',
  sideOffset = 4,
  ...props
}: React.ComponentProps<typeof PopoverPrimitive.Content>) {
  return (
    <PopoverPrimitive.Portal>
      <PopoverPrimitive.Content
        data-slot="popover-content"
        align={align}
        sideOffset={sideOffset}
        className={cn(
          'bg-popover text-popover-foreground data-[state=open]:animate-in data-[state=closed]:animate-out data-[state=closed]:fade-out-0 data-[state=open]:fade-in-0 data-[state=closed]:zoom-out-95 data-[state=open]:zoom-in-95 data-[side=bottom]:slide-in-from-top-2 data-[side=left]:slide-in-from-right-2 data-[side=right]:slide-in-from-left-2 data-[side=top]:slide-in-from-bottom-2 z-50 w-72 origin-(-radix-popover-content-transform-origin) rounded-md border p-4 shadow-md outline-hidden',
          className,
        )}
        {...props}
      />
    </PopoverPrimitive.Portal>
  )
}

function PopoverAnchor({
  ...props
}: React.ComponentProps<typeof PopoverPrimitive.Anchor>) {
  return <PopoverPrimitive.Anchor data-slot="popover-anchor" {...props} />
}

export { Popover, PopoverTrigger, PopoverContent, PopoverAnchor }
\end{DeskFlowCode}

\phantomsection
\addcontentsline{toc}{subsubsection}{deskflow\-frontend/components/ui/progress.tsx}
\DeskFlowListingCaption{deskflow-frontend/components/ui/progress.tsx}{lst:fd-ui-progress}
\begin{DeskFlowCode}
'use client'

import * as React from 'react'
import * as ProgressPrimitive from '@radix-ui/react-progress'

import { cn } from '@/lib/utils'

function Progress({
  className,
  value,
  ...props
}: React.ComponentProps<typeof ProgressPrimitive.Root>) {
  return (
    <ProgressPrimitive.Root
      data-slot="progress"
      className={cn(
        'bg-primary/20 relative h-2 w-full overflow-hidden rounded-full',
        className,
      )}
      {...props}
    >
      <ProgressPrimitive.Indicator
        data-slot="progress-indicator"
        className="bg-primary h-full w-full flex-1 transition-all"
        style={{ transform: `translateX(-${100 - (value || 0)}%)` }}
      />
    </ProgressPrimitive.Root>
  )
}

export { Progress }
\end{DeskFlowCode}

\phantomsection
\addcontentsline{toc}{subsubsection}{deskflow\-frontend/components/ui/radio\-group.tsx}
\DeskFlowListingCaption{deskflow-frontend/components/ui/radio-group.tsx}{lst:fd-ui-radio-group}
\begin{DeskFlowCode}
'use client'

import * as React from 'react'
import * as RadioGroupPrimitive from '@radix-ui/react-radio-group'
import { CircleIcon } from 'lucide-react'

import { cn } from '@/lib/utils'

function RadioGroup({
  className,
  ...props
}: React.ComponentProps<typeof RadioGroupPrimitive.Root>) {
  return (
    <RadioGroupPrimitive.Root
      data-slot="radio-group"
      className={cn('grid gap-3', className)}
      {...props}
    />
  )
}

function RadioGroupItem({
  className,
  ...props
}: React.ComponentProps<typeof RadioGroupPrimitive.Item>) {
  return (
    <RadioGroupPrimitive.Item
      data-slot="radio-group-item"
      className={cn(
        'border-input text-primary focus-visible:border-ring focus-visible:ring-ring/50 aria-invalid:ring-destructive/20 dark:aria-invalid:ring-destructive/40 aria-invalid:border-destructive dark:bg-input/30 aspect-square size-4 shrink-0 rounded-full border shadow-xs transition-[color,box-shadow] outline-none focus-visible:ring-[3px] disabled:cursor-not-allowed disabled:opacity-50',
        className,
      )}
      {...props}
    >
      <RadioGroupPrimitive.Indicator
        data-slot="radio-group-indicator"
        className="relative flex items-center justify-center"
      >
        <CircleIcon className="fill-primary absolute top-1/2 left-1/2 size-2 -translate-x-1/2 -translate-y-1/2" />
      </RadioGroupPrimitive.Indicator>
    </RadioGroupPrimitive.Item>
  )
}

export { RadioGroup, RadioGroupItem }
\end{DeskFlowCode}

\phantomsection
\addcontentsline{toc}{subsubsection}{deskflow\-frontend/components/ui/resizable.tsx}
\DeskFlowListingCaption{deskflow-frontend/components/ui/resizable.tsx}{lst:fd-ui-resizable}
\begin{DeskFlowCode}
'use client'

import * as React from 'react'
import { GripVerticalIcon } from 'lucide-react'
import * as ResizablePrimitive from 'react-resizable-panels'

import { cn } from '@/lib/utils'

function ResizablePanelGroup({
  className,
  ...props
}: React.ComponentProps<typeof ResizablePrimitive.PanelGroup>) {
  return (
    <ResizablePrimitive.PanelGroup
      data-slot="resizable-panel-group"
      className={cn(
        'flex h-full w-full data-[panel-group-direction=vertical]:flex-col',
        className,
      )}
      {...props}
    />
  )
}

function ResizablePanel({
  ...props
}: React.ComponentProps<typeof ResizablePrimitive.Panel>) {
  return <ResizablePrimitive.Panel data-slot="resizable-panel" {...props} />
}

function ResizableHandle({
  withHandle,
  className,
  ...props
}: React.ComponentProps<typeof ResizablePrimitive.PanelResizeHandle> & {
  withHandle?: boolean
}) {
  return (
    <ResizablePrimitive.PanelResizeHandle
      data-slot="resizable-handle"
      className={cn(
        'bg-border focus-visible:ring-ring relative flex w-px items-center justify-center after:absolute after:inset-y-0 after:left-1/2 after:w-1 after:-translate-x-1/2 focus-visible:ring-1 focus-visible:ring-offset-1 focus-visible:outline-hidden data-[panel-group-direction=vertical]:h-px data-[panel-group-direction=vertical]:w-full data-[panel-group-direction=vertical]:after:left-0 data-[panel-group-direction=vertical]:after:h-1 data-[panel-group-direction=vertical]:after:w-full data-[panel-group-direction=vertical]:after:translate-x-0 data-[panel-group-direction=vertical]:after:-translate-y-1/2 [&[data-panel-group-direction=vertical]>div]:rotate-90',
        className,
      )}
      {...props}
    >
      {withHandle && (
        <div className="bg-border z-10 flex h-4 w-3 items-center justify-center rounded-xs border">
          <GripVerticalIcon className="size-2.5" />
        </div>
      )}
    </ResizablePrimitive.PanelResizeHandle>
  )
}

export { ResizablePanelGroup, ResizablePanel, ResizableHandle }
\end{DeskFlowCode}

\phantomsection
\addcontentsline{toc}{subsubsection}{deskflow\-frontend/components/ui/scroll\-area.tsx}
\DeskFlowListingCaption{deskflow-frontend/components/ui/scroll-area.tsx}{lst:fd-ui-scroll-area}
\begin{DeskFlowCode}
'use client'

import * as React from 'react'
import * as ScrollAreaPrimitive from '@radix-ui/react-scroll-area'

import { cn } from '@/lib/utils'

function ScrollArea({
  className,
  children,
  ...props
}: React.ComponentProps<typeof ScrollAreaPrimitive.Root>) {
  return (
    <ScrollAreaPrimitive.Root
      data-slot="scroll-area"
      className={cn('relative', className)}
      {...props}
    >
      <ScrollAreaPrimitive.Viewport
        data-slot="scroll-area-viewport"
        className="focus-visible:ring-ring/50 size-full rounded-[inherit] transition-[color,box-shadow] outline-none focus-visible:ring-[3px] focus-visible:outline-1"
      >
        {children}
      </ScrollAreaPrimitive.Viewport>
      <ScrollBar />
      <ScrollAreaPrimitive.Corner />
    </ScrollAreaPrimitive.Root>
  )
}

function ScrollBar({
  className,
  orientation = 'vertical',
  ...props
}: React.ComponentProps<typeof ScrollAreaPrimitive.ScrollAreaScrollbar>) {
  return (
    <ScrollAreaPrimitive.ScrollAreaScrollbar
      data-slot="scroll-area-scrollbar"
      orientation={orientation}
      className={cn(
        'flex touch-none p-px transition-colors select-none',
        orientation === 'vertical' &&
          'h-full w-2.5 border-l border-l-transparent',
        orientation === 'horizontal' &&
          'h-2.5 flex-col border-t border-t-transparent',
        className,
      )}
      {...props}
    >
      <ScrollAreaPrimitive.ScrollAreaThumb
        data-slot="scroll-area-thumb"
        className="bg-border relative flex-1 rounded-full"
      />
    </ScrollAreaPrimitive.ScrollAreaScrollbar>
  )
}

export { ScrollArea, ScrollBar }
\end{DeskFlowCode}

\phantomsection
\addcontentsline{toc}{subsubsection}{deskflow\-frontend/components/ui/select.tsx}
\DeskFlowListingCaption{deskflow-frontend/components/ui/select.tsx}{lst:fd-ui-select}
\begin{DeskFlowCode}
'use client'

import * as React from 'react'
import * as SelectPrimitive from '@radix-ui/react-select'
import { CheckIcon, ChevronDownIcon, ChevronUpIcon } from 'lucide-react'

import { cn } from '@/lib/utils'

function Select({
  ...props
}: React.ComponentProps<typeof SelectPrimitive.Root>) {
  return <SelectPrimitive.Root data-slot="select" {...props} />
}

function SelectGroup({
  ...props
}: React.ComponentProps<typeof SelectPrimitive.Group>) {
  return <SelectPrimitive.Group data-slot="select-group" {...props} />
}

function SelectValue({
  ...props
}: React.ComponentProps<typeof SelectPrimitive.Value>) {
  return <SelectPrimitive.Value data-slot="select-value" {...props} />
}

function SelectTrigger({
  className,
  size = 'default',
  children,
  ...props
}: React.ComponentProps<typeof SelectPrimitive.Trigger> & {
  size?: 'sm' | 'default'
}) {
  return (
    <SelectPrimitive.Trigger
      data-slot="select-trigger"
      data-size={size}
      className={cn(
        "border-input data-[placeholder]:text-muted-foreground [&_svg:not([class*='text-'])]:text-muted-foreground focus-visible:border-ring focus-visible:ring-ring/50 aria-invalid:ring-destructive/20 dark:aria-invalid:ring-destructive/40 aria-invalid:border-destructive dark:bg-input/30 dark:hover:bg-input/50 flex w-fit items-center justify-between gap-2 rounded-md border bg-transparent px-3 py-2 text-sm whitespace-nowrap shadow-xs transition-[color,box-shadow] outline-none focus-visible:ring-[3px] disabled:cursor-not-allowed disabled:opacity-50 data-[size=default]:h-9 data-[size=sm]:h-8 *:data-[slot=select-value]:line-clamp-1 *:data-[slot=select-value]:flex *:data-[slot=select-value]:items-center *:data-[slot=select-value]:gap-2 [&_svg]:pointer-events-none [&_svg]:shrink-0 [&_svg:not([class*='size-'])]:size-4",
        className,
      )}
      {...props}
    >
      {children}
      <SelectPrimitive.Icon asChild>
        <ChevronDownIcon className="size-4 opacity-50" />
      </SelectPrimitive.Icon>
    </SelectPrimitive.Trigger>
  )
}

function SelectContent({
  className,
  children,
  position = 'popper',
  ...props
}: React.ComponentProps<typeof SelectPrimitive.Content>) {
  return (
    <SelectPrimitive.Portal>
      <SelectPrimitive.Content
        data-slot="select-content"
        className={cn(
          'bg-popover text-popover-foreground data-[state=open]:animate-in data-[state=closed]:animate-out data-[state=closed]:fade-out-0 data-[state=open]:fade-in-0 data-[state=closed]:zoom-out-95 data-[state=open]:zoom-in-95 data-[side=bottom]:slide-in-from-top-2 data-[side=left]:slide-in-from-right-2 data-[side=right]:slide-in-from-left-2 data-[side=top]:slide-in-from-bottom-2 relative z-50 max-h-(-radix-select-content-available-height) min-w-[8rem] origin-(-radix-select-content-transform-origin) overflow-x-hidden overflow-y-auto rounded-md border shadow-md',
          position === 'popper' &&
            'data-[side=bottom]:translate-y-1 data-[side=left]:-translate-x-1 data-[side=right]:translate-x-1 data-[side=top]:-translate-y-1',
          className,
        )}
        position={position}
        {...props}
      >
        <SelectScrollUpButton />
        <SelectPrimitive.Viewport
          className={cn(
            'p-1',
            position === 'popper' &&
              'h-[var(-radix-select-trigger-height)] w-full min-w-[var(-radix-select-trigger-width)] scroll-my-1',
          )}
        >
          {children}
        </SelectPrimitive.Viewport>
        <SelectScrollDownButton />
      </SelectPrimitive.Content>
    </SelectPrimitive.Portal>
  )
}

function SelectLabel({
  className,
  ...props
}: React.ComponentProps<typeof SelectPrimitive.Label>) {
  return (
    <SelectPrimitive.Label
      data-slot="select-label"
      className={cn('text-muted-foreground px-2 py-1.5 text-xs', className)}
      {...props}
    />
  )
}

function SelectItem({
  className,
  children,
  ...props
}: React.ComponentProps<typeof SelectPrimitive.Item>) {
  return (
    <SelectPrimitive.Item
      data-slot="select-item"
      className={cn(
        "focus:bg-accent focus:text-accent-foreground [&_svg:not([class*='text-'])]:text-muted-foreground relative flex w-full cursor-default items-center gap-2 rounded-sm py-1.5 pr-8 pl-2 text-sm outline-hidden select-none data-[disabled]:pointer-events-none data-[disabled]:opacity-50 [&_svg]:pointer-events-none [&_svg]:shrink-0 [&_svg:not([class*='size-'])]:size-4 *:[span]:last:flex *:[span]:last:items-center *:[span]:last:gap-2",
        className,
      )}
      {...props}
    >
      <span className="absolute right-2 flex size-3.5 items-center justify-center">
        <SelectPrimitive.ItemIndicator>
          <CheckIcon className="size-4" />
        </SelectPrimitive.ItemIndicator>
      </span>
      <SelectPrimitive.ItemText>{children}</SelectPrimitive.ItemText>
    </SelectPrimitive.Item>
  )
}

function SelectSeparator({
  className,
  ...props
}: React.ComponentProps<typeof SelectPrimitive.Separator>) {
  return (
    <SelectPrimitive.Separator
      data-slot="select-separator"
      className={cn('bg-border pointer-events-none -mx-1 my-1 h-px', className)}
      {...props}
    />
  )
}

function SelectScrollUpButton({
  className,
  ...props
}: React.ComponentProps<typeof SelectPrimitive.ScrollUpButton>) {
  return (
    <SelectPrimitive.ScrollUpButton
      data-slot="select-scroll-up-button"
      className={cn(
        'flex cursor-default items-center justify-center py-1',
        className,
      )}
      {...props}
    >
      <ChevronUpIcon className="size-4" />
    </SelectPrimitive.ScrollUpButton>
  )
}

function SelectScrollDownButton({
  className,
  ...props
}: React.ComponentProps<typeof SelectPrimitive.ScrollDownButton>) {
  return (
    <SelectPrimitive.ScrollDownButton
      data-slot="select-scroll-down-button"
      className={cn(
        'flex cursor-default items-center justify-center py-1',
        className,
      )}
      {...props}
    >
      <ChevronDownIcon className="size-4" />
    </SelectPrimitive.ScrollDownButton>
  )
}

export {
  Select,
  SelectContent,
  SelectGroup,
  SelectItem,
  SelectLabel,
  SelectScrollDownButton,
  SelectScrollUpButton,
  SelectSeparator,
  SelectTrigger,
  SelectValue,
}
\end{DeskFlowCode}

\phantomsection
\addcontentsline{toc}{subsubsection}{deskflow\-frontend/components/ui/separator.tsx}
\DeskFlowListingCaption{deskflow-frontend/components/ui/separator.tsx}{lst:fd-ui-separator}
\begin{DeskFlowCode}
'use client'

import * as React from 'react'
import * as SeparatorPrimitive from '@radix-ui/react-separator'

import { cn } from '@/lib/utils'

function Separator({
  className,
  orientation = 'horizontal',
  decorative = true,
  ...props
}: React.ComponentProps<typeof SeparatorPrimitive.Root>) {
  return (
    <SeparatorPrimitive.Root
      data-slot="separator"
      decorative={decorative}
      orientation={orientation}
      className={cn(
        'bg-border shrink-0 data-[orientation=horizontal]:h-px data-[orientation=horizontal]:w-full data-[orientation=vertical]:h-full data-[orientation=vertical]:w-px',
        className,
      )}
      {...props}
    />
  )
}

export { Separator }
\end{DeskFlowCode}

\phantomsection
\addcontentsline{toc}{subsubsection}{deskflow\-frontend/components/ui/sheet.tsx}
\DeskFlowListingCaption{deskflow-frontend/components/ui/sheet.tsx}{lst:fd-ui-sheet}
\begin{DeskFlowCode}
'use client'

import * as React from 'react'
import * as SheetPrimitive from '@radix-ui/react-dialog'
import { XIcon } from 'lucide-react'

import { cn } from '@/lib/utils'

function Sheet({ ...props }: React.ComponentProps<typeof SheetPrimitive.Root>) {
  return <SheetPrimitive.Root data-slot="sheet" {...props} />
}

function SheetTrigger({
  ...props
}: React.ComponentProps<typeof SheetPrimitive.Trigger>) {
  return <SheetPrimitive.Trigger data-slot="sheet-trigger" {...props} />
}

function SheetClose({
  ...props
}: React.ComponentProps<typeof SheetPrimitive.Close>) {
  return <SheetPrimitive.Close data-slot="sheet-close" {...props} />
}

function SheetPortal({
  ...props
}: React.ComponentProps<typeof SheetPrimitive.Portal>) {
  return <SheetPrimitive.Portal data-slot="sheet-portal" {...props} />
}

function SheetOverlay({
  className,
  ...props
}: React.ComponentProps<typeof SheetPrimitive.Overlay>) {
  return (
    <SheetPrimitive.Overlay
      data-slot="sheet-overlay"
      className={cn(
        'data-[state=open]:animate-in data-[state=closed]:animate-out data-[state=closed]:fade-out-0 data-[state=open]:fade-in-0 fixed inset-0 z-50 bg-black/50',
        className,
      )}
      {...props}
    />
  )
}

function SheetContent({
  className,
  children,
  side = 'right',
  ...props
}: React.ComponentProps<typeof SheetPrimitive.Content> & {
  side?: 'top' | 'right' | 'bottom' | 'left'
}) {
  return (
    <SheetPortal>
      <SheetOverlay />
      <SheetPrimitive.Content
        data-slot="sheet-content"
        className={cn(
          'bg-background data-[state=open]:animate-in data-[state=closed]:animate-out fixed z-50 flex flex-col gap-4 shadow-lg transition ease-in-out data-[state=closed]:duration-300 data-[state=open]:duration-500',
          side === 'right' &&
            'data-[state=closed]:slide-out-to-right data-[state=open]:slide-in-from-right inset-y-0 right-0 h-full w-3/4 border-l sm:max-w-sm',
          side === 'left' &&
            'data-[state=closed]:slide-out-to-left data-[state=open]:slide-in-from-left inset-y-0 left-0 h-full w-3/4 border-r sm:max-w-sm',
          side === 'top' &&
            'data-[state=closed]:slide-out-to-top data-[state=open]:slide-in-from-top inset-x-0 top-0 h-auto border-b',
          side === 'bottom' &&
            'data-[state=closed]:slide-out-to-bottom data-[state=open]:slide-in-from-bottom inset-x-0 bottom-0 h-auto border-t',
          className,
        )}
        {...props}
      >
        {children}
        <SheetPrimitive.Close className="ring-offset-background focus:ring-ring data-[state=open]:bg-secondary absolute top-4 right-4 rounded-xs opacity-70 transition-opacity hover:opacity-100 focus:ring-2 focus:ring-offset-2 focus:outline-hidden disabled:pointer-events-none">
          <XIcon className="size-4" />
          <span className="sr-only">Close</span>
        </SheetPrimitive.Close>
      </SheetPrimitive.Content>
    </SheetPortal>
  )
}

function SheetHeader({ className, ...props }: React.ComponentProps<'div'>) {
  return (
    <div
      data-slot="sheet-header"
      className={cn('flex flex-col gap-1.5 p-4', className)}
      {...props}
    />
  )
}

function SheetFooter({ className, ...props }: React.ComponentProps<'div'>) {
  return (
    <div
      data-slot="sheet-footer"
      className={cn('mt-auto flex flex-col gap-2 p-4', className)}
      {...props}
    />
  )
}

function SheetTitle({
  className,
  ...props
}: React.ComponentProps<typeof SheetPrimitive.Title>) {
  return (
    <SheetPrimitive.Title
      data-slot="sheet-title"
      className={cn('text-foreground font-semibold', className)}
      {...props}
    />
  )
}

function SheetDescription({
  className,
  ...props
}: React.ComponentProps<typeof SheetPrimitive.Description>) {
  return (
    <SheetPrimitive.Description
      data-slot="sheet-description"
      className={cn('text-muted-foreground text-sm', className)}
      {...props}
    />
  )
}

export {
  Sheet,
  SheetTrigger,
  SheetClose,
  SheetContent,
  SheetHeader,
  SheetFooter,
  SheetTitle,
  SheetDescription,
}
\end{DeskFlowCode}

\phantomsection
\addcontentsline{toc}{subsubsection}{deskflow\-frontend/components/ui/sidebar.tsx}
\DeskFlowListingCaption{deskflow-frontend/components/ui/sidebar.tsx}{lst:fd-ui-sidebar}
\begin{DeskFlowCode}
'use client'

import * as React from 'react'
import { Slot } from '@radix-ui/react-slot'
import { cva, VariantProps } from 'class-variance-authority'
import { PanelLeftIcon } from 'lucide-react'

import { useIsMobile } from '@/hooks/use-mobile'
import { cn } from '@/lib/utils'
import { Button } from '@/components/ui/button'
import { Input } from '@/components/ui/input'
import { Separator } from '@/components/ui/separator'
import {
  Sheet,
  SheetContent,
  SheetDescription,
  SheetHeader,
  SheetTitle,
} from '@/components/ui/sheet'
import { Skeleton } from '@/components/ui/skeleton'
import {
  Tooltip,
  TooltipContent,
  TooltipProvider,
  TooltipTrigger,
} from '@/components/ui/tooltip'

const SIDEBAR_COOKIE_NAME = 'sidebar_state'
const SIDEBAR_COOKIE_MAX_AGE = 60 * 60 * 24 * 7
const SIDEBAR_WIDTH = '16rem'
const SIDEBAR_WIDTH_MOBILE = '18rem'
const SIDEBAR_WIDTH_ICON = '3rem'
const SIDEBAR_KEYBOARD_SHORTCUT = 'b'

type SidebarContextProps = {
  state: 'expanded' | 'collapsed'
  open: boolean
  setOpen: (open: boolean) => void
  openMobile: boolean
  setOpenMobile: (open: boolean) => void
  isMobile: boolean
  toggleSidebar: () => void
}

const SidebarContext = React.createContext<SidebarContextProps | null>(null)

function useSidebar() {
  const context = React.useContext(SidebarContext)
  if (!context) {
    throw new Error('useSidebar must be used within a SidebarProvider.')
  }

  return context
}

function SidebarProvider({
  defaultOpen = true,
  open: openProp,
  onOpenChange: setOpenProp,
  className,
  style,
  children,
  ...props
}: React.ComponentProps<'div'> & {
  defaultOpen?: boolean
  open?: boolean
  onOpenChange?: (open: boolean) => void
}) {
  const isMobile = useIsMobile()
  const [openMobile, setOpenMobile] = React.useState(false)

  // This is the internal state of the sidebar.
  // We use openProp and setOpenProp for control from outside the component.
  const [_open, _setOpen] = React.useState(defaultOpen)
  const open = openProp ?? _open
  const setOpen = React.useCallback(
    (value: boolean | ((value: boolean) => boolean)) => {
      const openState = typeof value === 'function' ? value(open) : value
      if (setOpenProp) {
        setOpenProp(openState)
      } else {
        _setOpen(openState)
      }

      // This sets the cookie to keep the sidebar state.
      document.cookie = `${SIDEBAR_COOKIE_NAME}=${openState}; path=/; max-age=${SIDEBAR_COOKIE_MAX_AGE}`
    },
    [setOpenProp, open],
  )

  // Helper to toggle the sidebar.
  const toggleSidebar = React.useCallback(() => {
    return isMobile ? setOpenMobile((open) => !open) : setOpen((open) => !open)
  }, [isMobile, setOpen, setOpenMobile])

  // Adds a keyboard shortcut to toggle the sidebar.
  React.useEffect(() => {
    const handleKeyDown = (event: KeyboardEvent) => {
      if (
        event.key === SIDEBAR_KEYBOARD_SHORTCUT &&
        (event.metaKey || event.ctrlKey)
      ) {
        event.preventDefault()
        toggleSidebar()
      }
    }

    window.addEventListener('keydown', handleKeyDown)
    return () => window.removeEventListener('keydown', handleKeyDown)
  }, [toggleSidebar])

  // We add a state so that we can do data-state="expanded" or "collapsed".
  // This makes it easier to style the sidebar with Tailwind classes.
  const state = open ? 'expanded' : 'collapsed'

  const contextValue = React.useMemo<SidebarContextProps>(
    () => ({
      state,
      open,
      setOpen,
      isMobile,
      openMobile,
      setOpenMobile,
      toggleSidebar,
    }),
    [state, open, setOpen, isMobile, openMobile, setOpenMobile, toggleSidebar],
  )

  return (
    <SidebarContext.Provider value={contextValue}>
      <TooltipProvider delayDuration={0}>
        <div
          data-slot="sidebar-wrapper"
          style={
            {
              '-sidebar-width': SIDEBAR_WIDTH,
              '-sidebar-width-icon': SIDEBAR_WIDTH_ICON,
              ...style,
            } as React.CSSProperties
          }
          className={cn(
            'group/sidebar-wrapper has-data-[variant=inset]:bg-sidebar flex min-h-svh w-full',
            className,
          )}
          {...props}
        >
          {children}
        </div>
      </TooltipProvider>
    </SidebarContext.Provider>
  )
}

function Sidebar({
  side = 'left',
  variant = 'sidebar',
  collapsible = 'offcanvas',
  className,
  children,
  ...props
}: React.ComponentProps<'div'> & {
  side?: 'left' | 'right'
  variant?: 'sidebar' | 'floating' | 'inset'
  collapsible?: 'offcanvas' | 'icon' | 'none'
}) {
  const { isMobile, state, openMobile, setOpenMobile } = useSidebar()

  if (collapsible === 'none') {
    return (
      <div
        data-slot="sidebar"
        className={cn(
          'bg-sidebar text-sidebar-foreground flex h-full w-(-sidebar-width) flex-col',
          className,
        )}
        {...props}
      >
        {children}
      </div>
    )
  }

  if (isMobile) {
    return (
      <Sheet open={openMobile} onOpenChange={setOpenMobile} {...props}>
        <SheetContent
          data-sidebar="sidebar"
          data-slot="sidebar"
          data-mobile="true"
          className="bg-sidebar text-sidebar-foreground w-(-sidebar-width) p-0 [&>button]:hidden"
          style={
            {
              '-sidebar-width': SIDEBAR_WIDTH_MOBILE,
            } as React.CSSProperties
          }
          side={side}
        >
          <SheetHeader className="sr-only">
            <SheetTitle>Sidebar</SheetTitle>
            <SheetDescription>Displays the mobile sidebar.</SheetDescription>
          </SheetHeader>
          <div className="flex h-full w-full flex-col">{children}</div>
        </SheetContent>
      </Sheet>
    )
  }

  return (
    <div
      className="group peer text-sidebar-foreground hidden md:block"
      data-state={state}
      data-collapsible={state === 'collapsed' ? collapsible : ''}
      data-variant={variant}
      data-side={side}
      data-slot="sidebar"
    >
      {/* This is what handles the sidebar gap on desktop */}
      <div
        data-slot="sidebar-gap"
        className={cn(
          'relative w-(-sidebar-width) bg-transparent transition-[width] duration-200 ease-linear',
          'group-data-[collapsible=offcanvas]:w-0',
          'group-data-[side=right]:rotate-180',
          variant === 'floating' || variant === 'inset'
            ? 'group-data-[collapsible=icon]:w-[calc(var(-sidebar-width-icon)+(-spacing(4)))]'
            : 'group-data-[collapsible=icon]:w-(-sidebar-width-icon)',
        )}
      />
      <div
        data-slot="sidebar-container"
        className={cn(
          'fixed inset-y-0 z-10 hidden h-svh w-(-sidebar-width) transition-[left,right,width] duration-200 ease-linear md:flex',
          side === 'left'
            ? 'left-0 group-data-[collapsible=offcanvas]:left-[calc(var(-sidebar-width)*-1)]'
            : 'right-0 group-data-[collapsible=offcanvas]:right-[calc(var(-sidebar-width)*-1)]',
          // Adjust the padding for floating and inset variants.
          variant === 'floating' || variant === 'inset'
            ? 'p-2 group-data-[collapsible=icon]:w-[calc(var(-sidebar-width-icon)+(-spacing(4))+2px)]'
            : 'group-data-[collapsible=icon]:w-(-sidebar-width-icon) group-data-[side=left]:border-r group-data-[side=right]:border-l',
          className,
        )}
        {...props}
      >
        <div
          data-sidebar="sidebar"
          data-slot="sidebar-inner"
          className="bg-sidebar group-data-[variant=floating]:border-sidebar-border flex h-full w-full flex-col group-data-[variant=floating]:rounded-lg group-data-[variant=floating]:border group-data-[variant=floating]:shadow-sm"
        >
          {children}
        </div>
      </div>
    </div>
  )
}

function SidebarTrigger({
  className,
  onClick,
  ...props
}: React.ComponentProps<typeof Button>) {
  const { toggleSidebar } = useSidebar()

  return (
    <Button
      data-sidebar="trigger"
      data-slot="sidebar-trigger"
      variant="ghost"
      size="icon"
      className={cn('size-7', className)}
      onClick={(event) => {
        onClick?.(event)
        toggleSidebar()
      }}
      {...props}
    >
      <PanelLeftIcon />
      <span className="sr-only">Toggle Sidebar</span>
    </Button>
  )
}

function SidebarRail({ className, ...props }: React.ComponentProps<'button'>) {
  const { toggleSidebar } = useSidebar()

  return (
    <button
      data-sidebar="rail"
      data-slot="sidebar-rail"
      aria-label="Toggle Sidebar"
      tabIndex={-1}
      onClick={toggleSidebar}
      title="Toggle Sidebar"
      className={cn(
        'hover:after:bg-sidebar-border absolute inset-y-0 z-20 hidden w-4 -translate-x-1/2 transition-all ease-linear group-data-[side=left]:-right-4 group-data-[side=right]:left-0 after:absolute after:inset-y-0 after:left-1/2 after:w-[2px] sm:flex',
        'in-data-[side=left]:cursor-w-resize in-data-[side=right]:cursor-e-resize',
        '[[data-side=left][data-state=collapsed]_&]:cursor-e-resize [[data-side=right][data-state=collapsed]_&]:cursor-w-resize',
        'hover:group-data-[collapsible=offcanvas]:bg-sidebar group-data-[collapsible=offcanvas]:translate-x-0 group-data-[collapsible=offcanvas]:after:left-full',
        '[[data-side=left][data-collapsible=offcanvas]_&]:-right-2',
        '[[data-side=right][data-collapsible=offcanvas]_&]:-left-2',
        className,
      )}
      {...props}
    />
  )
}

function SidebarInset({ className, ...props }: React.ComponentProps<'main'>) {
  return (
    <main
      data-slot="sidebar-inset"
      className={cn(
        'bg-background relative flex w-full flex-1 flex-col',
        'md:peer-data-[variant=inset]:m-2 md:peer-data-[variant=inset]:ml-0 md:peer-data-[variant=inset]:rounded-xl md:peer-data-[variant=inset]:shadow-sm md:peer-data-[variant=inset]:peer-data-[state=collapsed]:ml-2',
        className,
      )}
      {...props}
    />
  )
}

function SidebarInput({
  className,
  ...props
}: React.ComponentProps<typeof Input>) {
  return (
    <Input
      data-slot="sidebar-input"
      data-sidebar="input"
      className={cn('bg-background h-8 w-full shadow-none', className)}
      {...props}
    />
  )
}

function SidebarHeader({ className, ...props }: React.ComponentProps<'div'>) {
  return (
    <div
      data-slot="sidebar-header"
      data-sidebar="header"
      className={cn('flex flex-col gap-2 p-2', className)}
      {...props}
    />
  )
}

function SidebarFooter({ className, ...props }: React.ComponentProps<'div'>) {
  return (
    <div
      data-slot="sidebar-footer"
      data-sidebar="footer"
      className={cn('flex flex-col gap-2 p-2', className)}
      {...props}
    />
  )
}

function SidebarSeparator({
  className,
  ...props
}: React.ComponentProps<typeof Separator>) {
  return (
    <Separator
      data-slot="sidebar-separator"
      data-sidebar="separator"
      className={cn('bg-sidebar-border mx-2 w-auto', className)}
      {...props}
    />
  )
}

function SidebarContent({ className, ...props }: React.ComponentProps<'div'>) {
  return (
    <div
      data-slot="sidebar-content"
      data-sidebar="content"
      className={cn(
        'flex min-h-0 flex-1 flex-col gap-2 overflow-auto group-data-[collapsible=icon]:overflow-hidden',
        className,
      )}
      {...props}
    />
  )
}

function SidebarGroup({ className, ...props }: React.ComponentProps<'div'>) {
  return (
    <div
      data-slot="sidebar-group"
      data-sidebar="group"
      className={cn('relative flex w-full min-w-0 flex-col p-2', className)}
      {...props}
    />
  )
}

function SidebarGroupLabel({
  className,
  asChild = false,
  ...props
}: React.ComponentProps<'div'> & { asChild?: boolean }) {
  const Comp = asChild ? Slot : 'div'

  return (
    <Comp
      data-slot="sidebar-group-label"
      data-sidebar="group-label"
      className={cn(
        'text-sidebar-foreground/70 ring-sidebar-ring flex h-8 shrink-0 items-center rounded-md px-2 text-xs font-medium outline-hidden transition-[margin,opacity] duration-200 ease-linear focus-visible:ring-2 [&>svg]:size-4 [&>svg]:shrink-0',
        'group-data-[collapsible=icon]:-mt-8 group-data-[collapsible=icon]:opacity-0',
        className,
      )}
      {...props}
    />
  )
}

function SidebarGroupAction({
  className,
  asChild = false,
  ...props
}: React.ComponentProps<'button'> & { asChild?: boolean }) {
  const Comp = asChild ? Slot : 'button'

  return (
    <Comp
      data-slot="sidebar-group-action"
      data-sidebar="group-action"
      className={cn(
        'text-sidebar-foreground ring-sidebar-ring hover:bg-sidebar-accent hover:text-sidebar-accent-foreground absolute top-3.5 right-3 flex aspect-square w-5 items-center justify-center rounded-md p-0 outline-hidden transition-transform focus-visible:ring-2 [&>svg]:size-4 [&>svg]:shrink-0',
        // Increases the hit area of the button on mobile.
        'after:absolute after:-inset-2 md:after:hidden',
        'group-data-[collapsible=icon]:hidden',
        className,
      )}
      {...props}
    />
  )
}

function SidebarGroupContent({
  className,
  ...props
}: React.ComponentProps<'div'>) {
  return (
    <div
      data-slot="sidebar-group-content"
      data-sidebar="group-content"
      className={cn('w-full text-sm', className)}
      {...props}
    />
  )
}

function SidebarMenu({ className, ...props }: React.ComponentProps<'ul'>) {
  return (
    <ul
      data-slot="sidebar-menu"
      data-sidebar="menu"
      className={cn('flex w-full min-w-0 flex-col gap-1', className)}
      {...props}
    />
  )
}

function SidebarMenuItem({ className, ...props }: React.ComponentProps<'li'>) {
  return (
    <li
      data-slot="sidebar-menu-item"
      data-sidebar="menu-item"
      className={cn('group/menu-item relative', className)}
      {...props}
    />
  )
}

const sidebarMenuButtonVariants = cva(
  'peer/menu-button flex w-full items-center gap-2 overflow-hidden rounded-md p-2 text-left text-sm outline-hidden ring-sidebar-ring transition-[width,height,padding] hover:bg-sidebar-accent hover:text-sidebar-accent-foreground focus-visible:ring-2 active:bg-sidebar-accent active:text-sidebar-accent-foreground disabled:pointer-events-none disabled:opacity-50 group-has-data-[sidebar=menu-action]/menu-item:pr-8 aria-disabled:pointer-events-none aria-disabled:opacity-50 data-[active=true]:bg-sidebar-accent data-[active=true]:font-medium data-[active=true]:text-sidebar-accent-foreground data-[state=open]:hover:bg-sidebar-accent data-[state=open]:hover:text-sidebar-accent-foreground group-data-[collapsible=icon]:size-8! group-data-[collapsible=icon]:p-2! [&>span:last-child]:truncate [&>svg]:size-4 [&>svg]:shrink-0',
  {
    variants: {
      variant: {
        default: 'hover:bg-sidebar-accent hover:text-sidebar-accent-foreground',
        outline:
          'bg-background shadow-[0_0_0_1px_hsl(var(-sidebar-border))] hover:bg-sidebar-accent hover:text-sidebar-accent-foreground hover:shadow-[0_0_0_1px_hsl(var(-sidebar-accent))]',
      },
      size: {
        default: 'h-8 text-sm',
        sm: 'h-7 text-xs',
        lg: 'h-12 text-sm group-data-[collapsible=icon]:p-0!',
      },
    },
    defaultVariants: {
      variant: 'default',
      size: 'default',
    },
  },
)

function SidebarMenuButton({
  asChild = false,
  isActive = false,
  variant = 'default',
  size = 'default',
  tooltip,
  className,
  ...props
}: React.ComponentProps<'button'> & {
  asChild?: boolean
  isActive?: boolean
  tooltip?: string | React.ComponentProps<typeof TooltipContent>
} & VariantProps<typeof sidebarMenuButtonVariants>) {
  const Comp = asChild ? Slot : 'button'
  const { isMobile, state } = useSidebar()

  const button = (
    <Comp
      data-slot="sidebar-menu-button"
      data-sidebar="menu-button"
      data-size={size}
      data-active={isActive}
      className={cn(sidebarMenuButtonVariants({ variant, size }), className)}
      {...props}
    />
  )

  if (!tooltip) {
    return button
  }

  if (typeof tooltip === 'string') {
    tooltip = {
      children: tooltip,
    }
  }

  return (
    <Tooltip>
      <TooltipTrigger asChild>{button}</TooltipTrigger>
      <TooltipContent
        side="right"
        align="center"
        hidden={state !== 'collapsed' || isMobile}
        {...tooltip}
      />
    </Tooltip>
  )
}

function SidebarMenuAction({
  className,
  asChild = false,
  showOnHover = false,
  ...props
}: React.ComponentProps<'button'> & {
  asChild?: boolean
  showOnHover?: boolean
}) {
  const Comp = asChild ? Slot : 'button'

  return (
    <Comp
      data-slot="sidebar-menu-action"
      data-sidebar="menu-action"
      className={cn(
        'text-sidebar-foreground ring-sidebar-ring hover:bg-sidebar-accent hover:text-sidebar-accent-foreground peer-hover/menu-button:text-sidebar-accent-foreground absolute top-1.5 right-1 flex aspect-square w-5 items-center justify-center rounded-md p-0 outline-hidden transition-transform focus-visible:ring-2 [&>svg]:size-4 [&>svg]:shrink-0',
        // Increases the hit area of the button on mobile.
        'after:absolute after:-inset-2 md:after:hidden',
        'peer-data-[size=sm]/menu-button:top-1',
        'peer-data-[size=default]/menu-button:top-1.5',
        'peer-data-[size=lg]/menu-button:top-2.5',
        'group-data-[collapsible=icon]:hidden',
        showOnHover &&
          'peer-data-[active=true]/menu-button:text-sidebar-accent-foreground group-focus-within/menu-item:opacity-100 group-hover/menu-item:opacity-100 data-[state=open]:opacity-100 md:opacity-0',
        className,
      )}
      {...props}
    />
  )
}

function SidebarMenuBadge({
  className,
  ...props
}: React.ComponentProps<'div'>) {
  return (
    <div
      data-slot="sidebar-menu-badge"
      data-sidebar="menu-badge"
      className={cn(
        'text-sidebar-foreground pointer-events-none absolute right-1 flex h-5 min-w-5 items-center justify-center rounded-md px-1 text-xs font-medium tabular-nums select-none',
        'peer-hover/menu-button:text-sidebar-accent-foreground peer-data-[active=true]/menu-button:text-sidebar-accent-foreground',
        'peer-data-[size=sm]/menu-button:top-1',
        'peer-data-[size=default]/menu-button:top-1.5',
        'peer-data-[size=lg]/menu-button:top-2.5',
        'group-data-[collapsible=icon]:hidden',
        className,
      )}
      {...props}
    />
  )
}

function SidebarMenuSkeleton({
  className,
  showIcon = false,
  ...props
}: React.ComponentProps<'div'> & {
  showIcon?: boolean
}) {
  // Random width between 50 to 90%.
  const width = React.useMemo(() => {
    return `${Math.floor(Math.random() * 40) + 50}%`
  }, [])

  return (
    <div
      data-slot="sidebar-menu-skeleton"
      data-sidebar="menu-skeleton"
      className={cn('flex h-8 items-center gap-2 rounded-md px-2', className)}
      {...props}
    >
      {showIcon && (
        <Skeleton
          className="size-4 rounded-md"
          data-sidebar="menu-skeleton-icon"
        />
      )}
      <Skeleton
        className="h-4 max-w-(-skeleton-width) flex-1"
        data-sidebar="menu-skeleton-text"
        style={
          {
            '-skeleton-width': width,
          } as React.CSSProperties
        }
      />
    </div>
  )
}

function SidebarMenuSub({ className, ...props }: React.ComponentProps<'ul'>) {
  return (
    <ul
      data-slot="sidebar-menu-sub"
      data-sidebar="menu-sub"
      className={cn(
        'border-sidebar-border mx-3.5 flex min-w-0 translate-x-px flex-col gap-1 border-l px-2.5 py-0.5',
        'group-data-[collapsible=icon]:hidden',
        className,
      )}
      {...props}
    />
  )
}

function SidebarMenuSubItem({
  className,
  ...props
}: React.ComponentProps<'li'>) {
  return (
    <li
      data-slot="sidebar-menu-sub-item"
      data-sidebar="menu-sub-item"
      className={cn('group/menu-sub-item relative', className)}
      {...props}
    />
  )
}

function SidebarMenuSubButton({
  asChild = false,
  size = 'md',
  isActive = false,
  className,
  ...props
}: React.ComponentProps<'a'> & {
  asChild?: boolean
  size?: 'sm' | 'md'
  isActive?: boolean
}) {
  const Comp = asChild ? Slot : 'a'

  return (
    <Comp
      data-slot="sidebar-menu-sub-button"
      data-sidebar="menu-sub-button"
      data-size={size}
      data-active={isActive}
      className={cn(
        'text-sidebar-foreground ring-sidebar-ring hover:bg-sidebar-accent hover:text-sidebar-accent-foreground active:bg-sidebar-accent active:text-sidebar-accent-foreground [&>svg]:text-sidebar-accent-foreground flex h-7 min-w-0 -translate-x-px items-center gap-2 overflow-hidden rounded-md px-2 outline-hidden focus-visible:ring-2 disabled:pointer-events-none disabled:opacity-50 aria-disabled:pointer-events-none aria-disabled:opacity-50 [&>span:last-child]:truncate [&>svg]:size-4 [&>svg]:shrink-0',
        'data-[active=true]:bg-sidebar-accent data-[active=true]:text-sidebar-accent-foreground',
        size === 'sm' && 'text-xs',
        size === 'md' && 'text-sm',
        'group-data-[collapsible=icon]:hidden',
        className,
      )}
      {...props}
    />
  )
}

export {
  Sidebar,
  SidebarContent,
  SidebarFooter,
  SidebarGroup,
  SidebarGroupAction,
  SidebarGroupContent,
  SidebarGroupLabel,
  SidebarHeader,
  SidebarInput,
  SidebarInset,
  SidebarMenu,
  SidebarMenuAction,
  SidebarMenuBadge,
  SidebarMenuButton,
  SidebarMenuItem,
  SidebarMenuSkeleton,
  SidebarMenuSub,
  SidebarMenuSubButton,
  SidebarMenuSubItem,
  SidebarProvider,
  SidebarRail,
  SidebarSeparator,
  SidebarTrigger,
  useSidebar,
}
\end{DeskFlowCode}

\phantomsection
\addcontentsline{toc}{subsubsection}{deskflow\-frontend/components/ui/skeleton.tsx}
\DeskFlowListingCaption{deskflow-frontend/components/ui/skeleton.tsx}{lst:fd-ui-skeleton}
\begin{DeskFlowCode}
import { cn } from '@/lib/utils'

function Skeleton({ className, ...props }: React.ComponentProps<'div'>) {
  return (
    <div
      data-slot="skeleton"
      className={cn('bg-accent animate-pulse rounded-md', className)}
      {...props}
    />
  )
}

export { Skeleton }
\end{DeskFlowCode}

\phantomsection
\addcontentsline{toc}{subsubsection}{deskflow\-frontend/components/ui/slider.tsx}
\DeskFlowListingCaption{deskflow-frontend/components/ui/slider.tsx}{lst:fd-ui-slider}
\begin{DeskFlowCode}
'use client'

import * as React from 'react'
import * as SliderPrimitive from '@radix-ui/react-slider'

import { cn } from '@/lib/utils'

function Slider({
  className,
  defaultValue,
  value,
  min = 0,
  max = 100,
  ...props
}: React.ComponentProps<typeof SliderPrimitive.Root>) {
  const _values = React.useMemo(
    () =>
      Array.isArray(value)
        ? value
        : Array.isArray(defaultValue)
          ? defaultValue
          : [min, max],
    [value, defaultValue, min, max],
  )

  return (
    <SliderPrimitive.Root
      data-slot="slider"
      defaultValue={defaultValue}
      value={value}
      min={min}
      max={max}
      className={cn(
        'relative flex w-full touch-none items-center select-none data-[disabled]:opacity-50 data-[orientation=vertical]:h-full data-[orientation=vertical]:min-h-44 data-[orientation=vertical]:w-auto data-[orientation=vertical]:flex-col',
        className,
      )}
      {...props}
    >
      <SliderPrimitive.Track
        data-slot="slider-track"
        className={
          'bg-muted relative grow overflow-hidden rounded-full data-[orientation=horizontal]:h-1.5 data-[orientation=horizontal]:w-full data-[orientation=vertical]:h-full data-[orientation=vertical]:w-1.5'
        }
      >
        <SliderPrimitive.Range
          data-slot="slider-range"
          className={
            'bg-primary absolute data-[orientation=horizontal]:h-full data-[orientation=vertical]:w-full'
          }
        />
      </SliderPrimitive.Track>
      {Array.from({ length: _values.length }, (_, index) => (
        <SliderPrimitive.Thumb
          data-slot="slider-thumb"
          key={index}
          className="border-primary ring-ring/50 block size-4 shrink-0 rounded-full border bg-white shadow-sm transition-[color,box-shadow] hover:ring-4 focus-visible:ring-4 focus-visible:outline-hidden disabled:pointer-events-none disabled:opacity-50"
        />
      ))}
    </SliderPrimitive.Root>
  )
}

export { Slider }
\end{DeskFlowCode}

\phantomsection
\addcontentsline{toc}{subsubsection}{deskflow\-frontend/components/ui/sonner.tsx}
\DeskFlowListingCaption{deskflow-frontend/components/ui/sonner.tsx}{lst:fd-ui-sonner}
\begin{DeskFlowCode}
'use client'

import { useTheme } from 'next-themes'
import { Toaster as Sonner, ToasterProps } from 'sonner'

const Toaster = ({ ...props }: ToasterProps) => {
  const { theme = 'system' } = useTheme()

  return (
    <Sonner
      theme={theme as ToasterProps['theme']}
      className="toaster group"
      style={
        {
          '-normal-bg': 'var(-popover)',
          '-normal-text': 'var(-popover-foreground)',
          '-normal-border': 'var(-border)',
        } as React.CSSProperties
      }
      {...props}
    />
  )
}

export { Toaster }
\end{DeskFlowCode}

\phantomsection
\addcontentsline{toc}{subsubsection}{deskflow\-frontend/components/ui/spinner.tsx}
\DeskFlowListingCaption{deskflow-frontend/components/ui/spinner.tsx}{lst:fd-ui-spinner}
\begin{DeskFlowCode}
import { Loader2Icon } from 'lucide-react'

import { cn } from '@/lib/utils'

function Spinner({ className, ...props }: React.ComponentProps<'svg'>) {
  return (
    <Loader2Icon
      role="status"
      aria-label="Loading"
      className={cn('size-4 animate-spin', className)}
      {...props}
    />
  )
}

export { Spinner }
\end{DeskFlowCode}

\phantomsection
\addcontentsline{toc}{subsubsection}{deskflow\-frontend/components/ui/switch.tsx}
\DeskFlowListingCaption{deskflow-frontend/components/ui/switch.tsx}{lst:fd-ui-switch}
\begin{DeskFlowCode}
'use client'

import * as React from 'react'
import * as SwitchPrimitive from '@radix-ui/react-switch'

import { cn } from '@/lib/utils'

function Switch({
  className,
  ...props
}: React.ComponentProps<typeof SwitchPrimitive.Root>) {
  return (
    <SwitchPrimitive.Root
      data-slot="switch"
      className={cn(
        'peer data-[state=checked]:bg-primary data-[state=unchecked]:bg-input focus-visible:border-ring focus-visible:ring-ring/50 dark:data-[state=unchecked]:bg-input/80 inline-flex h-[1.15rem] w-8 shrink-0 items-center rounded-full border border-transparent shadow-xs transition-all outline-none focus-visible:ring-[3px] disabled:cursor-not-allowed disabled:opacity-50',
        className,
      )}
      {...props}
    >
      <SwitchPrimitive.Thumb
        data-slot="switch-thumb"
        className={
          'bg-background dark:data-[state=unchecked]:bg-foreground dark:data-[state=checked]:bg-primary-foreground pointer-events-none block size-4 rounded-full ring-0 transition-transform data-[state=checked]:translate-x-[calc(100%-2px)] data-[state=unchecked]:translate-x-0'
        }
      />
    </SwitchPrimitive.Root>
  )
}

export { Switch }
\end{DeskFlowCode}

\phantomsection
\addcontentsline{toc}{subsubsection}{deskflow\-frontend/components/ui/table.tsx}
\DeskFlowListingCaption{deskflow-frontend/components/ui/table.tsx}{lst:fd-ui-table}
\begin{DeskFlowCode}
'use client'

import * as React from 'react'

import { cn } from '@/lib/utils'

function Table({ className, ...props }: React.ComponentProps<'table'>) {
  return (
    <div
      data-slot="table-container"
      className="relative w-full overflow-x-auto"
    >
      <table
        data-slot="table"
        className={cn('w-full caption-bottom text-sm', className)}
        {...props}
      />
    </div>
  )
}

function TableHeader({ className, ...props }: React.ComponentProps<'thead'>) {
  return (
    <thead
      data-slot="table-header"
      className={cn('[&_tr]:border-b', className)}
      {...props}
    />
  )
}

function TableBody({ className, ...props }: React.ComponentProps<'tbody'>) {
  return (
    <tbody
      data-slot="table-body"
      className={cn('[&_tr:last-child]:border-0', className)}
      {...props}
    />
  )
}

function TableFooter({ className, ...props }: React.ComponentProps<'tfoot'>) {
  return (
    <tfoot
      data-slot="table-footer"
      className={cn(
        'bg-muted/50 border-t font-medium [&>tr]:last:border-b-0',
        className,
      )}
      {...props}
    />
  )
}

function TableRow({ className, ...props }: React.ComponentProps<'tr'>) {
  return (
    <tr
      data-slot="table-row"
      className={cn(
        'hover:bg-muted/50 data-[state=selected]:bg-muted border-b transition-colors',
        className,
      )}
      {...props}
    />
  )
}

function TableHead({ className, ...props }: React.ComponentProps<'th'>) {
  return (
    <th
      data-slot="table-head"
      className={cn(
        'text-foreground h-10 px-2 text-left align-middle font-medium whitespace-nowrap [&:has([role=checkbox])]:pr-0 [&>[role=checkbox]]:translate-y-[2px]',
        className,
      )}
      {...props}
    />
  )
}

function TableCell({ className, ...props }: React.ComponentProps<'td'>) {
  return (
    <td
      data-slot="table-cell"
      className={cn(
        'p-2 align-middle whitespace-nowrap [&:has([role=checkbox])]:pr-0 [&>[role=checkbox]]:translate-y-[2px]',
        className,
      )}
      {...props}
    />
  )
}

function TableCaption({
  className,
  ...props
}: React.ComponentProps<'caption'>) {
  return (
    <caption
      data-slot="table-caption"
      className={cn('text-muted-foreground mt-4 text-sm', className)}
      {...props}
    />
  )
}

export {
  Table,
  TableHeader,
  TableBody,
  TableFooter,
  TableHead,
  TableRow,
  TableCell,
  TableCaption,
}
\end{DeskFlowCode}

\phantomsection
\addcontentsline{toc}{subsubsection}{deskflow\-frontend/components/ui/tabs.tsx}
\DeskFlowListingCaption{deskflow-frontend/components/ui/tabs.tsx}{lst:fd-ui-tabs}
\begin{DeskFlowCode}
'use client'

import * as React from 'react'
import * as TabsPrimitive from '@radix-ui/react-tabs'

import { cn } from '@/lib/utils'

function Tabs({
  className,
  ...props
}: React.ComponentProps<typeof TabsPrimitive.Root>) {
  return (
    <TabsPrimitive.Root
      data-slot="tabs"
      className={cn('flex flex-col gap-2', className)}
      {...props}
    />
  )
}

function TabsList({
  className,
  ...props
}: React.ComponentProps<typeof TabsPrimitive.List>) {
  return (
    <TabsPrimitive.List
      data-slot="tabs-list"
      className={cn(
        'bg-muted text-muted-foreground inline-flex h-9 w-fit items-center justify-center rounded-lg p-[3px]',
        className,
      )}
      {...props}
    />
  )
}

function TabsTrigger({
  className,
  ...props
}: React.ComponentProps<typeof TabsPrimitive.Trigger>) {
  return (
    <TabsPrimitive.Trigger
      data-slot="tabs-trigger"
      className={cn(
        "data-[state=active]:bg-background dark:data-[state=active]:text-foreground focus-visible:border-ring focus-visible:ring-ring/50 focus-visible:outline-ring dark:data-[state=active]:border-input dark:data-[state=active]:bg-input/30 text-foreground dark:text-muted-foreground inline-flex h-[calc(100%-1px)] flex-1 items-center justify-center gap-1.5 rounded-md border border-transparent px-2 py-1 text-sm font-medium whitespace-nowrap transition-[color,box-shadow] focus-visible:ring-[3px] focus-visible:outline-1 disabled:pointer-events-none disabled:opacity-50 data-[state=active]:shadow-sm [&_svg]:pointer-events-none [&_svg]:shrink-0 [&_svg:not([class*='size-'])]:size-4",
        className,
      )}
      {...props}
    />
  )
}

function TabsContent({
  className,
  ...props
}: React.ComponentProps<typeof TabsPrimitive.Content>) {
  return (
    <TabsPrimitive.Content
      data-slot="tabs-content"
      className={cn('flex-1 outline-none', className)}
      {...props}
    />
  )
}

export { Tabs, TabsList, TabsTrigger, TabsContent }
\end{DeskFlowCode}

\phantomsection
\addcontentsline{toc}{subsubsection}{deskflow\-frontend/components/ui/textarea.tsx}
\DeskFlowListingCaption{deskflow-frontend/components/ui/textarea.tsx}{lst:fd-ui-textarea}
\begin{DeskFlowCode}
import * as React from 'react'

import { cn } from '@/lib/utils'

function Textarea({ className, ...props }: React.ComponentProps<'textarea'>) {
  return (
    <textarea
      data-slot="textarea"
      className={cn(
        'border-input placeholder:text-muted-foreground focus-visible:border-ring focus-visible:ring-ring/50 aria-invalid:ring-destructive/20 dark:aria-invalid:ring-destructive/40 aria-invalid:border-destructive dark:bg-input/30 flex field-sizing-content min-h-16 w-full rounded-md border bg-transparent px-3 py-2 text-base shadow-xs transition-[color,box-shadow] outline-none focus-visible:ring-[3px] disabled:cursor-not-allowed disabled:opacity-50 md:text-sm',
        className,
      )}
      {...props}
    />
  )
}

export { Textarea }
\end{DeskFlowCode}

\phantomsection
\addcontentsline{toc}{subsubsection}{deskflow\-frontend/components/ui/toast.tsx}
\DeskFlowListingCaption{deskflow-frontend/components/ui/toast.tsx}{lst:fd-ui-toast}
\begin{DeskFlowCode}
'use client'

import * as React from 'react'
import * as ToastPrimitives from '@radix-ui/react-toast'
import { cva, type VariantProps } from 'class-variance-authority'
import { X } from 'lucide-react'

import { cn } from '@/lib/utils'

const ToastProvider = ToastPrimitives.Provider

const ToastViewport = React.forwardRef<
  React.ElementRef<typeof ToastPrimitives.Viewport>,
  React.ComponentPropsWithoutRef<typeof ToastPrimitives.Viewport>
>(({ className, ...props }, ref) => (
  <ToastPrimitives.Viewport
    ref={ref}
    className={cn(
      'fixed top-0 z-[100] flex max-h-screen w-full flex-col-reverse p-4 sm:bottom-0 sm:right-0 sm:top-auto sm:flex-col md:max-w-[420px]',
      className,
    )}
    {...props}
  />
))
ToastViewport.displayName = ToastPrimitives.Viewport.displayName

const toastVariants = cva(
  'group pointer-events-auto relative flex w-full items-center justify-between space-x-4 overflow-hidden rounded-md border p-6 pr-8 shadow-lg transition-all data-[swipe=cancel]:translate-x-0 data-[swipe=end]:translate-x-[var(-radix-toast-swipe-end-x)] data-[swipe=move]:translate-x-[var(-radix-toast-swipe-move-x)] data-[swipe=move]:transition-none data-[state=open]:animate-in data-[state=closed]:animate-out data-[swipe=end]:animate-out data-[state=closed]:fade-out-80 data-[state=closed]:slide-out-to-right-full data-[state=open]:slide-in-from-top-full data-[state=open]:sm:slide-in-from-bottom-full',
  {
    variants: {
      variant: {
        default: 'border bg-background text-foreground',
        destructive:
          'destructive group border-destructive bg-destructive text-destructive-foreground',
      },
    },
    defaultVariants: {
      variant: 'default',
    },
  },
)

const Toast = React.forwardRef<
  React.ElementRef<typeof ToastPrimitives.Root>,
  React.ComponentPropsWithoutRef<typeof ToastPrimitives.Root> &
    VariantProps<typeof toastVariants>
>(({ className, variant, ...props }, ref) => {
  return (
    <ToastPrimitives.Root
      ref={ref}
      className={cn(toastVariants({ variant }), className)}
      {...props}
    />
  )
})
Toast.displayName = ToastPrimitives.Root.displayName

const ToastAction = React.forwardRef<
  React.ElementRef<typeof ToastPrimitives.Action>,
  React.ComponentPropsWithoutRef<typeof ToastPrimitives.Action>
>(({ className, ...props }, ref) => (
  <ToastPrimitives.Action
    ref={ref}
    className={cn(
      'inline-flex h-8 shrink-0 items-center justify-center rounded-md border bg-transparent px-3 text-sm font-medium ring-offset-background transition-colors hover:bg-secondary focus:outline-none focus:ring-2 focus:ring-ring focus:ring-offset-2 disabled:pointer-events-none disabled:opacity-50 group-[.destructive]:border-muted/40 group-[.destructive]:hover:border-destructive/30 group-[.destructive]:hover:bg-destructive group-[.destructive]:hover:text-destructive-foreground group-[.destructive]:focus:ring-destructive',
      className,
    )}
    {...props}
  />
))
ToastAction.displayName = ToastPrimitives.Action.displayName

const ToastClose = React.forwardRef<
  React.ElementRef<typeof ToastPrimitives.Close>,
  React.ComponentPropsWithoutRef<typeof ToastPrimitives.Close>
>(({ className, ...props }, ref) => (
  <ToastPrimitives.Close
    ref={ref}
    className={cn(
      'absolute right-2 top-2 rounded-md p-1 text-foreground/50 opacity-0 transition-opacity hover:text-foreground focus:opacity-100 focus:outline-none focus:ring-2 group-hover:opacity-100 group-[.destructive]:text-red-300 group-[.destructive]:hover:text-red-50 group-[.destructive]:focus:ring-red-400 group-[.destructive]:focus:ring-offset-red-600',
      className,
    )}
    toast-close=""
    {...props}
  >
    <X className="h-4 w-4" />
  </ToastPrimitives.Close>
))
ToastClose.displayName = ToastPrimitives.Close.displayName

const ToastTitle = React.forwardRef<
  React.ElementRef<typeof ToastPrimitives.Title>,
  React.ComponentPropsWithoutRef<typeof ToastPrimitives.Title>
>(({ className, ...props }, ref) => (
  <ToastPrimitives.Title
    ref={ref}
    className={cn('text-sm font-semibold', className)}
    {...props}
  />
))
ToastTitle.displayName = ToastPrimitives.Title.displayName

const ToastDescription = React.forwardRef<
  React.ElementRef<typeof ToastPrimitives.Description>,
  React.ComponentPropsWithoutRef<typeof ToastPrimitives.Description>
>(({ className, ...props }, ref) => (
  <ToastPrimitives.Description
    ref={ref}
    className={cn('text-sm opacity-90', className)}
    {...props}
  />
))
ToastDescription.displayName = ToastPrimitives.Description.displayName

type ToastProps = React.ComponentPropsWithoutRef<typeof Toast>

type ToastActionElement = React.ReactElement<typeof ToastAction>

export {
  type ToastProps,
  type ToastActionElement,
  ToastProvider,
  ToastViewport,
  Toast,
  ToastTitle,
  ToastDescription,
  ToastClose,
  ToastAction,
}
\end{DeskFlowCode}

\phantomsection
\addcontentsline{toc}{subsubsection}{deskflow\-frontend/components/ui/toaster.tsx}
\DeskFlowListingCaption{deskflow-frontend/components/ui/toaster.tsx}{lst:fd-ui-toaster}
\begin{DeskFlowCode}
'use client'

import { useToast } from '@/hooks/use-toast'
import {
  Toast,
  ToastClose,
  ToastDescription,
  ToastProvider,
  ToastTitle,
  ToastViewport,
} from '@/components/ui/toast'

export function Toaster() {
  const { toasts } = useToast()

  return (
    <ToastProvider>
      {toasts.map(function ({ id, title, description, action, ...props }) {
        return (
          <Toast key={id} {...props}>
            <div className="grid gap-1">
              {title && <ToastTitle>{title}</ToastTitle>}
              {description && (
                <ToastDescription>{description}</ToastDescription>
              )}
            </div>
            {action}
            <ToastClose />
          </Toast>
        )
      })}
      <ToastViewport />
    </ToastProvider>
  )
}
\end{DeskFlowCode}

\phantomsection
\addcontentsline{toc}{subsubsection}{deskflow\-frontend/components/ui/toggle\-group.tsx}
\DeskFlowListingCaption{deskflow-frontend/components/ui/toggle-group.tsx}{lst:fd-ui-toggle-group}
\begin{DeskFlowCode}
'use client'

import * as React from 'react'
import * as ToggleGroupPrimitive from '@radix-ui/react-toggle-group'
import { type VariantProps } from 'class-variance-authority'

import { cn } from '@/lib/utils'
import { toggleVariants } from '@/components/ui/toggle'

const ToggleGroupContext = React.createContext<
  VariantProps<typeof toggleVariants>
>({
  size: 'default',
  variant: 'default',
})

function ToggleGroup({
  className,
  variant,
  size,
  children,
  ...props
}: React.ComponentProps<typeof ToggleGroupPrimitive.Root> &
  VariantProps<typeof toggleVariants>) {
  return (
    <ToggleGroupPrimitive.Root
      data-slot="toggle-group"
      data-variant={variant}
      data-size={size}
      className={cn(
        'group/toggle-group flex w-fit items-center rounded-md data-[variant=outline]:shadow-xs',
        className,
      )}
      {...props}
    >
      <ToggleGroupContext.Provider value={{ variant, size }}>
        {children}
      </ToggleGroupContext.Provider>
    </ToggleGroupPrimitive.Root>
  )
}

function ToggleGroupItem({
  className,
  children,
  variant,
  size,
  ...props
}: React.ComponentProps<typeof ToggleGroupPrimitive.Item> &
  VariantProps<typeof toggleVariants>) {
  const context = React.useContext(ToggleGroupContext)

  return (
    <ToggleGroupPrimitive.Item
      data-slot="toggle-group-item"
      data-variant={context.variant || variant}
      data-size={context.size || size}
      className={cn(
        toggleVariants({
          variant: context.variant || variant,
          size: context.size || size,
        }),
        'min-w-0 flex-1 shrink-0 rounded-none shadow-none first:rounded-l-md last:rounded-r-md focus:z-10 focus-visible:z-10 data-[variant=outline]:border-l-0 data-[variant=outline]:first:border-l',
        className,
      )}
      {...props}
    >
      {children}
    </ToggleGroupPrimitive.Item>
  )
}

export { ToggleGroup, ToggleGroupItem }
\end{DeskFlowCode}

\phantomsection
\addcontentsline{toc}{subsubsection}{deskflow\-frontend/components/ui/toggle.tsx}
\DeskFlowListingCaption{deskflow-frontend/components/ui/toggle.tsx}{lst:fd-ui-toggle}
\begin{DeskFlowCode}
'use client'

import * as React from 'react'
import * as TogglePrimitive from '@radix-ui/react-toggle'
import { cva, type VariantProps } from 'class-variance-authority'

import { cn } from '@/lib/utils'

const toggleVariants = cva(
  "inline-flex items-center justify-center gap-2 rounded-md text-sm font-medium hover:bg-muted hover:text-muted-foreground disabled:pointer-events-none disabled:opacity-50 data-[state=on]:bg-accent data-[state=on]:text-accent-foreground [&_svg]:pointer-events-none [&_svg:not([class*='size-'])]:size-4 [&_svg]:shrink-0 focus-visible:border-ring focus-visible:ring-ring/50 focus-visible:ring-[3px] outline-none transition-[color,box-shadow] aria-invalid:ring-destructive/20 dark:aria-invalid:ring-destructive/40 aria-invalid:border-destructive whitespace-nowrap",
  {
    variants: {
      variant: {
        default: 'bg-transparent',
        outline:
          'border border-input bg-transparent shadow-xs hover:bg-accent hover:text-accent-foreground',
      },
      size: {
        default: 'h-9 px-2 min-w-9',
        sm: 'h-8 px-1.5 min-w-8',
        lg: 'h-10 px-2.5 min-w-10',
      },
    },
    defaultVariants: {
      variant: 'default',
      size: 'default',
    },
  },
)

function Toggle({
  className,
  variant,
  size,
  ...props
}: React.ComponentProps<typeof TogglePrimitive.Root> &
  VariantProps<typeof toggleVariants>) {
  return (
    <TogglePrimitive.Root
      data-slot="toggle"
      className={cn(toggleVariants({ variant, size, className }))}
      {...props}
    />
  )
}

export { Toggle, toggleVariants }
\end{DeskFlowCode}

\phantomsection
\addcontentsline{toc}{subsubsection}{deskflow\-frontend/components/ui/tooltip.tsx}
\DeskFlowListingCaption{deskflow-frontend/components/ui/tooltip.tsx}{lst:fd-ui-tooltip}
\begin{DeskFlowCode}
'use client'

import * as React from 'react'
import * as TooltipPrimitive from '@radix-ui/react-tooltip'

import { cn } from '@/lib/utils'

function TooltipProvider({
  delayDuration = 0,
  ...props
}: React.ComponentProps<typeof TooltipPrimitive.Provider>) {
  return (
    <TooltipPrimitive.Provider
      data-slot="tooltip-provider"
      delayDuration={delayDuration}
      {...props}
    />
  )
}

function Tooltip({
  ...props
}: React.ComponentProps<typeof TooltipPrimitive.Root>) {
  return (
    <TooltipProvider>
      <TooltipPrimitive.Root data-slot="tooltip" {...props} />
    </TooltipProvider>
  )
}

function TooltipTrigger({
  ...props
}: React.ComponentProps<typeof TooltipPrimitive.Trigger>) {
  return <TooltipPrimitive.Trigger data-slot="tooltip-trigger" {...props} />
}

function TooltipContent({
  className,
  sideOffset = 0,
  children,
  ...props
}: React.ComponentProps<typeof TooltipPrimitive.Content>) {
  return (
    <TooltipPrimitive.Portal>
      <TooltipPrimitive.Content
        data-slot="tooltip-content"
        sideOffset={sideOffset}
        className={cn(
          'bg-foreground text-background animate-in fade-in-0 zoom-in-95 data-[state=closed]:animate-out data-[state=closed]:fade-out-0 data-[state=closed]:zoom-out-95 data-[side=bottom]:slide-in-from-top-2 data-[side=left]:slide-in-from-right-2 data-[side=right]:slide-in-from-left-2 data-[side=top]:slide-in-from-bottom-2 z-50 w-fit origin-(-radix-tooltip-content-transform-origin) rounded-md px-3 py-1.5 text-xs text-balance',
          className,
        )}
        {...props}
      >
        {children}
        <TooltipPrimitive.Arrow className="bg-foreground fill-foreground z-50 size-2.5 translate-y-[calc(-50%_-_2px)] rotate-45 rounded-[2px]" />
      </TooltipPrimitive.Content>
    </TooltipPrimitive.Portal>
  )
}

export { Tooltip, TooltipTrigger, TooltipContent, TooltipProvider }
\end{DeskFlowCode}

\phantomsection
\addcontentsline{toc}{subsubsection}{deskflow\-frontend/components/ui/use\-mobile.tsx}
\DeskFlowListingCaption{deskflow-frontend/components/ui/use-mobile.tsx}{lst:fd-ui-use-mobile}
\begin{DeskFlowCode}
import * as React from 'react'

const MOBILE_BREAKPOINT = 768

export function useIsMobile() {
  const [isMobile, setIsMobile] = React.useState<boolean | undefined>(undefined)

  React.useEffect(() => {
    const mql = window.matchMedia(`(max-width: ${MOBILE_BREAKPOINT - 1}px)`)
    const onChange = () => {
      setIsMobile(window.innerWidth < MOBILE_BREAKPOINT)
    }
    mql.addEventListener('change', onChange)
    setIsMobile(window.innerWidth < MOBILE_BREAKPOINT)
    return () => mql.removeEventListener('change', onChange)
  }, [])

  return !!isMobile
}
\end{DeskFlowCode}

\phantomsection
\addcontentsline{toc}{subsubsection}{deskflow\-frontend/components/ui/use\-toast.ts}
\DeskFlowListingCaption{deskflow-frontend/components/ui/use-toast.ts}{lst:fd-ui-use-toast}
\begin{DeskFlowCode}
'use client'

// Inspired by react-hot-toast library
import * as React from 'react'

import type { ToastActionElement, ToastProps } from '@/components/ui/toast'

const TOAST_LIMIT = 1
const TOAST_REMOVE_DELAY = 1000000

type ToasterToast = ToastProps & {
  id: string
  title?: React.ReactNode
  description?: React.ReactNode
  action?: ToastActionElement
}

const actionTypes = {
  ADD_TOAST: 'ADD_TOAST',
  UPDATE_TOAST: 'UPDATE_TOAST',
  DISMISS_TOAST: 'DISMISS_TOAST',
  REMOVE_TOAST: 'REMOVE_TOAST',
} as const

let count = 0

function genId() {
  count = (count + 1) % Number.MAX_SAFE_INTEGER
  return count.toString()
}

type ActionType = typeof actionTypes

type Action =
  | {
      type: ActionType['ADD_TOAST']
      toast: ToasterToast
    }
  | {
      type: ActionType['UPDATE_TOAST']
      toast: Partial<ToasterToast>
    }
  | {
      type: ActionType['DISMISS_TOAST']
      toastId?: ToasterToast['id']
    }
  | {
      type: ActionType['REMOVE_TOAST']
      toastId?: ToasterToast['id']
    }

interface State {
  toasts: ToasterToast[]
}

const toastTimeouts = new Map<string, ReturnType<typeof setTimeout>>()

const addToRemoveQueue = (toastId: string) => {
  if (toastTimeouts.has(toastId)) {
    return
  }

  const timeout = setTimeout(() => {
    toastTimeouts.delete(toastId)
    dispatch({
      type: 'REMOVE_TOAST',
      toastId: toastId,
    })
  }, TOAST_REMOVE_DELAY)

  toastTimeouts.set(toastId, timeout)
}

export const reducer = (state: State, action: Action): State => {
  switch (action.type) {
    case 'ADD_TOAST':
      return {
        ...state,
        toasts: [action.toast, ...state.toasts].slice(0, TOAST_LIMIT),
      }

    case 'UPDATE_TOAST':
      return {
        ...state,
        toasts: state.toasts.map((t) =>
          t.id === action.toast.id ? { ...t, ...action.toast } : t,
        ),
      }

    case 'DISMISS_TOAST': {
      const { toastId } = action

      // ! Side effects ! - This could be extracted into a dismissToast() action,
      // but I'll keep it here for simplicity
      if (toastId) {
        addToRemoveQueue(toastId)
      } else {
        state.toasts.forEach((toast) => {
          addToRemoveQueue(toast.id)
        })
      }

      return {
        ...state,
        toasts: state.toasts.map((t) =>
          t.id === toastId || toastId === undefined
            ? {
                ...t,
                open: false,
              }
            : t,
        ),
      }
    }
    case 'REMOVE_TOAST':
      if (action.toastId === undefined) {
        return {
          ...state,
          toasts: [],
        }
      }
      return {
        ...state,
        toasts: state.toasts.filter((t) => t.id !== action.toastId),
      }
  }
}

const listeners: Array<(state: State) => void> = []

let memoryState: State = { toasts: [] }

function dispatch(action: Action) {
  memoryState = reducer(memoryState, action)
  listeners.forEach((listener) => {
    listener(memoryState)
  })
}

type Toast = Omit<ToasterToast, 'id'>

function toast({ ...props }: Toast) {
  const id = genId()

  const update = (props: ToasterToast) =>
    dispatch({
      type: 'UPDATE_TOAST',
      toast: { ...props, id },
    })
  const dismiss = () => dispatch({ type: 'DISMISS_TOAST', toastId: id })

  dispatch({
    type: 'ADD_TOAST',
    toast: {
      ...props,
      id,
      open: true,
      onOpenChange: (open) => {
        if (!open) dismiss()
      },
    },
  })

  return {
    id: id,
    dismiss,
    update,
  }
}

function useToast() {
  const [state, setState] = React.useState<State>(memoryState)

  React.useEffect(() => {
    listeners.push(setState)
    return () => {
      const index = listeners.indexOf(setState)
      if (index > -1) {
        listeners.splice(index, 1)
      }
    }
  }, [state])

  return {
    ...state,
    toast,
    dismiss: (toastId?: string) => dispatch({ type: 'DISMISS_TOAST', toastId }),
  }
}

export { useToast, toast }
\end{DeskFlowCode}
\end{landscape}
